\documentclass[11pt]{article}

\usepackage[utf8]{inputenc}
\usepackage[T1]{fontenc}
\usepackage{lmodern}
\usepackage{booktabs}
\usepackage{enumitem}
\usepackage{amsmath, amssymb, amsthm}
\usepackage[margin=1in]{geometry}
\usepackage{microtype}
\PassOptionsToPackage{hyphens}{url}
\usepackage{xcolor}
\usepackage{hyperref}
\hypersetup{
  colorlinks=true, linkcolor=blue!50!black, urlcolor=blue!50!black, citecolor=blue!50!black,
  pdftitle={Transcendence of a planar reflection-group growth series and its relaxed companion},
  pdfauthor={Vico Bonfioli},
}

\theoremstyle{plain}
\newtheorem{theorem}{Theorem}
\newtheorem{lemma}[theorem]{Lemma}
\newtheorem{corollary}[theorem]{Corollary}
\newtheorem{proposition}[theorem]{Proposition}
\theoremstyle{definition}
\newtheorem{definition}[theorem]{Definition}
\newtheorem{hypothesis}[theorem]{Hypothesis}
\theoremstyle{remark}
\newtheorem{remark}[theorem]{Remark}

\newcommand{\Z}{\mathbb{Z}}
\newcommand{\Q}{\mathbb{Q}}
\newcommand{\R}{\mathbb{R}}
\newcommand{\C}{\mathbb{C}}
\newcommand{\F}{\mathbb{F}}
\renewcommand{\Re}{\operatorname{Re}}

\title{Transcendence of a planar reflection--group growth series\\ and its relaxed companion}
\author{Vico Bonfioli\thanks{Independent researcher. \texttt{vicobonfioli@gmail.com}.
Code, data, and the Lean~4 formalisation of \S\ref{sec:sitecost}, of
Proposition~\ref{prop:mobius} and of Lemma~\ref{lem:atomN}:
\url{https://github.com/ElVec1o/pythagorean-reflection-sequence}.}}
\date{2026}

\begin{document}
\maketitle

\begin{abstract}
\noindent Let $W_{\mathrm{univ}}=\langle R_0,R_1,R_2\rangle$ be the universal right-triangle
reflection group, the group generated by the reflections in the sides of a right triangle with
generic (transcendental) angle parameter, to which every right triangle group agrees up to a finite,
effectively computable depth \cite{Bonfioli2026a}. Let $U(x)=\sum_{n\ge0}u_nx^n$ be its geodesic
word-length growth series (the integer sequence \textup{OEIS~A396406} \cite{OEIS}), with
relaxed companion $V$. We prove:
\begin{enumerate}[leftmargin=1.9em,itemsep=0.12em,topsep=0.25em]
\item[\textup{(i)}] each catalytic building-block $q$-series $\Sigma_0,\Sigma_1,S_0,S_1$ is, by the
P\'olya--Carlson theorem, either a rational function whose poles are roots of unity or has
$|q|=1$ as a natural boundary and is transcendental over $\Q(q)$. For the two denominators
$\Sigma_1,S_1$ the first alternative is excluded with no growth hypothesis, by the infinitude of
their level sets $\Sigma_1=1$, $S_1=1$ in $(0,1)$ against the finiteness of the zero set of a
nonzero rational function (Corollary~\ref{cor:polepath}); that route consumes the same single
analytic input as everything else here, Lemma~\ref{lem:T2abs}. For the two numerators
$\Sigma_0,S_0$ the exclusion is instead under a coefficient-growth hypothesis, verified exactly for
the first $34$ square indices for $\Sigma_1$ and for all four blocks in the uniform form of
Hypothesis~\ref{hyp:growth} by exact integer arithmetic out to degree $1156$, and not proved
(Theorem~\ref{thm:blocks}, Remark~\ref{rem:growthstatus}); what the pole route reaches for the
numerators is the joint statement that $\Sigma_0$ and $S_0$ are not both rational;
\item[\textup{(ii)}] the common growth rate is $\beta_2=1.4916177871\ldots$, and
$\tfrac32$ is excluded;
\item[\textup{(iii)}] the travel resolvent $\Sigma_0/(1-\Sigma_1)$ is transcendental over $\Q(x)$,
unconditionally (Theorem~\ref{thm:travelres}), on three inputs all proved here: the pole set, from
the rectangle bound of Lemma~\ref{lem:T2abs}; the non-vanishing of the denominator block at the
travel poles, $|S_e(q_m)|\ge0.35\sqrt{\tau_m}$, by an elementary explicit chain with all constants
numerical and no cited asymptotic theorem (Appendix~\ref{app:annulus}); and the non-vanishing of
both numerators, from the exact identity $\Sigma_0(q_m)S_0(q_m)=2q_m/(1-q_m)$
(Proposition~\ref{prop:numnonvanish}), a consequence of the $q$-Pythagorean identity at a zero of
the $q$-cosine. The travel poles are the zeros of the Hahn--Exton $q$-Bessel function
$J^{(3)}_{-1/2}(Z(q);q^2)$, $Z=\sqrt{2(1-q)/q}$, which up to a positive factor is the Fredholm
determinant $\det(I-\mathcal T)$ of the travel transfer operator (Theorem~\ref{thm:fredholm}). The block assembly itself is developed in \S\ref{sec:assembly}: the travel and the
bulk block are one resolvent of the semiseparable kernel $q^{\max(a,b)}$ seen twice, one
three-term recursion covers both, and this proves the travel telescoping
$G^T_0=\Sigma_0/(1-\Sigma_1)$, the bulk dressing $b_0=S_0/(1-S_1)$, the reciprocity $t_0=b_1$
(equivalently $\det P=1$), the identity $t_1=P_{12}/S_e$, the M\"obius factorisation of the
gap-marked bulk resolvent, and the \emph{shape} of the local form at a travel pole: its leading
numerator is the square of one junction pairing $\langle\lambda,R\rangle$, which has an exact closed
form in $t_1$ and is positive at every travel pole (Theorem~\ref{thm:RJ}; analytic for
$\tau\le5\cdot10^{-3}$, a certificate at the first twelve poles). Hence $V$ and $U$ are
transcendental over $\Q(x)$ (Theorems~\ref{thm:V}, \ref{thm:U}). More holds: both are meromorphic on
$|x|<1$ with infinitely many poles there, so they are not $D$-finite ($u_n$ and $v_n$ are not
P-recursive), satisfy no linear $q$-difference equation with $Q$ not a root of unity and no Mahler
equation, and, by the P\'olya--Bertrandias theorem, have $|x|=1$ as a natural boundary
(Theorems~\ref{thm:nonDfinite}, \ref{thm:natboundary}); for the travel resolvent the same holds
unconditionally, with natural boundary $|q|=1$. The invariance of the travel block
\textup{(T)} is Proposition~\ref{prop:travelinv}, and the former hypothesis \textup{(L)} is a
theorem. All of this rests on the
strand-walk transfer model \textup{(M)}. Its local cost law is now a theorem inside the crossing
optimisation \textup{(}Lemma~\ref{lem:transport}, Corollary~\ref{cor:localcost}\textup{)}, and its
marker clause is corrected: with $d_L$ the last bulk deposit and $d_R$ the first travel deposit, the
junction cost is $\max(|d_L-1|,|d_R|)$, which depends on the sign of $d_L$ and not on its magnitude
alone; the form $\max(|d_L|-1,|d_R|)$ carried by earlier versions is false on every cell with
$d_L\le-2$ and $|d_R|\le|d_L|$ \textup{(}Remark~\ref{rem:markerfalse}\textup{)}. The transfer
model \textup{(M)} is a theorem (Theorem~\ref{thm:model}): the metric identity and the faithfulness
of the model to $W_{\mathrm{univ}}$ are formalised in Lean, and the block decomposition
\textup{(M3)}, corrected here, has a hand proof (Appendix~\ref{app:M3prime}) checked against exact
enumeration to $x^{26}$. The amplitude
estimate $(\star)$ ($|P_{12}(q_m)|w_m^3<2$ at every travel pole) is proved: the gate block is the symmetric second
difference of the $q$-cosine at its own zero, a two-interval uniform certificate bounds it by
$0.619\,\tau^{3/2}<\tau^{3/2}/\sqrt2$ for every pole with $\tau\le5\cdot10^{-3}$, and the six
poles above that threshold are checked directly at $120$-digit precision. An exact identity,
$S_e=\tfrac{qZ}2\mathfrak s(Z)-P_{12}$, turns the same certificate into the denominator bound
$|S_e|\ge0.63\sqrt\tau$, so the gate closes at every travel pole. The
turning-point
amplitude atom on the $U$-side, long the special-function obstacle, is an explicit rational inequality
\emph{formally verified in Lean~4 / Mathlib \cite{MathlibCommunity}} (Lemma~\ref{lem:atomN}), as is
the whole site cost chain of the model, Lemma~\ref{lem:transport} with
Corollaries~\ref{cor:localcost}, \ref{cor:lRclosed} and~\ref{cor:marker}
(\texttt{SiteCost.lean}, \texttt{MarkedSite.lean}, \texttt{PairingMatrix.lean},
\texttt{Realisation.lean}, \texttt{CutCross.lean});
\item[\textup{(iv)}] an orthogonal, analysis-free route via Christol's theorem would settle $U$
without any of the above: the $p$-kernel of $(u_n\bmod p)$ is computed to be maximally
non-automatic at $p=3,5$ and machine-checked at $p=3$, though a proof for dense-support series is
open.
\end{enumerate}
Every analytic input consumed by $V$ and by $U$ is proved here, and so is the algebra of the
block assembly; the combinatorial assembly \textup{(M3)} is a hand proof, not formalised, and two
finite ranges rest on computer-assisted interval certificates. On $U$ every coefficient is derived
elementarily
and machine-checked, the remainders are at the standard simple-saddle level, and the gate ($(\star)$ and the denominator bound) is proved, and with it the junction pairing at every travel pole
(analytically for $\tau\le5\cdot10^{-3}$, by certificate at the first twelve poles). The precise rigor standing, including
the arithmetic model of the committed certificates, is stated in the text. In sharp contrast with the rational growth of
hyperbolic groups and the algebraic growth series of Parry's wreath products, the geometric word
metric places this virtually-lamplighter group outside every algebraic class, and outside the
$D$-finite one.
\end{abstract}

\section{Introduction}\label{sec:intro}

A right triangle $T$ has three sides; the reflections across them generate a group
$W_T=\langle R_0,R_1,R_2\rangle\le\mathrm{Aff}(\R^2)$. For a right triangle whose angle
parameter $\zeta_T$ is generic, $W_T$ is the \emph{universal} group $W_{\mathrm{univ}}$ of
\cite{Bonfioli2026a}, and every other right triangle group agrees with $W_{\mathrm{univ}}$ in the
ball of some finite radius and deviates beyond it; the radius is effective
\cite[Theorem ``Effective universality'', \textup{thm:effective-universality}]{Bonfioli2026a}
\textup{(}for the $(1,2)$ triangle the first deviation is at depth $33$,
\cite[Theorem ``Sharp uniform universality at depth $32$'', \textup{thm:universality-sharp}]{Bonfioli2026a}\textup{)}.
Irrationality of the angles alone
\textup{(}Niven \cite{Niven1956}\textup{)} does not give genericity: a triangle in $\Q^2$ with
unequal legs has irrational angles, yet it is not universal. Throughout, $U(x)=\sum_nu_nx^n$ is the
growth series of $W_{\mathrm{univ}}$, which counts elements by minimal word length in
$\{R_0,R_1,R_2\}$. $W_{\mathrm{univ}}$ is not a Coxeter group \cite{IsolaPiergallini}, and, unlike in
the rational Coxeter case governed by Steinberg's formula \cite{StanleyEC2}, its growth is not a
priori rational. The sequence $u_n=1,3,5,8,13,21,34,55,89,\dots$ is \textup{OEIS~A396406}; it
coincides with the Fibonacci numbers $F_{n+3}$ for $1\le n\le9$ and first deviates at $n=10$. This
paper settles the arithmetic nature of the travel resolvent $\Sigma_0/(1-\Sigma_1)$, which is
transcendental over $\Q(x)$ with every input proved self-contained
(Theorem~\ref{thm:travelres}), and proves that both growth series are transcendental over $\Q(x)$
(Theorems~\ref{thm:V}, \ref{thm:U}): the junction non-vanishing \textup{(R-J)} is
Theorem~\ref{thm:RJ}, the travel-block invariance \textup{(T)} is
Proposition~\ref{prop:travelinv}, and the transfer model \textup{(M)} is Theorem~\ref{thm:model},
whose assembly step is a hand proof (Appendix~\ref{app:M3prime}) and is not formalised. The algebraic layer between the blocks and those two series, which
earlier versions of this paper carried as the hypotheses \textup{(R)} and \textup{(L)}, is proved
in \S\ref{sec:assembly}, and so is the local cost law of \textup{(M)}, together with the correction
of its marker clause (\S\ref{sec:sitecost}). The analytic side is closed on both: the amplitude estimate $(\star)$ and the
gate (numerator bound and denominator lower bound) are proved at every travel pole in
Appendix~\ref{app:star} (Theorem~\ref{thm:star}, Lemma~\ref{app:Se},
Corollary~\ref{app:gate}), and the two numerator non-vanishing statements the resolvents need are
proved in \S\ref{sec:numerators}.
Section~\ref{sec:beyond} goes past algebraicity. $U$ and $V$ continue meromorphically to $|x|<1$
(Lemma~\ref{lem:mero}), and their infinitely many poles there make them non-$D$-finite, so that
$u_n$ and $v_n$ are not P-recursive (Theorem~\ref{thm:nonDfinite}); they exclude linear
$q$-difference equations with $Q$ not a root of unity and Mahler equations
(Proposition~\ref{prop:noqdiff}); and, with the P\'olya--Bertrandias rationality theorem, they make
$|x|=1$ a natural boundary (Theorem~\ref{thm:natboundary}). For the travel resolvent these
statements are unconditional, with natural boundary $|q|=1$. $U$ also has a pole at $-x_m$ for
every travel pole (Proposition~\ref{prop:minusx}).

\paragraph{Context.} For word-growth series of groups, rationality is the norm in the classical
families: hyperbolic groups have rational growth series with respect to every finite generating set
(Cannon \cite{Cannon1984}), Coxeter systems are governed by Steinberg's rational formula \cite[\S5.12]{Humphreys1990}, and for
wreath products (including lamplighter-type groups) Parry computed the growth series for
natural generating sets and found them \emph{algebraic} \cite{Parry1992}. Provably transcendental
growth series are rare. The classical example is Stoll's \cite{Stoll1996}: the higher Heisenberg
group $H_5$ has a transcendental growth series with respect to its standard generating set, and a
rational one with respect to a specially chosen generating set; the three-dimensional Heisenberg
group, by contrast, has rational growth for every finite generating set (Duchin--Shapiro
\cite{DuchinShapiro2019}). Non-$D$-finite growth series are known as well. For groups of polynomial
growth, and for groups of intermediate growth such as Grigorchuk's \cite{Grigorchuk1984}, the growth
series has radius of convergence $1$ and integer coefficients, so by P\'olya--Carlson
\cite{Carlson1921} it is rational or has $|x|=1$ as a natural boundary; in the second case, which
includes Stoll's series for $H_5$ and every group of intermediate growth, it is not $D$-finite. And a
finitely presented group with unsolvable word problem has a non-computable growth function, while
a $D$-finite series with rational coefficients has computable coefficients, so its growth series is
not $D$-finite either. The group studied here is virtually the lamplighter
$\Z\wr\Z$ \cite{Bonfioli2026a}, and its generating set is not constructed but imposed by Euclidean
geometry; the three side-reflections of a generic right triangle. The results below show that for
this natural geometric system the growth series is transcendental and not $D$-finite, by an explicit
mechanism: infinitely many poles inside the disc $|x|<1$ of meromorphy, beyond the radius of
convergence $1/\beta_2$, together with a natural boundary on $|x|=1$ (\S\ref{sec:beyond}). For a
group generated by reflections this is the opposite of the Coxeter case, where Steinberg's formula
makes the growth series rational. It also places the example in sharp contrast with the
algebraicity of Parry's wreath-product series: the geometric word metric moves a
virtually-lamplighter group outside every algebraic class. The cogrowth series is settled by
general theory: $W_{\mathrm{univ}}$ is amenable, being virtually metabelian, and not virtually
nilpotent, so its cogrowth series is not $D$-finite by Bell--Mishna \cite{BellMishna2020}. The
higher-dimensional picture is subtler than a clean dimensional dichotomy \cite{Bonfioli2026c}: the
abstract right-angled Coxeter envelope of an $n$-orthoscheme has rational growth (rate
$r_n=1+2\cos\frac{2\pi}{n+3}$), but the geometric group is an amenable proper quotient that matches
this rational envelope only to a finite depth before deviating below it. The plane is the one
dimension in which the deviation has been computed past its onset; there it is the transcendental
series treated here; and whether every higher dimension is likewise eventually transcendental is
an open problem.

\paragraph{The relaxed/true dichotomy.} The word length $\ell_T(g)$ is realised by strand walks in
a finite transition system; permitting free isolated cycles yields a smaller, analytically tamer
\emph{relaxed} length $\ell_R(g)$, with $\ell_T=\ell_R+2c(g)$, $c(g)\in\Z_{\ge0}$ a cycle count.
The inequality $\ell_T\le\ell_R+2c$ is proved in \cite{Bonfioli2026a}. With $\ell_R$ and $c$ the
site-local closed forms of that paper, the identity $\ell_T=\ell_R+2c$ holds for every element and
is formalised in Lean (\texttt{metricAll}; \S\ref{sec:bridge}); that the closed form $\ell_R$ is
also the minimum over relaxed realisations is verified there, and bears only on reading $V$ as a
relaxed count.
The relaxed counting function is the companion $V$; one has $v_n\ge u_n$, equality for $n\le7$,
first strict at $n=8$. Both are governed by a linear $q$-difference (dilation $t\mapsto q^2t$)
\emph{catalytic functional equation} (\S\ref{sec:setup}) that is $q$-holonomic and not of
polynomial-kernel type.

\paragraph{Method.} The catalytic transfer is assembled from four integer $q$-series
$\Sigma_0,\Sigma_1,S_0,S_1$ with $j^2$-gapped coefficients. P\'olya--Carlson gives a dichotomy for
them (rational with roots-of-unity poles, or a natural boundary and transcendence over $\Q(q)$).
For the two denominators the rational alternative is excluded outright: their level sets
$\Sigma_1=1$ and $S_1=1$ in $(0,1)$ are infinite, and a rational function that is not identically
zero has finitely many zeros in $\C$ (Corollary~\ref{cor:polepath}). For the two numerators it is
excluded by super-polynomial coefficient growth along a sparse set of
indices, verified for the first $34$ square indices for $\Sigma_1$ and out to degree $1156$ for all
four blocks, and not proved (\S\ref{sec:blocks}).
Both $V$ and $U$ acquire poles at the \emph{travel poles} $q_m$, the roots of
$\Sigma_1^T(q)=1$; transcendence follows from accumulation of these poles at $q=1$ once a certain
auxiliary integral $T_2$ is shown to be small (\S\ref{sec:sd}). The travel poles are the zeros of
the Hahn--Exton $q$-Bessel function $J^{(3)}_{-1/2}(Z(q);q^2)$, $Z=\sqrt{2(1-q)/q}$, because that
function is, up to a positive factor, the Fredholm determinant $\det(I-\mathcal T)$ of the travel
transfer operator (Theorem~\ref{thm:fredholm}, Corollary~\ref{cor:qbessel}). What the theorems consume is only
$|T_2|\to0$, and that is proved self-contained in Lemma~\ref{lem:T2abs} by an absolute bound on a
rectangle: no saddle, no path through a saddle, no cited asymptotic theorem. Its quantitative part
is a verified enclosure of the contour majorant, in interval arithmetic and quantified over
parameter boxes, on $\tau\in[2/101^2,0.02]$, together with an explicit-constant estimate below
that; nothing there is a scan over sample values of $\tau$. It is the single
analytic input of this kind, shared by $V$ (\S\ref{sec:V}) and $U$ (\S\ref{sec:U}). The sharp size
of $T_2$, with the constant $\sqrt2/36$, is recorded separately in Remark~\ref{lem:cos} as a
numerically certain statement whose proof is incomplete (Remark~\ref{rem:cosgap}); nothing here
depends on it.

The $U$-side carries one extra ingredient. Writing the cocycle off-diagonal $P_{12}$ in exact
closed form as a Hahn--Exton $q$-Bessel function $J^{(3)}_{3/2}$, the transcendence gate reduces to
one amplitude bound on that series at the travel poles. This bound is \emph{proved}
(Appendix~\ref{app:star}): $P_{12}$ is the symmetric second difference of the $q$-cosine at its own
zero (Proposition~\ref{prop:P12second}), so the gate is a second-derivative estimate on a window of
width $\tau$, delivered by a two-interval uniform certificate; an exact identity then converts the
same certificate into the denominator lower bound, closing the gate
(Corollary~\ref{app:gate}). The assembled residue at a travel pole is the square of a junction
pairing $\langle\lambda,R\rangle$ (Proposition~\ref{prop:shape}), which has an exact closed form in
$t_1$; the gate makes it positive at every pole with $\tau\le5\cdot10^{-3}$, and a certificate covers
the first twelve (Theorem~\ref{thm:RJ}). The invariance of the travel block under the cycle marker,
\textup{(T)}, is proved (Proposition~\ref{prop:travelinv}). The sharp constants behind the picture are \emph{derived} in
Remark~\ref{rem:elemY3}: the pole--zero offset $\delta w=\tfrac{\sqrt2}8\tau^{3/2}$ and
$|P_{12}|/\tau^{3/2}\to1/(4\sqrt2)$, a fourfold margin below the gate threshold. (The
turning-point atom feeding that analysis, the rational envelope
$h(X)=(1+\tfrac3{2X^2})/(1+\tfrac1{X^2})\in[1,\tfrac32]$ of Lemma~\ref{lem:atomN}, is formally
verified in Lean~4.) $U$ and $V$ therefore rest on the same self-contained \emph{analysis}; no
special-function input separates them, and after \S\ref{sec:assembly} they share the junction
pairing \textup{(R-J)}, now Theorem~\ref{thm:RJ}, and the transfer model \textup{(M)}, now
Theorem~\ref{thm:model}.

Section~\ref{sec:modp} records an independent arithmetic route: by Christol's theorem
\cite{Christol1979,CKMR1980}, if
$(u_n\bmod p)$ is not $p$-automatic for one prime $p$ then $U$ is transcendental over $\Q(x)$; the
$p$-kernel is maximally non-automatic at $p=3,5$.

Appendix~\ref{app:beta2} treats a question separate from all of the above: whether the \emph{number}
$\beta_2$ is transcendental. That is open. What the appendix contributes is a reduction of it to a
value statement for a Hahn--Exton $q$-Bessel function at an algebraic base, and an exact location of
the obstruction, at the arithmetic ratio $\rho=\tfrac12$, one notch beyond the range in which
$q$-irrationality criteria are proved. No statement in the body depends on it.

\section{The catalytic equation and its building blocks}\label{sec:setup}

Write $q=x^2=e^{-\tau}$, $w=\sqrt{2/\tau}$. The relaxed length statistic satisfies a linear
$q$-difference equation with dilation $t\mapsto q^2t$; telescoping it produces four integer power
series in $q$,
\begin{equation}\label{eq:blocks}
  \Sigma_1(q)=\sum_{i\ge1}(-1)^{i-1}\frac{w^{2i}}{(2i)!}\,\eta_i,\quad
  S_e(q)=\sum_{j\ge0}\frac{(-2(1-q))^j q^{j(j+1)}}{(q;q)_{2j}},
\end{equation}
and their companions $\Sigma_0,S_0$, where $\eta_j=e^{-B_j}$ is the Gaussian-decaying form factor
of \S\ref{sec:sd} (equivalently $\eta_j=q^{j^2+j}[2(1-q)]^j(2j)!/[(q;q)_{2j}W^{2j}]$; the
identification of this display with $1-\Sigma_1=\cos(Z;q^2)$ of Proposition~\ref{prop:qtrig} is the
collapse used in Lemma~\ref{lem:infpoles}). We write $\Sigma_1^{\mathrm T}=\Sigma_1$ when the travel
r\^ole is to be emphasised. The relaxed travel resolvent is $G^T_0=\Sigma_0/(1-\Sigma_1)$
(Proposition~\ref{prop:travelsum}), and $U$ inherits its denominator $1-\Sigma_1$
(Proposition~\ref{prop:travelinv}); how the two blocks assemble into $V$ and $U$, and how far that
assembly is proved, is \S\ref{sec:assembly}. The \emph{travel poles}
\[
   q_1<q_2<\cdots\uparrow1,\qquad \Sigma_1^T(q_m)=1,
\]
form an infinite set accumulating at $q=1$ (Lemma~\ref{lem:infpoles}, proved in \S\ref{sec:sd}
once the collapse machinery is in place).

\begin{proposition}[Invariance of the travel block]\label{prop:travelinv}
The denominator $1-\Sigma_1$ is the identical, cycle-marker-free object in the generating functions
of $U$ and of $V$; hence $U$ and $V$ share the travel poles $q_m$.
\end{proposition}
\begin{proof}
A pure-travel element has no gap edge, since its span is $I_k$ and $f_j=\pm1$ there, so
Theorem~\ref{thm:nogap} gives $c=0$: some relaxed-optimal realisation has no isolated cycle. That
realisation is a single open walk, hence admissible for $\ell_T$, so $\ell_T\le\ell_R$; and
$\ell_T\ge\ell_R$ because the relaxed minimum ranges over a superset of the realisations. Therefore
$\ell_T=\ell_R$ on pure-travel elements and the two gradings of $T$ agree termwise. The travel
recursion is assembled from pure-travel runs alone, appending one $f=\pm1$ edge at a time, so it is
computed in the sector $c\equiv0$ and the travel block carries $y^0$ only.
\end{proof}

\noindent Earlier versions of this paper recorded the proposition as a verified conjecture resting
on two unproved statements, the localisation of cycles and the metric identity
$\ell_T=\ell_R+2c$. The first is now Corollary~\ref{cor:localzero}. The second is not needed: only
the inequality $\ell_T\le\ell_R$ is used, and it follows from the realisation exhibited by
Theorem~\ref{thm:nogap} together with the trivial reverse inequality, so the open lower bound of the
metric identity does not enter. Statements below that invoke this proposition no longer inherit a
conditional standing from it; see Theorem~\ref{thm:U} and Remark~\ref{rem:Ustatus}.

\section{Transcendence of the blocks}\label{sec:blocks}

The four blocks are the $k$-recursion sums
\begin{equation}\label{eq:krec}
  \Sigma_k=\sum_{j\ge0}A_{k+2j}\prod_{i<j}C_{k+2i},\qquad
  S_k=\sum_{j\ge0}\alpha_{k+2j}\prod_{i<j}\gamma_{k+2i},
\end{equation}
\[
  A_k=\frac{2q}{1-q^{k+1}},\quad
  C_k=\frac{2q^{k+3}}{1-q^{k+2}}-\frac{2q^{k+2}}{1-q^{k+1}},\quad
  \alpha_k=\frac{2q^{k+1}}{1-q^{k+1}},\quad
  \gamma_k=\alpha_{k+1}-\alpha_k ,
\]
$\Sigma_0,\Sigma_1$ carrying the travel r\^ole and $S_0,S_1$ the bulk r\^ole; $\Sigma_1$ and $S_1$
are the denominators of \eqref{eq:blocks} and $\Sigma_0,S_0$ their numerators. Every factor is a
Lambert series with integer coefficients. The factor $C_{1+2i}$ starts at order $2i+3$, so
$\prod_{i<j}C_{1+2i}$ starts at order $\sum_{i<j}(2i+3)=j^2+2j$ and, with $A_{1+2j}$ starting at
order $1$, the $j$-th summand of $\Sigma_1$ starts at order exactly $(j+1)^2$; only the finitely
many $j\le\sqrt n-1$ reach $[q^n]$. Each coefficient is therefore a finite integer sum, and
$\Sigma_1=2q+2q^3-4q^4+6q^5-4q^6+\cdots$. That integrality, with convergence on $|q|<1$, is the
whole input to P\'olya--Carlson.

\begin{theorem}\label{thm:blocks}
Let $\Sigma$ be any of $\Sigma_0,\Sigma_1,S_0,S_1$. Then $\Sigma$ has integer coefficients and
converges in $|q|<1$, and exactly one of the following holds:
\begin{enumerate}[label=\textup{(\alph*)},leftmargin=2.2em,itemsep=0.1em,topsep=0.2em]
\item $\Sigma$ is a rational function whose poles are roots of unity; or
\item $|q|=1$ is a natural boundary of $\Sigma$, and $\Sigma$ is transcendental over $\Q(q)$.
\end{enumerate}
Alternative \textup{(a)} is excluded, and hence \textup{(b)} holds, for any block whose
coefficients grow super-polynomially along a subsequence. For $\Sigma_1$ that growth is the
hypothesis
\begin{equation}\label{eq:growth}
   \bigl|[q^{(j+1)^2}]\Sigma_1\bigr|\ \ge\ 2^{\,j+1}\qquad\text{for infinitely many }j,
\end{equation}
which is \emph{not proved here}: see Remark~\textup{\ref{rem:growthstatus}}.
\end{theorem}

\begin{proof}
The coefficients are integers by the order count above: each is a finite sum of integers.
Convergence in $|q|<1$ is read off the closed forms: each block is a series whose $j$-th term carries the
Gaussian factor $q^{j^2+O(j)}$ against a denominator $(q;q)_{2j}$ or $(q;q)_{2j+1}$ bounded below
in modulus by $|(q;q)_\infty|>0$ on compact subsets of the punctured disc (for $\Sigma_1$ use
$1-\Sigma_1=\cos(Z;q^2)=\sum_{j\ge0}(-1)^jq^{j^2}(2(1-q))^j/(q;q)_{2j}$ of
Proposition~\ref{prop:qtrig}, in which the pole of $Z^2=2(1-q)/q$ at the origin is cancelled), so
each series converges absolutely and locally uniformly there. Hence
P\'olya--Carlson \cite{Carlson1921} applies and gives the dichotomy \textup{(a)}/\textup{(b)}: a
power series with integer coefficients convergent in the unit disc is either a rational function
whose poles are roots of unity, or has the unit circle as a natural boundary; a function with a
natural boundary is transcendental over $\Q(q)$. In case \textup{(a)} the coefficient sequence
satisfies a linear recurrence all of whose characteristic roots are roots of unity, hence
$|[q^n]\Sigma|\le Cn^d$ for some $C,d$; a subsequence growing faster than every polynomial in $n$
is therefore incompatible with \textup{(a)}. Under \eqref{eq:growth},
$|[q^n]\Sigma_1|\ge2^{\sqrt n}$ along $n=(j+1)^2$, which is such a subsequence.
\end{proof}

For the two denominators the growth hypothesis is not needed at all: alternative \textup{(a)} is
excluded by the pole set instead.

\begin{corollary}[the pole route to alternative \textup{(a)}]\label{cor:polepath}
Alternative \textup{(a)} of Theorem~\ref{thm:blocks} fails for $\Sigma_1$ and fails for $S_1$;
hence for each of the two denominators $|q|=1$ is a natural boundary and the block is transcendental
over $\Q(q)$. It also fails for at least one of the two numerators $\Sigma_0,S_0$. Neither
statement uses \eqref{eq:growth} or Hypothesis~\ref{hyp:growth}.
\end{corollary}

\begin{proof}
A rational function that is not identically zero has finitely many zeros in $\C$. The constraint is
global: unlike the identity theorem it requires no accumulation point inside a domain, so a zero set
that is merely infinite already contradicts it, wherever in $\C$ the zeros sit and wherever they
accumulate.

Suppose $\Sigma_1$ satisfies \textup{(a)}, so that on $|q|<1$ it agrees with a rational function
whose poles are roots of unity. Then $1-\Sigma_1$ is rational, and it is not identically zero,
because $\Sigma_1(0)=0$: every factor $A_k,C_k$ of \eqref{eq:krec} vanishes at $q=0$, equivalently
$1-\Sigma_1=\cos(Z;q^2)$ has constant term $1$ (Proposition~\ref{prop:qtrig}). A nonzero rational
function has finitely many zeros; but $\{q\in(0,1):\Sigma_1(q)=1\}$ is infinite by
Lemma~\ref{lem:infpoles}. Alternative \textup{(a)} is therefore impossible and
Theorem~\ref{thm:blocks} leaves \textup{(b)}. The same argument applies to $S_1$, with
$1-S_1=S_e$ and $S_1(0)=0$ (Proposition~\ref{prop:qtrig}) and with Lemma~\ref{lem:infzeros} in place
of Lemma~\ref{lem:infpoles}.

For the numerators, suppose $\Sigma_0$ and $S_0$ both satisfy \textup{(a)}. Then $\Sigma_0S_0$ is
rational and, by Proposition~\ref{prop:numnonvanish}, agrees with the rational function
$2q/(1-q)$ at every travel pole, hence at infinitely many points of $(0,1)$
(Lemma~\ref{lem:infpoles}); two rational functions
agreeing at infinitely many points are equal. They are not: $\Sigma_0S_0=4q^2+8q^3+4q^4+\cdots$
vanishes to order $2$ at $q=0$ and $2q/(1-q)$ to order $1$. This says nothing about either numerator
separately, \eqref{eq:numprod} constraining only the product.
\end{proof}

\noindent The two lemmas quoted are proved in \S\ref{sec:sd} and the two identities in
\S\ref{sec:qtrig}; of Theorem~\ref{thm:blocks} they use only its first clause, that the blocks
converge on $|q|<1$, so the dependency graph is acyclic. The series data
$\Sigma_1=2q+2q^3-4q^4+\cdots$, $S_1=2q^2+2q^4-2q^6+\cdots$, $\Sigma_0S_0=4q^2+8q^3+4q^4+\cdots$
and the two collapses $1-\Sigma_1=\cos(Z;q^2)$, $1-S_1=S_e$ are recomputed in exact integer
arithmetic at two truncation degrees in \texttt{f2\_pole\_route.py}.

\begin{remark}[what the pole route costs]\label{rem:polepathcost}
Corollary~\ref{cor:polepath} makes Theorem~\ref{thm:blocks} unconditional for $\Sigma_1$ and $S_1$,
and the two routes to it are of different kinds. \eqref{eq:growth} is checked in exact integer
arithmetic, with no precision to state, at $34$ indices, and remains a finite verification of a
statement with infinitely many instances. The pole route runs on Lemma~\ref{lem:infpoles}, hence on
Lemma~\ref{lem:T2abs}, and that lemma is proved on all of $\tau\in(0,0.02]$: a verified enclosure
of the contour majorant, in interval arithmetic and quantified over parameter boxes rather than
sampled, on $\tau\ge2/101^2$, and an explicit-constant estimate below that. The two infinitude
lemmas consume it along the infinite sequences $\tau^{(m)}=2/(m\pi)^2$ and $\tau^{[m]}$, $m\ge4$,
all of which lie in that range. So the pole route carries no unproved arithmetic input, while
\eqref{eq:growth} does; and for $\Sigma_0$ and for $S_0$ individually, \eqref{eq:growthG} remains
the only route. What the corollary buys besides the label is a shorter dependency list: the
exclusion of \textup{(a)} for the two denominators consumes the same single analytic input as
Theorems~\ref{thm:travelres}, \ref{thm:V}, \ref{thm:U} and Corollary~\ref{cor:beta}, instead of a
second input of its own. Proposition~\ref{prop:infpolesalt} gives both infinitude statements a
second proof with no $T_2$ input, so for the two denominators, and for the joint statement on the
numerators, the exclusion of \textup{(a)} does not need Lemma~\ref{lem:T2abs} at all.
\end{remark}

\noindent The hypothesis \eqref{eq:growth} concerns $\Sigma_1$ alone. Its counterpart for the other
three blocks is the following uniform statement, which is what alternative \textup{(a)} has to be
weighed against for $\Sigma_0$ and $S_0$, the two blocks Corollary~\ref{cor:polepath} does not
reach individually.

\begin{hypothesis}[coefficient growth]\label{hyp:growth}
Let $\Sigma$ be any of $\Sigma_0,\Sigma_1,S_0,S_1$. Then
\begin{equation}\label{eq:growthG}
   \bigl|[q^{n}]\Sigma\bigr|\ \ge\ 2^{\sqrt n-1}\qquad\text{for infinitely many }n .
\end{equation}
\end{hypothesis}

\noindent Statement \eqref{eq:growth} is the $\Sigma_1$ instance of \eqref{eq:growthG} in the
sharper form that names the witnessing indices $n=(j+1)^2$. Either form excludes alternative
\textup{(a)} for the block concerned, $2^{\sqrt n-1}$ outgrowing every polynomial in $n$.

\begin{remark}[why \eqref{eq:growth} is expected, and how far it is verified]\label{rem:growthstatus}
\eqref{eq:growth} is verified, not proved, and the reason it is expected is exact rather than
heuristic. By the order count above, the $j$-th summand of the telescoping
$\Sigma_1=\sum_{j\ge0}\bigl(\prod_{i<j}C_{1+2i}\bigr)A_{1+2j}$ has lowest degree exactly $(j+1)^2$,
with coefficient exactly $(-1)^j2^{\,j+1}$ (the leading $2$ of $A_{1+2j}$ against the leading $-2$
of each of the $j$ factors $C_{1+2i}$), and summands of index $>j$ start above that degree.
Hence
\[
   [q^{(j+1)^2}]\Sigma_1\;=\;(-1)^j2^{\,j+1}\;+\;\Bigl(\text{the degree-}(j+1)^2\text{ tails of the
   $j$ summands of index}<j\Bigr),
\]
a \emph{finite} correction to an exactly known leading term: \eqref{eq:growth} says precisely that
those $j$ tails do not cancel it. They do not merely fail to cancel; they reinforce. Exact integer
evaluation of \eqref{eq:krec} truncated at degree $1156$ gives
\[
   [q^{(j+1)^2}]\Sigma_1=2,\,-4,\,14,\,-52,\,178,\,-856,\,4626,\,-27524,\,150214,\,-816268,\dots
\]
and confirms \eqref{eq:growth} for every $j\le33$, with the sign $(-1)^j$ of the leading term at
every one of them and a margin that grows without bound: $|[q^{(j+1)^2}]\Sigma_1|/2^{\,j+1}$ is
$797$ at $j=9$, $1.2\cdot10^8$ at $j=20$ and $2.7\cdot10^{14}$ at $j=33$, where the value is
$-4649405079364022238518260$ against the bound $2^{34}$; the successive ratios
$|[q^{(j+2)^2}]\Sigma_1|/|[q^{(j+1)^2}]\Sigma_1|$ rise from $2$ to $6.22$ over that range.
Hypothesis~\ref{hyp:growth} holds for all four blocks with a witness in every window
$[(j+1)^2,(j+2)^2)$ for every $j\le32$, the range being cut off by the truncation degree and not by
a failure. The witnessing diagonals are $n=(j+1)(j+2)$ for $S_1$ and $n=(j+1)^2$ for $S_0$, each
carrying $|[q^n]\Sigma|\ge2^{\,j+1}$ throughout $4\le j\le32$ and $3\le j\le32$ respectively; for
$\Sigma_0$ the leading diagonal cancels: its values at $n=(j+1)^2$ are
$2,2,2,10,18,34,266,306,8170,330$, which drop below $2^{\,j+1}$ at $j=1,\dots,5$ and again at
$j=9$, so its witnesses are off-diagonal and \eqref{eq:growthG} is the only form available for it.

The computation is exact integer arithmetic in $\Z[[q]]/(q^{1161})$
(\texttt{blocks\_growth.py}); no floating point enters, so no precision is quoted. It is checked
against a second and independent derivation of the same coefficients from the $q$-trigonometric
closed forms of Proposition~\ref{prop:qtrig} and Lemma~\ref{lem:P12closed} (agreement to degree
$1160$ for each of $\Sigma_1,S_1,S_0$), and against the bulk block
$S_0/(1-S_1)=2q+2q^2+6q^3+2q^4+18q^5+6q^6+42q^7+\cdots$. For $j\le9$ the bound on $\Sigma_1$ is
also machine-checked in Lean~4 (\texttt{PolyaCarlson.lean}, theorem \texttt{coeff\_bound}); that
proof uses \texttt{native\_decide} and therefore depends, besides \texttt{propext}, on the
compiler-trust axiom \texttt{PolyaCarlson.coeff\_bound.\_native.native\_decide.ax\_1\_1}, which is
outside the three standard axioms.

No finite verification can establish a statement quantified over infinitely many indices.
\eqref{eq:growth} is numerically certain and unproved, and the transcendence conclusion
\textup{(b)} of Theorem~\ref{thm:blocks} inherits that standing wherever it is the route used,
which after Corollary~\ref{cor:polepath} is $\Sigma_0$ and $S_0$. What would close it is a lower
bound on the degree-$(j+1)^2$ coefficient of the $j$ preceding Lambert products, uniform in $j$:
the statement is finite for each $j$ and the obstruction is uniformity, not any single case.
\end{remark}

The dichotomy of Theorem~\ref{thm:blocks} uses no analytic gate. The exclusion of alternative
\textup{(a)} rests on \eqref{eq:growth} for $\Sigma_0$ and $S_0$, and on the pole route of
Corollary~\ref{cor:polepath}, hence on the analytic input of \S\ref{sec:sd}, for $\Sigma_1$ and
$S_1$. Neither transfers directly
to $V$ or $U$, which are \emph{ratios} of blocks in which the boundary singularities can cancel;
those require the pole-accumulation argument below, which is independent of \eqref{eq:growth}.

\section{The steepest-descent estimate}\label{sec:sd}

The one analytic input for both $V$ and $U$ is the smallness of an auxiliary integral. Let
$\varphi(y)=\log\frac{\sinh(y/2)}{y/2}=\sum_{n\ge1}\varphi_n y^{2n}$ with
$\varphi_n=\frac{(-1)^{n+1}\zeta(2n)}{n(2\pi)^{2n}}$. The form factor $B_s$ is \emph{defined} by
the closed form \eqref{eq:Bclosed} of Lemma~\ref{lem:Bbounded}; the familiar expansion
$B_s\sim\sum_{n\ge1}\varphi_n\tau^{2n}(P_n(2s)-s)$, $P_n(M)=\sum_{m=1}^{M}m^{2n}$ the Faulhaber
sum, is its Euler--Maclaurin \emph{asymptotic} series (equality at integer $2s$, divergent
otherwise; Lemma~\ref{lem:Bbounded}(v)) and is used only truncated. The block $S_e$ decomposes as $S_e=\cos W-T_2$ with $W=we^{-\tau/2}$ and
\begin{equation}\label{eq:T2}
   T_2=\sum_{n\ge1}(-1)^n\,g_n\,\frac{W^{2n}}{\Gamma(2n+1)},\qquad g_s=1-e^{-B_s}.
\end{equation}
What the theorems consume is only $T_2\to0$; its sharp size is recorded separately in
Remark~\ref{lem:cos} and is not used.

\begin{lemma}[amplitude regularity]\label{lem:Bbounded}
For $\Re s>-\tfrac12$ set $M=2s$ and
\begin{equation}\label{eq:Bclosed}
   B_s\;=\;\tfrac{\tau}4M(M{+}1)-M\log\tau-\log\Gamma(M{+}1)
   +\log(q;q)_\infty-\log(q^{M+1};q)_\infty-s\,\varphi(\tau),
\end{equation}
$\varphi(y)=\log\frac{\sinh(y/2)}{y/2}$ as above, each $q$-Pochhammer logarithm meaning the sum of
principal logarithms $\sum_{j\ge0}\operatorname{Log}(1-q^{M+1+j})$: every factor has
$|q^{M+1+j}|<1$, hence $\Re(1-q^{M+1+j})>0$, so this is the analytic branch on the half-plane
(the logarithm of the \emph{product} would differ by multiples of $2\pi i$). Then $B_s$ is analytic on $\Re s>-\tfrac12$, and
at every integer $n\ge0$
\begin{equation}\label{eq:Bsum}
   B_n\;=\;\sum_{i=1}^{2n}\varphi(i\tau)\;-\;n\,\varphi(\tau),
\end{equation}
so \eqref{eq:Bclosed} is the analytic interpolation of the collapse coefficients
$\eta_n=e^{-B_n}$ of \eqref{eq:blocks}. On the strip
\[
   \mathcal S=\{-\tfrac12<\Re s\le2w,\ |\Im s|\le2w\},\qquad \tau\le2/\pi^2,
\]
one has the explicit bounds
\begin{equation}\label{eq:Bbounds}
   |B_s|\ \le\ 30.3\,\sqrt\tau,\qquad |B_s'|\ \le\ 7.6\,\tau ,
\end{equation}
and, on the smaller region $|M|\le2w$ which contains the saddles, the truncation bound that
Remark~\ref{lem:cos} consumes:
\begin{equation}\label{eq:Btrunc}
   \bigl|B_s-\varphi_1\tau^2\bigl(P_1(M)-s\bigr)\bigr|\ \le\ 0.02\,\tau^{3/2},
   \qquad \varphi_1=\tfrac1{24}.
\end{equation}
Consequently $|e^{-B_s}|\le e^{30.3\sqrt\tau}$, the amplitude $A:=1-e^{-B_s}$ satisfies
$|A|\le1+e^{30.3\sqrt\tau}\to2$, and along any rectifiable $\Gamma\subset\mathcal S$ of length
$\ell(\Gamma)\le cw$,
\[
   V_\Gamma(A)\ \le\ \sup_\Gamma|e^{-B_s}|\int_\Gamma|B_s'|\,|ds|
   \ \le\ e^{30.3\sqrt\tau}\cdot7.6\,\tau\cdot cw\ =\ O(\sqrt\tau)\ \longrightarrow\ 0 ,
\]
using $\tau w=\sqrt{2\tau}$. The strip contains both saddles
$s^\ast=-\tfrac14\pm i(\tfrac w2+O(\sqrt\tau))$ of Remark~\ref{lem:cos} and the $O(w)$-length
portion of its steepest-descent contour through them (which is all the $\ell(\Gamma)\le cw$
clause above requires, the Gaussian width being $O(\sqrt w)$) and, being built from the larger
argument $w=We^{\tau/2}$, it contains the corresponding $W$-strip as well.
\end{lemma}

\begin{proof}
\emph{(i) Closed form, and \eqref{eq:Bsum}.} Unwinding
$\eta_n=q^{n^2+n}[2(1-q)]^n(2n)!/[(q;q)_{2n}W^{2n}]$ with $W^2=(2/\tau)e^{-\tau}$ gives
$\eta_n=e^{-\tau n^2}[(1-q)\tau]^n(2n)!/(q;q)_{2n}$. Now
$1-q^i=i\tau\,e^{-i\tau/2}\,\frac{\sinh(i\tau/2)}{i\tau/2}$, i.e.
$\log(1-q^i)=\log(i\tau)-\tfrac{i\tau}2+\varphi(i\tau)$; summing over $1\le i\le M$,
\begin{equation}\label{eq:qpochsum}
   \log(q;q)_M=\log\Gamma(M{+}1)+M\log\tau-\tfrac\tau4M(M{+}1)+\sum_{i=1}^{M}\varphi(i\tau).
\end{equation}
Since $\log[(1-q)\tau]=-\tfrac\tau2+2\log\tau+\varphi(\tau)$, substituting \eqref{eq:qpochsum} into
$-\log\eta_n$ collapses every elementary term and leaves exactly \eqref{eq:Bsum}; and rewriting
$\log(q;q)_M=\log(q;q)_\infty-\log(q^{M+1};q)_\infty$ in \eqref{eq:qpochsum} turns \eqref{eq:Bsum}
into \eqref{eq:Bclosed} for all $s$. Every ingredient of \eqref{eq:Bclosed} is analytic in $s$ for
$\Re s>-\tfrac12$: there $|q^{M+1}|<1$ keeps $(q^{M+1};q)_\infty$ zero-free and the branch above
canonical, and $\Re(M{+}1)>0$ keeps $M{+}1$ off the poles of $\Gamma$. (Both singular factors do in
fact cancel at $M=-1$, so $B_s$ continues past $\Re s=-\tfrac12$; we never need that.)

\emph{(ii) Splitting into Gamma blocks with exact cancellation.} By the Weierstrass product
$\frac{\sinh(y/2)}{y/2}=\prod_{k\ge1}\bigl(1+\frac{y^2}{4\pi^2k^2}\bigr)$ we have
$\varphi(y)=\sum_{k\ge1}\log\bigl(1+y^2/(4\pi^2k^2)\bigr)$, so with $R_k:=2\pi k/\tau$ and
$\prod_{n=1}^{M}(1+n^2a^2)=a^{2M}\,\Gamma(M{+}1{-}\tfrac ia)\Gamma(M{+}1{+}\tfrac ia)/
\bigl[\Gamma(1{-}\tfrac ia)\Gamma(1{+}\tfrac ia)\bigr]$ at $a=1/R_k$, \eqref{eq:Bsum} becomes
\begin{equation}\label{eq:Bblocks}
   B_s=\sum_{k\ge1}G(M;R_k)-s\,\varphi(\tau),\qquad
   G(M;R):=\int_0^M\Bigl[\psi(1{+}u{+}iR)+\psi(1{+}u{-}iR)-2\log R\Bigr]du,
\end{equation}
the integral along the segment $[0,M]$ fixing the branch (legitimate since $|M|<R_1$ puts the
segment far from the poles of $\psi(1{+}u\pm iR)$ at $u=-1-j\mp iR$). \emph{The $-2\log R$ is not a
normalisation: it is the exact cancellation of the two leading Stirling terms}, and it is what
makes the $k$-sum converge; without it each block would grow like $2M\log R$.

Identity \eqref{eq:Bblocks} is proved at integers, where the $i$-sum is finite and the interchange
of $\sum_i$ and $\sum_k$ is trivial; it extends to all $s$ as follows.

\emph{Domain.} Both sides are analytic on the half-plane $\Re M>-1$: the left by~(i), the right
because the segment $[0,M]$ then stays in $\Re u>-1$ while the poles of $\psi(1{+}u\pm iR)$ lie in
$\Re u\le-1$, and because $\sum_kG$ converges locally uniformly there, all but finitely many $k$
having $R_k>2|M|$ and hence $|G(M;R_k)|=O(|M|^3/R_k^2)$. (The disc $|M|<R_1$ of step~(iii) is where
the \emph{quantitative} majorant lives; it is not the domain of analyticity, and the recursion
below would walk out of it.)

\emph{Recursion.} Write $\widetilde B=B_s+s\varphi(\tau)$ for either side with the elementary term
restored. Both obey the \emph{same} unit recursion in $M$: on the left
$\tfrac\tau4[(M{+}1)(M{+}2)-M(M{+}1)]-\log\tau-\log(M{+}1)+\log(1-q^{M+1})=\varphi((M{+}1)\tau)$
term by term, and on the right $G(M{+}1;R)-G(M;R)=\log\bigl(1+(M{+}1)^2/R^2\bigr)$, which sums over
$k$ to the same $\varphi((M{+}1)\tau)$ by the Weierstrass product. (For $B_s$ itself both
increments are $\varphi((M{+}1)\tau)-\tfrac12\varphi(\tau)$, the extra term being common to the two
sides.)

\emph{Conclusion.} Hence the difference $D$ is $1$-periodic on $\Re M>-1$, so it continues to a
$1$-periodic \emph{entire} function. Its Fourier coefficients satisfy
$|c_m|\le e^{2\pi my}\sup_{0\le x\le1}|D(x{+}iy)|$ for every real $y$, and $D$ grows at most
polynomially: on the left by Stirling, on the right by splitting the $k$-sum at $R_k=2|M|$: the
$O(|M|\tau)$ low terms are each $O(|M|\log|M|)$ by Stirling and the remainder is
$O\bigl(|M|^3\sum_{R_k>2|M|}R_k^{-2}\bigr)=O(|M|^2\tau)$. Letting $y\to-\infty$ for $m>0$ and
$y\to+\infty$ for $m<0$ gives $c_m=0$ for $m\ne0$, so $D$ is constant; and $D(0)=0$ by
\eqref{eq:Bsum} at $n=0$, so $D\equiv0$.

\emph{(iii) The integrand.} For $|z|\ge1$ with $|\arg z|\le\tfrac{3\pi}4$ one has
$|\psi(1{+}z)-\log z-\tfrac1{2z}|\le|z|^{-2}$ (the constant is $0.13$ on that sector; the bound
fails without limit as $\arg z\to\pi$, which is why the sector, not a distance to the cut, is the
right hypothesis). On $\mathcal S$, $|M|\le4\sqrt2\,w$ and
$|M|/R_1=\tfrac{4\sqrt2\,w\tau}{2\pi}=\tfrac{4}{\pi}\sqrt{\tau}\le0.58$, so every $R_k>|M|$; and
along the segment $|\arg(u\pm iR_k)|\le120^\circ<135^\circ$ with $|u\pm iR_k|\ge R_1-|M|>13$, so
the estimate applies throughout. Adding the two conjugate terms, the logarithms
combine with $-2\log R$ to $\log\bigl(1+u^2/R^2\bigr)$ and the $\tfrac1{2z}$ terms to
$u/(u^2{+}R^2)$, whence for $|u|\le|M|$
\[
   \bigl|\psi(1{+}u{+}iR)+\psi(1{+}u{-}iR)-2\log R\bigr|
   \ \le\ \frac{|M|^2+|M|}{R^2-|M|^2}+\frac{2}{(R-|M|)^2}\ =:\ \gamma(R).
\]

\emph{(iv) Summation.} Hence $|G(M;R_k)|\le|M|\,\gamma(R_k)$ and, from \eqref{eq:Bblocks},
\[
   |B_s|\ \le\ |M|\sum_{k\ge1}\gamma(R_k)+|s|\,\varphi(\tau),
   \qquad
   |B_s'|\ \le\ 2\sum_{k\ge1}\gamma(R_k)+\varphi(\tau),
\]
the second because $\partial_MG=\psi(1{+}M{+}iR)+\psi(1{+}M{-}iR)-2\log R$ and $M=2s$. Both sums are
dominated by $\sum_k R_k^{-2}=\tau^2/24$; carrying the finite corrections explicitly at
$|M|=4\sqrt2\,w$, $\tau\le2/\pi^2$ gives $|M|\sum_k\gamma(R_k)+|s|\varphi(\tau)\le30.3\sqrt\tau$ and
$2\sum_k\gamma(R_k)+\varphi(\tau)\le7.6\,\tau$, which is \eqref{eq:Bbounds}. Both ratios are \emph{increasing} in $\tau$ on the range ($30.20,\ 25.14,\ 23.26,\ 22.21,\
21.64$ and $7.55,\ 6.28,\ 5.81,\ 5.55,\ 5.41$ at $\tau=2/\pi^2,\,0.1,\,0.05,\,0.02,\,0.005$,
decreasing to the limits $\tfrac{64}3$ and $\tfrac{16}3$ as $\tau\to0$), so each attains its
maximum at the right endpoint $\tau=2/\pi^2$, and the stated constants hold throughout.

For \eqref{eq:Btrunc}: expand the integrand of $G$ one order further, now \emph{keeping the third
Stirling term} $-\tfrac1{12z^2}$, which is what supplies the linear part of $P_1$: its integral
is $\tfrac{M}{6(M^2+R^2)}=\tfrac{M}{6R^2}+O(|M|^3/R^4)$. With
$\int_0^M\log(1+u^2/R^2)du=\tfrac{M^3}{3R^2}+O(|M|^5/R^4)$ and
$\int_0^Mu/(u^2{+}R^2)du=\tfrac{M^2}{2R^2}+O(|M|^4/R^4)$ this gives
$R^2G(M;R)\to P_1(M)=\tfrac{M^3}3+\tfrac{M^2}2+\tfrac M6$, so with $\sum_kR_k^{-2}=\tau^2/24$
exactly,
\[
   B_s=\varphi_1\tau^2\bigl(P_1(M)-s\bigr)+O(|M|^5\tau^4)+O(|s|\tau^4).
\]
On $|M|\le2w$ every remainder is $O(\tau^{3/2})$, and the constant is $0.02$ (worst $0.0177$ at
$\tau=2/\pi^2$, decreasing).

\emph{(v) What is \emph{not} used.} The $\tau$-expansion
$B_s=\sum_{n\ge1}\varphi_n\tau^{2n}(P_n(2s)-s)$, $P_n$ the Faulhaber sum, is the
Euler--Maclaurin \emph{asymptotic} expansion of \eqref{eq:Bclosed}: it is \emph{divergent} at
non-integer $2s$, since $P_n(2s)$ carries a Bernoulli part of size
$(2n)!\,e^{2\pi|\Im 2s|}/(2\pi)^{2n+1}$. Nothing above sums it; it is used only in the truncated
form ``leading term plus remainder'', as in Remark~\ref{lem:cos}, where the $n=1$ term at the saddle
supplies the constant $\sqrt2/36$.
\end{proof}

\noindent The algebraic skeleton of this proof is machine-checked in Lean~4
(\texttt{BboundedSkeleton.lean}, ten theorems, no \texttt{sorry}, axioms \texttt{propext},
\texttt{Classical.choice}, \texttt{Quot.sound}). Five carry the arithmetic content: the Gauss sum,
the summation \eqref{eq:qpochsum} built on it, the telescoping cancellation \eqref{eq:Bsum}, the
scaling identity $(\tau^2w^3)^2=8\tau$ (exact for $w$, only $e^{-3\tau/2}$-approximate for $W$),
and the strip condition $|M|/R_1<1$. The other five (the per-factor rewriting, the collapse of
$\log[(1-q)\tau]$, the cancellation $2M\log a+2M\log R=0$ at $aR=1$, the $w$-definition, and the
bounded-variation arithmetic) are \emph{rearrangements of their hypotheses}, verified as such:
they certify that the paper's algebra is what it claims, not the transcendental facts fed into
them. Those facts (the Weierstrass product, the digamma remainder, the convergence of the
$k$-sum, and the continuation argument of step~(ii)) are not formalised.

\begin{remark}[sharpness]\label{rem:Bsharp}
The majorant of step~(iv) is lossy by about a factor $6$: the true extremes on $\mathcal S$ are
$\Re B_s\ge-5.03\sqrt\tau$ and $|B_s'|\le5.46\,\tau$, and $\sup_{\mathcal S}|A|=9.02$ at
$\tau=2/\pi^2$ (\texttt{star\_gate\_certificate.py}, part~(7)). Only the proved bounds
\eqref{eq:Bbounds} are used anywhere.
\end{remark}

The sketch exhibits the saddles explicitly and derives the constants by Cauchy's theorem, Stirling's
formula and elementary Gaussian integration, and it would end by invoking a second-order remainder
estimate for a simple non-degenerate saddle, \cite[Ch.~4, Thm.~7.1]{Olver}. Its hypotheses are
\emph{not} verified here (the saddle is flat rather than non-degenerate, and the amplitude's
bounded variation is available only on a bounded piece of the path), which is one of the reasons
the sketch is not a proof (Remark~\ref{rem:cosgap}). No theorem of this paper rests on it.

\begin{remark}[the sharp size of $T_2$; \textsc{heuristic}]\label{lem:cos}
Numerically, $T_2=\tfrac{\sqrt2}{36}\sqrt\tau\,\sin X-\tfrac{161}{1296}\,\tau\cos X+O(\tau^{3/2})$
uniformly in the phase, for any argument $X$ with $\tau X^2\in[2e^{-\tau},2]$; the constant and the
exponent of the remainder are certain to eight digits over $m\le1024$. The sketch below indicates
why, by steepest descent, but it does \emph{not} constitute a proof, for the reasons set out in
Remark~\ref{rem:cosgap}. Accordingly this statement carries no weight anywhere in this paper:
the only property of $T_2$ that the theorems consume is the crude bound
$|T_2|\to0$ of Lemma~\ref{lem:T2abs}, which is proved, and which is what
Lemma~\ref{lem:infpoles} cites.
\end{remark}

\begin{proof}[Sketch, not a proof; see Remark~\ref{rem:cosgap}]
Since $\cos W=\sum_{n\ge0}(-1)^nW^{2n}/(2n)!$, \eqref{eq:T2} gives the exact alternating sum
$S_e=\cos W-T_2=\sum_{n\ge0}(-1)^n e^{-B_n}W^{2n}/(2n)!$ (with $B_0=0$). Let
$b(z)=e^{-B_z}W^{2z}/\Gamma(2z+1)$, analytic for $\Re z>-\tfrac12$ (by \eqref{eq:Bclosed} and
Lemma~\ref{lem:Bbounded}; \emph{not} by termwise polynomiality, the expansion being divergent off
$z$). As $\pi\csc\pi z$ has residue $(-1)^n$ at each integer, Cauchy's theorem gives
$S_e=\tfrac1{2\pi i}\oint b(z)\,\pi\csc(\pi z)\,dz$ over a contour enclosing $\{0,1,2,\dots\}$. Write the
integrand as $e^{h(z)}\pi\csc\pi z$ with $h(z)=2z\log W-\log\Gamma(2z+1)$; including the $\csc$ phase, the
exponent is stationary where $2\log W-2\psi(2z+1)\pm i\pi=0$, so by $\psi(w)=\log w+O(1/w)$ there are two
simple, non-degenerate saddles
\[
   z^\ast_\pm=\pm\tfrac{i}{2}W\bigl(1+O(1/W)\bigr),\qquad h''(z^\ast_\pm)=\pm\tfrac{4i}{W}+O(W^{-2})\neq0.
\]
Deform the contour to the steepest-descent paths through $z^\ast_\pm$; this is legitimate because $b$
decays super-exponentially as $\Re z\to+\infty$ and $\csc$ suppresses $|\Im z|\to\infty$. A direct Stirling
evaluation shows that at the modulation-free level $B\equiv0$ the two saddles sum to $\cos W$ \emph{exactly}: the divergent factors $e^{\pm\pi W/2}$ from $\csc(\pi z^\ast_\pm)$ and from $1/\Gamma(\pm iW+1)$ cancel
identically, each saddle contributing $\tfrac12 e^{\pm iW}$. Restoring the modulation scales each saddle by
$e^{-B_{z^\ast_\pm}}$; since $P_1(2z)=\tfrac{2z(2z+1)(4z+1)}6$ and $\varphi_1=\tfrac1{24}$,
\[
   B_{z^\ast_+}=\varphi_1\tau^2\bigl(P_1(iW)-\tfrac{iW}{2}\bigr)+O(\tau)
   \qquad\text{(the $\tau^{3/2}$ part of the remainder is \eqref{eq:Btrunc}; the $O(\tau)$
   also absorbs $B_{z^\ast_+}-B_{iW/2}$ and the dropped $\varphi_1\tau^2W^2/2$)}
   =\tfrac{\tau^2}{24}\Bigl(-\tfrac{i}{3}W^3\Bigr)+O(\tau)=-i\tfrac{\sqrt2}{36}\sqrt\tau+O(\tau)
\]
(using $\tau^2W^3=2\sqrt2\,\sqrt\tau\,e^{-3\tau/2}$, the factor $e^{-3\tau/2}$ being absorbed in the $O(\tau)$), and $B_{z^\ast_-}=\overline{B_{z^\ast_+}}$. Hence
\[
   S_e=\tfrac12 e^{-B_{z^\ast_+}}e^{iW}+\tfrac12 e^{-B_{z^\ast_-}}e^{-iW}+E
   =\cos W-\tfrac{\sqrt2}{36}\sqrt\tau\,\sin W+E .
\]
The remainder $E$ is $O(\tau)$ by the classical second-order estimate for a simple saddle
\cite[Ch.~4, Thm.~7.1]{Olver}: parametrising the descent path by arclength $s$, the integrand equals
$e^{h(z^\ast)}e^{-s^2}\bigl(1+O(sW^{-1/2})+O(\tau)\bigr)$, the next saddle coefficient is
$\propto h'''(z^\ast)=O(W^{-2})=O(\tau)$, and the bounded-variation hypothesis on the amplitude
$e^{-B_z}$ along the path is supplied by Lemma~\ref{lem:Bbounded}. This citation is the proof's one
non-elementary step. The
saddle prediction matches $S_e$ to $O(\tau)$ over $\tau=0.02,0.01,0.005$ numerically
(\texttt{saddle\_verify.py}). Thus $T_2=\tfrac{\sqrt2}{36}\sqrt\tau\,\sin W+O(\tau)$, with $\sin W=\sin
w+O(\sqrt\tau)$.
\end{proof}

\begin{lemma}[absolute bound on $T_2$]\label{lem:T2abs}
For $\tau\le0.02$, $\ |T_2|\le0.17\,\tau^{1/4}$; in particular $|T_2|\le0.064<1$ there, and
$|T_2|\to0$ as $\tau\to0$. The argument uses no saddle, no path through a saddle, and no cited
asymptotic theorem: only the Mellin--Barnes residues, a rectangle contained in the strip
$\mathcal S$ of Lemma~\ref{lem:Bbounded}, and the amplitude bounds proved there. The same holds at
any argument $X$ with $\tau X^2\in[2e^{-\tau},2]$; the constant above is the one for $X=w$, which
is the case Lemma~\ref{lem:infpoles} uses and the worse of the two.
\end{lemma}

\begin{proof}
Since $\operatorname{Res}_{s=n}\pi/\sin\pi s=(-1)^n$ and $h(s)=X^{2s}g_s/\Gamma(2s+1)$ is analytic
on $\Re s>-\tfrac12$ (Lemma~\ref{lem:Bbounded}(i)), the terms of \eqref{eq:T2} are the residues of
$h(s)\pi/\sin\pi s$ at $s=1,2,\dots$. Split the series at $n=\lfloor2w\rfloor$.

\emph{The tail.} At integers, $B_n=\sum_{i\le2n}\varphi(i\tau)-n\varphi(\tau)\ge n\varphi(\tau)\ge0$
by \eqref{eq:Bsum} and the monotonicity of $\varphi$, so $0\le g_n=1-e^{-B_n}\le1$ and
$\bigl|\sum_{n>\sigma_\ast}(-1)^ng_nX^{2n}/(2n)!\bigr|\le\sum_{n>\sigma_\ast}X^{2n}/(2n)!$. This
is a sawtooth in $\tau$, jumping each time $2w$ crosses an integer, with a superexponentially
decreasing envelope: $7.5\cdot10^{-10}$ at $\tau=0.02$ and $1.9\cdot10^{-11}$ at $\tau=0.0127$.
It is bounded together with the contour integral below and is nowhere the binding term.

\emph{The head.} The remaining residues are enclosed by the rectangle
$R=\{\tfrac12\le\Re s\le\sigma_\ast,\ |\Im s|\le X/2\}$, where $\sigma_\ast$ is the largest
half-odd-integer below $2w$ (so $|\sin\pi s|=\cosh\pi t\ge1$ on the right edge; taking the edge at
$2w$ itself would put it through a pole whenever $2w\in\Z$, e.g.\ at $\tau=0.02$), which lies \emph{wholly inside} $\mathcal S$, so the
amplitude bounds of Lemma~\ref{lem:Bbounded} apply at every point of $\partial R$; no bound on
$B_s$ outside $\mathcal S$ is needed anywhere. On $\Re s=\tfrac12$ we have $|\sin\pi s|=\cosh\pi t\ge1$,
so no pole is grazed. On the horizontal edges the two exponentials cancel:
$|\pi/\sin\pi s|\le\pi/\sinh(\pi X/2)\le2.01\pi e^{-\pi X/2}$ (valid once $X\ge1.69$, hence
throughout $\tau\le0.02$, where $X\ge9.9$), while $1/|\Gamma(2\sigma+1+iX)|$ carries the
compensating $e^{+\pi X/2}$; the resulting kernel is $O(\sqrt{2\pi/X})$, attaining $0.7887$ at $\sigma=\tfrac12$ against $\sqrt{2\pi/X}=0.7927$ when $\tau=0.02$.

For the amplitude: where $|2s|\le2w$, \eqref{eq:Btrunc} gives
$|B_s|\le\tfrac{\tau^2}{24}\bigl|\tfrac{M^3}3+\tfrac{M^2}2-\tfrac M3\bigr|+0.02\,\tau^{3/2}$ with
$M=2s$, hence $|g_s|=|1-e^{-B_s}|\le|B_s|e^{|B_s|}$; elsewhere on $\partial R$,
$|g_s|\le1+e^{30.3\sqrt\tau}$ by \eqref{eq:Bbounds}. Write $\mathcal B(\tau,X)$ for the resulting
majorant: the boundary integral of those bounds against
$\bigl|X^{2s}/\Gamma(2s{+}1)\bigr|\,\bigl|\pi/\sin\pi s\bigr|$, divided by $2\pi$, plus the tail.
The prose above fixes the geometry; what remains is to bound $\mathcal B$ on the whole region
$\tau\in(0,0.02]$, $\tau X^2\in[2e^{-\tau},2]$, and that is done in two ranges.

\emph{(1) $\tau\ge2/101^2$: verified enclosure.} On $w=\sqrt{2/\tau}\in[10,101]$,
$X\in[we^{-\tau/2},w]$, the majorant is enclosed in MPFR interval arithmetic with outward directed
rounding at every operation, by a branch and bound on the parameter box: a box is accepted when the
enclosed supremum of $\mathcal B$ over the whole box does not exceed the infimum of
$0.17\,\tau^{1/4}$ on it, and is bisected otherwise. The contour integrals are enclosed cellwise,
each cell contributing its integrand enclosure times its width, so what is returned bounds the
integral rather than estimating it; the cells tile each edge exactly, and the initial boxes are cut
at the half-integer steps of $\sigma_\ast$, so no box straddles a jump of the contour. The scan
terminates with $232$ accepted boxes covering the range and no failures, and returns
$\mathcal B\le0.16978\,\tau^{1/4}$, a box-level figure comparing the largest bound on a box with
the smallest $\tau^{1/4}$ on it. Run at $96$ and again at $160$ bits it produces the same box
decomposition and the same figure. Pointwise at $\tau=0.02$, $X=w=10$, where the cross-check below locates the maximum of the ratio,
the enclosure is $\mathcal B\le0.061685=0.164031\,\tau^{1/4}$, and the scan repeated against the constant $0.166$ in
place of $0.17$ also terminates without failures. (\texttt{tools/t2abs\_iv}; the earlier
adaptive-quadrature certificate \texttt{t2abs\_certificate.py} shares no code with it and is kept as
a cross-check, its seven tabulated ratios $0.1623,\ 0.0428,\ 0.0308,\ 0.0226,\ 0.0173,\ 0.0136,\
0.0115$ at $\tau=2\cdot10^{-2},\dots,10^{-5}$ being reproduced from above.)

\emph{(2) $\tau<2/101^2$: explicit constants.} No finite scan reaches $\tau\to0$. Here
$X\ge101e^{-1/101^2}>100$, $w\le Xe^{1/X^2}\le1.0001X$ and $\tau\le2/X^2$, so
$0.17\,\tau^{1/4}=0.17\cdot2^{1/4}w^{-1/2}\ge0.2021\,X^{-1/2}$, and it suffices to bound
$\mathcal B$ by that.

On the left edge $2s+1=2+2it$, and $|\Gamma(1+iy)|^2=\pi y/\sinh\pi y$ makes the kernel exact:
$\bigl|X^{2s}/\Gamma(2s{+}1)\bigr|\,\bigl|\pi/\sin\pi s\bigr|=X\pi\sqrt{g(\pi t)}/\sqrt{1+4t^2}$
with $g(u)=\tanh u/u\le\min(1,1/u)$, hence at most $X\pi$ for $t\le1/\pi$ and at most
$\tfrac12X\sqrt\pi\,t^{-3/2}$ beyond. There $|M|=\sqrt{1+4t^2}\le1+2t\le1+X<2w$, so
\eqref{eq:Btrunc} applies along the whole edge; with
$\tfrac{(1+2t)^3}3+\tfrac{(1+2t)^2}2+\tfrac{1+2t}3=\tfrac83t^3+6t^2+\tfrac{14}3t+\tfrac76$
integrated against $t^{-3/2}$, and collecting the range $t\le1/\pi$ and the
$0.02\,\tau^{3/2}$ remainder on the same power of $X$, the edge contributes
$L\le0.03033\,X^{-1/2}$ for $X\ge100$.

On a horizontal edge put $a=2\sigma+1$, $u=a/X$ and
$\psi(u)=u-\arctan u-\tfrac u2\log(1+u^2)$, so $\psi'(u)=-\tfrac12\log(1+u^2)$ and $\psi$ is
negative, decreasing and concave on $u>0$. Stirling with Binet's remainder,
$\log\Gamma(z)=(z-\tfrac12)\operatorname{Log}z-z+\tfrac12\log2\pi+\mu(z)$ with
$|\mu(z)|\le\mu(\Re z)\le1/(12\Re z)$ (the kernel of Binet's first integral is positive on
$(0,\infty)$), cancels every term of size $a\log X$ and leaves
\begin{equation}\label{eq:hker}
   \frac{X^{2\sigma}}{|\Gamma(a+iX)|}\cdot\frac{\pi}{\sinh(\pi X/2)}
   \;=\;\frac{\sqrt{2\pi}}{1-e^{-\pi X}}\;X^{-1/2}(1+u^2)^{1/4}e^{X\psi(u)}e^{-\Re\mu},
   \qquad|\Re\mu|\le\tfrac1{12a}\le\tfrac1{24}.
\end{equation}
With $d\sigma=\tfrac X2du$ and $|M|\le X\sqrt{1+u^2}$, \eqref{eq:Btrunc} applies for $u\le\sqrt3$,
where $|M|\le2X\le2w$, and gives $|g_s|\le1.00449\bigl[(1+u^2)^{3/2}/(18X)+0.33502\,X^{-2}\bigr]$;
for $u>\sqrt3$ only $|g_s|\le1+e^{30.3\sqrt\tau}\le2.5352$ is used. Two elementary inequalities
close the $u$-integral: $\psi(u)\le-u^3/12$ on $[0,1]$, since $\log(1+x)\ge x/2$ there; and
$\psi(u)\le\psi(1)+\psi'(1)(u-1)$ for $u\ge1$, by concavity. With
$\int_0^\infty e^{-Xu^3/12}du=\Gamma(\tfrac43)(12/X)^{1/3}$ and $(1+u^2)^{7/4}\le2^{7/4}$ on
$[0,1]$, one edge contributes
$H\le0.50142\,X^{-5/6}+1.06902\,X^{-11/6}+4.5\cdot10^{-9}e^{-0.131(X-100)}$.

The right edge and the tail are exponentially small. On $\Re s=\sigma_\ast$ one has
$|\pi/\sin\pi s|\le\pi$ and $a=2\sigma_\ast+1>4X-1$, while
$\log|\Gamma(a+iY)|\ge(a-\tfrac12)\log a-Y^2/a-a+\tfrac12\log2\pi-\tfrac1{12a}$ for $|Y|\le X$;
since $c\mapsto c(1-\log c)$ decreases and $a/X\in[3.99,4.011]$, the kernel there is at most
$\pi e^{-1.2805X}$. And $n_0=\lfloor\sigma_\ast\rfloor+1>2X-1$ with $(2n)!\ge(2n/e)^{2n}$ makes the
tail at most $1.07\,e^{-1.51X}$. Both are below $10^{-50}$ throughout. Collecting, for $X\ge100$,
\[
   \mathcal B\;\le\;0.00966\,X^{-1/2}+0.15961\,X^{-5/6}+0.34028\,X^{-11/6},
\]
which is at most $0.2021\,X^{-1/2}$ as soon as
$0.15961\,X^{-1/3}+0.34028\,X^{-4/3}\le0.2021-0.00966=0.19244$; the left side decreases and
equals $0.03512$ at $X=100$. The constants of this range are regenerated, and the two $\psi$ inequalities and the
rearrangement \eqref{eq:hker} independently checked against \texttt{mpmath}'s complex $\Gamma$, in
\texttt{t2abs\_smalltau.py} at $40$ and again at $60$ digits.

At $\tau=\tau_4=0.012665$, the largest value Lemma~\ref{lem:infpoles} uses, the bound is $0.0144$,
and what Lemma~\ref{lem:infpoles} consumes there is only $|T_2|<1$.
\end{proof}

\begin{remark}[why Remark~\ref{lem:cos} is not used]\label{rem:cosgap}
The steepest-descent sketch is recorded because it produces the sharp constant, but it is not a
proof, on three counts. (i) It asserts that the two saddles sum to $\cos W$ \emph{exactly}; they do
not: the leading saddle term carries a phase $e^{-i/(24W)}$, so the two sum to
$\cos W+\tfrac1{24W}\sin W+O(W^{-2})$. (ii) It sizes the first saddle correction by
$h'''(z^\ast)=O(W^{-2})$; but the saddle is \emph{flat} ($h''=4i/W\to0$), so the expansion parameter
is $W^{-1}$ and the correction is $c_1=i/(24W)+O(W^{-2})$. Errors (i) and (ii) are each about
three quarters of the quantity being computed and cancel exactly, which is why the stated constant
is nevertheless right; taken literally the chain yields $\sqrt2/144$, four times too small.
(iii) The deformation is justified by the claim that $\csc$ suppresses $|\Im z|\to\infty$, which is
false for the exact $B$, as the numbers in the previous proof show; and the uniform remainder is
attributed to a steepest-descent error estimate whose hypotheses (a parameter-independent phase and
amplitude, an explicit path, an explicit truncation order) are not verified here. A complete proof
needs the two-term flat-saddle expansion together with an error estimate on the curved path; that is
not written, and nothing in this paper waits on it.
\end{remark}

\begin{lemma}[infinitude and accumulation of the travel poles]\label{lem:infpoles}
The set $\{q\in(0,1):\Sigma_1^T(q)=1\}$ is infinite, and together with
Lemma~\ref{lem:polefinite} it accumulates exactly at $q=1$. In particular there is a travel pole
with $w\in(m\pi,(m+1)\pi)$ for every large $m$.
\end{lemma}

\begin{proof}
By Proposition~\ref{prop:qtrig}, $1-\Sigma_1^T=\cos(Z;q^2)$ with $Z^2=z_0^2e^{\tau}$. The collapse
above defines the radius-free coefficients $e^{-B_n}=q^{n^2+n}[2(1-q)]^n(2n)!/[(q;q)_{2n}W^{2n}]$
through $S_e=\cos(z_0;q^2)=\sum_{n\ge0}(-1)^ne^{-B_n}W^{2n}/(2n)!$; replacing $z_0^2$ by
$Z^2=z_0^2e^\tau$ multiplies the $n$-th coefficient by $e^{n\tau}$, i.e.\ replaces $W$ by
$We^{\tau/2}=w$ exactly:
\[
 1-\Sigma_1^T\;=\;\sum_{n\ge0}(-1)^n e^{-B_n}\,\frac{w^{2n}}{(2n)!}\;=\;\cos w-T_2(w),
\]
with $T_2(\cdot)$ the series \eqref{eq:T2} evaluated at argument $w$ in place of $W$; the
coefficients $e^{-B_n}$ are the \emph{same} ones, so no re-derivation of the amplitude is involved.
Only the argument moves, and it moves inside the region where the two estimates are already
proved: Lemma~\ref{lem:Bbounded} is stated on a strip built from the larger argument $w$ (and
containing the $W$-strip), so it applies verbatim, and Lemma~\ref{lem:T2abs} is stated for any
argument with $\tau X^2\in[2e^{-\tau},2]$, hence at $X=w$, where $\tau w^2=2$ identically. Thus
$|T_2(w)|\le0.064$ once $\tau\le0.02$, which along the extreme-phase sequence $w=m\pi$,
$\tau^{(m)}=2/(m\pi)^2$, holds for every $m\ge4$ (there $\tau\le\tau_4=0.012665$ and the bound is
$0.0144$). Along that sequence
\[
 1-\Sigma_1^T\big|_{w=m\pi}\;=\;(-1)^m-T_2(m\pi),
\]
which has the sign of $(-1)^m$ for all $m\ge m_0$, with $m_0=4$ from
Lemma~\ref{lem:T2abs} (only $|T_2(m\pi)|<1$ is used, a far weaker input than the sharp size
recorded in Remark~\ref{lem:cos}). By continuity of $q\mapsto\Sigma_1^T(q)$ the
intermediate value theorem yields a zero of $1-\Sigma_1^T$ with $w\in(m\pi,(m+1)\pi)$ for every
$m\ge m_0$: infinitely many travel poles, with $\tau\downarrow0$, i.e.\ $q\uparrow1$. Since every
travel pole has $q>\tfrac1{10}$ and the pole set is locally finite in $(0,1)$
(Lemma~\ref{lem:polefinite}), its only accumulation point is $q=1$.
\end{proof}

\noindent The alternation is in fact immediate: $|T_2(\pi)|\le0.027$ by direct evaluation and
$|T_2(m\pi)|=O(\tau_m)$ thereafter, so $m_0=1$ and there is a travel pole in every window
$(m\pi,(m+1)\pi)$, $m\ge1$. The dominant pole $q_1$ is the one below this ladder, in
$w\in(0,\pi)$. The transport identity and the sign alternation (checked for $m\le40$, worst
$|T_2(m\pi)|=0.0267$) are verified in \texttt{star\_gate\_certificate.py}.

\emph{Alternative proof of the first sentence, without $T_2$.} Proposition~\ref{prop:infpolesalt}
proves that the travel poles are infinite from the Fredholm determinant of
Theorem~\ref{thm:fredholm} alone: $\lambda_1(q)>(1-q^2)/(2q)$, the uniform domination
$|t_k|\le(2|y|)^k/(2k-1)!!$ with Hurwitz's theorem, which gives
$\lambda_k(q)\sim(1-q)(2k-1)^2\pi^2/8$ at fixed $k$, and the fact that distinct branches cross the
level $1$ at distinct $q$; no single-crossing property is needed. With Lemma~\ref{lem:polefinite}
this also gives the accumulation at $q=1$. It does not localise the poles: the window statement
above, a travel pole with $w\in(m\pi,(m+1)\pi)$ for every large $m$, still uses
Lemma~\ref{lem:T2abs}.

\begin{lemma}[infinitude of the zeros of the bulk block]\label{lem:infzeros}
The zero set of $S_e$ in $(0,1)$, equivalently $\{q\in(0,1):S_1(q)=1\}$, is infinite.
\end{lemma}

\begin{proof}
Here no transport is needed: $S_e=\cos W-T_2$ at $W=we^{-\tau/2}$ is \eqref{eq:T2} itself, the
series for which both estimates were set up. The map
$\tau\mapsto W(\tau)=\sqrt{2/\tau}\,e^{-\tau/2}$ is continuous and strictly decreasing on
$(0,\infty)$ with $W\to\infty$ as $\tau\downarrow0$, so each integer $m\ge4$ gives a unique
$\tau^{[m]}$ with $W(\tau^{[m]})=m\pi$; since $W<w$, $\tau^{[m]}<2/(m\pi)^2\le\tau_4=0.012665$, and
$\tau W^2=2e^{-\tau}$ identically, so Lemma~\ref{lem:T2abs} applies at every one of them and only
$|T_2|<1$ is used. Then $S_e=(-1)^m-T_2$ has the sign of $(-1)^m$ along the sequence, and the
$\tau^{[m]}$ decrease in $m$, so the intermediate value theorem applied to the continuous function
$q\mapsto S_e(q)$ (Theorem~\ref{thm:blocks}, first clause) gives a zero in each of the disjoint
intervals $(e^{-\tau^{[m]}},e^{-\tau^{[m+1]}})$, $m\ge4$. Finally $S_e=1-S_1$ by
Proposition~\ref{prop:qtrig}.
\end{proof}

\noindent The zeros of $S_e$ are not the travel poles: those are the zeros of $\cos(Z;q^2)$, which
is the same function a half $q$-step away (Proposition~\ref{prop:qtrig}), and $S_e$ is bounded away
from $0$ at every travel pole by Proposition~\ref{prop:selower}. Lemma~\ref{lem:infzeros} is used
only in Corollary~\ref{cor:polepath}. The alternation it asserts is checked directly for
$4\le m\le40$ in \texttt{f2\_pole\_route.py}, at $60$ and again at $100$ digits: the series for
$S_e$ cancels heavily against intermediate terms far larger than its value, and $30$ digits already
returns wrong signs there.

\begin{remark}[the Hubbard--Stratonovich integral representation]\label{rem:gaussint}
Writing $q^{j(j+1)}=q^{j}e^{-j^2\tau}$ and linearising the Gaussian through
$e^{-j^2\tau}=(4\pi\tau)^{-1/2}\int_{\R}e^{-u^2/(4\tau)}e^{iju}\,du$, the denominator block admits the
\emph{exact} real integral
\begin{equation}\label{eq:HS}
   S_e=\frac1{\sqrt{4\pi\tau}}\int_{-\infty}^{\infty}e^{-u^2/(4\tau)}\,\Psi(u,q)\,du,
   \qquad
   \Psi(u,q)=\sum_{j\ge0}\frac{\bigl(-2(1-q)q\,e^{iu}\bigr)^{j}}{(q;q)_{2j}},
\end{equation}
the interchange of sum and integral being justified by absolute convergence
($\sum_j|Z|^{j}/(q;q)_{2j}<\infty$, $Z=-2(1-q)q\,e^{iu}$, against the integrable Gaussian). Because the
coefficients of $\Psi$ in $\zeta=e^{iu}$ are real, $\Psi(-u,q)=\overline{\Psi(u,q)}$, so \eqref{eq:HS}
is real and may be written over a single variable, $S_e=(4\pi\tau)^{-1/2}\!\int_\R
e^{-u^2/(4\tau)}\Re\Psi\,du$; the exact Gaussian moment underlying it,
$\int_{\R}e^{-u^2/(4\tau)}e^{iju}\,du=\sqrt{4\pi\tau}\,e^{-\tau j^2}$, is machine-checked in Lean~4
(\texttt{GaussHS.lean}, theorem \texttt{gaussian\_moment}; axioms \texttt{propext},
\texttt{Classical.choice}, \texttt{Quot.sound}, no \texttt{sorry}).

This representation, though exact and convenient, is \emph{not} an analytic shortcut. The amplitude $\Psi(\cdot,q)$ is entire in $\zeta$, but it is
\emph{not} uniformly bounded as $q\to1^-$: on any fixed window it grows like $\cosh\!\bigl(w\sin(u/2)\bigr)$
(numerically $|\Psi|$ already exceeds $300$ at $u=1,\ \tau=10^{-2}$), and its termwise limit
$\cos(w\,e^{iu/2})$ \emph{differs} from $\Psi$ by an $O(1)$ factor at the dominant scale $j\sim\tau^{-1/2}$
(where the subleading part of $(q;q)_{2j}$ is itself $O(1)$; numerically $\Psi/\cos(we^{iu/2})\approx0.62$).
The size of $S_e$ is therefore governed by the balance between this exponential growth and the Gaussian; a
simple \emph{complex} saddle at $u^\ast\approx-\sqrt{2\tau}$, of the same order of difficulty as the direct
sum, \emph{not} a bounded-amplitude stationary-phase integral with trivial hypotheses
\cite[Ch.~II]{Wong}. In particular the
bounded-variation input of Lemma~\ref{lem:Bbounded} is \emph{not} eliminated by passing to \eqref{eq:HS};
it, together with the rectangle bound of Lemma~\ref{lem:T2abs}, remains the analytic core. The value
of \eqref{eq:HS} is representational (one real variable, an entire integrand, an exact Lean-checked moment),
not a reduction of hypotheses; the numerical agreement of its saddle value with $S_e$ is recorded in
\texttt{gaussint\_verify.py}.
\end{remark}

\begin{remark}[an elementary derivation of the coefficients of Remark~\ref{lem:cos}]\label{rem:elemcos}
The leading term and the constant $\sqrt2/36$ of Remark~\ref{lem:cos} can be obtained without the
steepest-descent theorem, by an exact algebraic reduction. Writing $1-e^{-k\tau}=2e^{-k\tau/2}\sinh(k\tau/2)$
throughout, the $j$-th term of $S_e=\sum_{j\ge0}(-2(1-q))^jq^{j(j+1)}/(q;q)_{2j}$ collapses to
\[
   S_e=\sum_{j\ge0}(-1)^j e^{-j\tau}\,\frac{\sinh^j(\tau/2)}{\prod_{k=1}^{2j}\sinh(k\tau/2)}
      =\sum_{j\ge0}(-1)^j\frac{(2/\tau)^j}{(2j)!}\,e^{-j\tau+E_j},
   \quad E_j=\tfrac{j\tau^2}{24}-\tfrac{\tau^2}{144}(2j)(2j{+}1)(4j{+}1)+O(\tau^4j^5),
\]
$E_j$ being the exact $\sinh$-curvature defect. Two elementary facts finish the extraction. First, the
factor $e^{-j\tau}$ resums \emph{in closed form}: $\sum_j(-1)^j(2e^{-\tau}/\tau)^j/(2j)!=\cos W$ with
$W=\sqrt{2/\tau}\,e^{-\tau/2}$, which is exactly the frequency of Remark~\ref{lem:cos}; so the phase
correction needs no estimate at all. Second, the entire-function moment identity
$\sum_j(-1)^j y^j j^m/(2j)!=(y\,d/dy)^m\cos\sqrt y$ evaluates the curvature term through its leading part
$-\tfrac{\tau^2}{9}(y\,d/dy)^3\cos\sqrt y=-\tfrac{\tau^2}{9}\cdot\tfrac{W^3}{8}\sin W=-\tfrac{\sqrt2}{36}\sqrt\tau\,\sin W$
(using $\tau^2W^3=2\sqrt2\,\sqrt\tau\,e^{-3\tau/2}$, the factor $e^{-3\tau/2}$ being absorbed in the $O(\tau)$), giving
\[
   S_e=\cos W-\tfrac{\sqrt2}{36}\sqrt\tau\,\sin W+O(\tau),\qquad W=\sqrt{2/\tau}\,e^{-\tau/2},
\]
all constants explicit and verified to ten digits (\texttt{elemcos\_verify.py}). Carried to higher moments
the same reduction produces every subsequent coefficient as a finite combination of $W^k\sin W,\,W^k\cos W$,
so it serves the $O(\tau^2)$-relative demand of the $U$-numerator in the same elementary terms. The one input
it does \emph{not} render elementary is the remainder estimate; that the alternating expansion is
asymptotic with $O(\tau)$ error. That step is a \emph{cancellation} statement: not only does the triangle
inequality on $\sum_j(2/\tau)^j/(2j)!\,|\cdot|$ return $\cosh W\asymp e^{\sqrt{2/\tau}}$, but the partial
sums of the $\cos W$ series themselves swing up to $\asymp e^{c\sqrt{2/\tau}}$ before cancelling to $O(1)$
(numerically $\max_J|S_J|=9,\,48,\,1459,\,7.6\!\times\!10^4$ at $\tau=0.1,0.05,0.02,0.01$). Hence \emph{every}
partial-summation method (Abel, summation-by-parts, Euler transformation) fails, since each is
controlled by these exponentially large partial sums; only a global resummation captures the cancellation.
This global resummation is exactly the two-saddle steepest-descent computation carried out
self-contained in the proof of Remark~\ref{lem:cos}: the saddles $z^\ast_\pm=\pm iW/2$ supply the
cancellation the partial sums cannot. Together, the elementary extraction here and that saddle argument reduce
Remark~\ref{lem:cos} to Cauchy's theorem, Stirling's formula, elementary Gaussian integration, the
bounded-variation bound of Lemma~\ref{lem:Bbounded}, and \emph{one} cited estimate: the classical
remainder bound for a simple saddle \cite[Ch.~4, Thm.~7.1]{Olver}, used once in the proof of
Remark~\ref{lem:cos}. It is not independent of that citation, and we do not claim otherwise;
Appendix~\ref{app:annulus} provides the route that is.
\end{remark}

\begin{remark}[derivation of the $U$-numerator amplitude]\label{rem:elemY3}
The sharp picture behind Theorem~\ref{thm:star}: the asymptotic value
$|\sum_kd_k|=\tfrac3{8\sqrt2}\tau^{5/2}(1+O(\sqrt\tau))$ at the travel poles, on the one explicit
series $P_{12}=\tfrac{2q^3}{1-q^3}\sum_kd_k$,
$d_k=(-2)^k(1-q)^kq^{k^2+3k}/[(q^2;q^2)_k(q^5;q^2)_k]$, which submits to the identical reduction. Writing
$1-q^m=2e^{-m\tau/2}\sinh(m\tau/2)$ throughout collapses the terms \emph{exactly} to
\[
   d_k=(-1)^k e^{-k\tau}\,
   \frac{\sinh^{k+1}(\tfrac\tau2)\,\sinh(\tfrac{3\tau}2)\,\sinh((k{+}1)\tau)}
        {\prod_{m=1}^{2k+3}\sinh(\tfrac{m\tau}2)},
\]
and the elementary expansion (exact linear absorption into the argument, Faulhaber sums for the
$\sinh$-curvature, the moment identities $k^m\mapsto(y\,d/dy)^m$) yields, with
$\Phi(y)=(\sin W-W\cos W)/(2W^3)$, $W=\sqrt y$, $D=y\,d/dy$,
\[
   \sum_kd_k=6\Bigl[\Phi+c_3D^3\Phi+c_2D^2\Phi+\tfrac{c_3^2}2D^6\Phi
   +(c_5+c_2c_3)D^5\Phi+\tfrac{c_3^3}6D^9\Phi\Bigr]\bigl(y_*\bigr)+O(\tau^3),
\]
$y_*=\tfrac2\tau e^{-\tau-23\tau^2/36}$, $c_2=-\tfrac5{12}\tau^2$, $c_3=-\tfrac19\tau^2$,
$c_5=\tfrac1{450}\tau^4$; \emph{every constant an exact rational, derived, not fitted}
(\texttt{elemY3\_verify.py}; the error is verified $O(\tau^3)$ over $\tau\in[0.005,0.08]$ and the formula
reproduces the exact series at the poles $m=7,\dots,33$ to the displayed accuracy). At the travel poles the
expansion gives $|\sum_kd_k|/\tau^{5/2}\to3/(8\sqrt2)=0.2652$, i.e.\ exactly the observed gate value
$|P_{12}|/\tau^{3/2}\to1/(4\sqrt2)$ with its fourfold margin: structurally, the pole sits at offset
$\delta w$ from the zero of the displayed expansion, and the constant $3/(8\sqrt2)$ arises as
$(3\tau/2)\,\delta w/\tau^{5/2}$.

The \emph{phase-matching step} that fixes $\delta w$ is carried out by the same machinery. The pole
equation $\Sigma_1(q_m)=1$ collapses exactly to
$1-\Sigma_1=\sum_{j\ge0}(-1)^j\sinh^j(\tfrac\tau2)/\prod_{m=1}^{2j}\sinh(\tfrac{m\tau}2)$; the
$S_e$-series stripped of its $e^{-j\tau}$ phase; so the identical expansion applies with
$y_*=\tfrac2\tau e^{\tau^2/36}$, $c_2=-\tfrac1{12}\tau^2$ and the \emph{universal} $c_3=-\tfrac19\tau^2$,
$c_5=\tfrac1{450}\tau^4$. Solving the two phase conditions perturbatively in the common variable
$w=\sqrt{2/\tau}$ near each grid point $(M+\tfrac12)\pi$ gives
\[
   \chi_{\mathrm{pole}}=-\tfrac{\sqrt2}{36}\sqrt\tau+\tfrac{41\sqrt2}{2400}\tau^{3/2}+O(\tau^2),
   \qquad
   \chi_{\mathrm{zero}}=-\tfrac{\sqrt2}{36}\sqrt\tau-\tfrac{259\sqrt2}{2400}\tau^{3/2}+O(\tau^2).
\]
The $\sqrt\tau$ phases agree \emph{identically}; both are produced by the universal cubic curvature
$c_3$, and the constant $\sqrt2/36$ is that of Remark~\ref{lem:cos}; this is precisely why the travel
poles track the zeros of $Y_3(1)$. The $\tau^{3/2}$ terms differ by exactly
\[
   \delta w=\chi_{\mathrm{pole}}-\chi_{\mathrm{zero}}=\tfrac{\sqrt2}8\,\tau^{3/2}+O(\tau^2),
   \qquad\text{whence}\qquad
   |\!\sum_kd_k|=\tfrac{3\tau}2\,\delta w\,(1+O(\sqrt\tau))=\tfrac3{8\sqrt2}\tau^{5/2}(1+O(\sqrt\tau)),
\]
and therefore $|P_{12}|/\tau^{3/2}\to1/(4\sqrt2)=0.17678<1/\sqrt2$: \emph{the numerator gate, derived,
with its fourfold margin}. Verification: the two phase series match the exact pole and zero locations to
seven digits at $m=11,17,23,29$, and $\delta w/\tau^{3/2}=0.17699,\,0.17687,\,0.17683,\,0.17681\to
\sqrt2/8=0.176777$ (\texttt{phase\_match\_verify.py}); the pole-side series independently reproduces the
McMahon-type coefficients of the travel-pole expansion. The remainder discipline throughout is that of
Remark~\ref{lem:cos}: every coefficient extracted elementarily, the $O(\cdot)$ terms controlled by the same
two-saddle argument.
\end{remark}

\subsection{The travel-pole equation as a $q$-trigonometric identity}\label{sec:qtrig}

The two series that govern the poles admit an exact identification with the $q$-cosine of
Koornwinder and Swarttouw \cite{KS1992}, $\cos(z;q^2):={}_1\phi_1(0;q;q^2,q^2z^2)
=\sum_{j\ge0}(-1)^j q^{j^2+j}z^{2j}/(q;q)_{2j}$.

\begin{proposition}\label{prop:qtrig}
Let $z_0=\sqrt{2(1-q)}$ and $Z=z_0/\sqrt q$. Then
\[
   1-S_1=S_e=\cos(z_0;q^2),\qquad 1-\Sigma_1=\cos(Z;q^2).
\]
In particular the travel poles $q_m$ are exactly the solutions of $\cos(Z;q^2)=0$, and $S_e$ is the
same function evaluated one half $q$-step below its own zero.
\end{proposition}

\begin{proof}
Collapse the $k$-recursion \eqref{eq:krec} on its two odd ladders, the travel one carrying
$\Sigma_1$ and the bulk one carrying $S_1$. Write $1-q^m=2e^{-m\tau/2}\sinh(m\tau/2)$ in the
$\sinh$-forms of the two series
(\S\ref{sec:sd}): the products collapse to
$1-\Sigma_1=\sum_j(-1)^j\bigl(2(1-q)\bigr)^jq^{j^2}/(q;q)_{2j}$ and
$1-S_1=\sum_j(-1)^j\bigl(2(1-q)\bigr)^jq^{j^2+j}/(q;q)_{2j}$, the second being $S_e$ as defined in
\eqref{eq:blocks}. Comparing with the series for
$\cos(z;q^2)$ term by term forces $q^2z^2=2(1-q)q$, i.e.\ $z=Z$, in the first case and
$q^2z^2=2(1-q)q^2$, i.e.\ $z=z_0$, in the second; both series converge for $|q|<1$, the factor
$q^{j^2}$ dominating. Both collapses are recomputed as identities of integer power series, to
degree $1160$ in \texttt{blocks\_growth.py} and at two truncation degrees in
\texttt{f2\_pole\_route.py}. The per-term identity is machine-checked in Lean~4
(\texttt{GramCasoratian.lean}, theorem \texttt{sigma1\_term}; axioms \texttt{propext},
\texttt{Classical.choice}, \texttt{Quot.sound}, no \texttt{sorry}).
\end{proof}

\noindent\emph{The Hahn--Exton normalisation.} Write
$J^{(3)}_\nu(x;q^2)=x^\nu\frac{(q^{2\nu+2};q^2)_\infty}{(q^2;q^2)_\infty}\,
{}_1\phi_1(0;q^{2\nu+2};q^2,q^2x^2)$ for the Hahn--Exton $q$-Bessel function, in the
normalisation of \cite[eq.~(2.4)]{KoelinkSwarttouw1994} with base $q^2$, and put
\[
   c_q:=\frac{(q^2;q^2)_\infty}{(q;q^2)_\infty}>0 .
\]
Since $(q^2;q^2)_k(q;q^2)_k=(q;q)_{2k}$ and $(q^2;q^2)_k(q^3;q^2)_k=(q;q)_{2k+1}/(1-q)$, comparing
coefficients gives, for $x>0$,
\[
   \cos(x;q^2)=c_q\,x^{1/2}J^{(3)}_{-1/2}(x;q^2),\qquad
   \sin(x;q^2)=c_q\,x^{1/2}J^{(3)}_{1/2}(x;q^2).
\]
So the travel poles are the zeros in $(0,1)$ of $q\mapsto J^{(3)}_{-1/2}(Z(q);q^2)$
(Corollary~\ref{cor:qbessel}); Proposition~\ref{prop:zerostructure}(i) records the same identity in
the variable of Appendix~\ref{app:beta2}.

\begin{remark}\label{rem:qtrig}
The identification is consistent with the pole data: at all $58$ tabulated poles the computed value
of $\cos(z_0;q^2)$ agrees with the stored $S_e$ to at least thirty digits, and $|S_e|/\sqrt\tau$
lies in $[0.698,0.708]$, approaching $1/\sqrt2$ (\texttt{qtrig\_verify.py}). The half-step
structure accounts for the size of the gate quantity: $Z-z_0=z_0(1-\sqrt q)/\sqrt q
=\tau^{3/2}/\sqrt2+O(\tau^{5/2})$, the pole condition annihilates $\cos(\,\cdot\,;q^2)$ at
$Z$, and the $q$-derivative rule $(1-q)D_{q,z}\sin(z;q^2)=\cos(z;q^2)$ of \cite{KS1992} gives the
function slope $1/(1-q)\sim1/\tau$ near the zero, so the sampled value is of size
$(\tau^{3/2}/\sqrt2)(1/\tau)=\sqrt{\tau/2}$, the observed constant. Turning this observation into a proved lower bound on
$|S_e(q_m)|$ independent of Remark~\ref{lem:cos} requires uniform control of
$\cos'(\,\cdot\,;q^2)$ between $z_0$ and $Z$ as $q\to1^-$. That control is supplied by
Appendix~\ref{app:annulus} (Lemma~\ref{app:M2} bounds $\max|f''|$ on $[qZ,Z]$, and
Proposition~\ref{prop:selower} converts it into the lower bound), independently of
Remark~\ref{lem:cos}.
\end{remark}

\subsection{The pole equation as a Fredholm determinant}\label{sec:fredholmdet}

The $q$-cosine of Proposition~\ref{prop:qtrig} is the Fredholm determinant of the travel operator.
Throughout this subsection $0<q<1$, $D_e=\operatorname{diag}(2q^{1+s})_{s\ge0}$ and
$K=[q^{\max(s,s')}]_{s,s'\ge0}$, so that $\mathcal T=D_eK$ is the travel operator of
Definition~\ref{def:model}. Put
\[
   \mathbf S:=D_e^{1/2}KD_e^{1/2},\quad \mathbf S[s,s']=2q^{1+(s+s')/2+\max(s,s')},\quad
   F(\lambda,q):=\sum_{k\ge0}\frac{(-2\lambda(1-q))^k\,q^{k^2}}{(q;q)_{2k}}\ \ (\lambda\in\C).
\]
The spectral parameter $\lambda$ is not to be confused with the junction vectors
$\lambda^{(y)}$ of \S\ref{sec:modelassembly}.

\begin{theorem}[the travel determinant]\label{thm:fredholm}
Let $0<q<1$ and $Z=\sqrt{2(1-q)/q}$.
\begin{enumerate}[label=\textup{(\roman*)},leftmargin=2.2em,itemsep=0.1em,topsep=0.2em]
\item $\mathbf S$ is a positive definite trace-class operator on $\ell^2(\Z_{\ge0})$, with
$\operatorname{tr}\mathbf S=2q/(1-q^2)$.
\item $\mathcal T$ and $\mathbf S$ have the same principal minors,
$\det\mathcal T[I,I]=\det\mathbf S[I,I]$ for every finite $I\subset\Z_{\ge0}$, and the same non-zero
eigenvalues, the eigenvectors corresponding by $R=D_e^{1/2}\phi$: if $\mathbf S\phi=\kappa\phi$ with
$\phi\in\ell^2$ and $\kappa\ne0$, then $R\in\ell^1$ and $\mathcal TR=\kappa R$; conversely, if
$R\in\ell^1$ and $\mathcal TR=\kappa R$ with $\kappa\ne0$, then $\phi=D_e^{-1/2}R\in\ell^2$ and
$\mathbf S\phi=\kappa\phi$.
\item For every $\lambda\in\C$, with absolutely convergent series,
\[
   \det(I-\lambda\mathbf S)=\sum_{k\ge0}(-\lambda)^k\!\!\sum_{|I|=k}\det\mathbf S[I,I]
   =F(\lambda,q),
\]
\[
   F(\lambda,q)=\cos(\sqrt\lambda\,Z;q^2)
   =c_q\,(\sqrt\lambda\,Z)^{1/2}J^{(3)}_{-1/2}(\sqrt\lambda\,Z;q^2),
\]
the last two expressions being even in $\sqrt\lambda$. We write $\det(I-\lambda\mathcal T)$ for the
same Fredholm series formed with the minors of $\mathcal T$; by \textup{(ii)} it is the same
function.
\item The eigenvalues of $\mathbf S$ are $1/\lambda_k(q)$, $k\ge1$, with
$0<\lambda_1(q)<\lambda_2(q)<\cdots\to\infty$, all simple; $\sum_k1/\lambda_k(q)=2q/(1-q^2)$ and
$F(\lambda,q)=\prod_{k\ge1}\bigl(1-\lambda/\lambda_k(q)\bigr)$.
\item $1-\Sigma_1=F(1,q)=\det(I-\mathcal T)$, and $S_e=F(q,q)=\det(I-q\mathcal T)=\det(I-M_0)$: the
bulk operator $M_0$ of Definition~\ref{def:model} is $q\mathcal T$ after the index shift $b=s+1$.
\end{enumerate}
\end{theorem}

\begin{proof}
(i) For $x\in\ell^2$ put $X=D_e^{1/2}x$, which lies in $\ell^1$ by Cauchy--Schwarz. Since
$q^{\max(s,s')}=(1-q)\sum_{t\ge\max(s,s')}q^t$,
\[
   \langle x,\mathbf Sx\rangle=\langle X,KX\rangle
   =(1-q)\sum_{t\ge0}q^t\Bigl|\sum_{s\le t}X_s\Bigr|^2\ \ge\ 0,
\]
the double sums converging absolutely, with equality only if every partial sum of $X$ vanishes,
that is $X=0$, that is $x=0$. $\mathbf S$ is Hilbert--Schmidt, hence bounded, and it is positive
definite; a positive operator with summable diagonal is trace class, with trace
$\sum_s2q^{1+s}q^s=2q/(1-q^2)$.

(ii) $\mathcal T[I,I]=D_e[I,I]\,K[I,I]$ and $\mathbf S[I,I]=D_e[I,I]^{1/2}K[I,I]D_e[I,I]^{1/2}$
with $D_e[I,I]$ diagonal, so the two determinants agree. If $\mathbf S\phi=\kappa\phi$, then
$R=D_e^{1/2}\phi\in\ell^1$ and $\mathcal TR=D_e^{1/2}\mathbf S\phi=\kappa R$. Conversely,
$\mathcal TR=\kappa R$ gives $R=\kappa^{-1}D_e(KR)$ with $|(KR)_s|\le\|R\|_{\ell^1}$, so
$\phi=D_e^{-1/2}R=\kappa^{-1}D_e^{1/2}KR\in\ell^2$ and
$\mathbf S\phi=D_e^{1/2}KR=\kappa\phi$. The operators are not boundedly similar, $D_e^{-1/2}$ being
unbounded; only \textup{(ii)} is used.

(iii) For trace-class $A$, $\det(I-\lambda A)=\sum_k(-\lambda)^k\operatorname{tr}\Lambda^k(A)$
\cite[Ch.~3]{Simon} (the von Koch form; compare Fredholm's formula for continuous kernels,
\cite[Thm.~3.10]{Simon}), and computing the trace of $\Lambda^k(\mathbf
S)$ in the orthonormal basis $e_{s_1}\wedge\cdots\wedge e_{s_k}$, $s_1<\cdots<s_k$, whose diagonal
entries are the principal minors, gives the first equality. The principal minors of the positive
$\mathbf S$ are $\ge0$, so the inner sums converge absolutely, to the values computed next, and the
outer series then converges absolutely for every $\lambda$. For $I=\{s_1<\cdots<s_k\}$ the matrix $K[I,I]$ is $(a_{\max(i,j)})$ with
$a_i=q^{s_i}$ decreasing; subtracting row $i+1$ from row $i$ for $i<k$ leaves a lower triangular
matrix with diagonal $a_1-a_2,\dots,a_{k-1}-a_k,a_k$, so
$\det K[I,I]=q^{s_k}\prod_{i<k}(q^{s_i}-q^{s_{i+1}})$ and, with the gaps
$g_i=s_{i+1}-s_i\ge1$,
\[
   \det\mathbf S[I,I]=2^kq^k\,q^{2(s_1+\cdots+s_k)}\prod_{i<k}\bigl(1-q^{g_i}\bigr).
\]
Since $s_1+\cdots+s_k=ks_1+\sum_{j<k}(k-j)g_j$, the sum over $s_1\ge0$ and $g_1,\dots,g_{k-1}\ge1$
factorises; with $\sum_{g\ge1}q^{2rg}(1-q^g)=q^{2r}(1-q)/\bigl((1-q^{2r})(1-q^{2r+1})\bigr)$,
\[
   \sum_{|I|=k}\det\mathbf S[I,I]
   =\frac{(2q)^k}{1-q^{2k}}\prod_{r=1}^{k-1}\frac{q^{2r}(1-q)}{(1-q^{2r})(1-q^{2r+1})}
   =\frac{(2(1-q))^k\,q^{k^2}}{(q;q)_{2k}} .
\]
Hence $\det(I-\lambda\mathbf S)=F(\lambda,q)$; the case $k=1$ is the trace. Finally
$(2\lambda(1-q))^kq^{k^2}=q^{k^2+k}(\sqrt\lambda\,Z)^{2k}$, since $Z^2=2(1-q)/q$, so
$F(\lambda,q)=\cos(\sqrt\lambda\,Z;q^2)$, and the Hahn--Exton form is the normalisation after
Proposition~\ref{prop:qtrig}.

(iv) For positive trace-class $\mathbf S$ the determinant is the product
$\prod_k(1-\lambda\kappa_k)$ over its eigenvalues $\kappa_k>0$ counted with multiplicity
\cite[Ch.~3]{Simon}, so the $1/\kappa_k$ are the zeros of $F(\cdot,q)$, the multiplicity of
$\kappa_k$ being the order of the zero. By \textup{(iii)} these are the $\lambda=x^2/Z^2$ with $x>0$
a zero of $\cos(x;q^2)=c_qx^{1/2}J^{(3)}_{-1/2}(x;q^2)$ ($\cos(\cdot;q^2)$ is even with
$\cos(0;q^2)=1$). For $\nu>-1$ the zeros of $J^{(3)}_\nu(\cdot;q^2)$ are real and simple and form
an infinite sequence \cite[Cor.~3.2, Lemma~3.3, Thm.~3.4]{KoelinkSwarttouw1994}; here
$\nu=-\tfrac12$ with base $q^2$. So the eigenvalues are simple and infinitely many, and the sum of
the $1/\lambda_k$ is the trace.

(v) At $\lambda=1$, $F(1,q)=\sum_k(-1)^k(2(1-q))^kq^{k^2}/(q;q)_{2k}=1-\Sigma_1$
(Proposition~\ref{prop:qtrig}); at $\lambda=q$, $F(q,q)=\sum_k(-2(1-q))^kq^{k^2+k}/(q;q)_{2k}=S_e$
by \eqref{eq:blocks}. And $M_0[s+1,s'+1]=2q^{s+1}q^{\max(s,s')+1}=q\,\mathcal T[s,s']$, so $M_0$
has the principal minors of $q\mathcal T$.
\end{proof}

\noindent The identity $\det(I-\lambda\mathbf S)=F(\lambda,q)$ was checked in double precision
against truncated matrices at $q=0.3,0.6,0.8$ and four values of $\lambda$ (one of them non-real),
to within $3.2\cdot10^{-15}$, and the gap sum against a brute-force sum of minors for $k\le3$ at
$q=0.3$.

\begin{remark}[a $q$-string]\label{rem:qstring}
With the Jackson integral $\int_0^1f\,d_qt=(1-q)\sum_{s\ge0}q^sf(q^s)$, let
$K_qf(x)=\int_0^1\min(x,t)f(t)\,d_qt$ on the lattice $x\in\{q^s\}$. At $x=q^s$,
$(KD_eu)_s=\sum_{s'}q^{\max(s,s')}2q^{1+s'}u_{s'}=\tfrac{2q}{1-q}(K_qu)(q^s)$, that is,
\[
   \mathcal T^{\!\top}=KD_e=\frac{2q}{1-q}\,K_q ,
\]
so up to the transposition, which changes neither the principal minors nor the eigenvalues,
$\mathcal T=\tfrac{2q}{1-q}K_q$. The kernel $\min(x,t)$ is the Green function of $-d^2/dx^2$ on
$[0,1]$ with a Dirichlet end at $0$ and a Neumann end at $1$: $\mathcal T$ is the compliance operator
of a $q$-string with uniform Jackson mass, fixed at $0$ and free at $1$, and $u=KR$ is its
displacement ($u_s\to0$ as $x=q^s\to0$). As $q\to1$ the kernel tends to the classical one, with
eigenvalues $4/((2k-1)^2\pi^2)$; this is the heuristic behind
$\lambda_k(q)\sim(1-q)(2k-1)^2\pi^2/8$, proved for fixed $k$ in
Proposition~\ref{prop:infpolesalt}.
\end{remark}

\begin{corollary}[the travel poles are Hahn--Exton zeros]\label{cor:qbessel}
A point $q\in(0,1)$ is a travel pole if and only if
\[
   J^{(3)}_{-1/2}\bigl(Z(q);q^2\bigr)=0,\qquad Z(q)=\sqrt{2(1-q)/q},
\]
if and only if $1$ is an eigenvalue of $\mathcal T(q)$, that is $\lambda_k(q)=1$ for some $k$. That
$k$ is unique, the eigenvalue is simple, its eigenvector is the travel null vector $R$ of
Proposition~\ref{prop:shape}, and $\phi=D_e^{-1/2}R$ is the corresponding eigenvector of $\mathbf S$.
\end{corollary}

\begin{proof}
Theorem~\ref{thm:fredholm}\textup{(iii)--(v)}, with $c_q>0$ and $Z>0$, and
Proposition~\ref{prop:qtrig}; uniqueness of $k$ is the simplicity in
Theorem~\ref{thm:fredholm}\textup{(iv)}, and the eigenvector is Theorem~\ref{thm:fredholm}\textup{(ii)}
with Lemma~\ref{lem:nullspace}.
\end{proof}

\noindent The zero theory used is that of Koelink--Swarttouw \cite{KoelinkSwarttouw1994}, at
$\nu=-\tfrac12$ and base $q^2$, where its hypothesis $\nu>-1$ holds. The eigenvalues also agree with
the Jacobi-matrix description of the zeros of $J^{(3)}_\nu$ by \v{S}tampach and \v{S}\v{t}ov\'{\i}\v{c}ek
\cite{StampachStovicek2014}, in its Friedrichs extension; the positive form of $\mathbf S$ is the
Friedrichs realisation, on which $u_s=(KR)_s\to0$. At $q_1,\dots,q_{10}$, located at $90$-digit
precision, $|J^{(3)}_{-1/2}(Z;q^2)|\le2\cdot10^{-79}$, and exactly $m-1$ zeros of
$\cos(\cdot;q_m^2)$ lie below $Z(q_m)$ (a numerical check, not a certificate).

\begin{proposition}[infinitely many travel poles, without $T_2$]\label{prop:infpolesalt}
\begin{enumerate}[label=\textup{(\alph*)},leftmargin=2.2em,itemsep=0.1em,topsep=0.2em]
\item $\lambda_1(q)>(1-q^2)/(2q)$ on $(0,1)$. In particular every $\lambda_k(q)>1$ for
$q\le\sqrt2-1$, and there is no travel pole in $(0,\sqrt2-1]$.
\item For each fixed $k$, $\lambda_k(q)/(1-q)\to(2k-1)^2\pi^2/8$ as $q\to1^-$.
\item Each $\lambda_k$ is continuous on $(0,1)$, and there is $c_k\in(\sqrt2-1,1)$ with
$\lambda_k(c_k)=1$. The $c_k$ are pairwise distinct travel poles; in particular there are infinitely
many travel poles.
\end{enumerate}
None of this uses $T_2$, Lemma~\ref{lem:T2abs}, or a single-crossing property.
\end{proposition}

\begin{proof}
(a) The largest eigenvalue of the positive trace-class $\mathbf S$ is strictly less than its trace,
since $\mathbf S$ has infinitely many positive eigenvalues: $1/\lambda_1<2q/(1-q^2)$. And
$(1-q^2)/(2q)\ge1$ exactly for $q\le\sqrt2-1$. At such $q$ no $\lambda_k$ equals $1$, so
$1-\Sigma_1=F(1,q)\ne0$ by Theorem~\ref{thm:fredholm}\textup{(iv)}.

(b) Put $\varepsilon=1-q$ and $\lambda=\varepsilon y$. The $k$-th term of $F(\varepsilon y,q)$ is
$t_k(y,\varepsilon)=(-1)^k(2y\varepsilon^2)^kq^{k^2}/(q;q)_{2k}$. By the arithmetic--geometric mean
inequality, $1-q^n=\varepsilon\sum_{i<n}q^i\ge n\varepsilon q^{(n-1)/2}$, so
$(q;q)_{2k}\ge\varepsilon^{2k}(2k)!\,q^{k(2k-1)/2}$ and
\[
   |t_k(y,\varepsilon)|\ \le\ \frac{(2|y|)^k\,q^{k/2}}{(2k)!}\ \le\ \frac{(2|y|)^k}{(2k-1)!!},
\]
uniformly in $\varepsilon\in(0,1)$. Each $t_k(y,\varepsilon)$ tends to $(-1)^k(2y)^k/(2k)!$ as
$\varepsilon\to0$, because $q^{k^2}\to1$ and $(q;q)_{2k}/\varepsilon^{2k}\to(2k)!$. The domination
is summable locally uniformly in $y\in\C$, so $F(\varepsilon y,1-\varepsilon)\to\cos\sqrt{2y}$
uniformly on compact sets. The limit has the simple zeros $y_j=(2j-1)^2\pi^2/8$, and the zeros of
$y\mapsto F(\varepsilon y,1-\varepsilon)$ are the $\lambda_j(q)/\varepsilon$, positive and ordered.
By Hurwitz's theorem, for $\varepsilon$ small each disc $|y-y_j|<\delta$, $j\le k$, contains exactly
one of them and $|y|\le y_k+\delta$ contains no other; hence $\lambda_j(q)/\varepsilon\to y_j$ for
$j\le k$.

(c) $q\mapsto\mathbf S(q)$ is continuous in trace norm, hence in operator norm, on $(0,1)$: its
entries are continuous and $\sum_{s,s'}|\mathbf S[s,s']|$ converges locally uniformly. The ordered
eigenvalues of positive compact operators are $1$-Lipschitz in operator norm (min--max), so each
$1/\lambda_k$, and so each $\lambda_k$, is continuous. By (a) $\lambda_k>1$ on $(0,\sqrt2-1]$, and by
(b) $\lambda_k(q)\to0$ as $q\to1$, so the intermediate value theorem gives $c_k$; it is a travel pole
by Corollary~\ref{cor:qbessel}. If $c_j=c_k$ with $j\ne k$, then $\lambda_j(c_k)=\lambda_k(c_k)=1$,
against simplicity (Theorem~\ref{thm:fredholm}\textup{(iv)}).

The same argument with $\lambda_k(q)/q$ in place of $\lambda_k(q)$ gives Lemma~\ref{lem:infzeros}
without $T_2$ too: $S_e=F(q,q)=\prod_k(1-q/\lambda_k(q))$ vanishes exactly where some
$\lambda_k(q)=q$; $\lambda_k(q)/q>(1-q^2)/(2q^2)>1$ for $q<1/\sqrt3$ by (a); $\lambda_k(q)/q\to0$ by
(b); and two branches cannot both equal $q$ at the same point.
\end{proof}

\subsection{The numerators, and why they do not vanish at the poles}\label{sec:numerators}

Proposition~\ref{prop:qtrig} identifies the two \emph{denominators}. The two numerators are the
companion $q$-sine at the same two arguments, and that is what settles their non-vanishing at the
travel poles. Throughout, $\mathfrak s(u)=q^{-1/2}\sin(q^{1/2}u;q^2)$ is the half-step $q$-sine of
Lemma~\ref{lem:P12closed}, $\sin(z;q^2)=\sum_{j\ge0}(-1)^jq^{j^2+j}z^{2j+1}/(q;q)_{2j+1}$.

\begin{lemma}[the numerators in closed form]\label{lem:qsinenum}
Identically in $0<q<1$,
\[
   \Sigma_0=\frac{2q}{Z}\,\mathfrak s(Z),\qquad S_0=\frac{2q}{z_0}\,\mathfrak s(z_0).
\]
With Proposition~\ref{prop:qtrig} this identifies all four blocks as one $q$-trigonometric pair
sampled at two arguments a half $q$-step apart: $(\Sigma_0,\,1-\Sigma_1)$ is
$\bigl(\tfrac{2q}u\mathfrak s(u),\,\cos(u;q^2)\bigr)$ at $u=Z$, and $(S_0,\,S_e)$ is the same pair
at $u=z_0=\sqrt q\,Z$. In particular $S_0=\tfrac{2q}{1-q}S_o$ with $S_o$ the odd bulk block of
Lemma~\ref{lem:P12closed}, which is the second identity of the dictionary \eqref{eq:dict} read as a
statement about the explicit series, and is therefore independent of the transfer model.
\end{lemma}

\begin{proof}
Collapse the $k$-recursion \eqref{eq:krec} on the even ladders, exactly as
Proposition~\ref{prop:qtrig} does on the odd ones. From $\alpha_k=2q^{k+1}/(1-q^{k+1})$,
\[
  |\gamma_k|=\alpha_k-\alpha_{k+1}=\frac{2(1-q)q^{k+1}}{(1-q^{k+1})(1-q^{k+2})},\qquad
  |C_k|=\frac{2(1-q)q^{k+2}}{(1-q^{k+1})(1-q^{k+2})} ,
\]
$\gamma_k$ and $C_k$ themselves being negative on $(0,1)$, so both even ladders alternate. For the bulk one,
$\sum_{i<j}(2i+1)=j^2$ and the denominators telescope to $(q;q)_{2j}$, giving
$\prod_{i<j}|\gamma_{2i}|=[2(1-q)]^jq^{j^2}/(q;q)_{2j}$ and hence
\[
  S_0=\sum_{j\ge0}(-1)^j\alpha_{2j}\prod_{i<j}|\gamma_{2i}|
     =2\sum_{j\ge0}(-1)^j\frac{[2(1-q)]^j\,q^{(j+1)^2}}{(q;q)_{2j+1}} .
\]
Write $q^{(j+1)^2}=q\,q^{j^2+j}q^{j}$ and $[2(1-q)]^jq^{j}=\bigl(q^{1/2}z_0\bigr)^{2j}$: the sum is
$(2q/(q^{1/2}z_0))\sin(q^{1/2}z_0;q^2)=(2q/z_0)\mathfrak s(z_0)$. For the travel ladder,
$\sum_{i<j}(2i+2)=j^2+j$ and the denominators telescope to $(q;q)_{2j}$ again, so
$\prod_{i<j}|C_{2i}|=[2(1-q)]^jq^{j^2+j}/(q;q)_{2j}$ and, with $A_{2j}=2q/(1-q^{2j+1})$,
\[
  \Sigma_0=2q\sum_{j\ge0}(-1)^j\frac{[2(1-q)]^j\,q^{j^2+j}}{(q;q)_{2j+1}}
          =\frac{2q}{z_0}\sin(z_0;q^2)=\frac{2q^{3/2}}{z_0}\,\mathfrak s(Z)=\frac{2q}{Z}\,\mathfrak s(Z),
\]
using $[2(1-q)]^j=z_0^{2j}$ and $z_0=q^{1/2}Z$. The $S_o$ statement is Lemma~\ref{lem:P12closed}'s
$S_o=\tfrac{z_0}2\mathfrak s(z_0)$ combined with $z_0^2=2(1-q)$. Both collapses are the exponent
bookkeeping machine-checked in \texttt{QTrigIdentities.lean} and are verified numerically off the
poles, hence as identities, in \texttt{qnumerator\_certificate.py}.
\end{proof}

\begin{lemma}[half-step rules and the pole invariant]\label{lem:halfstep}
Identically in $0<q<1$, with $c(u)=\cos(u;q^2)$:
\begin{enumerate}[label=\textup{(\roman*)},leftmargin=2.2em,itemsep=0.1em,topsep=0.2em]
\item $c(qu)=c(u)+q^{3/2}u\,\mathfrak s(q^{1/2}u)$;
\item $\mathfrak s(qu)=\mathfrak s(u)-u\,c(q^{1/2}u)$;
\item $c(u)\,c(q^{1/2}u)+q^{3/2}\,\mathfrak s(u)\,\mathfrak s(q^{1/2}u)=1$.
\end{enumerate}
Consequently, at every travel pole, where $c(Z)=0$:
\[
   q^{3/2}\,\mathfrak s(z_0)\,\mathfrak s(Z)=1 ,
\]
so neither $\mathfrak s(Z)$ nor $\mathfrak s(z_0)$ vanishes there.
\end{lemma}

\begin{proof}
(i) and (ii) are termwise. From $(q;q)_{2j}=(q;q)_{2j-1}(1-q^{2j})$,
$c(u)-c(qu)=\sum_{j\ge1}(-1)^jq^{j^2+j}u^{2j}/(q;q)_{2j-1}$; shifting $j\mapsto j+1$ and factoring
out $-qu$ leaves $-qu\sin(qu;q^2)=-q^{3/2}u\,\mathfrak s(q^{1/2}u)$, which is (i). From
$(q;q)_{2j+1}=(q;q)_{2j}(1-q^{2j+1})$,
$\mathfrak s(u)-\mathfrak s(qu)=q^{-1/2}\sum_{j\ge0}(-1)^jq^{j^2+j}(q^{1/2}u)^{2j+1}/(q;q)_{2j}$;
factoring out $q^{1/2}u$ leaves $u\,c(q^{1/2}u)$, which is (ii). This is the exponent
matching machine-checked as \texttt{cshift\_exponent} and \texttt{sshift\_exponent}. For (iii),
substitute (i) and (ii) into the Hahn--Exton $q$-Wronskian
$\mathfrak s(u)c(qu)-c(u)\mathfrak s(qu)=u$ (\cite[eq.~(4.3.8)]{Swarttouw1992};
\cite[eq.~(4.21)]{Daalhuis1994}): the $\mathfrak s(u)c(u)$ terms cancel and what is left is
$u\bigl[c(u)c(q^{1/2}u)+q^{3/2}\mathfrak s(u)\mathfrak s(q^{1/2}u)\bigr]$, so the $q$-Wronskian and
(iii) are equivalent for $u\ne0$. Identity (iii) is the $q$-Pythagorean identity of
\cite{KS1992}. Taking $u=Z$ and $c(Z)=0$ gives the pole invariant, which is the step already used
inside the proof of Proposition~\ref{prop:P12second}.
\end{proof}

\begin{proposition}[the numerators do not vanish at the travel poles]\label{prop:numnonvanish}
At every travel pole $q_m$,
\begin{equation}\label{eq:numprod}
   \Sigma_0(q_m)\,S_0(q_m)\;=\;\frac{2q_m}{1-q_m}\;\ne\;0 .
\end{equation}
In particular $\Sigma_0(q_m)\ne0$ and $S_0(q_m)\ne0$, at \emph{every} travel pole, with no
asymptotic input and no exceptional set. Moreover $|\Sigma_0(q_m)|\,|S_0(q_m)|\ge q_mw_m^2$, so at
least one of the two numerators exceeds $\sqrt{q_m}\,w_m$ in modulus.
\end{proposition}

\begin{proof}
By Lemma~\ref{lem:qsinenum}, $\Sigma_0S_0=\bigl(4q^2/(z_0Z)\bigr)\mathfrak s(Z)\mathfrak s(z_0)$,
and by Lemma~\ref{lem:halfstep} the product $\mathfrak s(Z)\mathfrak s(z_0)$ equals $q^{-3/2}$ at a
travel pole. Since $z_0Z=z_0^2/\sqrt q=2(1-q)/\sqrt q$, the two powers of $q$ combine to
$4q^2\cdot q^{-3/2}\cdot\sqrt q/(2(1-q))=2q/(1-q)$. The size statement follows from
$1-q\le\tau=2/w^2$.
\end{proof}

\begin{remark}\label{rem:numsize}
The two numerators are individually of size $w_m$: $|\Sigma_0(q_m)|/w_m$ and $|S_0(q_m)|/w_m$ rise
monotonically to $1$ over the first $32$ travel poles, from $0.731$ and $0.893$ at the dominant
pole $q_1=0.449453\ldots$ to $0.99992$ and $0.99997$ at the last, with $\Sigma_0$ strictly
alternating in sign
(\texttt{qnumerator\_certificate.py}, mpmath at two precisions, worst deviation in
\eqref{eq:numprod} $4.5\cdot10^{-28}$; the dominant pole is included, it lying below the
extreme-phase ladder and being bracketed separately). This matches the asymptotic
$\Sigma_0\sim w\sin w$, the
amplitude companion of $1-\Sigma_1\sim1-\cos w$; but \eqref{eq:numprod} is an identity and the
asymptotic is not used anywhere.
\end{remark}

\subsection{Spectral consequences at the travel poles}\label{sec:spectral}

Theorem~\ref{thm:fredholm} turns questions about the travel poles into questions about the
eigenvalue branches $\lambda_k(q)$. This subsection records what that gives, and at what standing.
Write $\mathbf B(q):=1-\Sigma_1(q)=\det(I-\mathcal T(q))$.

\begin{lemma}[Hellmann--Feynman]\label{lem:HF}
The map $q\mapsto\mathbf S(q)$ is real-analytic from $(0,1)$ into the trace-class operators, and
each $\lambda_k$ is real-analytic on $(0,1)$. Let $\mathcal TR=\lambda_k(q)^{-1}R$ with $R\ne0$ real
(Theorem~\ref{thm:fredholm}\textup{(ii)}), and put $u=KR$. Then
\[
   \frac{\lambda_k'(q)}{\lambda_k(q)}=-\frac{\mathfrak e_1+\mathfrak e_2}{q}-\frac{1-2q}{q(1-q)},
   \qquad
   \mathfrak e_1=\frac{\sum_ss\,q^{-s}(u_{s+1}-u_s)^2}{\sum_sq^{-s}(u_{s+1}-u_s)^2},\quad
   \mathfrak e_2=\frac{\sum_ss\,q^su_s^2}{\sum_sq^su_s^2} .
\]
In particular $\lambda_k'(q)<0$ if and only if $\mathfrak e_1+\mathfrak e_2>(2q-1)/(1-q)$.
\end{lemma}

\begin{proof}
With the principal branch of $q^{1/2}$, the entries $2q^{1+(s+s')/2+\max(s,s')}$ are holomorphic
on $\{|q|<1,\ \Re q>0\}$ and the sum of their moduli converges locally uniformly there, so
$\mathbf S$ is holomorphic in trace norm, and self-adjoint for real $q$. Its non-zero eigenvalues are
isolated and simple (Theorem~\ref{thm:fredholm}\textup{(iv)}), so by analytic perturbation theory for
a bounded holomorphic self-adjoint family \cite[Ch.~II and Ch.~VII]{Kato} each eigenvalue
$\kappa(q)=1/\lambda_k(q)$ and its eigenvector $\phi(q)$ are real-analytic, and
$\kappa'=\langle\phi,\mathbf S'\phi\rangle/\langle\phi,\phi\rangle$. No $q$-dependent domain
enters: $\mathbf S(q)$ is a bounded family.

It remains to evaluate $\langle\phi,\mathbf S'\phi\rangle$. Put $R=D_e^{1/2}\phi$ and $u=KR$, so
that $D_eu=\mathcal TR=\kappa R$ and $\phi_s^2=\kappa^{-2}2q^{1+s}u_s^2$; put
$N=\sum_sq^su_s^2$ and $E=\sum_sq^{-s}(u_{s+1}-u_s)^2$. Since
$q\partial_q\mathbf S[s,s']=\bigl(1+\tfrac{s+s'}2+\max(s,s')\bigr)\mathbf S[s,s']$,
\[
   q\langle\phi,\mathbf S'\phi\rangle=\kappa\langle\phi,\phi\rangle+\kappa\sum_ss\,\phi_s^2
   +\sum_{s,s'}R_sR_{s'}\max(s,s')\,q^{\max(s,s')},
\]
and the first two terms are $(2qN/\kappa)(1+\mathfrak e_2)$. For the third apply $q\partial_q$ to
$q^{\max(s,s')}=(1-q)\sum_{t\ge\max(s,s')}q^t$: with $\rho_t=\sum_{s\le t}R_s$, bounded by
$\|R\|_{\ell^1}$, it equals $\sum_t\bigl(t(1-q)-q\bigr)q^t\rho_t^2$, all sums converging absolutely.
Since $u_{t+1}-u_t=\sum_{s'}(q^{\max(t+1,s')}-q^{\max(t,s')})R_{s'}=-(1-q)q^t\rho_t$, one has
$q^t\rho_t^2=q^{-t}(u_{t+1}-u_t)^2/(1-q)^2$, so the third term is
$E\bigl((1-q)\mathfrak e_1-q\bigr)/(1-q)^2$. Finally $\langle R,KR\rangle=(1-q)\sum_tq^t\rho_t^2
=E/(1-q)$ and $\langle R,KR\rangle=\sum_sR_su_s=2qN/\kappa$, so $E=2q(1-q)N/\kappa$. Collecting,
$q\langle\phi,\mathbf S'\phi\rangle=(2qN/\kappa)\bigl(1+\mathfrak e_1+\mathfrak e_2-q/(1-q)\bigr)$
and $\langle\phi,\phi\rangle=2qN/\kappa^2$, so
$\kappa'/\kappa=(\mathfrak e_1+\mathfrak e_2)/q+(1-2q)/(q(1-q))$, and
$\lambda_k'/\lambda_k=-\kappa'/\kappa$.
\end{proof}

\noindent At $q_1,\dots,q_6$ the right-hand side agrees to $12$ digits with
$\lambda_m'(q_m)=-\partial_qF/\partial_\lambda F$ computed from the explicit $F$ of
Theorem~\ref{thm:fredholm}. At a travel pole the eigenvalue $1$ has eigenvector the travel null vector
(Corollary~\ref{cor:qbessel}), so the following is a statement about $R$ alone.

\begin{hypothesis}[\textup{(T$^\ast_{\mathrm{poles}}$)}]\label{hyp:Tstar}
At every travel pole $q_m>\tfrac12$, the travel null vector $R$, with $u=KR$ and
$\mathfrak e_1,\mathfrak e_2$ as in Lemma~\ref{lem:HF}, satisfies
\[
   \mathfrak e_1+\mathfrak e_2\ >\ \frac{2q_m-1}{1-q_m}.
\]
\end{hypothesis}

\noindent At a pole $q_m\le\tfrac12$ (only $q_1=0.4494\ldots$) the inequality holds automatically:
the right-hand side is $\le0$ and $\mathfrak e_2>0$, since $u$ is not supported at $s=0$ (if it were,
$R=\kappa^{-1}D_eu$ would be a multiple of $\delta_0$ and $(KR)_1=qR_0\ne0$).

\begin{proposition}[simple poles and the sign of $\mathbf B'$; conditional on
\textup{(T$^\ast_{\mathrm{poles}}$)}]\label{prop:simplepoles}
Assume Hypothesis~\ref{hyp:Tstar}. Then:
\begin{enumerate}[label=\textup{(\roman*)},leftmargin=2.2em,itemsep=0.1em,topsep=0.2em]
\item each $\lambda_k$ takes the value $1$ at exactly one point $c_k\in(0,1)$, and
$c_1<c_2<\cdots$; the $m$-th travel pole is $q_m=c_m$, and $\lambda_j(q_m)<1$ for $j<m$,
$\lambda_m(q_m)=1$, $\lambda_j(q_m)>1$ for $j>m$;
\item $\lambda_m'(q_m)<0$, and $\mathbf B$ has a simple zero at $q_m$:
$\mathbf B'(q_m)=\lambda_m'(q_m)\prod_{j\ne m}\bigl(1-1/\lambda_j(q_m)\bigr)\ne0$. Hence in
Proposition~\ref{prop:shape} $p=\pi=1$: every travel pole is a simple pole of $(I-\mathcal T)^{-1}$;
\item $\operatorname{sign}\mathbf B'(q_m)=(-1)^m$.
\end{enumerate}
\end{proposition}

\begin{proof}
(i) At a point $c$ with $\lambda_k(c)=1$, $c$ is a travel pole with eigenvector $R$
(Corollary~\ref{cor:qbessel}), so Hypothesis~\ref{hyp:Tstar} and the note after it, with
Lemma~\ref{lem:HF}, give $\lambda_k'(c)<0$. A real-analytic function that is strictly decreasing at
every point where it equals $1$ takes the value $1$ at most once: between two consecutive such points
it would be $<1$ just after the first and $>1$ just before the second, and the intermediate value
theorem would give a third point between them. With Proposition~\ref{prop:infpolesalt}, $\lambda_k$
takes the value $1$ exactly once, at $c_k$, and $\lambda_k<1$ exactly on $(c_k,1)$. Since
$\lambda_k(c_{k+1})<\lambda_{k+1}(c_{k+1})=1$, $c_k<c_{k+1}$. The travel poles are the points where
some $\lambda_k$ equals $1$ (Corollary~\ref{cor:qbessel}), so their increasing enumeration is
$q_m=c_m$, and $\lambda_j(q_m)<1$ exactly for $j<m$.

(ii) Near $q_m$ let $P(q)$ be the analytic rank-one eigenprojection of $\kappa_m=1/\lambda_m$
\cite[Ch.~VII]{Kato}. Since $\mathbf S$ commutes with $P$,
$\mathbf B=\det(I-\mathbf S)=(1-\kappa_m)\det\bigl(I-\mathbf S(I-P)\bigr)$, and the second factor
is real-analytic near $q_m$, equal to $\prod_{j\ne m}(1-1/\lambda_j(q_m))\ne0$ at $q_m$. So
$\mathbf B'(q_m)=-\kappa_m'(q_m)\prod_{j\ne m}(1-1/\lambda_j(q_m))$ with
$-\kappa_m'(q_m)=\lambda_m'(q_m)<0$, and $p=1$. For $\pi$, on $\ell^1$
\[
   (I-\mathcal T)^{-1}=I+D_e^{1/2}(I-\mathbf S)^{-1}D_e^{1/2}K,
\]
as follows from $(I-D_eK)D_e^{1/2}=D_e^{1/2}(I-\mathbf S)$ and $D_e^{1/2}KD_eK=\mathbf SD_e^{1/2}K$;
here $D_e^{1/2}K:\ell^1\to\ell^2$ and $D_e^{1/2}:\ell^2\to\ell^1$ are bounded and analytic in $q$.
Near $q_m$, $(I-\mathbf S)^{-1}=P/(1-\kappa_m)$ plus the reduced resolvent, which is holomorphic
there because the other eigenvalues stay away from $1$; and $1-\kappa_m$ has a simple zero. So
$(I-\mathcal T)^{-1}$ has a pole of order at most $1$, and $\pi\ge p\ge1$
(Proposition~\ref{prop:shape}) gives $\pi=p=1$.

(iii) By (i) the product in (ii) has exactly $m-1$ negative factors, and $\lambda_m'(q_m)<0$.
\end{proof}

\noindent The alternation $\operatorname{sign}\mathbf B'(q_m)=(-1)^m$ is checked numerically for
$m\le10$ ($\mathbf B'(q_1)=-2.630$, $\mathbf B'(q_2)=28.71$, \dots, $\mathbf B'(q_{10})=6661$).

\begin{proposition}[the sign of $S_e$, and half-interlacing]\label{prop:signlaws}
Put $\mathbf v=D_e^{1/2}\mathbf 1=(\sqrt{2q^{1+s}})_{s\ge0}\in\ell^2$.
\begin{enumerate}[label=\textup{(\alph*)},leftmargin=2.2em,itemsep=0.1em,topsep=0.2em]
\item Off the travel poles in $(0,1)$,
$\Sigma_0/\mathbf B=\langle\mathbf v,(I-\mathbf S)^{-1}\mathbf v\rangle$. At a travel pole $q_m$, with
$\lambda_k(q_m)=1$ and $\phi$ the normalised eigenvector,
\[
   \Sigma_0(q_m)=w_m\prod_{j\ne k}\bigl(1-1/\lambda_j(q_m)\bigr),\qquad
   w_m=\langle\mathbf v,\phi\rangle^2>0 .
\]
\item Under Hypothesis~\ref{hyp:Tstar}, $\operatorname{sign}\Sigma_0(q_m)=(-1)^{m-1}$; hence
$\Sigma_0$ has at least one zero in each interval $(q_m,q_{m+1})$.
\item Under Hypothesis~\ref{hyp:Tstar}, if moreover $\lambda_{m-1}(q_m)<q_m$, then
$\operatorname{sign}S_e(q_m)=(-1)^{m-1}$.
\end{enumerate}
\end{proposition}

\begin{proof}
(a) By Proposition~\ref{prop:travelsum}, $\Sigma_0/\mathbf B=\mathbf 1^{\!\top}(I-\mathcal
T)^{-1}D_e\mathbf 1$ off the poles, and by the resolvent identity in the proof of
Proposition~\ref{prop:simplepoles},
$(I-\mathcal T)^{-1}D_e=D_e+D_e^{1/2}(I-\mathbf S)^{-1}\mathbf SD_e^{1/2}
=D_e^{1/2}(I-\mathbf S)^{-1}D_e^{1/2}$; and $\mathbf 1^{\!\top}D_e^{1/2}g=\langle\mathbf v,g\rangle$
for $g\in\ell^2$. For the value at a pole fix $q$ and put
$N_q(\lambda)=\det(I-\lambda\mathbf S+\mathbf v\mathbf v^{\!\top})-\det(I-\lambda\mathbf S)$, entire
in $\lambda$. Where $I-\lambda\mathbf S$ is invertible, the rank-one identity
$\det(A+\mathbf v\mathbf v^{\!\top})=\det A\,(1+\langle\mathbf v,A^{-1}\mathbf v\rangle)$ and the
expansion of $\mathbf S$ in its orthonormal eigenbasis $(\phi_j)$ (complete, $\mathbf S$ being
compact, self-adjoint and injective) give
\[
   N_q(\lambda)=\det(I-\lambda\mathbf S)\,\langle\mathbf v,(I-\lambda\mathbf S)^{-1}\mathbf v\rangle
   =\sum_jw_j\prod_{i\ne j}\bigl(1-\lambda/\lambda_i(q)\bigr),\qquad
   w_j=\langle\mathbf v,\phi_j\rangle^2 ,
\]
and since $|\prod_{i\ne j}(1-\lambda/\lambda_i)|\le e^{|\lambda|\operatorname{tr}\mathbf S}$ and
$\sum_jw_j=\|\mathbf v\|^2$, the series converges locally uniformly and the identity holds for all
$\lambda$. At $\lambda=1$ and $q$ off the poles, $N_q(1)=\Sigma_0(q)$; both sides are continuous in
$q$, $N_q(1)$ as a difference of Fredholm determinants of trace-norm continuous families, so
$N_q(1)=\Sigma_0(q)$ on all of $(0,1)$. At $q=q_m$ every term $j\ne k$ of the series contains the
factor $1-1/\lambda_k(q_m)=0$. This gives the formula, and $w_m>0$ because $\Sigma_0(q_m)\ne0$
(Proposition~\ref{prop:numnonvanish}).

(b) By Proposition~\ref{prop:simplepoles}\textup{(i)}, $k=m$ and exactly $m-1$ factors are negative.
So $\Sigma_0$ alternates in sign along the poles, and, being continuous, vanishes between
consecutive ones.

(c) By Theorem~\ref{thm:fredholm}\textup{(iv)},\textup{(v)} at $\lambda=q$,
$S_e(q_m)=\prod_j\bigl(1-q_m/\lambda_j(q_m)\bigr)$. For $j>m$, $\lambda_j>1>q_m$; for $j=m$ the
factor is $1-q_m>0$; for $j<m$, $\lambda_j\le\lambda_{m-1}<q_m$ and the factor is negative.
\end{proof}

\noindent The extra condition in (c) would follow from a shift inequality
$\lambda_n(q)\le q^2\lambda_{n+1}(q)$, which gives $\lambda_{m-1}(q_m)\le q_m^2<q_m$; that
inequality is \emph{not} proved in this paper, so (c) carries it as a hypothesis. Numerically, at
$m=2,\dots,6$, $\lambda_{m-1}(q_m)=0.112,\ 0.360,\ 0.510,\ 0.605,\ 0.670$, well below $q_m^2$, and
$S_e(q_m)=0.520,\ -0.210,\ 0.127,\ -0.091,\ 0.071,\ -0.058$ for $m=1,\dots,6$ ($50$ digits). The
sign alternation of $\Sigma_0$ is the one recorded for the first $32$ poles in
Remark~\ref{rem:numsize}.

\begin{remark}[standing of Propositions~\ref{prop:simplepoles} and~\ref{prop:signlaws}]\label{rem:spectralstanding}
Earlier versions of this paper stated neither Hypothesis~\ref{hyp:Tstar} nor the three inputs that
our working notes carried beside it for the same purpose: \textup{(I1)} a Sturm oscillation theorem
for the Friedrichs realisation of the eigenvalue problem; \textup{(I2)} Hellmann--Feynman on a
$q$-dependent form domain; \textup{(I3)} Proposition~\ref{prop:qtrig} with
Lemma~\ref{lem:infpoles}. Here \textup{(I2)} is Lemma~\ref{lem:HF}, proved because $\mathbf S(q)$ is
a bounded norm-analytic family, so Kato's theory applies with no domain question;
\textup{(I3)} is covered by Proposition~\ref{prop:qtrig}, Theorem~\ref{thm:fredholm} and
Proposition~\ref{prop:infpolesalt}, and Lemma~\ref{lem:infpoles} is not used; \textup{(I1)} is not used at all: it would turn the index
statement of Proposition~\ref{prop:simplepoles}\textup{(i)} into the count of $m-1$ sign changes of
$R_s$, which nothing here needs. What remains is Hypothesis~\ref{hyp:Tstar} itself, which is
\emph{not} proved here. Numerically, the ratio of its two sides exceeds $2$ at every pole
$2\le m\le50$ and decreases towards $2$ ($2.145$ at $m=2$, $2.014$ at $m=5$). A proof with the
constant $1.25$ in front of the right-hand side, from a flux identity expressing
$\mathfrak e_1+\mathfrak e_2$ as twice a mean position, an amplitude bound, and an interval
certificate on $q\in[\tfrac12,1-10^{-6}]$, has been written and independently reviewed, but it is
not part of this paper. So Proposition~\ref{prop:simplepoles} and
Proposition~\ref{prop:signlaws}\textup{(b)} are proved given Hypothesis~\ref{hyp:Tstar},
Proposition~\ref{prop:signlaws}\textup{(c)} given also the shift inequality, and
Proposition~\ref{prop:signlaws}\textup{(a)} unconditionally. No theorem of this paper depends on
them: Proposition~\ref{prop:shape} needs only $\pi\ge p\ge1$, and Theorem~\ref{thm:RJ} gives the
non-vanishing of the residue directly.
\end{remark}

\begin{remark}[what the known theory does not give]\label{rem:knowntheory}
Two statements about the travel poles lie outside what Theorem~\ref{thm:fredholm} and the
literature on Hahn--Exton zeros deliver.

\emph{The accumulation law.} The zero theory of \cite{KoelinkSwarttouw1994,StampachStovicek2014}
concerns the zeros in $x$ at fixed $q$; the travel poles sit on the diagonal $x=Z(q)$. With
Hurwitz's theorem it gives each branch at fixed $k$,
$\lambda_k(q)\sim\varepsilon(2k-1)^2\pi^2/8$ with $\varepsilon=1-q$
(Proposition~\ref{prop:infpolesalt}\textup{(b)}). The $m$-th pole is where $\lambda_m=1$, that is
in the regime $k=m\asymp\varepsilon^{-1/2}$, and no fixed-$k$ statement reaches it. So the law
$\varepsilon_m=1-q_m\sim8/(\pi^2(2m-1)^2)$, equivalently $w_m\sim(m-\tfrac12)\pi$, recorded
numerically in Remark~\ref{rem:family}, is not a consequence. It needs a count lemma: an asymptotic
for $\#\{k:\lambda_k(q)<1\}$, the number of zeros of $\cos(\cdot;q^2)$ in $(0,Z(q))$, uniform as
$q\to1$. Lemma~\ref{lem:T2abs} puts at least one travel pole in each window of $w$ (the proof of
Lemma~\ref{lem:infpoles}), and Proposition~\ref{prop:simplepoles}\textup{(i)} identifies the index
under Hypothesis~\ref{hyp:Tstar}, but neither bounds the number of poles in a window. The count
lemma is open.

\emph{Complex $q$.} The identity $\det(I-\mathcal T)=\cos(Z;q^2)$ continues analytically to complex
$q$, but positivity and self-adjointness do not, and the zero theory cited is for $0<q<1$. We know of
no result on zeros of $q$-Bessel functions at complex base that would locate the non-real zeros of
Remark~\ref{rem:polesdensity}; their density on the circle remains a conjecture.
\end{remark}

\section{The transfer model and the block assembly}\label{sec:assembly}

The travel block and the bulk block are one object seen twice. Both are the resolvent of the
semiseparable kernel $q^{\max(a,b)}$ against a geometric diagonal weight, and a single three-term
recursion in the section index $k$ covers both: its coefficients are $\alpha_k,\gamma_k$ in the
bulk case and $A_k,C_k$ in the travel case, which is why the four series of \eqref{eq:krec} are
four ladders on one recursion and not four separate constructions. This section develops that into
the block assembly the rest of the paper consumes.

Sections~\ref{sec:sect}--\ref{sec:mob} are unconditional statements about explicitly given
operators, and \S\ref{sec:sitecost} is an unconditional statement about the crossing optimisation
that those operators encode. Section~\ref{sec:modelassembly} states the strand-walk input
\textup{(M)} that identifies the one with the other and proves it (Theorem~\ref{thm:model}, with
the assembly step proved in Appendix~\ref{app:M3prime}). Section~\ref{sec:RJ} evaluates the junction
pairing at the travel poles in closed form.

Throughout, $|q|<1$ and $x=q^{1/2}$. Besides the four blocks of \eqref{eq:krec} we use the third
ladder on the same recursion,
\begin{equation}\label{eq:Tladder}
   T_k=\sum_{j\ge0}\alpha_{k+2j+1}\prod_{i<j}\gamma_{k+2i}\qquad(k\ge0).
\end{equation}
Since $\gamma_k=-2q^{k+1}(1-q)/[(1-q^{k+1})(1-q^{k+2})]$ and $C_k=q\gamma_k$, on every compact
subset of $|q|<1$ there is a $\Xi$ with $|\gamma_k|,|C_k|\le\Xi|q|^{k+1}$, so
$\bigl|\prod_{i<j}\gamma_{k+2i}\bigr|\le\Xi^{j}|q|^{j^2+jk}$ and likewise for $C$; as
$\alpha_k$ and $A_k$ are bounded there, $S_k,\Sigma_k,T_k$ converge absolutely and locally
uniformly and are holomorphic on $|q|<1$ for every $k\ge0$.

\subsection{The section identity}\label{sec:sect}

\begin{definition}[section identity]\label{def:model}
Fix $s_0\in\{0,1\}$ and $\vartheta\in\{0,1\}$ with $s_0+\vartheta=1$, and put
$e_s=2q^{\vartheta+s}$ for $s\ge s_0$. Given a bounded \emph{source} $\mathbf c=(c_s)_{s\ge s_0}$,
a family $\Phi=(\Phi_s)_{s\ge s_0}$ satisfies the \emph{section identity} with data
$(s_0,\vartheta,\mathbf c)$ if
\begin{equation}\label{eq:sectionid}
   \Phi_s=e_s\Bigl(c_s+\sum_{s'\ge s_0}q^{\max(s,s')}\Phi_{s'}\Bigr),\qquad s\ge s_0 .
\end{equation}
Its \emph{sections} are $G_k=\sum_{s\ge s_0}q^{ks}\Phi_s$, $k\ge0$. In operator form
$(I-M)\Phi=E$ with
\begin{equation}\label{eq:transfer}
   M[s,s']=e_s\,q^{\max(s,s')},\qquad E_s=e_sc_s .
\end{equation}
The \emph{bulk} instance is $(s_0,\vartheta)=(1,0)$ and the \emph{travel} instance is
$(s_0,\vartheta)=(0,1)$; we write $M_0$ for the bulk operator and $\mathcal T$ for the travel one.
\end{definition}

\begin{lemma}[existence, uniqueness, analyticity]\label{lem:modelexist}
For $|q|<\tfrac13$ and any bounded $\mathbf c$, equation \eqref{eq:sectionid} has a unique solution
in $\ell^1$, namely $\Phi=\sum_{n\ge0}M^nE$, the series converging in $\ell^1$ and coefficientwise
in $q$, one extra power of $q$ per term. One has $|\Phi_s|\le C|q|^{\vartheta+s}$ with $C$
independent of $s$, and each section $G_k$ is holomorphic at $q=0$.
\end{lemma}

\begin{proof}
$\|M\|_{\ell^1\to\ell^1}=\sup_{s'}\sum_s|M[s,s']|\le\sup_{s'}\sum_{s\ge s_0}2|q|^{\vartheta+s}
=2|q|/(1-|q|)<1$ for $|q|<\tfrac13$, using $\vartheta+s_0=1$ and $|q|^{\max(s,s')}\le1$. Hence
$I-M$ is invertible on $\ell^1$ and the Neumann series converges. Every entry of $M$ is divisible
by $q^{\vartheta+s}q^{\max(s,s')}$, hence by $q$, so $M^nE$ is divisible by $q^{n+1}$ and the
series converges coefficientwise. Reading the $s$-th row of $\Phi=E+M\Phi$ gives
$|\Phi_s|\le2|q|^{\vartheta+s}(\|\mathbf c\|_\infty+\|\Phi\|_{\ell^1})$. Since
$|\Phi_s|=O(|q|^{s})$, each $G_k$ is a locally uniformly convergent sum of holomorphic functions.
\end{proof}

\subsection{One telescoping lemma for both blocks}\label{sec:tel}

\begin{lemma}[catalytic telescoping]\label{lem:tel}
Let $0<|q|<1$ and let $\Phi\in\ell^1$ satisfy \eqref{eq:sectionid} with data
$(s_0,\vartheta,\mathbf c)$. Put
\[
  a_k=\frac{2q^{\vartheta+(k+1)s_0}}{1-q^{k+1}},\qquad
  b_k=\frac{2q^{\vartheta+k+2}}{1-q^{k+2}}-\frac{2q^{\vartheta+k+1}}{1-q^{k+1}},\qquad
  \mathcal C_k=\sum_{s\ge s_0}e_sc_sq^{ks},
\]
and $\mathfrak a_k=\sum_{j\ge0}a_{k+2j}\prod_{i<j}b_{k+2i}$,
$\mathfrak c_k=\sum_{j\ge0}\mathcal C_{k+2j}\prod_{i<j}b_{k+2i}$. Then for every $k\ge0$
\begin{equation}\label{eq:threeterm}
   G_k=\mathcal C_k+a_k\,G_1+b_k\,G_{k+2},
\end{equation}
the series $\mathfrak a_k,\mathfrak c_k$ converge, $G_k=\mathfrak c_k+\mathfrak a_kG_1$, and if
$\mathfrak a_1\ne1$ then
\begin{equation}\label{eq:solved}
   G_1=\frac{\mathfrak c_1}{1-\mathfrak a_1},\qquad
   G_k=\mathfrak c_k+\mathfrak a_k\,\frac{\mathfrak c_1}{1-\mathfrak a_1}\quad(k\ge0).
\end{equation}
In the bulk instance $a_k=\alpha_k$, $b_k=\gamma_k$ and $\mathfrak a_k=S_k$; in the travel
instance $a_k=A_k$, $b_k=C_k$ and $\mathfrak a_k=\Sigma_k$.
\end{lemma}

\begin{proof}
Split the inner sum of \eqref{eq:sectionid} at $s'=s$, where $\max(s,s')$ switches:
\[
  \Phi_s=e_sc_s+e_s\!\!\sum_{s'\ge s}\!q^{s'}\Phi_{s'}+e_sq^{s}\!\!\sum_{s_0\le s'<s}\!\Phi_{s'} .
\]
Multiply by $q^{ks}$ and sum over $s\ge s_0$. All three double sums converge absolutely, since
$\Phi\in\ell^1$ and $|q|<1$, so the orders of summation may be exchanged. The first gives
$\mathcal C_k$. In the second, with $e_s=2q^{\vartheta+s}$,
\[
  \sum_{s\ge s_0}2q^{\vartheta+(k+1)s}\sum_{s'\ge s}q^{s'}\Phi_{s'}
  =2q^{\vartheta}\sum_{s'\ge s_0}q^{s'}\Phi_{s'}\sum_{s=s_0}^{s'}q^{(k+1)s}
  =\frac{2q^{\vartheta}}{1-q^{k+1}}\sum_{s'\ge s_0}q^{s'}\Phi_{s'}
      \bigl(q^{(k+1)s_0}-q^{(k+1)(s'+1)}\bigr),
\]
which is $a_kG_1-\frac{2q^{\vartheta+k+1}}{1-q^{k+1}}G_{k+2}$, because
$\sum_{s'}q^{s'}\Phi_{s'}=G_1$ and $\sum_{s'}q^{(k+2)s'}\Phi_{s'}=G_{k+2}$. In the third,
\[
  \sum_{s\ge s_0}2q^{\vartheta+(k+2)s}\sum_{s_0\le s'<s}\Phi_{s'}
  =2q^{\vartheta}\sum_{s'\ge s_0}\Phi_{s'}\sum_{s>s'}q^{(k+2)s}
  =\frac{2q^{\vartheta+k+2}}{1-q^{k+2}}\,G_{k+2}.
\]
Adding the three gives \eqref{eq:threeterm}. Iterating it $J$ times,
\[
  G_k=\sum_{j<J}\bigl(\mathcal C_{k+2j}+a_{k+2j}G_1\bigr)\prod_{i<j}b_{k+2i}
      +G_{k+2J}\prod_{i<J}b_{k+2i}.
\]
For fixed $q$ one has $|b_m|\le\Xi|q|^{m+1}$ with $\Xi=4/(1-|q|)$, so
$\bigl|\prod_{i<J}b_{k+2i}\bigr|\le\Xi^{J}|q|^{J^2+Jk}\to0$; and $|G_m|\le\|\Phi\|_{\ell^1}$
uniformly in $m$, so the remainder vanishes as $J\to\infty$. The same bound, against
$\sup_k|a_k|<\infty$ and $\sup_k|\mathcal C_k|\le2\|\mathbf c\|_\infty/(1-|q|)$, makes
$\mathfrak a_k$ and $\mathfrak c_k$ converge. This gives $G_k=\mathfrak c_k+\mathfrak a_kG_1$;
taking $k=1$ gives $G_1(1-\mathfrak a_1)=\mathfrak c_1$, whence \eqref{eq:solved}. The
identifications of $a_k,b_k$ are read off the displays: $(s_0,\vartheta)=(1,0)$ gives
$a_k=2q^{k+1}/(1-q^{k+1})=\alpha_k$ and $b_k=\alpha_{k+1}-\alpha_k=\gamma_k$, while
$(s_0,\vartheta)=(0,1)$ gives $a_k=2q/(1-q^{k+1})=A_k$ and
$b_k=2q^{k+3}/(1-q^{k+2})-2q^{k+2}/(1-q^{k+1})=C_k$.
\end{proof}

\begin{remark}[where each feature of the model is used]\label{rem:telhyp}
Three features enter \eqref{eq:threeterm}, each exactly once. That the coupling is
$q^{\max(s,s')}$ is what makes the split at $s'=s$ produce two \emph{geometric} sums; a symmetric
kernel not of the form $\mu^{\max}$ destroys the three-term relation. That $e_s$ is geometric in
$s$ is what makes those sums summable in closed form; a non-geometric edge weight leaves $G_k$
coupled to a continuum of sections. That the source is geometric is used only to evaluate
$\mathcal C_k$; both \eqref{eq:threeterm} and \eqref{eq:solved} hold for any bounded source. The
step $k\mapsto k+2$, and with it the two-ladder structure of $S_k$ and $\Sigma_k$, comes from the
two occurrences of $q^{(k+2)s'}$ above.
\end{remark}

\begin{proposition}[the travel $k$-sum telescopes]\label{prop:travelsum}
Let $\Phi$ solve the travel section identity, $(s_0,\vartheta,\mathbf c)=(0,1,\mathbf 1)$. Then, as
meromorphic functions on $|q|<1$,
\[
  G^{T}_k=\frac{\Sigma_k}{1-\Sigma_1}\quad(k\ge0),\qquad\text{in particular}\quad
  G^{T}_0=\sum_{s\ge0}\Phi_s=\frac{\Sigma_0}{1-\Sigma_1} .
\]
\end{proposition}

\begin{proof}
Lemma~\ref{lem:tel} with $(s_0,\vartheta)=(0,1)$ and $c_s\equiv1$ gives
$\mathcal C_k=\sum_{s\ge0}2q^{1+s}q^{ks}=2q/(1-q^{k+1})=A_k=a_k$, so
$\mathfrak c_k=\mathfrak a_k=\Sigma_k$ and \eqref{eq:solved} reads
$G^T_k=\Sigma_k+\Sigma_k\Sigma_1/(1-\Sigma_1)=\Sigma_k/(1-\Sigma_1)$. Both sides are holomorphic
near $q=0$ and the right-hand side is meromorphic on $|q|<1$, so the identity extends by analytic
continuation. Combinatorially this is the Neumann series of Lemma~\ref{lem:modelexist}:
$\sum_s\Phi_s=\sum_{n\ge0}\mathbf 1^{\!\top}\mathcal T^nE$, whose $n$-th term is the generating
function of the travel runs with exactly $n+1$ edges.
\end{proof}

\begin{proposition}[the bulk dressing]\label{prop:bulkdress}
Let $\Phi$ solve the bulk section identity with $\mathbf c=\mathbf 1$ and let $\Psi$ solve it with
$c_s=q^s$. Set
\[
  b_0=\sum_{s\ge1}\Phi_s,\qquad t_0=\sum_{s\ge1}q^s\Phi_s,\qquad
  b_1=\sum_{s\ge1}\Psi_s,\qquad t_1=\sum_{s\ge1}q^s\Psi_s .
\]
Then, as meromorphic functions on $|q|<1$,
\[
  b_0=\frac{S_0}{1-S_1},\qquad t_0=\frac{S_1}{1-S_1},\qquad
  t_1=\frac{T_1}{1-S_1},\qquad b_1=T_0+\frac{S_0T_1}{1-S_1}.
\]
\end{proposition}

\begin{proof}
Lemma~\ref{lem:tel} with $(s_0,\vartheta)=(1,0)$. For $\mathbf c=\mathbf 1$,
$\mathcal C_k=\sum_{s\ge1}2q^{s(k+1)}=\alpha_k=a_k$, so $\mathfrak c_k=\mathfrak a_k=S_k$ and
$G_k=S_k/(1-S_1)$; take $k=0,1$. For $c_s=q^s$,
$\mathcal C_k=\sum_{s\ge1}2q^{s(k+2)}=\alpha_{k+1}$, so
$\mathfrak c_k=\sum_j\alpha_{k+2j+1}\prod_{i<j}\gamma_{k+2i}=T_k$ by \eqref{eq:Tladder}, and
\eqref{eq:solved} gives $t_1=G_1=T_1/(1-S_1)$ and $b_1=G_0=T_0+S_0T_1/(1-S_1)$. Analytic
continuation as before.
\end{proof}

\noindent With $1-S_1=S_e$ (Proposition~\ref{prop:qtrig}) the first of these is $b_0=S_0/S_e$, the
third line of the dictionary \eqref{eq:dict}; and $b_0$ is holomorphic and non-zero at every travel
pole, by Propositions~\ref{prop:selower} and \ref{prop:numnonvanish} respectively.

\subsection{The cocycle, the reciprocity, and the dictionary}\label{sec:cocycleid}

The bulk section identity has a second, equivalent, description as a two-point problem for a
first-order system. That system is the cocycle of \S\ref{sec:gateblock}, and comparing the two
descriptions identifies $t_1$ with $P_{12}/S_e$.

\begin{corollary}[the cocycle is the same object]\label{cor:cocycle}
Let $\Phi$ solve the bulk section identity with source $\mathbf c$, and put
$L_b=\sum_{a\le b}\Phi_a$ and $\mathcal E_b=-\sum_{a\ge b}q^a\Phi_a$. Then
$\Phi_b=2q^bc_b+2q^{2b}L_{b-1}-2q^b\mathcal E_b$ and
\begin{equation}\label{eq:cocyclerec}
  \begin{pmatrix}L_b\\ \mathcal E_{b+1}\end{pmatrix}
  =\mathcal M_b\begin{pmatrix}L_{b-1}\\ \mathcal E_{b}\end{pmatrix}
   +\begin{pmatrix}2q^{b}c_b\\ 2q^{2b}c_b\end{pmatrix},\qquad
  \mathcal M_b=\begin{pmatrix}1+2q^{2b}&-2q^{b}\\ 2q^{3b}&1-2q^{2b}\end{pmatrix},
\end{equation}
with $L_0=0$, $\mathcal E_1=-t_0$ or $-t_1$ according to the source, and $\mathcal E_b\to0$. Let
$(\ell,r)$ and $(\ell',r')$ be the solutions of the homogeneous recursion
\textup{(}$\mathbf c=0$\textup{)} with data $(\ell_0,r_1)=(1,0)$ and $(\ell'_0,r'_1)=(0,1)$, and
put $(P_{11},P_{12})=\lim_b(\ell_b,r_{b+1})$ and $(P_{21},P_{22})=\lim_b(\ell'_b,r'_{b+1})$, the
limits existing because $\sum_b\|\mathcal M_b-I\|<\infty$. Eliminating $L$ between the two lines
of \eqref{eq:cocyclerec} at $\mathbf c=0$ gives
\[
   \mathcal E_{b+2}=\bigl(1+q^3-2(1-q)q^{2b+2}\bigr)\mathcal E_{b+1}-q^3\mathcal E_b ,
\]
which is the scalar three-term recursion of \S\ref{sec:gateblock}; so $r$ and $r'$ are its two
standard solutions and $P_{12},P_{22}$ their limits, $P_{12}$ being the entry whose closed form is
\eqref{eq:p12closed}. Then, identically in $0<q<1$,
\[
   P_{11}=1+T_0,\qquad P_{12}=T_1,\qquad P_{21}=-S_0,\qquad P_{22}=S_e=1-S_1 .
\]
In particular $t_1=P_{12}/S_e$ is an identity between independently defined objects, not a
definition.
\end{corollary}

\begin{proof}
Substituting $L_b=L_{b-1}+\Phi_b$ and $\mathcal E_{b+1}=\mathcal E_b+q^b\Phi_b$ into
$\Phi_b=2q^b(c_b+q^bL_b-\mathcal E_{b+1})$, which is \eqref{eq:sectionid} after the split at
$s'=b$, the two terms $\pm2q^{2b}\Phi_b$ cancel and leave
$\Phi_b=2q^bc_b+2q^{2b}L_{b-1}-2q^b\mathcal E_b$; feeding this back gives \eqref{eq:cocyclerec}.
The homogeneous solutions form a two-dimensional space, since \eqref{eq:cocyclerec} determines
$(L_b,\mathcal E_{b+1})$ from $(L_{b-1},\mathcal E_b)$.

Take $\mathbf c=\mathbf 1$, so $L_0=0$, $L_b\to b_0$, $\mathcal E_1=-t_0$ and
$\mathcal E_{b+1}\to0$. Then $(L,\mathcal E-1)$ solves the homogeneous recursion: in the first line
the source $2q^bc_b=2q^b$ is cancelled by $-2q^b\cdot(-1)$, and in the second the source
$2q^{2b}c_b=2q^{2b}$ is cancelled by $(1-2q^{2b})(-1)-(-1)$. Its data are $(0,-1-t_0)$, so
$(L,\mathcal E-1)=-(1+t_0)(\ell',r')$, and letting $b\to\infty$,
\[
   b_0=-(1+t_0)P_{21},\qquad -1=-(1+t_0)P_{22} .
\]
By Proposition~\ref{prop:bulkdress} and $S_e=1-S_1$ one has $1+t_0=1+S_1/S_e=1/S_e$, whence
$P_{22}=S_e$ and $P_{21}=-b_0S_e=-S_0$.

Take now $c_s=q^s$, with data $L'_0=0$, $L'_b\to b_1$, $\mathcal E'_1=-t_1$,
$\mathcal E'_{b+1}\to0$. Here $(L'+1,\mathcal E')$ solves the homogeneous recursion: in the first
line $(1+2q^{2b})\cdot1$ exceeds $1$ by $2q^{2b}$, which is the source $2q^bc_b=2q^{2b}$, and in
the second $2q^{3b}\cdot1$ is the source $2q^{2b}c_b=2q^{3b}$. Its data are $(1,-t_1)$, so
$(L'+1,\mathcal E')=(\ell,r)-t_1(\ell',r')$, and letting $b\to\infty$,
\[
   b_1+1=P_{11}-t_1P_{21}=P_{11}+t_1S_0,\qquad 0=P_{12}-t_1P_{22}=P_{12}-t_1S_e .
\]
The second gives $P_{12}=t_1S_e=T_1$ by Proposition~\ref{prop:bulkdress}, and the first then gives
$P_{11}=b_1+1-t_1S_0=T_0+S_0T_1/S_e+1-(T_1/S_e)S_0=1+T_0$. All four limits are taken for
$|q|<\tfrac13$, where Lemma~\ref{lem:modelexist} applies, and extend to $0<q<1$ by analytic
continuation, both sides being meromorphic there.
\end{proof}

\begin{proposition}[reciprocity]\label{prop:recip}
$t_0=b_1$ identically in $q$. Equivalently $T_0S_e+T_1S_0=S_1$, equivalently
$P_{11}P_{22}-P_{12}P_{21}=1$.
\end{proposition}

\begin{proof}
Two independent proofs.

\emph{Kernel symmetry.} Write $D=\operatorname{diag}(2q^b)_{b\ge1}$ and $K[b,a]=q^{\max(a,b)}$, so
that $M_0=DK$, $E=D\mathbf 1$, and $u:=(2q^{2b})_b=Dv$ with $v=(q^a)_a$. The two sources of
Proposition~\ref{prop:bulkdress} are $E$ and $u$, so $\Phi=(I-M_0)^{-1}E$, $\Psi=(I-M_0)^{-1}u$,
and
\[
   t_0=v^{\!\top}(I-M_0)^{-1}E,\qquad b_1=\mathbf 1^{\!\top}(I-M_0)^{-1}u .
\]
For $|q|<\tfrac13$, $X:=(I-M_0)^{-1}D=\sum_{n\ge0}(DK)^nD$, and the $n$-th summand is
$DKD\cdots KD$, a palindromic product of the symmetric matrices $D$ and $K$, hence symmetric; the
series converges entrywise. Therefore
$t_0=v^{\!\top}X\mathbf 1=\sum_{a,b}v_aX[a,b]=\sum_{a,b}X[b,a]v_a=\mathbf 1^{\!\top}Xv=b_1$, every
sum being absolutely convergent. Both sides are meromorphic on $|q|<1$ by
Proposition~\ref{prop:bulkdress}, so the identity extends.

\emph{Casoratian.} $\det\mathcal M_b=(1+2q^{2b})(1-2q^{2b})+4q^{3b}q^b=1$, so
$\ell_br'_{b+1}-\ell'_br_{b+1}$ does not depend on $b$; at $b=0$ it is $1$, whence
$P_{11}P_{22}-P_{21}P_{12}=1$. Substituting the four entries of Corollary~\ref{cor:cocycle} gives
$(1+T_0)S_e+T_1S_0=1$, that is $T_0S_e+T_1S_0=1-S_e=S_1$; dividing by $S_e$ and using
Proposition~\ref{prop:bulkdress} turns this into $b_1=T_0+S_0T_1/S_e=S_1/S_e=t_0$.
\end{proof}

\noindent The symmetry of $K$ is used and is necessary: replacing $q^{\max(a,b)}$ by the
asymmetric $q^{\max(a,b)}+q^{a+2b}$ makes $t_0-b_1=4q^{8}+O(q^{9})$
(\texttt{assembly\_certificate.py}, Part~2). The determinant identity
$P_{11}P_{22}-P_{12}P_{21}=1$ is the Casoratian of the cocycle and is automatic; what carries
content is the identification of its four entries in Corollary~\ref{cor:cocycle}, which turns it
into the reciprocity.

\begin{corollary}[the dictionary]\label{cor:dict}
Identically in $0<q<1$,
\[
  S_e=1-S_1,\qquad S_o=\tfrac{1-q}{2q}S_0,\qquad b_0=\frac{S_0}{S_e},\qquad
  t_0=b_1=\frac{S_1}{S_e},\qquad t_1=\frac{T_1}{S_e}=\frac{P_{12}}{S_e},
\]
and consequently $\kappa:=t_0b_1-b_0t_1=(S_1^2-S_0P_{12})/S_e^2$. This is \eqref{eq:dict} together
with its two remaining entries.
\end{corollary}

\begin{proof}
The first line is Proposition~\ref{prop:qtrig} and Lemma~\ref{lem:qsinenum}; the rest is
Proposition~\ref{prop:bulkdress} with Proposition~\ref{prop:recip} and
Corollary~\ref{cor:cocycle}, and $\kappa$ is substitution.
\end{proof}

\subsection{Compactness, the null space, and the singularities}\label{sec:fred}

\begin{lemma}[compactness and meromorphy]\label{lem:fred}
For $|q|<1$ the operator $M$ of \eqref{eq:transfer}, and its gap-marked version with kernel
\eqref{eq:gapkernel} below, are compact on $\ell^1$. Hence $q\mapsto(I-M)^{-1}$ is meromorphic on
$|q|<1$, with poles of finite order, and at each pole $\ker(I-M)\ne0$.
\end{lemma}

\begin{proof}
Let $M_N$ agree with $M$ on rows $s\le N$ and vanish above. Then
$\|M-M_N\|_{\ell^1\to\ell^1}\le\sup_{s'}\sum_{s>N}|e_s||q|^{\max(s,s')}
\le\sum_{s>N}2|q|^{\vartheta+s}\to0$, so $M$ is a norm limit of finite-rank operators and is
compact; the rank-one gap term is compact a fortiori. The family is holomorphic in $q$ on the
connected domain $|q|<1$ and $I-M$ is invertible near $q=0$ (Lemma~\ref{lem:modelexist}), so the
analytic Fredholm theorem applies.
\end{proof}

\begin{lemma}[the null space is at most one-dimensional]\label{lem:nullspace}
Let $0<|q|<1$ and let $R\in\ell^1$ satisfy $R_s=e_s\sum_{s'}q^{\max(s,s')}R_{s'}$ for all
$s\ge s_0$. Then its sections satisfy $G^R_k=G^R_1\mathfrak a_k$ for all $k\ge0$, and $G^R_1=0$
forces $R=0$. In particular $\dim\ker(I-M)\le1$, and a non-zero kernel element exists only where
$\mathfrak a_1=1$.
\end{lemma}

\begin{proof}
Lemma~\ref{lem:tel} with $\mathbf c=0$ gives $\mathcal C_k=0$, hence $\mathfrak c_k=0$ and
$G^R_k=\mathfrak a_kG^R_1$; its derivation used only $R\in\ell^1$ and $|q|<1$. Taking $k=1$ gives
$G^R_1=\mathfrak a_1G^R_1$, so $\mathfrak a_1=1$ unless $G^R_1=0$. If $G^R_1=0$ then $G^R_k=0$ for
every $k$. The function $F(t)=\sum_{s\ge s_0}R_st^s$ is holomorphic on $|t|<1$ because
$R\in\ell^1$, and $F(q^k)=G^R_k=0$ for all $k\ge0$; the points $q^k$ accumulate at the interior
point $t=0$, so $F\equiv0$ and $R=0$. Hence the kernel is spanned by the class normalised by
$G^R_1=1$.
\end{proof}

\begin{corollary}[singularities of the two gapless resolvents]\label{cor:sing}
On $|q|<1$, the poles of $(I-M_0)^{-1}$ lie in the zero set of $S_e$ and the poles of
$(I-\mathcal T)^{-1}$ lie in the zero set of $1-\Sigma_1$; conversely every travel pole is a pole
of $(I-\mathcal T)^{-1}$.
In particular $b_0,b_1,t_0,t_1$ are holomorphic at every travel pole $q_m$, where $S_e(q_m)\ne0$
by Proposition~\ref{prop:selower}, and $b_0(q_m)=S_0(q_m)/S_e(q_m)\ne0$ by
Proposition~\ref{prop:numnonvanish}.
\end{corollary}

\begin{proof}
By Lemma~\ref{lem:fred} a pole forces a non-zero kernel element, and by Lemma~\ref{lem:nullspace}
that forces $\mathfrak a_1=1$, which is $S_1=1$, i.e.\ $S_e=0$, in the bulk instance and
$\Sigma_1=1$ in the travel one. Conversely, at a travel pole $\Sigma_0/(1-\Sigma_1)$ is unbounded
because $\Sigma_0(q_m)\ne0$ (Proposition~\ref{prop:numnonvanish}), so by
Proposition~\ref{prop:travelsum} the travel resolvent is singular there.
\end{proof}

\subsection{The gap term is rank one: the M\"obius factorisation}\label{sec:mob}

Marking the gap bridges changes the bulk kernel by one term. With
$\mathfrak g=\mathfrak g(q,y)=\sum_{L\ge1}q^Ly^{L-1}=q/(1-qy)$, the gap-marked bulk kernel and
system are
\begin{equation}\label{eq:gapkernel}
   K_{\mathfrak g}(a,b)=q^{\max(a,b)}+q^{a+b}\,\mathfrak g,\qquad
   P_b=2q^b\Bigl(q^b+\sum_{a\ge1}K_{\mathfrak g}(a,b)P_a\Bigr),\qquad B(q,y)=\sum_{b\ge1}P_b ,
\end{equation}
the source $q^b$ being the cost of the outer span-end site of a bulk run (Lemma~\ref{M3:lem:bulk}). In operator form
$(I-M_0-\mathfrak g\,uv^{\!\top})P=u$ with
\begin{equation}\label{eq:rankone}
   M_0[b,a]=2q^bq^{\max(a,b)},\qquad E_b=2q^b,\qquad u_b=2q^{2b},\qquad v_a=q^a ,
\end{equation}
since $2q^b\cdot q^{a+b}=(2q^{2b})(q^a)$: \emph{the gap term is exactly rank one.} With the four
scalars of Proposition~\ref{prop:bulkdress},
\[
  b_0=\mathbf 1^{\!\top}(I-M_0)^{-1}E,\quad b_1=\mathbf 1^{\!\top}(I-M_0)^{-1}u,\quad
  t_0=v^{\!\top}(I-M_0)^{-1}E,\quad t_1=v^{\!\top}(I-M_0)^{-1}u .
\]

\begin{proposition}[M\"obius factorisation]\label{prop:mobius}
Let $q$ be such that $I-M_0$ is invertible on $\ell^1$ and $1-\mathfrak g\,t_1\ne0$. Then
$I-M_0-\mathfrak guv^{\!\top}$ is invertible on $\ell^1$,
\begin{equation}\label{eq:smresolvent}
   (I-M_0-\mathfrak guv^{\!\top})^{-1}
   =R_0+\frac{\mathfrak g}{1-\mathfrak gt_1}\,R_0uv^{\!\top}R_0,\qquad R_0=(I-M_0)^{-1},
\end{equation}
and the $E$-sourced sum $B^E(q,y):=\mathbf 1^{\!\top}(I-M_0-\mathfrak guv^{\!\top})^{-1}E$ is
$(b_0+\mathfrak g\kappa)/(1-\mathfrak g t_1)$ with $\kappa=t_0b_1-b_0t_1$, which is
\eqref{eq:mobius}. \textup{(}The bulk sum of the model has source $u$; see
Remark~\ref{rem:bulksource}.\textup{)} In particular $B^E$ is a M\"obius function of $\mathfrak g$, with a single zero,
at $\mathfrak g_0=-b_0/\kappa$, and a single pole, at $\mathfrak g=1/t_1$.
\end{proposition}

\begin{proof}
Sherman--Morrison. Put $\delta=1-\mathfrak g\,v^{\!\top}R_0u=1-\mathfrak g t_1\ne0$. Then the
right-hand side of \eqref{eq:smresolvent} is a two-sided inverse of
$I-M_0-\mathfrak g uv^{\!\top}$:
\[
  (I-M_0-\mathfrak g uv^{\!\top})\bigl(R_0+\mathfrak g\delta^{-1}R_0uv^{\!\top}R_0\bigr)
  =I+\mathfrak g\bigl(\delta^{-1}-1-\mathfrak g\delta^{-1}t_1\bigr)uv^{\!\top}R_0
  =I ,
\]
since $\delta^{-1}(1-\delta-\mathfrak gt_1)=0$; the other side is identical. All pairings converge
absolutely: $u,E\in\ell^1$, $v,\mathbf 1\in\ell^\infty$, and $R_0$ is bounded on $\ell^1$.
Applying $\mathbf 1^{\!\top}(\,\cdot\,)E$,
\[
  B^E=\mathbf 1^{\!\top}R_0E+\mathfrak g\delta^{-1}(\mathbf 1^{\!\top}R_0u)(v^{\!\top}R_0E)
   =b_0+\frac{\mathfrak g\,b_1t_0}{1-\mathfrak g t_1}
   =\frac{b_0+\mathfrak g(t_0b_1-b_0t_1)}{1-\mathfrak g t_1}. \qedhere
\]
\end{proof}

\begin{remark}[the bulk sum of the model]\label{rem:bulksource}
The system \eqref{eq:gapkernel} has right-hand side $u$, not $E$. By \eqref{eq:smresolvent},
\[
   B(q,y)=\mathbf 1^{\!\top}(I-M_0-\mathfrak guv^{\!\top})^{-1}u=\frac{b_1}{1-\mathfrak g\,t_1},
   \qquad
   B_0(q,y):=1+\mathfrak g\sum_{b\ge1}q^bP_b=\frac1{1-\mathfrak g\,t_1},
\]
where $B_0$ is the weight of the junction-adjacent bulk chains that are empty or begin with a gap
run. The M\"obius form $(b_0+\mathfrak g\kappa)/(1-\mathfrak gt_1)$ of
Proposition~\ref{prop:mobius} is the resolvent of $E$. It enters below only through $b_0$, $t_1$ and
the gate, and not as the bulk sum.
\end{remark}

\begin{remark}[the two hypotheses are the two the paper already carries]\label{rem:mobhyp}
Invertibility of $I-M_0$ is $S_e\ne0$ (Corollary~\ref{cor:sing}), which at the travel poles is
Proposition~\ref{prop:selower}. And $1-\mathfrak g_V t_1\ne0$ is $s\ne1$, the gate half $|s|<1$ of
\eqref{eq:gate}, which is Theorem~\ref{thm:star} with Lemma~\ref{app:Se};
$1-\mathfrak g_Ut_1\ne0$ follows, since $\mathfrak g_U=\mathfrak g_V/(1+q)<\mathfrak g_V$ and
$t_1>0$ at the poles. So Proposition~\ref{prop:mobius} consumes no hypothesis that
\S\ref{sec:inputL} was not already carrying.
\end{remark}

\subsection*{Machine verification of the telescoping, the factorisation and the reciprocity}

Proposition~\ref{prop:mobius} is formally verified in Lean~4 against Mathlib
\cite{MathlibCommunity} in \texttt{Mobius.lean} and \texttt{MobiusL1.lean}. Its hypotheses are its
Lean hypotheses: `$I-M_0$ is invertible' is an operator $R_0$ with $R_0(I-M_0)=(I-M_0)R_0=I$, and
$1-\mathfrak gt_1\ne0$ is carried as stated. With those in hand nothing analytic remains.
\texttt{Mobius.lean} proves the factorisation for an arbitrary module over an arbitrary field:
\texttt{opN\_opS} and \texttt{opS\_opN} are the two halves of \eqref{eq:smresolvent},
\texttt{solution\_unique} is the unique solvability of \eqref{eq:gapkernel},
\texttt{B\_mobius} is \eqref{eq:mobius}, and \texttt{B\_eq\_zero\_iff} and \texttt{B\_pole} are the
single zero at $\mathfrak g_0=-b_0/\kappa$ and the single pole at $\mathfrak g=1/t_1$.

\texttt{MobiusL1.lean} instantiates that statement at the case the proposition is about, $K=\C$
and $E=\ell^1$ with the concrete data of \eqref{eq:rankone}. For $|q|<1$ a kernel with rows
dominated by a summable sequence defines a $\C$-linear operator on $\ell^1$ (\texttt{opOfKer}) and
a bounded sequence defines a functional (\texttt{funcOfSeq}); \texttt{M0}, \texttt{uElt},
\texttt{EElt}, \texttt{vFun} and \texttt{oneFun} are $M_0$, $u$, $E$, $v^{\!\top}$ and
$\mathbf 1^{\!\top}$ of \eqref{eq:rankone}, with the index $n$ standing for site $n+1$.
\texttt{rank\_one} is the identity $2q^b\,q^{a+b}=(2q^{2b})(q^a)$, \texttt{gapSystem\_iff}
identifies the operator equation $(I-M_0-\mathfrak guv^{\!\top})P=E$ with the display
\eqref{eq:gapkernel} coordinatewise, \texttt{l1\_solution\_unique} is the unique $\ell^1$
solvability of \eqref{eq:gapkernel}, \texttt{l1\_smresolvent\_left} and
\texttt{l1\_smresolvent\_right} are \eqref{eq:smresolvent} for those operators, and
\texttt{l1\_B\_of\_solution} is the conclusion: an $\ell^1$ solution of \eqref{eq:gapkernel} has
$B=\sum_bP_b=(b_0+\mathfrak g\kappa)/(1-\mathfrak gt_1)$.

The first sentence of Proposition~\ref{prop:recip}, $t_0=b_1$, is verified in
\texttt{Reciprocity.lean} in the operator form its first proof uses, with $E$ and $F$ kept apart:
$D\colon F\to E$ and $K\colon E\to F$ symmetric for the pairing, $M_0=DK$, and the two sources
$D\mathbf 1$ and $Dv$. \texttt{pushThrough} is $(I-DK)^{-1}D=D(I-KD)^{-1}$, which replaces the
palindromic Neumann series of the paper's proof by an algebraic identity;
\texttt{X\_symm} is the symmetry of $X$ and \texttt{t0\_eq\_b1} the conclusion. The
necessity of the symmetry of $K$ is recorded as \texttt{K\_symm\_needed}.

Lemma~\ref{lem:tel} is verified in \texttt{Telescope.lean}. The site $s=n+s_0$ is indexed by
$n\in\Z_{\ge0}$, so $e_s$ is \texttt{eCo} and $s_0+\vartheta=1$ is the hypothesis \texttt{hs}; the source
is bounded and $\Phi\in\ell^1$ is carried as summability of $\|\Phi_n\|$, which is what
\texttt{MobiusL1.summable\_norm} supplies for a $\Phi$ given as an element of the $\ell^1$ space of
\texttt{MobiusL1.lean} (\texttt{l1\_three\_term}, \texttt{l1\_G\_eq}). \texttt{row\_sum} is the
split of the inner sum at $s'=s$ together with the finite and the infinite geometric sum it
produces; \texttt{three\_term} is \eqref{eq:threeterm}, the interchange of the two summations being
Fubini against the majorant $2|q|^{n+1}\|\Phi_m\|$, which is summable on $\Z_{\ge0}^2$ as a product;
\texttt{norm\_bTel\_le} is the bound $|b_k|\le\Xi|q|^{k+1}$ with $\Xi=4/(1-|q|)$ and
\texttt{summable\_bProd\_norm} is the paper's $\Xi^J|q|^{J^2+Jk}\to0$, obtained by the ratio test;
\texttt{summable\_aFrak} and \texttt{summable\_cFrak} are the convergence of $\mathfrak a_k$ and
$\mathfrak c_k$; \texttt{G\_eq} is $G_k=\mathfrak c_k+\mathfrak a_kG_1$, proved by telescoping
$R_j=G_{k+2j}\prod_{i<j}b_{k+2i}$ and letting $R_J\to0$; \texttt{G\_one\_solved} and
\texttt{G\_solved} are \eqref{eq:solved}. The last sentence of the lemma, the identification of the
two instances, is \texttt{aTel\_bulk}, \texttt{bTel\_bulk}, \texttt{aTel\_travel} and
\texttt{bTel\_travel}.

Proposition~\ref{prop:travelsum} and Proposition~\ref{prop:bulkdress} are verified in
\texttt{SectionLadders.lean}, which evaluates $\mathcal C_k$ at the three geometric sources:
\texttt{CC\_travel} is $\mathcal C_k=2q/(1-q^{k+1})=A_k$, \texttt{CC\_bulk} is
$\mathcal C_k=\alpha_k$ and \texttt{CC\_bulk\_geo} is $\mathcal C_k=\alpha_{k+1}$, whence
\texttt{cFrak\_travel}, \texttt{cFrak\_bulk} and \texttt{cFrak\_bulk\_geo} identify $\mathfrak c_k$
with $\Sigma_k$, with $S_k$ and with $T_k$. \texttt{travelsum} is
Proposition~\ref{prop:travelsum} and \texttt{bulkdress\_b0\_t0}, \texttt{bulkdress\_t1\_b1} are the
four scalars of Proposition~\ref{prop:bulkdress}. \texttt{Sladder\_eq}, \texttt{SigmaLadder\_eq}
and \texttt{Tladder\_eq} are the three ladders written out as in \eqref{eq:krec} and
\eqref{eq:Tladder}. What is verified is the implication, not the existence of the solution: the
$\ell^1$ solution is a hypothesis, which is where the paper's $|q|<\tfrac13$ lives, and no analytic
continuation step is needed, \texttt{G\_eq} holding on the whole of $|q|<1$ for any $\ell^1$
solution. Lemma~\ref{lem:modelexist} is not formalised.

The first of the two clauses of Proposition~\ref{prop:recip} that begin ``equivalently'' is
verified as \texttt{SectionLadders.recip\_ladder}: given the two bulk solutions, $t_0=b_1$ and
$T_0S_e+T_1S_0=S_1$ are the same statement, the equivalence being
\texttt{recip\_ladder\_iff}. The reciprocity itself enters there as a hypothesis; it is proved
separately in \texttt{Reciprocity.lean}, in the operator form recorded above, and the two are not
joined, that joining being the instantiation of \texttt{Reciprocity.lean} at the concrete
$\ell^1$ resolvent.

Not formalised in this section: Corollary~\ref{cor:cocycle}, and with it the second
``equivalently'' clause of Proposition~\ref{prop:recip}, the one in terms of $P_{ij}$. The
obstruction is the existence of the four limits $\lim_b(\ell_b,r_{b+1})$ and
$\lim_b(\ell'_b,r'_{b+1})$: Mathlib carries no convergence theorem for an infinite product of
matrices $\prod_b\mathcal M_b$ with $\sum_b\|\mathcal M_b-I\|<\infty$, and supplying one is a
statement about products, not about this cocycle. The rest of the corollary, the passage from
\eqref{eq:sectionid} to \eqref{eq:cocyclerec} and the identification of the four entries once the
limits exist, needs only the split already carried out in \texttt{Telescope.row\_sum}.

\begin{lemma}[a positivity segment]\label{lem:posseg}
For $0<q\le0.45$ one has $S_e(q)\ge0.32$ and $|s(q)|=\tfrac{q}{1-q}|P_{12}|/S_e\le0.78$. Since
$q_1=0.4494536\ldots<0.45$, both hold on $(0,q_1]$.
\end{lemma}

\begin{proof}
By Proposition~\ref{prop:qtrig}, $S_e=\sum_{j\ge0}(-1)^j[2(1-q)]^jq^{j^2+j}/(q;q)_{2j}$, and by
\eqref{eq:p12banked},
$P_{12}=\sum_{k\ge1}(-1)^{k-1}[2(1-q)]^kq^{k^2+k+1}(1-q^{2k})/(q;q)_{2k+1}$. For every $n$,
$(q;q)_n\ge(q;q)_\infty\ge1-q-q^2$, the last step by truncating Euler's pentagonal series
$1-q-q^2+q^5+q^7-\cdots$, whose tail is positive for $q\le\tfrac12$; on $q\le0.45$ this is
$\ge0.3475$. Each term $[2(1-q)]^jq^{j^2+j}$ with $j\ge1$ increases on $(0,0.45]$, since
$(j^2+j)/q-j/(1-q)>0$ there, so it is at most its value at $q=0.45$: successively
$0.222750$, $0.010048$, $0.000092$, the remaining terms summing to less than $10^{-5}$. Hence
$S_e\ge1-0.23290/0.3475\ge0.32$. The same estimate applied to $[2(1-q)]^kq^{k^2+k+1}$, likewise
increasing, and to $(1-q^{2k})\le1$, gives $0.100238$, $0.004521$, $0.000042$ and a tail below
$10^{-6}$, so $|P_{12}|\le0.104801/0.3475\le0.3016$ and
$|s|\le\tfrac{0.45}{0.55}\cdot\tfrac{0.3016}{0.32}\le0.78$.
\end{proof}

\subsection{The site cost}\label{sec:sitecost}

The subsections above concern the operators of Definition~\ref{def:model}. This one concerns the
combinatorial quantity those operators are meant to encode: the per-site cost of the crossing
optimisation of \cite{Bonfioli2026a}. Its value is computed here in closed form, which settles the
first half of the model input \textup{(M)} below and corrects the marker clause of the form in
which that input was previously stated.

\begin{definition}[the pairing optimisation]\label{def:pairing}
Fix marker data $(\varepsilon^\ast,\delta^\ast)\in\{\pm1\}\times\{0,1\}$, a displacement
$k^\ast\in\Z$, and deposits $d_j\in\Z$ on the integer edges $j$, edge $j$ joining site $j$ to site
$j+1$. The travel indicator is $f_j=1$ for $0\le j<k^\ast$, $f_j=-1$ for $k^\ast\le j<0$ and
$f_j=0$ otherwise, and $d_j\equiv f_j\pmod2$. A \emph{realisation} consists of
\begin{itemize}[leftmargin=1.6em,itemsep=0.1em,topsep=0.2em]
\item crossing counts $m_j\ge\max(|d_j|,|f_j|)$ with $m_j\equiv d_j\pmod2$, whose support is an
      interval containing $0$; edge $j$ is crossed $u_j=(m_j+f_j)/2$ times upward and
      $\mathrm{dn}_j=(m_j-f_j)/2$ times downward;
\item a sign in $\{\pm\}$ on each crossing, with $p^u_j$ of the up-crossings and $p^d_j$ of the
      down-crossings carrying $+$, subject to
      $d_j=2p^d_j-\mathrm{dn}_j+u_j-2p^u_j$;
\item at every site $s$, a bijection between the arrivals and the departures at $s$.
\end{itemize}
The arrivals at $s$ are the up-crossings of edge $s-1$, of class $(\mathrm L,\pm)$ by their sign,
and the down-crossings of edge $s$, of class $(\mathrm R,\pm)$; the departures are the
down-crossings of edge $s-1$, of class $(\mathrm L,\pm)$, and the up-crossings of edge $s$, of
class $(\mathrm R,\pm)$. In addition there is a virtual arrival of class $(\mathrm L,+)$ at $s=0$,
and a virtual departure at $s=k^\ast$, of class $(\mathrm R,\varepsilon^\ast)$ if $\delta^\ast=1$
and $(\mathrm L,\varepsilon^\ast)$ if $\delta^\ast=0$. A matched pair costs $0$ on the same side
with equal signs, $2$ on the same side with opposite signs, and $1$ on opposite sides. The cost of
a realisation is $\sum_jm_j$ plus the sum of the pairing costs; $\ell_R$ is the minimum cost, and
$c$ is the least number of components other than the one carrying the two virtual events, over the
realisations of minimum cost.
\end{definition}

\begin{remark}[standing of Definition~\ref{def:pairing}]\label{rem:pairingstatus}
This is the optimisation of \cite[Metric formula, as verified]{Bonfioli2026a}, restated. That it computes the relaxed
word length is verified there and is not proved. Everything in this subsection is a statement about
the optimisation itself and is proved outright. The growth-series statements of this paper do not
consume it: they go through the site-local closed forms $\ell_R,c$ and the Lean identity
\texttt{metricAll} (Theorem~\ref{thm:model}), and the optimisation enters only the interpretation
of $V$ as a relaxed count.
\end{remark}

\begin{lemma}[the transportation value]\label{lem:transport}
Let $A^\pm,B^\pm,C^\pm,D^\pm$ be non-negative integers with
$A^++A^-+B^++B^-=C^++C^-+D^++D^-$, and consider the transportation problem with supplies
$A^+,A^-,B^+,B^-$ and demands $C^+,C^-,D^+,D^-$ in the four classes
$(\mathrm L,+),(\mathrm L,-),(\mathrm R,+),(\mathrm R,-)$ under the cost of
Definition~\textup{\ref{def:pairing}}. Put
\[
   \alpha=(C^+-C^-)-(A^+-A^-),\qquad
   \beta =(B^+-B^-)-(D^+-D^-),\qquad
   \Phi  =(A^++A^-)-(C^++C^-).
\]
Then the minimum cost equals $\max\bigl(|\alpha|,|\beta|,|\Phi|\bigr)$.
\end{lemma}

\begin{proof}
Write $N^a_{\mathrm L}=A^++A^-$, $N^d_{\mathrm L}=C^++C^-$, $N^a_{\mathrm R}=B^++B^-$,
$N^d_{\mathrm R}=D^++D^-$, so $\Phi=N^a_{\mathrm L}-N^d_{\mathrm L}=N^d_{\mathrm R}-N^a_{\mathrm R}$,
and $\sigma^a_{\mathrm L}=A^+-A^-$, $\sigma^d_{\mathrm L}=C^+-C^-$, and similarly on the right, so
that $\alpha=\sigma^d_{\mathrm L}-\sigma^a_{\mathrm L}$ and
$\beta=\sigma^a_{\mathrm R}-\sigma^d_{\mathrm R}$.

Let $\pi$ be the mass sent from left supplies to right demands and $\pi'$ that from right supplies
to left demands. The mass staying on the left is $g=N^a_{\mathrm L}-\pi=N^d_{\mathrm L}-\pi'$, so
$\pi'=\pi-\Phi$; put $P=\pi+\pi'=2\pi-\Phi$. Cross transports cost $1$ each and their cost does not
depend on the signs, so for fixed $\pi,\pi'$ the cost is $P+2\,\mathrm{mis}_{\mathrm L}
+2\,\mathrm{mis}_{\mathrm R}$, where $\mathrm{mis}_{\mathrm L}$ is the least number of opposite-sign
pairs among the $g$ same-side pairs on the left, minimised over which supplies stay. If $a^+$ of the
staying supplies and $c^+$ of the served demands carry $+$, that number is $|a^+-c^+|$, since
$\min(a^+,c^+)+\min(g-a^+,g-c^+)=g-|a^+-c^+|$ pairs can be matched with equal signs and no more.
Here $a^+$ ranges over the integer interval $[\max(0,g-A^-),\min(A^+,g)]$, which is non-empty
because $0\le g\le A^++A^-$, and $c^+$ over $[\max(0,g-C^-),\min(C^+,g)]$, likewise non-empty. The
least $|a^+-c^+|$ over two non-empty integer intervals $[l_1,r_1],[l_2,r_2]$ is
$\max(0,l_1-r_2,l_2-r_1)$. Here $l_1-r_2=\max(0,g-A^-)-\min(C^+,g)$, which is $g-A^--C^+$ when
$g\ge A^-$ and $C^+\le g$, and is at most $0$ otherwise, so
$\max(0,l_1-r_2)=\max(0,g-A^--C^+)$; symmetrically $\max(0,l_2-r_1)=\max(0,g-C^--A^+)$, and
\[
   \mathrm{mis}_{\mathrm L}=\max\bigl(0,\ g-A^--C^+,\ g-C^--A^+\bigr).
\]
Substituting $A^\pm=(N^a_{\mathrm L}\pm\sigma^a_{\mathrm L})/2$ and
$C^\pm=(N^d_{\mathrm L}\pm\sigma^d_{\mathrm L})/2$ gives
$g-A^--C^+=\tfrac{\Phi-\alpha}2-\pi$ and $g-C^--A^+=\tfrac{\Phi+\alpha}2-\pi$, hence
$2\,\mathrm{mis}_{\mathrm L}=\max(0,|\alpha|+\Phi-2\pi)=\max(0,|\alpha|-P)$. The same computation on
the right gives $2\,\mathrm{mis}_{\mathrm R}=\max(0,|\beta|-\Phi-2\pi')=\max(0,|\beta|-P)$.
Therefore the cost is
\[
   \phi(P)=P+(|\alpha|-P)^++(|\beta|-P)^+ ,
\]
which for $P\ge0$ is at least $\max(|\alpha|,|\beta|)$, with equality exactly on
$\min(|\alpha|,|\beta|)\le P\le\max(|\alpha|,|\beta|)$, and which equals $P$ above that range.

It remains to identify the admissible $P$. From $\pi\ge0$ and $\pi'\ge0$ one gets $P\ge|\Phi|$;
from $\pi\le\min(N^a_{\mathrm L},N^d_{\mathrm R})$ and $\pi'\le\min(N^a_{\mathrm R},N^d_{\mathrm L})$
one gets, using $\Phi=N^a_{\mathrm L}-N^d_{\mathrm L}=N^d_{\mathrm R}-N^a_{\mathrm R}$, the single
constraint $P\le\min(N^a_{\mathrm L}+N^d_{\mathrm L},\,N^a_{\mathrm R}+N^d_{\mathrm R})$. Finally
$P\equiv\Phi\pmod2$, and $\alpha\equiv\beta\equiv\Phi\pmod2$ because $\sigma\equiv N\pmod2$ on each
side. Now $|\alpha|\le|\sigma^d_{\mathrm L}|+|\sigma^a_{\mathrm L}|\le
N^d_{\mathrm L}+N^a_{\mathrm L}$ and $|\Phi|\le N^a_{\mathrm L}+N^d_{\mathrm L}$, and symmetrically
$|\beta|,|\Phi|\le N^a_{\mathrm R}+N^d_{\mathrm R}$. Hence
$P^\ast:=\max\bigl(|\Phi|,\min(|\alpha|,|\beta|)\bigr)$ is admissible: it has the right parity, it
is at least $|\Phi|$, and it is at most $\max(|\Phi|,|\alpha|)\le N^a_{\mathrm L}+N^d_{\mathrm L}$
and at most $\max(|\Phi|,|\beta|)\le N^a_{\mathrm R}+N^d_{\mathrm R}$. If
$|\Phi|\le\max(|\alpha|,|\beta|)$ then $P^\ast$ lies in the equality range and
$\phi(P^\ast)=\max(|\alpha|,|\beta|)$, which is optimal since $\phi\ge\max(|\alpha|,|\beta|)$
throughout. If $|\Phi|>\max(|\alpha|,|\beta|)$ then every admissible $P$ satisfies $P\ge|\Phi|$ and
$\phi(P)=P$, so the minimum is $|\Phi|$, attained at $P^\ast=|\Phi|$. In both cases the minimum is
$\max(|\alpha|,|\beta|,|\Phi|)$.
\end{proof}

\begin{corollary}[the local cost law]\label{cor:localcost}
In Definition~\textup{\ref{def:pairing}}, let $\alpha_s,\beta_s,\Phi_s$ be the three quantities of
Lemma~\textup{\ref{lem:transport}} formed from the arrivals and departures at site $s$. Then
\[
  \alpha_s=d_{s-1}-[s=0]+\varepsilon^\ast[s=k^\ast][\delta^\ast=0],\qquad
  \beta_s=d_{s}-\varepsilon^\ast[s=k^\ast][\delta^\ast=1],
\]
\[
  \Phi_s=f_{s-1}+[s=0]-[s=k^\ast][\delta^\ast=0],
\]
one has $|\Phi_s|\le\min(|\alpha_s|,|\beta_s|)$ in every case, and the minimal pairing cost at $s$
is
\[
   \mathrm{Site}(s)=\max\bigl(|\alpha_s|,|\beta_s|\bigr),
\]
independently of the crossing counts $m_j$, of the sign splits $p^u_j,p^d_j$, and of the travel
indicators. In particular at a site carrying no virtual event
$\mathrm{Site}(s)=\max(|d_{s-1}|,|d_s|)$.
\end{corollary}

\begin{proof}
The displays are read off Definition~\ref{def:pairing}: the up-crossings of edge $s-1$ contribute
$\sigma^a_{\mathrm L}$, the down-crossings contribute $\sigma^d_{\mathrm L}$, and
$d_{s-1}=2p^d_{s-1}-\mathrm{dn}_{s-1}+u_{s-1}-2p^u_{s-1}=\sigma^d_{\mathrm L}-\sigma^a_{\mathrm L}$,
which is $\alpha_s$ before the virtual events; the virtual arrival adds one to $A^+$ and the virtual
departure adds one to $C^{\varepsilon^\ast}$ or to $D^{\varepsilon^\ast}$. Likewise
$d_s=\sigma^a_{\mathrm R}-\sigma^d_{\mathrm R}$ and
$\Phi_s=u_{s-1}-\mathrm{dn}_{s-1}=f_{s-1}$ before the virtual events. Since the sign splits enter
only through $\sigma^a,\sigma^d$, and those enter $\alpha_s,\beta_s,\Phi_s$ only through the fixed
combinations $d_{s-1},d_s,f_{s-1}$, none of the three depends on $m_j$ or on $p^u_j,p^d_j$.

For $|\Phi_s|\le\min(|\alpha_s|,|\beta_s|)$: at a site with no virtual event $\Phi_s=f_{s-1}=f_s$,
since arrivals and departures balance, and $|d_j|\ge|f_j|$, since a bulk deposit is even with $f=0$
and a travel deposit is odd with $|f|=1$. The virtual events sit at $s=0$ and $s=k^\ast$, which for
$k^\ast\ne0$ are exactly the two sites where $f$ jumps, and there are eight remaining cases. For $s=0<k^\ast$: $\Phi_0=f_{-1}+1=1$, while $\alpha_0=d_{-1}-1$ is odd
and $\beta_0=d_0$ is odd. For $s=0$ and $k^\ast<0$: $\Phi_0=-1+1=0$. For $s=k^\ast>0$: $\Phi=1-1=0$
if $\delta^\ast=0$, and $\Phi=1$ with $\alpha=d_{k^\ast-1}$ odd, $\beta=d_{k^\ast}-\varepsilon^\ast$
odd, if $\delta^\ast=1$. For $s=k^\ast<0$: $\Phi=-1$ with $\alpha=d_{k^\ast-1}+\varepsilon^\ast$ odd
and $\beta=d_{k^\ast}$ odd if $\delta^\ast=0$, and $\Phi=0$ if $\delta^\ast=1$. For $k^\ast=0$,
where the two virtual events share the site: $\Phi_0=0$ if $\delta^\ast=0$, and $\Phi_0=1$ with
$\alpha_0=d_{-1}-1$ and $\beta_0=d_0-\varepsilon^\ast$ both odd if $\delta^\ast=1$. In each case
$|\Phi_s|\le\min(|\alpha_s|,|\beta_s|)$, so Lemma~\ref{lem:transport} gives
$\mathrm{Site}(s)=\max(|\alpha_s|,|\beta_s|)$.
\end{proof}

\begin{corollary}[the closed form for $\ell_R$, and rigidity of the crossing counts]\label{cor:lRclosed}
Let the \emph{span} be the least interval of edges containing $0$, all $j$ with $d_j\ne0$ and all
$j$ with $f_j\ne0$. Every realisation of minimum cost has support exactly the span and crossing
counts
\[
   m_j=\max(|d_j|,|f_j|), \quad\text{forced to } 2 \text{ when that value is }0,
\]
and
\[
   \ell_R=\sum_{j\in\mathrm{span}}m_j+\sum_{s}\max\bigl(|\alpha_s|,|\beta_s|\bigr),
\]
the second sum over the sites of the span.
\end{corollary}

\begin{proof}
By Corollary~\ref{cor:localcost} the pairing cost at every site is independent of the $m_j$, so the
total cost is $\sum_jm_j$ plus a quantity that does not depend on the $m_j$; hence a minimum-cost
realisation minimises $\sum_jm_j$ alone, subject to $m_j\ge\max(|d_j|,|f_j|)$, $m_j\equiv d_j$, and
the support being an interval containing $0$. Since $d_j\equiv f_j$, the value
$\max(|d_j|,|f_j|)$ has the parity of $d_j$ and is the least admissible value on the span except
where it vanishes, where the support condition forces $m_j\ge2$. Enlarging the support beyond the
span adds an edge with $m\ge2$ and creates one further site, of cost $\max(0,0)=0$, without
changing any other site cost, so it adds $2$ and never helps.
\end{proof}

\noindent Corollary~\ref{cor:localcost} is the first clause of \textup{(M1)} below, proved. The
marker clause needs the sign, and the form previously used does not carry it.

\begin{corollary}[the marker junctions]\label{cor:marker}
Let $k^\ast>0$, let $d_L=d_{-1}$ be the last bulk deposit \textup{(}even\textup{)} and $d_R=d_0$ the
first travel deposit \textup{(}odd\textup{)}. Then the cost of the junction at site $0$ is
\[
   \mathrm{Site}_0(d_L,d_R)=\max\bigl(|d_L-1|,\ |d_R|\bigr),
\]
for all four marker data. At the far junction, site $k^\ast$, with $d_L=d_{k^\ast-1}$ odd and
$d_R=d_{k^\ast}$ even, the cost is $\max(|d_L+\varepsilon^\ast|,|d_R|)$ if $\delta^\ast=0$ and
$\max(|d_L|,|d_R-\varepsilon^\ast|)$ if $\delta^\ast=1$; it is not independent of
$(\varepsilon^\ast,\delta^\ast)$ and it is not the mirror image of $\mathrm{Site}_0$.
\end{corollary}

\begin{proof}
Corollary~\ref{cor:localcost}. At site $0$ neither $\varepsilon^\ast$ nor $\delta^\ast$ occurs.
\end{proof}

\begin{remark}[the earlier marker form is false]\label{rem:markerfalse}
Earlier versions of this paper, and \cite{supplement}, carried
$\mathrm{Site}_0(d_L,d_R)=\max(|d_L|-1,|d_R|)$. That agrees with
Corollary~\ref{cor:marker} exactly on $d_L\ge0$ and fails on every cell with $d_L\le-2$ and
$|d_R|\le|d_L|$: the smallest is $d_L=-2$, $d_R=1$, where the true cost is $3$ and the earlier form
gives $1$. The verification on which the earlier form rested covered $0\le d_L\le12$ only. The
junction cost therefore depends on the \emph{sign} of the last bulk deposit, not on its magnitude
alone.

The correction is local. Every bulk edge cost and every interior bulk site cost depends on the
deposit only through its magnitude (Corollary~\ref{cor:localcost}), so reversing the sign of the
junction-adjacent deposit of a bulk run is a weight-preserving involution on bulk runs; writing
$B_\sigma$ for the generating function of the bulk runs whose junction-adjacent deposit has
magnitude $\sigma$ \textup{(}so $B_{2b}=P_b$ of \eqref{eq:gapkernel}, and
$B_0=1/(1-\mathfrak gt_1)\ne0$ of Remark~\ref{rem:bulksource} collects the empty chain and the chains
opening with a gap run\textup{)}, the two signs carry $B_\sigma/2$ each, and the near-junction vector
of \textup{(M3)} is
\begin{equation}\label{eq:junctionsym}
   \mu^{(y)}_s=B_0\,x^{2s+1}+\sum_{b\ge1}\frac{P_b}{2}
   \Bigl(x^{\max(2b-1,\,2s+1)}+x^{\max(2b+1,\,2s+1)}\Bigr),
\end{equation}
in place of $\sum_\sigma B_\sigma x^{\max(\sigma-1,2s+1)}$. The far-junction vector is
\begin{equation}\label{eq:farjunction}
   E^{(y)}_s=B_0^{[s]}x^{2s}+B_0\,x^{2s+2}+\sum_{b\ge1}P_b
   \Bigl(x^{\max(2s,\,2b)}+x^{\max(2s+2,\,2b)}\Bigr),\qquad
   B_0^{[s]}=\begin{cases}B_{0z}:=1+y\,\mathfrak g\sum_bq^bP_b,&s=0,\\ B_0,&s\ge1,\end{cases}
\end{equation}
and $\lambda^{(y)}:=E^{(y)}+2\mu^{(y)}$. The slot $B_{0z}$ carries the junction cut: at $y=1$ it
equals $B_0$, and at $y=q$ it equals $B_0\bigl(1-t_1q/(1+q)\bigr)$. This changes no rank statement below:
Proposition~\ref{prop:junction} concludes that the junction is not of rank one, and the symmetrised
matrix \eqref{eq:junctionsym} has rank $6$ on the block $\sigma\in\{0,2,4,6,8,10\}$,
$2s+1\in\{1,3,5,7,9,11\}$, by exact elimination over $\Q$ at $x=\tfrac12$, against rank $5$ for the
earlier form on the same block.
\end{remark}

\begin{proposition}[the cut criterion and the shield lower bound]\label{prop:cut}
Call a site $s$ of the span \emph{cut} if $\alpha_s=\beta_s=\Phi_s=0$, and let $Z$ be the set of
cut sites interior to the span. At a cut site every minimum-cost pairing matches each arrival with a
departure on its own side, so no strand crosses $s$; at any other interior site $s$ some
minimum-cost pairing lets a strand cross. Consequently every minimum-cost realisation has at least
$|Z|+1$ components and
\[
   c\ \ge\ |Z| .
\]
A maximal run of $L$ gap edges of the bulk \textup{(}$d_j=f_j=0$\textup{)} flanked by bulk deposits
and containing no virtual event contributes exactly its $L-1$ interior sites to $Z$.
\end{proposition}

\begin{proof}
At a cut site $\phi(P)=P$ in the notation of Lemma~\ref{lem:transport}, so the unique minimiser is
$P=0$, which is the assertion. At any other interior site $\max(|\alpha_s|,|\beta_s|)>0$, since
$|\Phi_s|\le\min(|\alpha_s|,|\beta_s|)$; the equality range of $\phi$ is
$[\min(|\alpha_s|,|\beta_s|),\max(|\alpha_s|,|\beta_s|)]$ and contains a value $P>0$ admissible for
the caps of Lemma~\ref{lem:transport}: namely $\min(|\alpha_s|,|\beta_s|)$ when that is positive,
which is the choice $P^\ast$ of Lemma~\ref{lem:transport}; and $2$ otherwise, since
$\min(|\alpha_s|,|\beta_s|)=0$ forces $\Phi_s=0$, hence $\alpha_s\equiv\beta_s\equiv0\pmod2$ and
$\max(|\alpha_s|,|\beta_s|)\ge2$. The two caps of Lemma~\ref{lem:transport} are $m_{s-1}$ and $m_s$,
each augmented by $1$ for a virtual event of $s$ on that side, and both are at least $2$ when
$\Phi_s=0$: an edge with $f=0$ has $m\ge2$ on the span by Corollary~\ref{cor:lRclosed}, while
$\Phi_s=f_{s-1}+[s=0]-[s=k^\ast][\delta^\ast=0]=f_s+[s=k^\ast][\delta^\ast=1]$ shows that an adjacent
edge with $|f|=1$ forces a virtual event of $s$ on its side, which supplies the missing unit. Removing the $|Z|$ cut sites disconnects the crossings of the span
into at least $|Z|+1$ classes, no two of which can lie in the same component; one component carries
the virtual events, so at least $|Z|$ are isolated cycles. For the last sentence, an interior site
of such a run has $d_{s-1}=d_s=0$ and $f_{s-1}=0$ and no virtual event, hence is cut, while the two
sites at the ends of the run have $|\alpha|$ resp.\ $|\beta|$ equal to the flanking bulk deposit,
which is even and non-zero.
\end{proof}

\begin{remark}[what the shield law still owes]\label{rem:shieldowes}
Proposition~\ref{prop:cut} gives $c\ge|Z|$. The reverse inequality $c\le|Z|$ asks for one
minimum-cost realisation whose components are exactly the $|Z|+1$ classes cut out by $Z$, that is,
for pairings that connect all crossings within each class. It is not proved here. It is verified by
exhaustive enumeration, over every sign split and every site bijection, on all bulk configurations
with $k^\ast=0$, up to $8$ edges, even deposits of magnitude at most $4$ and both signs, and all
four marker data: $1\,048\,544$ configurations, gap runs up to length $6$, $0$ exceptions to both
$c=|Z|$ and the closed form of Corollary~\ref{cor:lRclosed}
(\texttt{code/zeta\_probe/tools/sitecost}, mode \texttt{shield}). The closed form $|Z|$ also recovers, by inspection, the boundary shield
$\mathrm{sh}_R=\mathbf 1[\varepsilon^\ast=-1\text{ or }\delta^\ast=1]$ that
\cite[the explicit transducer]{Bonfioli2026a} records at $k^\ast=0$: by
Corollary~\ref{cor:localcost} the marker site is cut exactly when $d_{-1}=0$, $\delta^\ast=0$ and
$\varepsilon^\ast=1$, which is exactly where that indicator vanishes.
\end{remark}

\begin{remark}[the boundary term in the cut count, and the check at $k^\ast\ne0$]\label{rem:boundaryterm}
The set $Z$ of Proposition~\ref{prop:cut} consists of cut sites \emph{interior} to the span. That
count is short by one whenever both markers sit at an endpoint of the span on the side opposite the
edges, since cutting there isolates the marker component and leaves every crossing in closed walks.
The configuration is exactly
\[
   k^\ast=0,\qquad \delta^\ast=0,\qquad \mathrm{lo}=0<\mathrm{hi},
\]
and it is the boundary shield of Remark~\ref{rem:shieldowes}: the marker site is cut precisely where
$\mathrm{sh}_R$ vanishes. The clause $\mathrm{lo}=0<\mathrm{hi}$ is needed as well, since without it
the identity is miscounted, having $c=0$ and one cut site.

Writing $Z^{+}$ for $Z$ together with that site when the displayed condition holds, the shield law
$c=|Z^{+}|$ has no exceptions on the $3\,336\,511$ elements of word length at most $31$, over all
$k^\ast$, with $c=(\ell_T-\ell_R)/2$ computed from breadth-first word length and the closed form of
Corollary~\ref{cor:lRclosed} (\texttt{code/zeta\_probe/tools/nogap}). This extends the check of
Remark~\ref{rem:shieldowes}, which covers the bulk $k^\ast=0$ only. Enlarging $Z$ by every marker
site instead, rather than by the boundary configuration above, produces $8555$ violations of
Proposition~\ref{prop:cut} at word length $19$, which is how the condition was located.

The count $|Z^{+}|$ is the site-local cut count \texttt{cTrue} of \texttt{CorrectedSpan.lean}, and
$\ell_T=\ell_R+2|Z^{+}|$, with $\ell_R$ the closed form of Corollary~\ref{cor:lRclosed}, is now the
Lean theorem \texttt{metricAll} (Theorem~\ref{thm:model}). What remains unproved is only the
reverse inequality $c\le|Z^{+}|$ for the component count $c$ of the optimisation of
Definition~\ref{def:pairing}; the growth-series results do not use it.
\end{remark}

\subsection{The gap-free case: the reverse inequality holds when $Z$ is empty}\label{ss:nogap}

Call an edge $j$ of the span a \emph{gap edge} when $d_j=0$ and $f_j=0$; by
Corollary~\ref{cor:lRclosed} such an edge is forced to $m_j=2$ by reachability. The reverse
inequality is out of reach in general (Remark~\ref{rem:shieldowes}), but it holds, with a short
proof, on the elements that carry no gap edge at all.

\begin{lemma}[the sign split is forced at minimum multiplicity]\label{lem:splitforced}
Let $j$ be a span edge with no gap, so $d_j\ne0$. Then $m_j^\ast=|d_j|$, and the split
$p^d_j=p^u_j+(d_j-f_j)/2$ subject to $0\le p^u_j\le u_j$ and $0\le p^d_j\le dn_j$ has exactly one
solution: $p^u_j=0$ when $d_j>0$ and $p^u_j=u_j$ when $d_j<0$. Consequently every up-crossing of
edge $j$ carries the sign $-\operatorname{sgn}(d_j)$ and every down-crossing the sign
$+\operatorname{sgn}(d_j)$.
\end{lemma}

\begin{proof}
$d_j\equiv f_j\pmod 2$ and $d_j\ne0$ give $|d_j|\ge|f_j|$ in both parities, so $m_j^\ast=|d_j|$ and
$2\,dn_j=d_j-f_j$ when $d_j>0$. Then $p^d_j=p^u_j+dn_j\le dn_j$ forces $p^u_j\le0$, hence
$p^u_j=0$. The case $d_j<0$ is the mirror image.
\end{proof}

\begin{lemma}[the pair cost sees only the side pattern]\label{lem:sidesonly}
At a site both of whose edges are gap-free, the left arrivals all carry one sign and the left
departures the opposite sign, and likewise on the right. Hence
$\operatorname{pc}=2$ for a bounce and $1$ for a pass, independently of which strands are matched.
\end{lemma}

\begin{proof}
Left arrivals are the up-crossings of the left edge and left departures its down-crossings, so
Lemma~\ref{lem:splitforced} gives them opposite signs; same on the right.
\end{proof}

\begin{lemma}[free swaps]\label{lem:freeswap}
Under Lemma~\ref{lem:sidesonly}, the $2$-swap exchanging the departures of two pairs changes the
site cost by $0$ exactly when the two arrivals lie on the same side or the two departures do; the
only cost-lowering patterns are two bounces on opposite sides. In particular a pair that is a bounce
never blocks.
\end{lemma}

\begin{proof}
Sixteen side patterns, checked exhaustively; formalised as
\texttt{NoGapMerge.swap\_free\_iff}, \texttt{swap\_neg\_iff} and \texttt{bounce\_never\_blocks}.
\end{proof}

\begin{theorem}[no gap edge forces $c=0$]\label{thm:nogap}
If $g$ has no gap edge then $c(g)=0$.
\end{theorem}

\begin{proof}
Since the site cost of Lemma~\ref{lem:transport} depends on neither the crossing counts nor the
sign splits, every relaxed-optimal realisation has $m_j=m_j^\ast$, so
Lemmas~\ref{lem:splitforced}--\ref{lem:freeswap} apply at every site. Take a relaxed-optimal
realisation with the fewest components and suppose there are at least two. If two components hold
strands of the leftmost span edge, that site is shared and both have $(R\to R)$ bounces there, so by
Lemma~\ref{lem:freeswap} they merge at no cost. Otherwise one component owns that edge; let $B$ be
another component whose leftmost site $r$ is minimal, so $r$ exceeds the left end of the span. Edge
$r-1$ then lies in the span with $m_{r-1}\ge1$, and its strands are not $B$'s, so site $r$ is
shared; $B$'s ends at $r$ come only from edge $r$, so $B$ bounces there and again the merge is free.
A $2$-swap across two components merges them, so the component count drops at no cost, contradicting
minimality. Hence there is one component and no isolated cycle.

One configuration escapes the last step: if $r$ carries the virtual arrival, $B$ may have a single
$(L\to R)$ pass and no bounce, and a neighbour with an opposing $(R\to L)$ pass and $d_{r-1}<0$ would
block. That neighbour then has no left arrival at $r$, so its strands on edge $r-1$ are all
down-crossings, and any up-crossing there would belong to a third component which merges with $B$
freely; so the blocker also needs $u_{r-1}=0$, that is $d_{r-1}=f_{r-1}=-1$ and a single strand. But
$f_{r-1}=-1$ forces $k^\ast<0$, and $r=0$ since $r$ carries the virtual arrival, so the marker
component runs from site $0$ to site $k^\ast<0$ and must traverse edge $r-1$; its one strand is
therefore the marker component's, contradicting the assumption that it belongs to another. The
configuration does not occur.
\end{proof}

\begin{corollary}\label{cor:localzero}
If the lamp support of $g$ lies inside its travel interval then $c(g)=0$.
\end{corollary}

\begin{proof}
Support inside $I_k$ makes the span equal to $I_k$, on which $f_j=\pm1$, so no edge is a gap edge;
apply Theorem~\ref{thm:nogap}.
\end{proof}

Corollary~\ref{cor:localzero} is the statement used in the lifting argument for the true series, where it
had been verified but not proved. Verification of Theorem~\ref{thm:nogap} independent of the proof:
$2\,505\,271$ elements of word length at most $31$, no exceptions
(\texttt{code/zeta\_probe/tools/nogap}).

\begin{remark}[the argument does not extend past $Z=\emptyset$]\label{rem:nogapstops}
Theorem~\ref{thm:nogap} rests on Lemma~\ref{lem:splitforced}, which needs $d_j\ne0$. A gap edge has
$m=2$ and $d=0$, so $p^d=p^u$ with $u=dn=1$ admits two splits and the homogeneity of
Lemma~\ref{lem:sidesonly} fails. Enumerating all $256$ sign-and-side configurations of a $2$-swap,
$52$ block, of which $36$ share an arrival or a departure side, and a cost-$0$ bounce is blocked in
$28$; both criteria of Lemma~\ref{lem:freeswap} are therefore false in general. The reverse shield
inequality is not reachable by this route.
\end{remark}

\begin{remark}[the relaxed length decomposes at every $k^\ast$]\label{rem:decomp}
Corollary~\ref{cor:lRclosed} is verified by direct enumeration only for $k^\ast=0$
(Remark~\ref{rem:shieldowes}), but it holds at every $k^\ast$. The site cost depends on neither
$m_j$ nor the sign splits, so $\sum_j m_j$ is minimised separately at $m_j=m_j^\ast$; edge $j$
presents a departure at site $j$ and an arrival at site $j+1$, so the matchings at distinct sites act
on disjoint sets of ends and are chosen independently; and in the relaxed model every family of
per-site matchings is admissible, isolated cycles being free. No term falls below its local minimum
and all local minima are attained together. The true metric behaves differently precisely because the
single-walk requirement couples the sites, which is what Theorem~\ref{thm:nogap} has to handle.
\end{remark}

\subsection*{Machine verification of the site cost}

Lemma~\ref{lem:transport}, Corollary~\ref{cor:localcost}, Corollary~\ref{cor:lRclosed},
Corollary~\ref{cor:marker}, the counterexample of Remark~\ref{rem:markerfalse} and all of
Proposition~\ref{prop:cut} except its component count are formally verified in Lean~4 against
Mathlib \cite{MathlibCommunity}, in six files: \texttt{SiteCost.lean} (the transportation problem
at one site), \texttt{MarkedSite.lean} (the virtual events and the local cost law),
\texttt{PairingMatrix.lean} (bijections against multiplicity matrices),
\texttt{Realisation.lean} (the whole edge path), \texttt{CutCross.lean} (the middle sentence of
Proposition~\ref{prop:cut}) and \texttt{CutComponents.lean} (the counting step of the component
bound, over an abstract graph).

A pairing at a site is the structure \texttt{Plan}: sixteen non-negative integers with the eight
matching equations of Definition~\ref{def:pairing}; \texttt{Plan.cost} is its cost and
\texttt{Plan.cross} the number of pairs on opposite sides. The minimum of
Lemma~\ref{lem:transport} is formalised as \texttt{IsLeast}, so both halves are proved separately:
\texttt{cost\_lower\_bound} holds for every pairing and \texttt{exists\_plan\_cost\_eq} exhibits
one of exactly that cost. Definition~\ref{def:pairing} asks for a \emph{bijection} between the
arrivals and the departures rather than for a matrix, and the two are reconciled:
\texttt{planOfBij} sends a bijection to its multiplicity matrix without changing the cost, and
\texttt{bijOfMatrix} realises every matrix with the right row and column sums by a bijection, so
\texttt{bij\_min} is Lemma~\ref{lem:transport} stated over bijections and
\texttt{MarkedSite.site\_cost\_bij} is Corollary~\ref{cor:localcost} stated over bijections.
A site is the structure \texttt{MarkedSite}: the displacement $k^\ast$, the site $s$, the marker
data, and the crossing data of the two adjacent edges, constrained to realise the travel
indicator \texttt{travel}. A realisation of a whole path is the structure \texttt{Realisation}
over a \texttt{PathData}, which carries $k^\ast$, the marker data, the deposits and the span
$[A,B]$; \texttt{PathData.emptyPath} is a witness that the hypotheses are consistent.

\begin{center}\small
\begin{tabular}{lll}
\hline
Paper step & Lean declaration & axioms \\
\hline
$\mathrm{cost}\ge\max(|\alpha|,|\beta|,|\Phi|)$, six duals & \texttt{cost\_lower\_bound} & standard\\
a pairing attaining it & \texttt{exists\_plan\_cost\_eq} & standard\\
Lemma~\ref{lem:transport}, the minimum & \texttt{transport\_min} & standard\\
a bijection gives a matrix of equal cost & \texttt{planOfBij\_cost} & standard\\
a matrix gives a bijection of equal cost & \texttt{bijOfMatrix\_cost} & standard\\
Lemma~\ref{lem:transport} over bijections & \texttt{bij\_min} & standard\\
collapse to $\max(|\alpha|,|\beta|)$ & \texttt{siteCost\_eq} & standard\\
$\alpha_s=d_{s-1}-[s{=}0]+\varepsilon^\ast[s{=}k^\ast][\delta^\ast{=}0]$ & \texttt{MarkedSite.alpha\_eq} & standard\\
$\beta_s=d_s-\varepsilon^\ast[s{=}k^\ast][\delta^\ast{=}1]$ & \texttt{MarkedSite.beta\_eq} & standard\\
$\Phi_s=f_{s-1}+[s{=}0]-[s{=}k^\ast][\delta^\ast{=}0]$ & \texttt{MarkedSite.Phi\_eq} & standard\\
$|\Phi_s|\le\min(|\alpha_s|,|\beta_s|)$, all site types & \texttt{MarkedSite.Phi\_le\_min} & standard\\
Corollary~\ref{cor:localcost}, the minimum at any site & \texttt{MarkedSite.site\_cost} & standard\\
Corollary~\ref{cor:localcost} over bijections & \texttt{MarkedSite.site\_cost\_bij} & standard\\
independence of $m_j$ and of $p^u_j,p^d_j$ & \texttt{MarkedSite.cost\_indep} & \texttt{propext}\\
interior site cost $\max(|d_{s-1}|,|d_s|)$ & \texttt{MarkedSite.cost\_interior} & \texttt{propext}\\
Corollary~\ref{cor:marker}, near junction & \texttt{PathData.marker\_near} & \texttt{propext}\\
Corollary~\ref{cor:marker}, far junction, $\delta^\ast{=}0$ & \texttt{PathData.marker\_far\_left} & \texttt{propext}\\
Corollary~\ref{cor:marker}, far junction, $\delta^\ast{=}1$ & \texttt{PathData.marker\_far\_right} & \texttt{propext}\\
\quad and it is neither $(\varepsilon^\ast,\delta^\ast)$-free & \texttt{PathData.marker\_far\_depends} & \emph{none}\\
\quad nor the mirror of $\mathrm{Site}_0$ & \quad(the same declaration) & \\
off the span every site is empty & \texttt{Realisation.pair\_cost\_zero\_outside} & standard\\
Corollary~\ref{cor:lRclosed}, the minimum & \texttt{lR\_closed} & standard\\
Corollary~\ref{cor:lRclosed}, rigidity of $m_j$ & \texttt{Realisation.rigidity} & standard\\
Proposition~\ref{prop:cut}, cut site: no crossing & \texttt{cut\_forces\_no\_cross} & standard\\
\quad the same over a realisation & \texttt{Realisation.cut\_no\_cross} & standard\\
the transportation value at a prescribed cross mass & \texttt{exists\_plan\_of\_cross} & standard\\
the two caps $m_{s-1},m_s\ge2$ where $\Phi_s=0$ & \texttt{Realisation.capL\_two} & standard\\
& \texttt{Realisation.capR\_two} & standard\\
Proposition~\ref{prop:cut}, a non-cut site: a crossing & \texttt{Realisation.exists\_cross\_of\_not\_cut} & standard\\
Proposition~\ref{prop:cut}, a gap run's cut sites & \texttt{PathData.gap\_run\_cut} & standard\\
$Z$ interior to the span & \texttt{PathData.cutSet\_lt}, \texttt{cutSet\_le} & standard\\
every edge of the span carries a crossing & \texttt{Realisation.m\_pos\_of\_min} & standard\\
$\mathrm{blk}$ constant on components & \texttt{CutComponents.blk\_reachable} & \texttt{propext}\\
$\mathrm{blk}$ attains $0,\dots,|Z|$ on the span & \texttt{CutComponents.exists\_pos\_with\_gz} & standard\\
at least $|Z|+1$ components, given the graph & \texttt{CutComponents.exists\_injective\_components} & standard\\
\quad and $c\ge|Z|$ from it & \texttt{\dots\_avoiding} & standard\\
the two marker forms agree on $d_L\ge0$ & \texttt{marker\_forms\_agree} & standard\\
and differ at $(d_L,d_R)=(-2,1)$ & \texttt{marker\_forms\_differ} & \emph{none}\\
\hline
\end{tabular}
\end{center}

\noindent ``standard'' means \texttt{propext}, \texttt{Classical.choice},
\texttt{Quot.sound}, or a subset of them; several of the rows above in fact report only
\texttt{propext} and \texttt{Quot.sound}. The last row is closed by \texttt{decide} and depends
on no axiom at all, as do \texttt{PathData.marker\_far\_depends} and
\texttt{Realisation.site\_cost\_eq}. The cost of a realisation is summed over the sites of its
support interval, which is legitimate because both virtual events lie there
(\texttt{Realisation.virtual\_in\_window}) and every site outside it is empty.

The middle sentence of Proposition~\ref{prop:cut}, that at an interior site which is not cut some
minimum-cost pairing lets a strand cross, is verified in the form the proof uses.
\texttt{exists\_plan\_cost\_eq} settles every case except $\min(|\alpha_s|,|\beta_s|)=0$ with
$\Phi_s=0$; there \texttt{exists\_plan\_of\_cross} runs the construction of
Lemma~\ref{lem:transport} at the prescribed cross mass $P=2$, which is admissible because
$\Phi_s=0$ forces both caps to be at least $2$, and that is a path-level fact, proved as
\texttt{Realisation.capL\_two} and \texttt{Realisation.capR\_two} from
$m_j\ge\max(|d_j|,|f_j|)$ on the span together with
$\Phi_s=f_{s-1}+[s{=}0]-[s{=}k^\ast][\delta^\ast{=}0]=f_s+[s{=}k^\ast][\delta^\ast{=}1]$.

The counting step of the component bound is verified in \texttt{CutComponents.lean}, over an
abstract graph. Write $\mathrm{blk}(v)$ for the number of cut sites at or below the edge carrying
the crossing $v$. \texttt{blk\_adj} and \texttt{blk\_reachable} are that $\mathrm{blk}$ is constant
on connected components, given the locality hypothesis \texttt{Local}: an edge of the graph joins
crossings on the two edges adjacent to one site, and joins crossings on \emph{different} edges only
at a site that is not cut. \texttt{exists\_pos\_with\_gz} is the discrete intermediate value
argument, that $\mathrm{blk}$ takes every value in $0,\dots,|Z|$ on the span, and
\texttt{exists\_injective\_components} concludes with an injection of $\{0,\dots,|Z|\}$ into the
connected components, which is ``at least $|Z|+1$ components'' with no finiteness hypothesis.
\texttt{exists\_injective\_components\_avoiding} is the passage to $c\ge|Z|$ once the two virtual
events are known to share a component. Two of the three inputs are supplied against the structures
of \texttt{Realisation.lean}: \texttt{PathData.cutSet} is $Z$ with its two bounds
(\texttt{cutSet\_lt}, \texttt{cutSet\_le}), and \texttt{Realisation.m\_pos\_of\_min} is the
occupancy, every edge of the span of a minimum-cost realisation carrying a crossing.

Two things in this subsection are \emph{not} formalised. First, that the optimisation of
Definition~\ref{def:pairing} computes the relaxed word length (Remark~\ref{rem:pairingstatus}, an
input of \cite{Bonfioli2026a}, verified there and not proved). Second, the strand graph itself, and
with it the component count $c\ge|Z|$. \texttt{Realisation.pair} stores the multiplicity matrix of
the site pairing rather than the bijection, so the vertex set and the incidence relation that
\texttt{CutComponents.lean} quantifies over are not definable against the present structure and
would need a further definition. \texttt{Realisation.cut\_no\_cross} already carries the
mathematical content of the locality hypothesis at the level of the matrix; what is missing is the
passage from a matrix with vanishing cross mass to an incidence relation on individual crossings.

\subsection{The model \textup{(M)}, and the assembly}\label{sec:modelassembly}

Everything above is a statement about the explicit operators of Definition~\ref{def:model} and
\eqref{eq:gapkernel}, or about the optimisation of Definition~\ref{def:pairing}. What connects
those to the growth series is the strand-walk input of Route~B, which we state once and name.

\begin{theorem}[the transfer model \textup{(M)}]\label{hyp:model}\label{thm:model}
\textup{(M1)} \emph{Local cost law.} The site cost between adjacent active edges carrying deposits
$d,d'$ is $x^{\max(|d|,|d'|)}$, independent of the crossing counts and of the travel indicator.
Bulk deposits are even, $|d|=2s$, and travel deposits odd, $|d|=2s+1$, so this cost is
$q^{\max(s,s')}$ in the bulk and $x^{2\max(s,s')+1}$ on the travel side, and the edge weights are
$2q^s$ and $2x^{2s+1}$ respectively: Definition~\ref{def:model} in its two instances.
At the near marker site the cost is $x^{\mathrm{Site}_0(d_L,d_R)}$ with
$\mathrm{Site}_0(d_L,d_R)=\max(|d_L-1|,|d_R|)$, $d_L$ the last bulk deposit and $d_R$ the first
travel deposit, for all four marker data $(\varepsilon^\ast,\delta^\ast)$; at the far marker site
it is the $(\varepsilon^\ast,\delta^\ast)$-dependent cost of Corollary~\ref{cor:marker}.
\textup{(M2)} \emph{Shield law.} A maximal interior run of $L$ forced gap edges costs $q^L$ and
carries $L-1$ isolated cycles; hence the gap-bridge kernel is $\mathfrak g(q,y)=q/(1-qy)$ and the
gap-marked bulk kernel is the $K_{\mathfrak g}$ of \eqref{eq:gapkernel}.
\textup{(M3)} \emph{Decomposition.} The bivariate series of \S\ref{sec:bridge} is
$\mathcal W=\mathcal W_++\mathcal W_-+\mathcal W_0$, the sectors $k^\ast>0$, $k^\ast<0$, $k^\ast=0$,
with
\begin{equation}\label{eq:assembly}
   \mathcal W_+=x^{-1}\bigl\langle\mu^{(y)},(I-\mathcal T)^{-1}D_e\lambda^{(y)}\bigr\rangle,\quad
   \mathcal W_-=x^{-1}\bigl\langle\tfrac12E^{(y)},(I-\mathcal T)^{-1}D_e\lambda^{(y)}\bigr\rangle,\quad
   \mathcal W_++\mathcal W_-=\frac1{2x}\bigl\langle\lambda^{(y)},(I-\mathcal T)^{-1}D_e\lambda^{(y)}\bigr\rangle,
\end{equation}
where $\mathcal T=D_eK$ is the travel operator of Definition~\ref{def:model},
$D_e=\operatorname{diag}(e_s)$ with $e_s=2q^{1+s}$, and $\mu^{(y)},E^{(y)},\lambda^{(y)}$ are as in
\eqref{eq:junctionsym}--\eqref{eq:farjunction}. The sum over the marker data
$(\varepsilon^\ast,\delta^\ast)$ is already carried out in these vectors: $\delta^\ast=1$ at the far
site gives $2\mu$ and $\delta^\ast=0$ gives $E$ when $k^\ast>0$, and the roles are exchanged when
$k^\ast<0$. The $k^\ast=0$ sector is
$\mathcal W_0=A_0+B_0A_1+B_0^2A_2$, with $A_i$ built from bulk quantities alone; it contains no
travel resolvent.
\end{theorem}

\begin{proof}[Proof and standing]
Let $(k^\ast,\varepsilon^\ast,\delta^\ast,d)$ parametrise the elements of $W_{\mathrm{univ}}$ as
in \cite[Theorem ``Normal form and lamp model'', \textup{thm:normal-form}]{Bonfioli2026a}. Two facts
are formalised in Lean~4 with the standard axioms only:
\begin{itemize}
\item the affine realisation intertwines the three generators and is injective, and every element
of the model is reachable, so the word length of the model is the Cayley length in
$W_{\mathrm{univ}}$ (\texttt{RJPhi.lean}: \texttt{intertwine\_s1}, \texttt{intertwine\_s2},
\texttt{intertwine\_s3}, \texttt{injective}; reachability is \texttt{reaches\_of\_phiZ} in
\texttt{PhiLipschitz.lean});
\item $\mathrm{wordLength}=\ell_R+2c$ for every element, with $\ell_R$ and $c$ the site-local
closed forms \texttt{lRTrue}, \texttt{cTrue} of \texttt{CorrectedSpan.lean}
(\texttt{RJMetricAll.lean}, theorem \texttt{metricAll}).
\end{itemize}
Given these, $\ell_R$ and $c$ are sums of site and edge costs whose only non-bulk sites are $0$ and
$k^\ast$ (Corollaries~\ref{cor:localcost}, \ref{cor:marker}; a gap run of length $L$ contributes
$q^Ly^{L-1}$ and a travel-interior site is never a cut). Summing these site costs chain by chain
gives \eqref{eq:gapkernel} with source $u$, the vectors
\eqref{eq:junctionsym}--\eqref{eq:farjunction}, and the three sectors of \eqref{eq:assembly}; this
is Theorem~\ref{thm:M3prime}, proved in Appendix~\ref{app:M3prime}, and the regrouping of
$\mathcal W_0$ is Remark~\ref{M3:rem:W0regroup}. That last step is a hand proof, with two
independent adversarial reviews, and it is not formalised. Every identity holds in
$\Z[y][[x]]$, and the series agree with enumeration as recorded in \S\ref{sec:bridge}.
\end{proof}

\noindent Earlier versions carried \textup{(M)} as a hypothesis, with three named weak links: the
reverse shield inequality $c\le L-1$ (Remark~\ref{rem:shieldowes}), the metric formula, and
\textup{(M3)}. The first two are no longer separate inputs. The cut count $c$ that enters
$\mathcal W$ is the site-local count \texttt{cTrue}, and \texttt{metricAll} proves
$\mathrm{wordLength}=\ell_R+2c$ for it directly, so neither the reverse inequality nor the
optimisation of Definition~\ref{def:pairing} is consumed. One caveat: $\ell_R$ in $\mathcal W$ is the
closed-form relaxed length of \cite{Bonfioli2026a}. That it also equals the minimum over relaxed
realisations (\cite{Bonfioli2026a}, \textup{rem:metric-formula-status}) affects only the
\emph{interpretation} of $V$ as a relaxed count, not $U$ and not the series identities. The third,
\textup{(M3)}, is proved in Appendix~\ref{app:M3prime}, in the corrected form stated above. Nothing in \S\S\ref{sec:setup}--\ref{sec:numerators}, \S\ref{sec:sd},
\S\S\ref{sec:sect}--\ref{sec:sitecost} above, or Appendices~\ref{app:annulus}--\ref{app:star} uses
it: those speak only about the explicit series of \eqref{eq:krec}, about $P_{12}$, about the
operators just defined, and about the optimisation of Definition~\ref{def:pairing}.

$V=\mathcal W(x,1)$ and $U=\mathcal W(x,q)$ (\S\ref{sec:bridge}), the two evaluations differing
only through $\mathfrak g_V=\tfrac q{1-q}$, $\mathfrak g_U=\tfrac q{1-q^2}$ and the slot $B_{0z}$;
and $\lambda^{(y)},\mu^{(y)},\mathcal W_0$ are holomorphic at a travel pole $q_m$, by
\eqref{eq:smresolvent}, as soon as $S_e(q_m)\ne0$ and $1-\mathfrak gt_1\ne0$ there, which is
Proposition~\ref{prop:selower} and Corollary~\ref{app:gate} ($|\mathfrak g_Vt_1|\le0.983$, and
$|\mathfrak g_Ut_1|<|\mathfrak g_Vt_1|$). In particular $\mathcal W_0$ has no travel pole: by
Remark~\ref{M3:rem:W0regroup} it is a polynomial in $B_0$ with coefficients built from the bulk
vector $P$ and monomials, so its only possible singularities are at $S_e=0$ and at
$1-\mathfrak gt_1=0$.

The growth series are defined as power series, with radius of convergence
$x_1=\sqrt{q_1}=1/\beta_2=0.6704\ldots$ (Corollary~\ref{cor:beta}), and the local forms below are
statements about functions on a larger disc. The following lemma supplies that continuation and
identifies it with the series.

\begin{lemma}[meromorphic continuation]\label{lem:mero}
Let $y$ be either a complex constant with $|y|\le1$ or $y=x^2$. Then the specialisation
$x\mapsto\mathcal W(x,y)$ of the series of Theorem~\ref{thm:M3prime} converges near $x=0$ and is
the Taylor series at $0$ of a single-valued meromorphic function on the disc $|x|<1$, given there
by \eqref{eq:assembly} and $\mathcal W_0=A_0+B_0A_1+B_0^2A_2$ read as analytic expressions. Its
poles in $0<|x|<1$ lie over points $q=x^2$ at which $1-\Sigma_1(q)=0$, $S_e(q)=0$ or
$1-\mathfrak g(q,y)\,t_1(q)=0$. In particular $V=\mathcal W(x,1)$ and $U=\mathcal W(x,x^2)$ are
meromorphic on $|x|<1$, and $G^T_0=\Sigma_0/(1-\Sigma_1)$ is meromorphic on $|q|<1$.
\textup{(}The continuation concerns the explicit operators of \S\S\ref{sec:sect}--\ref{sec:mob}; its
identification with $U$ and $V$ is Theorem~\ref{thm:model}.\textup{)}
\end{lemma}

\begin{proof}
Put $q=x^2$, so $|q|<1$; $\mathfrak g=q/(1-qy)$ is holomorphic on the disc, since $|qy|<1$ in both
cases.

\emph{Operators.} $q\mapsto\mathcal T(q)$ and $q\mapsto M_0(q)$ are holomorphic on $|q|<1$ as maps
into the bounded operators on $\ell^1$, in operator norm: the row truncations $M_N$ in the proof of
Lemma~\ref{lem:fred} have finitely many non-zero rows, each a norm-convergent $\ell^\infty$-valued
power series in $q$ (the row $s$ is $e_s\sum_{s'}q^{\max(s,s')}\delta_{s'}$), hence are
holomorphic in norm, and $\|M-M_N\|\le\sum_{s>N}2|q|^{\vartheta+s}\to0$ uniformly on $|q|\le r<1$; a locally
uniform norm limit of holomorphic operator-valued maps is holomorphic. With compactness
(Lemma~\ref{lem:fred}) and invertibility near $q=0$ (Lemma~\ref{lem:modelexist}), the analytic
Fredholm theorem makes $(I-\mathcal T)^{-1}$ and $R_0=(I-M_0)^{-1}$ meromorphic on $|q|<1$ in
operator norm, with poles only where $1-\Sigma_1=0$, resp.\ $S_e=0$ (Corollary~\ref{cor:sing}).

\emph{Bulk quantities.} $u=(2q^{2b})_b$ is a norm-convergent, hence holomorphic, $\ell^1$-valued
function, so $\Psi=R_0u$ is meromorphic $\ell^1$-valued and $t_1=v^{\!\top}\Psi$ is meromorphic.
$1-\mathfrak gt_1$ tends to $1$ as $x\to0$, so it is not identically zero, and $B_0=1/(1-\mathfrak
gt_1)$, $P=B_0\Psi$, $\mathfrak gS=B_0-1$ and $B_{0z}=1+y(B_0-1)$ are meromorphic
(Remark~\ref{M3:rem:W0regroup}).

\emph{Boundary vectors and the sectors $k^\ast\ne0$.} Every entry of
\eqref{eq:app-mu}--\eqref{eq:app-lambda} is a combination of $B_0$, $B_{0z}$ and sums
$\sum_bP_bx^{m_b}$ with $m_b\ge0$, so $|\mu^{(y)}_s|,|E^{(y)}_s|,|\lambda^{(y)}_s|\le
C\,(|B_0|+|B_{0z}|+\|P\|_{\ell^1})$ uniformly in $s$, locally uniformly in $x$ away from the poles
of $B_0$ and $\Psi$. Hence $D_e\lambda^{(y)}$ has $\ell^1$-tails
$\sum_{s>N}|e_s\lambda^{(y)}_s|=O(|q|^{N+2})$, locally uniformly, and is meromorphic
$\ell^1$-valued; so is $(I-\mathcal T)^{-1}D_e\lambda^{(y)}$. Its pairing with the bounded vectors
$\mu^{(y)}$ and $\tfrac12E^{(y)}$ is an absolutely convergent sum, locally uniformly convergent
away from the poles (an $\ell^1$-valued continuous map has uniformly small tails on compact sets),
hence meromorphic. This gives $\mathcal W_\pm$; the prefactor $x^{-1}$ is at most a pole at $x=0$,
and is removed below.

\emph{The sector $k^\ast=0$.} By Remark~\ref{M3:rem:W0regroup},
$\mathcal W_0=A_0+B_0A_1+B_0^2A_2$. The $A_i$ are not finite sums: each is an absolutely convergent
double sum, over the left and the right chain kinds $\mathsf E,\mathsf Z,\mathsf N_{\pm,b}$ of
\eqref{eq:app-W0} ($b\ge1$), of monomials $x^jy^\chi$ with $j\ge0$, $\chi\in\{0,1\}$, times at most
two entries of $\Psi$. Since $|x^jy^\chi|\le1$, the double sum is dominated by a constant times
$(1+\|\Psi\|_{\ell^1})^2$, converges locally uniformly where $\Psi$ is holomorphic, and is
meromorphic.

\emph{Identification with the $x$-adic series.} Choose $r_0>0$ so small that on $|x|<r_0$ the
operators $\mathcal T$, $M_0$ and $M_0+\mathfrak guv^{\!\top}$ have $\ell^1$-norm below $\tfrac12$
(Lemma~\ref{lem:modelexist}; the rank-one term has norm at most $|\mathfrak g|\,\|u\|_{\ell^1}$) and
$|\mathfrak gt_1|<\tfrac12$. There every object above is the sum
of a norm-convergent Neumann series ($(I-\mathcal T)^{-1}$, $R_0$, $P=(I-M_0-\mathfrak
guv^{\!\top})^{-1}u$, and $B_0=\sum_n(\mathfrak gt_1)^n$), whose terms are holomorphic on
$|x|<r_0$ and whose Taylor series at $0$ are the terms of the $x$-adic Neumann series of
Lemma~\ref{M3:lem:conv}, of $x$-adic order tending to infinity. If $\sum_nh_n$ converges uniformly
on a disc about $0$, the Taylor coefficients of the sum are the sums of the Taylor coefficients of
the $h_n$ (Weierstrass); when $\operatorname{ord}h_n\to\infty$ each of these sums is finite and is
the $x$-adic sum. Applying this to each Neumann series, to the absolutely convergent pairings and
to the double sums $A_i$, the Taylor series at $0$ of the analytic $\mathcal W(x,y)$ is the
$x$-adic series of Theorem~\ref{thm:M3prime}, specialised at $y$ (the specialisation
$\Z[y][[x]]\to\C[[x]]$ at a constant $y$, or at $y=x^2$, is $x$-adically continuous and commutes
with the sums). In particular that function is holomorphic at $x=0$. At $y=1$ the series is $V$,
and at $y=x^2$ it is $U$ (Corollary~\ref{cor:M3prime-growth}). A power series that agrees near $0$
with a function meromorphic on the disc is continued by it, uniquely. For $G^T_0$, $\Sigma_0$ and
$1-\Sigma_1\not\equiv0$ are holomorphic on $|q|<1$ (Theorem~\ref{thm:blocks}, first clause).
\end{proof}

\begin{proposition}[the shape of the local form]\label{prop:shape}
Let $q_m$ be a travel pole, let $p\ge1$ be the order of the zero of $1-\Sigma_1$ there
\textup{(}finite, Lemma~\ref{lem:polefinite}\textup{)}, and let $\pi$ be the order of the pole of
$(I-\mathcal T)^{-1}$ there \textup{(}finite, Lemma~\ref{lem:fred}\textup{)}. Then $\pi\ge p\ge1$
and, for each $y$,
\[
   \mathcal W(x,y)=\frac{\mathcal N_y}{(q-q_m)^{\pi}}+\bigl(\text{poles of order}<\pi\bigr)+H_m ,
   \qquad H_m\ \text{holomorphic at }q_m,
\]
with the leading numerator a perfect square
\begin{equation}\label{eq:junctionpairing}
   \mathcal N_y=\frac{c_\star}{2x_m}\,\Pi_y(q_m)^2,\qquad
   \Pi_y(q_m):=\langle\lambda^{(y)},R\rangle ,\qquad c_\star\ne0,\ x_m=\sqrt{q_m},
\end{equation}
where $R$ spans $\ker(I-\mathcal T(q_m))$, normalised by its sections,
\[
   \sum_{s\ge0}q^{ks}R_s=\Sigma_k(q_m)\quad(k\ge0),\qquad\text{in particular }
   \sum_sR_s=\Sigma_0(q_m),\quad\sum_sq^sR_s=\Sigma_1(q_m)=1 ,
\]
and $L$ spans $\ker(I-\mathcal T(q_m)^{\!\top})$, explicitly $L_s=R_s/e_s$. There are no cross
terms between the sectors. In particular $\Pi_y(q_m)\ne0$ implies that $q_m$ is a pole of $\mathcal W(\,\cdot\,,y)$, of order exactly
$\pi$.
\end{proposition}

\begin{proof}
Write $\mathcal R(q)=(I-\mathcal T(q))^{-1}=\sum_{j=1}^{\pi}A_j(q-q_m)^{-j}+\text{hol}$, with
$A_\pi\ne0$. Multiplying $(I-\mathcal T(q))\mathcal R(q)=I$ by $(q-q_m)^{\pi}$ and letting
$q\to q_m$ gives $(I-\mathcal T(q_m))A_{\pi}=0$, so
$\operatorname{ran}A_{\pi}\subseteq\ker(I-\mathcal T(q_m))$, which is one-dimensional by
Lemma~\ref{lem:nullspace}; the same step on $\mathcal R(q)(I-\mathcal T(q))=I$ gives
$A_{\pi}(I-\mathcal T(q_m))=0$. Hence $A_\pi x=c_\star\langle L,x\rangle R$ with $c_\star\ne0$,
$R$ spanning the kernel and $L$ the cokernel. The normalisation of $R$ is
Lemma~\ref{lem:nullspace}: $G^R_k=G^R_1\Sigma_k$ and $\Sigma_1(q_m)=1$, so $G^R_1=1$ gives
$G^R_k=\Sigma_k(q_m)$. That $L_s=R_s/e_s$ spans the cokernel follows from
$\mathcal T=D_eK$ with $D_e=\operatorname{diag}(e_s)$ and $K$ symmetric: from $D_eKR=R$ one gets
$KD_e(D_e^{-1}R)=KR=D_e^{-1}R$, that is $\mathcal T^{\!\top}L=L$, and $L\in\ell^\infty$ because
$|R_s|\le|e_s|\,\|R\|_{\ell^1}$; the cokernel is one-dimensional by Fredholm. Substituting the
Laurent expansion into $\mathcal W_++\mathcal W_-$ of \eqref{eq:assembly}, and using that
$\lambda^{(y)},\mathcal W_0$ are holomorphic at $q_m$, the leading coefficient is
$(2x_m)^{-1}\langle\lambda,A_\pi D_e\lambda\rangle=(2x_m)^{-1}c_\star\langle L,D_e\lambda\rangle
\langle\lambda,R\rangle=(2x_m)^{-1}c_\star\langle\lambda,R\rangle^2$, since $D_eL=R$. It is
non-zero exactly when $\Pi_y(q_m)\ne0$, and then $q_m$ is a pole of order $\pi$. Finally
$\langle\mathbf 1,\mathcal RE\rangle=\Sigma_0/(1-\Sigma_1)$ (Proposition~\ref{prop:travelsum}) has
a pole of order exactly $p$ at $q_m$, since $\Sigma_0(q_m)\ne0$
(Proposition~\ref{prop:numnonvanish}), while a matrix element of $\mathcal R$ has pole order at
most $\pi$; hence $\pi\ge p$.
\end{proof}

\begin{proposition}[the marker junction is not rank one]\label{prop:junction}
The junction matrix $J[\sigma,s]$ of \eqref{eq:junctionsym} has rank at least $3$. Already the
one-sign matrix $\bigl[x^{\max(\sigma-1,2s+1)}\bigr]$ has rank at least $3$: on
$\sigma\in\{0,4,8\}$, $2s+1\in\{1,5,9\}$ it is
\[
   \begin{pmatrix}x&x^5&x^9\\ x^3&x^5&x^9\\ x^7&x^7&x^9\end{pmatrix},\qquad
   \det=x^{9}\,(x-x^{3})(x^{5}-x^{7})\ \ne\ 0\quad(0<x<1),
\]
its rank on $\sigma\in\{0,2,4,6,8\}$, $2s+1\in\{1,3,5,7,9\}$ is $4$, and the symmetrised matrix
\eqref{eq:junctionsym} has rank $6$ on $\sigma\in\{0,2,4,6,8,10\}$, $2s+1\in\{1,3,5,7,9,11\}$. By
contrast the gap-bridge term $q^{a+b}$ of \eqref{eq:gapkernel} is rank one, which is why
\eqref{eq:mobius} factorises. Consequently $\mu^{(y)}$ is not a function of the total bulk sum
$B(q,y)$ alone, and $\Pi_y(q_m)\ne0$ does not follow from $\Sigma_0(q_m)\ne0$ together with
$B(q_m,y)\ne0$. At a travel pole, however, the pairing against the null vector $R$ collapses to a
function of $t_1$ alone (Theorem~\ref{thm:RJ}).
\end{proposition}

\begin{proof}
The displayed block is the case $\sigma-1\in\{-1,3,7\}$, $2s+1\in\{1,5,9\}$, which interleaves as
$-1<1<3<5<7<9$; subtracting the second row from the first gives $(x-x^3,0,0)$, and expanding along
it leaves $(x-x^3)(x^5x^9-x^9x^7)=x^9(x-x^3)(x^5-x^7)$. The two rank statements are exact Gaussian
elimination over $\Q$, at $q=\tfrac13$ and at $x=\tfrac12$ respectively
(\texttt{assembly\_certificate.py}, Part~4). For the second-last sentence: by \textup{(M3)}, $\mu^{(y)}$
is the image of $(B_\sigma)_\sigma$ under the junction matrix, whose rank is at least $3$.
\end{proof}

\begin{theorem}[the bulk dictionary at the travel poles]\label{thm:L}
At every travel pole $q_m$ the factorisation \eqref{eq:mobius} holds with
$\kappa=t_0b_1-b_0t_1$ and $t_0=b_1$, and the dictionary \eqref{eq:dict} holds identically in
$0<q<1$, in the completed form of Corollary~\ref{cor:dict}. Consequently the bulk quantities
$B_0=1/(1-\mathfrak gt_1)$ and $B=b_1/(1-\mathfrak gt_1)$ of Remark~\ref{rem:bulksource} are finite at
every travel pole, for $y=1$ and $y=q$.
\end{theorem}

\begin{proof}
Proposition~\ref{prop:mobius}, whose two hypotheses are supplied by Corollary~\ref{cor:sing} with
Proposition~\ref{prop:selower} and by Remark~\ref{rem:mobhyp}; Proposition~\ref{prop:recip};
Corollary~\ref{cor:dict}. Finiteness is $|\mathfrak g_Vt_1|\le0.983$ (Corollary~\ref{app:gate}) and
$\mathfrak g_U<\mathfrak g_V$.
\end{proof}

\begin{remark}[what \textup{(L)} consumed]\label{rem:Lgap}
That the bulk transfer is the operator of \eqref{eq:rankone}, that is, \textup{(M1)} and
\textup{(M2)}. Of these, \textup{(M1)} is now Corollary~\ref{cor:localcost}, the exponent
$q^{\max(s,s')}$ of the bulk kernel being $x^{\max(|d|,|d'|)}$ at even deposits; and of
\textup{(M2)} the cost $q^L$ of a gap run and the inequality $c\ge L-1$ are
Corollary~\ref{cor:lRclosed} and Proposition~\ref{prop:cut}, and $c\le L-1$ follows from the metric
identity \texttt{metricAll} with the site-local cut count (Theorem~\ref{thm:model}), which no
longer consumes the optimisation of Definition~\ref{def:pairing}. Everything the paper previously listed as
unproved \emph{downstream} of that, namely the factorisation \eqref{eq:mobius}, the reciprocity
$t_0=b_1$, the entry $t_1=P_{12}/S_e$, and the resolvent lines of \eqref{eq:dict}, is proved above.
The residue at $q_m$ is governed not by the bulk sum but by the pairing $\Pi_y(q_m)$
(Proposition~\ref{prop:shape}), which Theorem~\ref{thm:RJ} evaluates.
\end{remark}

\subsection{The junction pairing \textup{(R-J)}}\label{sec:RJ}

\textup{(R-J)} is the statement that $\Pi_y(q_m)=\langle\lambda^{(y)},R\rangle\ne0$ at infinitely
many travel poles, for $y=1$ and $y=q$. Earlier versions carried it as a hypothesis. It is now a
theorem.

\begin{theorem}[closed form of the junction pairing]\label{thm:RJ}
Let $q_m$ be a travel pole, $x=\sqrt{q_m}$, and $R$ normalised as in
Proposition~\ref{prop:shape}. Then
\begin{align}
   (1-\mathfrak g_Vt_1)\,\Pi_1(q_m)
     &=\frac{2q(1+x)}{1-q}\Bigl(t_1+\frac{1+x}{2x}\Bigr),\label{eq:RJclosed1}\\
   \Pi_q(q_m)
     &=\frac{2q\,B_0(1+x)}{1-x}\Bigl(t_1\frac{1-x+q}{1+q}+\frac1{2x}\Bigr),
     \qquad B_0=\frac1{1-\mathfrak g_Ut_1}. \label{eq:RJclosedq}
\end{align}
Consequently $\Pi_1(q_m)>0$ and $\Pi_q(q_m)>0$ at every travel pole $q_m$.
\end{theorem}

\begin{proof}
Throughout $q=q_m$.

\emph{Reduction to the ungapped bulk vector.} Put $\Psi_0=1$ and let $\Psi=(\Psi_b)_{b\ge1}=R_0u$ be
the solution of the bulk system with $\mathfrak g=0$ and source $u$, that is
\[
   \Psi_b=2q^bF_b,\qquad F_t:=\sum_{b\ge0}q^{\max(b,t)}\Psi_b\qquad(b\ge1),
\]
the term $b=0$ of $F_t$ being the source $q^t$. Then $t_1=v^{\!\top}R_0u=\sum_{b\ge1}q^b\Psi_b$ and,
by \eqref{eq:smresolvent}, $P=R_0u/(1-\mathfrak gt_1)=B_0\Psi$, where $B_0=1/(1-\mathfrak gt_1)$
(Remark~\ref{rem:bulksource}). So every bulk weight in \eqref{eq:junctionsym}--\eqref{eq:farjunction}
is $B_0\Psi_b$, with $b=0$ standing for the empty and gap-opening chains. In the variable $q=x^2$,
\[
   x^{\max(2s,2b)}=q^{\max(s,b)},\quad x^{\max(2s+2,2b)}=q\,q^{\max(s,b-1)},\quad
   x^{\max(2s+1,2b-1)}=x\,q^{\max(s,b-1)},\quad x^{\max(2s+1,2b+1)}=x\,q^{\max(s,b)},
\]
with $\max(s,-1)=s$. At $y=1$ the slot $B_{0z}$ equals $B_0$, and \eqref{eq:junctionsym},
\eqref{eq:farjunction} combine into
\begin{equation}\label{eq:lambdaPsi}
   \lambda^{(1)}_s=B_0\sum_{b\ge0}\Psi_b\Bigl((1+x)\,q^{\max(s,b)}+(q+x)\,q^{\max(s,b-1)}\Bigr),
\end{equation}
hence $\langle\lambda^{(1)},R\rangle=B_0\bigl((1+x)X+(q+x)Y\bigr)$ with
\[
   X=\sum_{b\ge0}\Psi_b\,(KR)_b,\qquad Y=\sum_{b\ge0}\Psi_b\,G_b,\qquad
   (KR)_b=\sum_{t\ge0}q^{\max(b,t)}R_t,\quad G_b=\sum_{s\ge0}q^{\max(s,b-1)}R_s .
\]
All double sums converge absolutely ($\Psi\in\ell^1$ with $|\Psi_b|=O(q^{2b}|\Psi|_{\ell^1})$,
$R\in\ell^1$), so they may be summed in either order.

\emph{Two facts about $R$.} From $\mathcal TR=R$, $R_s=2q^{s+1}(KR)_s$ for every $s\ge0$; and the
normalisation $\sum_sq^sR_s=\Sigma_1(q_m)=1$ of Proposition~\ref{prop:shape} says $(KR)_0=G_0=1$. In
particular $R_0=2q$.

\emph{The $X$ shift identity.} Summing first over $t$ and splitting off $b=0$,
$X=1+S$ with $S=\sum_{b\ge1}\Psi_b(KR)_b$. Summing first over $b$ instead, $X=\sum_tR_tF_t$. The
term $t=0$ is $R_0F_0=2q(1+t_1)$, since $F_0=\sum_bq^b\Psi_b=1+t_1$; for $t\ge1$,
$R_tF_t=2q^{t+1}(KR)_t\cdot\Psi_t/(2q^t)=q\,\Psi_t(KR)_t$. Hence $X=2q(1+t_1)+qS$, and eliminating
$S$,
\[
   X=1+\frac{X-(1+t_1)R_0}{q},\qquad\text{so}\qquad X=\frac{q(1+2t_1)}{1-q}.
\]

\emph{The $Y$ shift identity.} Splitting off $b=0$ and using $G_b=(KR)_{b-1}=R_{b-1}/(2q^b)$ for
$b\ge1$, $Y=1+\sum_{s\ge0}\Psi_{s+1}R_s/(2q^{s+1})$. Summing first over $b$ instead, and using
$q^{\max(s,b-1)}=q^{-1}q^{\max(s+1,b)}$, $Y=\sum_sR_sF_{s+1}/q=\sum_sR_s\Psi_{s+1}/(2q^{s+2})$. The
two expressions give
\[
   Y=1+qY,\qquad\text{so}\qquad Y=\frac1{1-q}.
\]

\emph{The closed forms.} With $q=x^2$,
$(1+x)X+(q+x)Y=\bigl((1+x)q(1+2t_1)+q+x\bigr)/(1-q)$, and since
$(1+x)q+q+x=x(1+x)^2$ this is \eqref{eq:RJclosed1}. At $y=q$ only the slot $s=0$ of $E^{(q)}$
changes, from $B_0$ to $B_{0z}=1+q\,\mathfrak g_US=B_0\bigl(1-t_1q/(1+q)\bigr)$ (using
$\mathfrak g_US=B_0-1=\mathfrak g_Ut_1B_0$). So
$\Pi_q=B_0\bigl((1+x)X+(q+x)Y\bigr)+(B_{0z}-B_0)R_0$, now with $B_0=1/(1-\mathfrak g_Ut_1)$, and
collecting terms with $(1+x)^2(1-x+q)=(1+x)(1+q)-q(1-q)$ (both equal $1+x+x^3+x^4$) gives
\eqref{eq:RJclosedq}.

\emph{Positivity for $\tau_m\le5\cdot10^{-3}$.} By Corollary~\ref{app:gate}, $|\mathfrak g_Vt_1|\le0.983$,
and $\mathfrak g_U<\mathfrak g_V$. Hence $1-\mathfrak gt_1\ge0.017$ for both $y$, and
$q|t_1|\le0.983(1-q)$, so $|t_1|<0.12$ once $q\ge0.9$ (here $q\ge e^{-0.005}$). Since
$(1+x)/(2x)\ge1$, $1/(2x)>1/2$ and $0<(1-x+q)/(1+q)<1$, both brackets are positive, and so are the
prefactors.

\emph{Positivity at the poles above that threshold.} This part is a computer-assisted certificate:
at $q_1,\dots,q_{12}$, $\Pi_1$ and $\Pi_q$ are enclosed in MPFR interval arithmetic and found positive,
together with $1-\mathfrak gt_1>0$ (\texttt{code/zeta\_probe/rj\_certificates/r35rust/src/bin/r54b.rs}). The pole
index comes from \texttt{code/zeta\_probe/rj\_certificates/r39pole}, an interval-arithmetic zero count
of $1-\Sigma_1$ which isolates exactly $13$ zeros on $[0,0.9988]$ \textup{(}the right endpoint
taken as its binary64 value\textup{)}, each in its own cell. Its tiling of the interval was repaired
after a first review found gaps of order $10^{-77}$ between cells, and the repaired code was rebuilt
and re-run from current source by a second reviewer; the trusted base is MPFR and the tool's own code.
\end{proof}

\noindent Only infinitely many poles are needed for Theorems~\ref{thm:V} and~\ref{thm:U}, and the
analytic part ($\tau_m\le5\cdot10^{-3}$) supplies them; the certificate is needed only for the
statement ``at every travel pole''. The algebraic core is formalised in Lean~4
(\S\ref{sec:RJlean}); the eigen-relations of $R$, the gate and the analytic estimates behind it enter
the formal statements as hypotheses and are not formalised.

\begin{corollary}[\textup{(R-J)} at the dominant pole]\label{prop:RJdominant}
$\Pi_1(q_1)>0$ and $\Pi_q(q_1)>0$; this is the case $m=1$ of Theorem~\ref{thm:RJ}.
\end{corollary}

\subsection{Formalisation of the junction pairing and of the model}\label{sec:RJlean}

The following files in \texttt{lean/with\_mathlib} are checked by Lean~4 against Mathlib
\cite{MathlibCommunity}. Each carries \texttt{\#print axioms} checks on its main statements, and the
recorded output lists only the standard axioms; none uses \texttt{sorry} or
\texttt{native\_decide}.
\begin{itemize}
\item \texttt{RJMain.lean} and \texttt{RJShift.lean}: the $X$ and $Y$ shift identities
(\texttt{X\_shift}, \texttt{Y\_shift}) and their solutions (\texttt{X\_closed}, \texttt{Y\_closed}),
from the eigen-recursion of $R$, the bulk recursion of $\Psi$ and the normalisation, including the
interchange of the double sums from $\ell^1$ summability (\texttt{summable\_swap}).
\item \texttt{RJClosedForm.lean}: the algebra of \eqref{eq:RJclosed1}--\eqref{eq:RJclosedq}
(\texttt{closed\_form}, \texttt{closed\_form\_yq}) and the positivity step from the gate bound
(\texttt{gate\_bound}, \texttt{bracket\_pos}, \texttt{bracket\_yq\_pos}).
\item \texttt{RJIdentities.lean}: auxiliary finite identities for the travel recursion.
\item \texttt{RJLemmaJ.lean}: the cut criterion at the junction sites used in
Appendix~\ref{app:M3prime}.
\item \texttt{RJPhi.lean}: the faithfulness of the model, i.e.\ the affine realisation intertwines
the three generators and is injective.
\item \texttt{RJMetricAll.lean}: \texttt{metricAll}, $\mathrm{wordLength}=\ell_R+2c$ for every
element, with $\ell_R,c$ the closed forms \texttt{lRTrue}, \texttt{cTrue}.
\end{itemize}
What is \emph{not} formalised: the assembly \textup{(M3)} of Appendix~\ref{app:M3prime}; the gate
(Appendix~\ref{app:star}) and every analytic estimate behind it; the identification of the Lean
eigen-recursion hypotheses with the operators $\mathcal T$ and $M_0$ at an actual travel pole; and
the interval certificates. The formal statements take these as hypotheses.

\subsection{Certificate}\label{sec:assemblycert}

\begin{remark}[the computational certificate]\label{rem:assemblycert}
\texttt{code/zeta\_probe/assembly\_certificate.py} re-derives the statements of this section in
two independent arithmetic models and exits $0$ with every check passing. Numerical agreement is
evidence and is not a proof; the proofs above are what carry the statements. What the certificate
adds is an independent arithmetic model and the hypothesis deletions.

\emph{Part 1, exact arithmetic in $\Z[[q]]/(q^{N+1})$ at the two truncation orders $N=28$ and
$N=40$, no floating point}: Propositions~\ref{prop:travelsum}, \ref{prop:bulkdress} and
\ref{prop:recip} and Corollaries~\ref{cor:cocycle} and \ref{cor:dict}, coefficient by
coefficient; the bulk block $b_0$ has coefficients
$2,2,6,2,18,6,42,18,118,50,282,190,706,594$ and the travel block the travel-run counts
$2,6,10,26,54,114,274,582,1298$; $P_{11}=1+T_0$, $P_{12}=T_1$, $P_{21}=-S_0$, $P_{22}=S_e=1-S_1$,
$P_{11}S_e+P_{12}S_0=1$, the Hahn--Exton form \eqref{eq:p12closed}, and $S_o=\frac{1-q}{2q}S_0$.

\emph{Part 2, exact arithmetic in $\Q[\mathfrak g][[q]]/(q^{19})$}: the factorisation of
Proposition~\ref{prop:mobius} holds identically in the free parameter $\mathfrak g$, and three
hypothesis deletions each break it, with the first surviving order recorded. A rank-two gap term
fails at $q^{6}$; a gap term $q^{a+2b}$ in place of $q^{a+b}$ fails at $q^{4}$; and the asymmetric
gapless kernel $q^{\max(a,b)}+q^{a+2b}$ breaks $t_0=b_1$ at $q^{8}$.

\emph{Part 3, two precisions, $40$ and $70$ significant decimal digits}: at the first ten travel
poles, $q_1=0.4494536306$ through $q_{10}=0.997756895$, the dictionary $b_0=S_0/S_e$,
$t_0=b_1=S_1/S_e$, $t_1=P_{12}/S_e$ holds to relative error $2.5\cdot10^{-26}$ at $40$ digits and
$8.8\cdot10^{-57}$ at $70$, the loss of about $14$ digits being the cancellation in the
alternating ladders at the highest pole; the two precisions agree on the pole locations to
$1.6\cdot10^{-31}$. The resolvent is evaluated there by the $O(N)$ two-point solve of
Corollary~\ref{cor:cocycle}, that is by the analytic continuation, the defining Neumann series not
converging at those $q$. The M\"obius value is cross-checked against an independent dense solve of
the truncated system at $q=0.5,0.7,q_1$ and $y=1,q$, agreeing to $6.2\cdot10^{-28}$. The gate
signs $b_0>0$ and $s<1$ hold at all ten, with $b_0\tau\to2$ and $s\to\tfrac14$. (Earlier versions
also checked a ``lifting increment'' $(1-s)B_V+qb_0$; it refers to the $E$-sourced M\"obius form and
not to the model's bulk sum, and is withdrawn.)

\emph{Part 4, exact rational arithmetic}: the one-sign junction matrix
$[x^{\max(\sigma-1,2s+1)}]$ has rank $4$ over $\Q$ at $q=\tfrac13$ on the $5\times5$ block, its
interleaved $3\times3$ minor equals $9/524288$ at $x=\tfrac12$, the symmetrised junction matrix
\eqref{eq:junctionsym} has rank $6$ on the $6\times6$ block at $x=\tfrac12$, and the gap-bridge
matrix $[q^{a+b}]$ has rank $1$.

\emph{Part 5, exact integer arithmetic, \texttt{code/zeta\_probe/tools/sitecost}}: the site cost is
computed by two independent exact solvers, a min-cost flow and a transportation dynamic program,
which agree on all $1\,012\,664$ supply-demand pairs with entries at most $7$; on the same range the
closed form $\max(|\alpha|,|\beta|,|\Phi|)$ of Lemma~\ref{lem:transport} is exact, with $0$ misses.
Corollary~\ref{cor:localcost} is confirmed on $4\,532\,157$ site configurations, each evaluated by
both solvers, namely every crossing count from the minimum to the minimum plus $10$, every sign
split, both signs of both deposits with $|d|\le24$, every travel indicator, all four marker data and
all eight site types
(bulk interior, travel interior in both directions, the near and far junctions for $k^\ast>0$ and
$k^\ast<0$, and the coincident junction at $k^\ast=0$): $0$ exceptions, with the cost independent of
the crossing counts and of the sign splits in every case. Deleting any one of the three pairing
costs is witnessed necessary: replacing the sign-flip bounce cost $2$ by $0,1,3,4$, the pass cost
$1$ by $0,2,3$, or the bounce cost $0$ by $1,2$ each breaks
Corollary~\ref{cor:localcost}, with $2544$ to $7392$ exceptions out of $7524$ configurations and the
smallest counterexample recorded. The marker closed forms of Corollary~\ref{cor:marker} hold with
$0$ exceptions on the full signed range, and the form $\max(|d_L|-1,|d_R|)$ fails on $42$ of $156$
cells at the near junction and on $36$ of $156$ at its $k^\ast<0$ counterpart, all with $d_L<0$.
Corollaries~\ref{cor:lRclosed} and the shield count are confirmed as recorded in
Remark~\ref{rem:shieldowes}.
\end{remark}

\section{The relaxed series \texorpdfstring{$V$}{V}}\label{sec:V}

The relaxed series is the travel resolvent dressed by the bulk. The resolvent itself is settled
unconditionally; passing from it to $V$ uses the transfer model (Theorem~\ref{thm:model}) and the
junction pairing (Theorem~\ref{thm:RJ}).

\begin{theorem}[the travel resolvent]\label{thm:travelres}
The travel resolvent $G^T_0=\Sigma_0/(1-\Sigma_1)$ is transcendental over $\Q(x)$. This is
unconditional: it uses no growth hypothesis, no asymptotic theorem, and no statement about the
strand-walk model. Its one numerical input is the contour bound of Lemma~\ref{lem:T2abs}, of which
it consumes only $|T_2(m\pi)|<1$, against a computed worst value $0.027$; that bound is a verified
enclosure over the parameter region, not a sampled one, and its arithmetic model is stated with the
lemma.
\end{theorem}

\begin{proof}
$\Sigma_0$ and $1-\Sigma_1$ are holomorphic on $|q|<1$ (Theorem~\ref{thm:blocks}, first clause) and
$1-\Sigma_1\not\equiv0$. The zero set of $1-\Sigma_1$ in $(0,1)$ is infinite and accumulates
exactly at $q=1$ (Lemma~\ref{lem:infpoles}, which consumes from Lemma~\ref{lem:T2abs} only
$|T_2|<1$, together with Lemma~\ref{lem:polefinite}). At each such $q_m$ the numerator does not
vanish, by Proposition~\ref{prop:numnonvanish}. Hence every $q_m$ is a pole of $G^T_0$, of order
the vanishing order of $1-\Sigma_1$ there, which is finite. So $G^T_0$ has infinitely many poles in
$|q|<1$ accumulating at $|q|=1$; under $q=x^2$ these are infinitely many poles in $|x|<1$
accumulating at $|x|=1$. An algebraic function over $\Q(x)$ has only finitely many poles in
$|x|<1$, so no algebraic function has such a pole set; hence $G^T_0$ is transcendental. It is not
$D$-finite either, in $q$ or in $x$, and $|q|=1$ is its natural boundary
(Theorems~\ref{thm:nonDfinite} and~\ref{thm:natboundary}): for $D$-finiteness the finiteness of the
pole set of an algebraic function is replaced by Lemma~\ref{lem:odeposes}, which confines the poles
in the disc of a solution of a linear differential equation to the finitely many zeros of its
leading coefficient.
\end{proof}

\begin{theorem}\label{thm:V}
The relaxed growth series $V$ is transcendental over $\Q(x)$.
\textup{(}Standing: Theorems~\ref{thm:model} and~\ref{thm:RJ}.\textup{)}
\end{theorem}

\begin{proof}
By Lemma~\ref{lem:mero}, $V$ is meromorphic on $|x|<1$, and the local forms below are those of
this continuation. By Proposition~\ref{prop:shape}, at a travel pole $q_m$ the relaxed series has the local form
$V=\mathcal N_1(q-q_m)^{-\pi}+(\text{poles of order}<\pi)+H_m$ with $H_m$ holomorphic, $\pi\ge1$,
and $\mathcal N_1=(2x_m)^{-1}c_\star\Pi_1(q_m)^2$ with $c_\star\ne0$. That proposition consumes two
non-vanishing statements, and they are separate: $S_e(q_m)\ne0$, in the quantitative form
$|S_e(q_m)|\ge0.35\sqrt{\tau_m}$ of Proposition~\ref{prop:selower}, which makes the bulk boundary
vectors $\lambda^{(1)},\mu^{(1)}$ holomorphic at $q_m$ and says nothing about vanishing; and
$\Sigma_0(q_m)\ne0$, from Proposition~\ref{prop:numnonvanish}, which pins the pole order
$\pi\ge p\ge1$; and $\Pi_1(q_m)\ne0$ at every travel pole, by Theorem~\ref{thm:RJ}.

So the leading Laurent coefficient is non-zero at every travel pole, and $V$ has a pole at each of
them. The travel poles are infinite in number and accumulate at $q=1$, that is at $|x|=1$
(Lemma~\ref{lem:infpoles}); this is the boundary of the disc $|x|<1$ of meromorphy, not of the disc
of convergence $|x|<x_1$, which only the dominant pole bounds. They are locally finite in $(0,1)$
(Lemma~\ref{lem:polefinite}), so an infinite subset of them accumulates at $q=1$ and nowhere else.
Hence $V$ has infinitely many poles in $|q|<1$ accumulating at $|q|=1$; under $q=x^2$ these are
infinitely many poles in $|x|<1$ accumulating at $|x|=1$. An algebraic function over $\Q(x)$ has
only finitely many poles in $|x|<1$, so an infinite pole set accumulating at the boundary is
impossible for one; hence $V$ is transcendental. No comparison of $T_2$ against $\cos W$ is made or
needed, and none is available at this strength, the proved bound on $|T_2|$ being
$0.17\,\tau^{1/4}$ (Lemma~\ref{lem:T2abs}), which exceeds $\sqrt{\tau/2}$ as $\tau\to0$;
Remark~\ref{lem:cos} sharpens the size of $T_2$ and is not used.
\end{proof}

\begin{remark}[what became of \textup{(R)}]\label{rem:Rstatus}
\emph{What \textup{(R)} was.} Route~B decomposes a relaxed element into two bulk runs, one travel
interval of length $|k|$ and four marker junctions; \textup{(R)} was that decomposition read as a
local form for $V$ at a travel pole, \emph{together with} a product form for its numerator,
$\mathcal N_V=\Sigma_0\cdot S_0/(1-S_1)\cdot J$ with $J$ a junction factor.

\emph{What is now proved.} The $k$-sum telescoping (Proposition~\ref{prop:travelsum}), the bulk
dressing (Proposition~\ref{prop:bulkdress}), and the shape of the local form, with a rank-one
numerator paired against the null vector of the travel transfer
(Proposition~\ref{prop:shape}). These are the whole of \textup{(R)} except the product form.

\emph{What is refuted.} The product form. The marker junction cost is
$\max(|d_L-1|,|d_R|)$ (Corollary~\ref{cor:marker}), a max-coupling of rank at least $3$ with an explicit non-zero $3\times3$
minor (Proposition~\ref{prop:junction}), so the numerator is not a multiple of the total bulk sum
and the step ``$\mathcal N_V(q_m)\ne0$ because $\Sigma_0(q_m)\ne0$ and $S_0(q_m)\ne0$'' is not
available. Proposition~\ref{prop:numnonvanish} is unaffected and is used, in
Theorem~\ref{thm:travelres} and inside Proposition~\ref{prop:shape}; what fails is the deduction
that leant on it here.

\emph{What replaces it.} The junction pairing \textup{(R-J)}, that $\Pi_1$ is non-zero at
infinitely many travel poles, which is strictly weaker than \textup{(R)} and is
Theorem~\ref{thm:RJ}: at a travel pole $\Pi_1$ collapses to an explicit function of $t_1$, and the
gate makes it positive.

\emph{Where it is used.} Only in Theorem~\ref{thm:V} and, at $y=q$, in Theorem~\ref{thm:U}.
Theorem~\ref{thm:travelres}, Propositions~\ref{prop:numnonvanish} and \ref{prop:selower},
\S\ref{sec:sd} and Appendices~\ref{app:annulus}--\ref{app:star} do not touch it. All of
\S\ref{sec:assembly} rests on the transfer model \textup{(M)}, Theorem~\ref{thm:model}.

\emph{Evidence.} The blocks reproduce their own combinatorial series exactly: $S_0/(1-S_1)$
against the bulk-run counts $2,2,6,2,18,6,42,18,118,50,282,190,706,594$ and
$\Sigma_0/(1-\Sigma_1)$ against the travel-run counts $2,6,10,26,54,114,274,582,1298$
(\texttt{blocks\_growth.py}, and \texttt{assembly\_certificate.py} Part~1 in exact integer
arithmetic at two truncation orders). $V$ itself is computed to $131$ exact terms by an
independent catalytic transfer, not from the blocks.
\end{remark}

\begin{corollary}\label{cor:beta}
The exponential growth rate of $v_n$ is $\beta_2=1.4916177871\ldots$, with
$\beta_2^2$ the reciprocal of the smallest travel pole $q_1$; the value $\tfrac32$ is excluded.
It needs only the location of the dominant travel pole: not Remark~\ref{lem:cos}, not
Appendix~\ref{app:annulus}, and not Appendix~\ref{app:star}. It does share with
Theorem~\ref{thm:V} the transfer model \textup{(M)}, which is what places $q_1$ on $V$ rather than
on $G^T_0$ alone, and at the dominant pole the junction pairing is positive
(Corollary~\ref{prop:RJdominant}). The
lower bound $\beta_2$ for $u_n$ and $v_n$ is free even of \textup{(M)}, coming from the
pure-travel subcount \cite{Bonfioli2026a}. The same rate for $u_n$
follows from Proposition~\ref{prop:travelinv} (the former hypothesis \textup{(T)} of
Theorem~\ref{thm:U}, now proved)
together with the positivity of $\Pi_q(q_1)$; $u_n\le v_n$ gives the upper bound
unconditionally.
\end{corollary}

Whether the \emph{number} $\beta_2$ is transcendental is a separate question, open, and with a
clean reduction to a value statement for a Hahn--Exton $q$-Bessel function; it is treated in
Appendix~\ref{app:beta2}, which is independent of everything else in this paper.

\section{The true series \texorpdfstring{$U$}{U}}\label{sec:U}

Transcendence of $U$ needs two things that the analysis of \S\ref{sec:sd} and
Appendix~\ref{app:star} does not supply: that $U$ has the travel poles at all, and that the residue
of $U$ at each of them is non-zero. The first was the combinatorial hypothesis \textup{(T)}; it is
now Proposition~\ref{prop:travelinv} and is proved (\S\ref{sec:inputT}). The second was previously
the hypothesis \textup{(L)}; its algebraic content is now a theorem (\S\ref{sec:mob},
Theorem~\ref{thm:L}), and the residue itself is the junction pairing \textup{(R-J)} at $y=q$,
Theorem~\ref{thm:RJ}, the same statement that Theorem~\ref{thm:V} consumes at $y=1$. This section
states \textup{(T)} in full, with its dependencies, with every place it is used, and with the proof
that closes it, and records what the bulk dictionary gives, before turning to the analytic content. A reader
should not have to open \cite{supplement} to see what Theorem~\ref{thm:U} rests on.

\subsection{The bivariate bridge}\label{sec:bridge}

With $c(g)\in\Z_{\ge0}$ the cycle count of \S\ref{sec:intro}, set
\[
   \mathcal W(x,y)=\sum_{n,c\ge0}w_{n,c}\,x^ny^c,\qquad
   w_{n,c}=\#\{g:\ell_R(g)=n,\ c(g)=c\},
\]
the grading being by \emph{relaxed} length. Then $V(x)=\mathcal W(x,1)$ and, since
$\ell_T=\ell_R+2c$ and $q=x^2$, $U(x)=\mathcal W(x,x^2)=\mathcal W(x,q)$: the two series are two
evaluations of one bivariate object, at $y=1$ and at $y=q$, and everything separating them is
carried by the marker $y$. The second identity is the one place where the metric identity
$\ell_T=\ell_R+2c$ is used. With $\ell_R,c$ the site-local closed forms of \cite{Bonfioli2026a},
$\mathrm{wordLength}=\ell_R+2c$ holds for every element (Lean, \texttt{metricAll} in
\texttt{RJMetricAll.lean}, resting on \texttt{wordLength\_eq\_lRTrue\_add\_two\_cTrue} and
\texttt{reaches\_of\_phiZ} in \texttt{PhiLipschitz.lean}), and the word length of the model is the
Cayley length in $W_{\mathrm{univ}}$ (Lean, \texttt{RJPhi.lean}); so $U=\mathcal W(x,q)$
unconditionally. That the closed form $\ell_R$ is also the minimum over relaxed realisations is
\cite{Bonfioli2026a}'s metric formula (its \textup{rem:metric-formula-status}); it bears on reading
$V$ as a relaxed count, and on nothing proved here.
The corrected assembly \eqref{eq:assembly} agrees with a breadth-first enumeration of
$W_{\mathrm{univ}}$ as exact series: at $y=q$ to $x^{26}$ in all three sectors; at $y=q^2$ to
$x^{22}$; and at symbolic $y$ in all $816$ coefficients with $\ell_T\le31$ ($5.03\cdot10^6$
elements), where a deliberately corrupted model is detected. $v_0,\dots,v_{19}$ are certified by the
same enumeration together with the bound $\ell_R\ge2c+6$. Earlier versions cited agreement to
$v_{14}$; that check (\texttt{cycle\_weighted\_gf.py}) was a sweep of the bridge
$\ell_T=\ell_R+2c$ and did not assemble the blocks, and the model as then printed fails it.

\subsection{The former input \textup{(T)}: the travel block carries no marker}\label{sec:inputT}

\textup{(T)} is Proposition~\ref{prop:travelinv}: the denominator $1-\Sigma_1$ is $y$-free, hence
literally the same series in $\mathcal W(x,1)$ and in $\mathcal W(x,q)$, so $U$ has the travel
poles $q_m$ of \S\ref{sec:setup}.

\emph{Where it is used.} In the proof of Theorem~\ref{thm:U}, to transport the pole set of
Lemma~\ref{lem:infpoles} from $V$ to $U$, and in the same way in Corollary~\ref{cor:beta} for the
growth rate of $u_n$. Nothing in \S\ref{sec:sd}, Appendix~\ref{app:annulus} or
Appendix~\ref{app:star} touches it.

\emph{What it rests on.} Nothing further: it is proved. The localisation of cycles, that a
pure-travel element has $c=0$, is Corollary~\ref{cor:localzero}, a case of
Theorem~\ref{thm:nogap}. The metric identity is not needed. Theorem~\ref{thm:nogap} exhibits a
relaxed-optimal realisation with no isolated cycle; that realisation is a single open walk, so it is
admissible for $\ell_T$ and gives $\ell_T\le\ell_R$, while $\ell_T\ge\ell_R$ holds because the
relaxed minimum ranges over a superset of the realisations. Hence $\ell_T=\ell_R$ on pure-travel
elements, the two gradings agree termwise, and the travel recursion \eqref{eq:krec}, assembled from
pure-travel runs alone, is computed inside a $c\equiv0$ sector, so the marker cannot enter it.

In particular the reverse inequality of $\ell_T=\ell_R+2c$ -- open when earlier versions of this
section listed it as a dependency, and since proved in \cite{Bonfioli2026a} -- plays no role here
regardless: only the direction that follows from an exhibited realisation is used.

\emph{Evidence.} The pure-travel generating function computed in both gradings to radius $18$ is
identical coefficient by coefficient,
$1,3,5,8,12,17,24,35,52,78,116,173,256,383,572,858,1276,1907,2832$, with ratio tending to
$\beta_2$.

\subsection{The former input \textup{(L)}: the bulk dictionary and the gate}\label{sec:inputL}

The marker enters the block assembly through exactly one factor, the gap-bridge kernel
$\mathfrak g(q,y)=q/(1-qy)$ of \textup{(M2)}, summing a maximal run of $L$ forced gap edges, each
of weight $q$ and each an isolated cycle except the one shielded by the inner block's turn-around.
Its two evaluations, $\mathfrak g_V=\mathfrak g(q,1)=\tfrac{q}{1-q}$ and
$\mathfrak g_U=\mathfrak g(q,q)=\tfrac{q}{1-q^2}=\tfrac{\mathfrak g_V}{1+q}$, are holomorphic and
non-zero on $(0,1)$, hence at every travel pole.

The bulk resolvent is semiseparable, $M_0[b,a]=2q^bq^{\max(a,b)}$ plus a rank-one gap term of
weight $\mathfrak g$, and the resolvent of $E$ is a M\"obius function of $\mathfrak g$:
\begin{equation}\label{eq:mobius}
   B^E(q,y)=\frac{b_0+\mathfrak g\,\kappa}{1-\mathfrak g\,t_1},\qquad \kappa=t_0b_1-b_0t_1,
\end{equation}
with $b_0,b_1,t_0,t_1$ the $y$-free gapless-resolvent sums of
Proposition~\ref{prop:bulkdress}, which satisfy the reciprocity $t_0=b_1$
(Proposition~\ref{prop:recip}). The factorisation \eqref{eq:mobius} is proved as
Proposition~\ref{prop:mobius}; its two hypotheses are $S_e\ne0$ and $s\ne1$, which the paper was
carrying already (Remark~\ref{rem:mobhyp}). Write $s=\mathfrak g_Vt_1$, so that $qt_1=s(1-q)$.
What the assembly consumes is the \emph{gate}
\begin{equation}\label{eq:gate}
   b_0>0\qquad\text{and}\qquad s:=\tfrac{q}{1-q}\,t_1<1,\qquad t_1=\frac{P_{12}}{S_e},
\end{equation}
on the off-diagonal cocycle entry $P_{12}$, together with the dictionary identifying the entries of
\eqref{eq:mobius} with the blocks of \eqref{eq:blocks},
\begin{equation}\label{eq:dict}
   S_e=1-S_1,\qquad S_o=\tfrac{1-q}{2q}\,S_0,\qquad
   b_0=\frac{S_0}{1-S_1}=\frac{2q}{1-q}\cdot\frac{S_o}{S_e},
\end{equation}
$S_o$ being identified in Lemma~\ref{lem:P12closed} and the two resolvent lines in
Corollary~\ref{cor:dict}. The remaining entry $t_1=P_{12}/S_e$ is an identity between
independently defined objects (Corollary~\ref{cor:cocycle}), not a definition, so all four lines
of \eqref{eq:dict} now stand on the same footing.

\emph{What is proved.} The factorisation \eqref{eq:mobius} and the dictionary \eqref{eq:dict} are
Theorem~\ref{thm:L}. Three ingredients previously listed as
verified rather than proved are now proved: the reciprocity $t_0=b_1$
(Proposition~\ref{prop:recip}), the entry $t_1=P_{12}/S_e$ (Corollary~\ref{cor:cocycle}), and the
identification of the resolvent quantities $b_0,b_1,t_0$ with the blocks
(Proposition~\ref{prop:bulkdress}). That the bulk transfer is $M_0$ with the rank-one gap term and
source $u$ is part of \textup{(M)}, now Theorem~\ref{thm:model}. On the numerical side, the kernel
$K(a,b)=q^{\max(a,b)}$ on active half-sizes reproduces
$S_0/(1-S_1)=2q+2q^2+6q^3+2q^4+18q^5+6q^6+42q^7+\cdots$ (\texttt{blocks\_growth.py}), and the
whole dictionary is re-derived in exact integer arithmetic and at two precisions in
\texttt{assembly\_certificate.py}.

\emph{Where it is used.} In Theorem~\ref{thm:RJ}, through $t_1$ and $1-\mathfrak gt_1$; and inside
Corollary~\ref{app:gate}, where the half $b_0>0$ is proved from a sign chain running through the
closed form in \eqref{eq:dict}. The other half of the gate, $|s|<1$, uses neither: it is
Theorem~\ref{thm:star} with Lemma~\ref{app:Se}, and it is unconditional.

\subsection{The gate block in closed form}\label{sec:gateblock}

Both columns of the cocycle satisfy the scalar three-term recursion
$y_{n+1}=(1+q^3-2(1-q)q^{2n+2})y_n-q^3y_{n-1}$; setting $x=q^n$ turns it into the linear
$q$-difference equation
\begin{equation}\label{eq:qdiff}
  Y(qx)+q^3Y(x/q)=\bigl(1+q^3-2(1-q)q^2x^2\bigr)Y(x),
\end{equation}
whose regular solution is, in closed form, a Hahn--Exton $q$-Bessel function
\cite{Hahn1953,Exton1978,KS1992},
\begin{equation}\label{eq:p12closed}
  Y_3(x)=\sum_{k\ge0}d_kx^{2k+3},\qquad
  d_k=\frac{(-2)^k(1-q)^kq^{k^2+3k}}{(q^2;q^2)_k(q^5;q^2)_k}=J^{(3)}_{3/2}\ \text{normalised},
\end{equation}
and $P_{12}=\tfrac{2q^3}{1-q^3}Y_3(1)$. The denominator $S_e$ is the block of \eqref{eq:blocks},
controlled by Lemma~\ref{app:Se}. It remains to bound the numerator residual against its confluent
limit; this is where the two-Bessel identity enters.

\subsection{The amplitude atom, formally verified}

As $q\to1$, equation \eqref{eq:qdiff} degenerates to $\tau Y''-(2\tau/x)Y'+2Y=0$, solved by
$x^{3/2}J_{3/2}(Wx)$. The $q$-corrected connection coefficient across the turning point $X=\nu=\tfrac32$
is governed by the following identity, whose point is that the confluence is \emph{entire}; it
passes \emph{through} the turning point, bypassing the Liouville--Green singularity where asymptotic
methods stall.

\begin{lemma}[two-Bessel confluence; amplitude normalisation]\label{lem:atomN}
The regular solution of \eqref{eq:qdiff} has the confluent form
\[
   Y_3(x)=\gamma_{\mathrm{cl}}\,x^{3/2}\Bigl[J_{3/2}(X)+\tfrac{\tau X}2J_{5/2}(X)\Bigr]+O(\tau^2),
   \qquad X=wx,\ \gamma_{\mathrm{cl}}=\tfrac{3\sqrt\pi}4\bigl(\tfrac w2\bigr)^{-3/2},
\]
and its single-solution envelope satisfies $A/A_{\mathrm{cl}}=1+\tau\,h(X)+O(\tau^2)$ with the
\emph{explicit rational}
\[
   h(X)=\frac{1+3/(2X^2)}{1+1/X^2}\ \in\ [1,\tfrac32]\qquad(X>0).
\]
Consequently $\gamma/\gamma_{\mathrm{cl}}=1+O(\tau)$ with explicit constant $\tfrac32$.
\end{lemma}

\begin{proof}
From the $q$-Pochhammer confluence $d_k=\tau^{-k}D_k(1-k\tau+O(\tau^2))$ with
$D_k=(-1)^k/[2^kk!\,(\tfrac52)_k]$, the leading sum is $\gamma_{\mathrm{cl}}x^{3/2}J_{3/2}(X)$ by the
Bessel representation of ${}_0F_1$, and the $-k\tau$ correction contributes
$\tfrac{\tau X}2\gamma_{\mathrm{cl}}x^{3/2}J_{5/2}(X)$ via ${}_0F_1(;\tfrac72;\cdot)$. Both series are
entire in $X$. The envelope of $J_{3/2}+\tfrac{\tau X}2J_{5/2}$ follows from
$A^2=M_{3/2}^2+(\tfrac{\tau X}2)^2M_{5/2}^2+\tau X(J_{3/2}J_{5/2}+Y_{3/2}Y_{5/2})$ and the closed
half-integer cross-term identity
\[
   J_{3/2}(X)J_{5/2}(X)+Y_{3/2}(X)Y_{5/2}(X)=\frac{4}{\pi X^2}\Bigl(1+\frac3{2X^2}\Bigr),
\]
which with $M_{3/2}^2=\tfrac2{\pi X}(1+\tfrac1{X^2})$ gives $h(X)=(1+\tfrac3{2X^2})/(1+\tfrac1{X^2})$;
this is monotone decreasing from $h(0^+)=\tfrac32$ to $h(\infty)=1$, so $1\le h\le\tfrac32$.
\end{proof}

\begin{remark}[machine verification]\label{rem:lean}
The cross-term identity and the bound $1\le h(X)\le\tfrac32$ have been formalised and checked in
Lean~4 against Mathlib (file \texttt{AtomN.lean}); the proof terms depend only on the standard
axioms \texttt{propext}, \texttt{Classical.choice}, \texttt{Quot.sound}, with no \texttt{sorry}.
The identity is proved by a single \texttt{linear\_combination} off $\sin^2+\cos^2=1$, and the bound
by \texttt{nlinarith}.
\end{remark}

\begin{theorem}\label{thm:U}
The growth series $U$ of $W_{\mathrm{univ}}$ is transcendental over $\Q(x)$.
\textup{(}Standing: Theorems~\ref{thm:model} and~\ref{thm:RJ}.\textup{)}

The input \textup{(T)} that earlier versions of this theorem carried is gone:
Proposition~\ref{prop:travelinv} is proved, via Theorem~\ref{thm:nogap} and
Corollary~\ref{cor:localzero}, and the proof does not use the metric identity.

The hypothesis \textup{(L)} that earlier versions of this theorem carried is gone: its content,
the factorisation \eqref{eq:mobius} with the dictionary \eqref{eq:dict} and the reciprocity
$t_0=b_1$, is Theorem~\ref{thm:L}. The residue is the square of $\Pi_q(q_m)$, which is
Theorem~\ref{thm:RJ}.

Independently of \textup{(M)}, \textup{(T)} and \textup{(R-J)}, the half
$|s|<1$ of the gate \eqref{eq:gate} holds at every travel pole. The numerator amplitude estimate
$(\star)$ (at every travel pole
$q_m$, $|P_{12}(q_m)|\,w_m^3<2$) is Theorem~\ref{thm:star} (Appendix~\ref{app:star}):
$|P_{12}|\le0.619\,\tau^{3/2}$ for $\tau\le5\cdot10^{-3}$ and $|P_{12}|w^3\le0.75$ at the six
poles above that threshold. The passage to $|s|<1$ is
Corollary~\ref{app:gate}: $|s|\le0.983$ on the uniform range and $|s|\le0.296$ at the six
exceptional poles, with the denominator lower bound $|S_e|\ge0.63\sqrt\tau$ supplied by
Lemma~\ref{app:Se}. The remaining half $b_0>0$ is proved in the same corollary from a sign chain
through the closed form for $b_0$, which is Corollary~\ref{cor:dict}. The
leading-order picture of Remark~\ref{rem:elemY3} (the phase-matched zero giving
$|P_{12}|/\tau^{3/2}\to1/(4\sqrt2)$, a fourfold margin below the threshold $1/\sqrt2$) is
consistent with, and explained by, the second-difference identity of
Proposition~\ref{prop:P12second}.
\end{theorem}

\begin{proof}
By Lemma~\ref{lem:mero} at $y=x^2$, $U$ is meromorphic on $|x|<1$, and the local forms below are
those of this continuation. By \textup{(T)}, $U$ inherits the travel poles. By Proposition~\ref{prop:shape} at $y=q$, at each
travel pole $U$ has the local form $\mathcal N_q(q-q_m)^{-\pi}+(\text{lower order})+H_m$ with
$\mathcal N_q=(2x_m)^{-1}c_\star\Pi_q(q_m)^2$ and $c_\star\ne0$; the boundary vectors are holomorphic there
because the gate \eqref{eq:gate} holds at \emph{every} travel pole. Indeed, on the uniform range
$\tau\le5\cdot10^{-3}$, $|P_{12}|\le0.619\,\tau^{3/2}$ (Theorem~\ref{thm:star}) and
$|S_e|\ge0.63\sqrt\tau$ (Lemma~\ref{app:Se}) give
$|s|=\tfrac{q}{1-q}\,|P_{12}|/|S_e|\le\tfrac{\tau}{e^\tau-1}\cdot\tfrac{0.619}{0.63}\le0.983<1$,
and $b_0>0$ follows from the sign chain of the same lemmas through the closed form of
Corollary~\ref{cor:dict}; the six exceptional poles are checked directly
(Corollary~\ref{app:gate}). By Theorem~\ref{thm:RJ} the leading coefficient $\mathcal N_q$ is
non-zero at every $q_m$, so $U$ has an uncancelled pole at each of them; these accumulate at
$q=1$ (Lemmas~\ref{lem:infpoles} and \ref{lem:polefinite}), so $U$ has infinitely many poles in
the unit disc accumulating at its boundary and is transcendental over $\Q(x)$, as in
Theorem~\ref{thm:V}. Transcendence does not consume the exceptional-pole enumeration of
Theorem~\ref{thm:star}: the uniform range $\tau\le5\cdot10^{-3}$ alone carries infinitely many
poles (Lemma~\ref{lem:infpoles} produces them at $\tau\downarrow0$); the enumeration only upgrades
``at infinitely many travel poles'' to ``at every travel pole'' in the gate. For orientation: the
computed values over $m=5,\dots,71$ are $|P_{12}|/\tau^{3/2}\to1/(4\sqrt2)=0.1768$ and
$s\to\tfrac14$, well inside every bound above.
\end{proof}

\begin{remark}[the assembly is machine-checked]\label{rem:uassembly}
The composition just performed (the gate arithmetic, the two halves of \eqref{eq:gate}, the
passage to the pole set, and the step from infinitely many poles to non-algebraicity) is
verified in Lean~4 (\texttt{UAssembly.lean}, seven theorems, no \texttt{sorry}, axioms
\texttt{propext}, \texttt{Classical.choice}, \texttt{Quot.sound}), each analytic input entering as
an explicit hypothesis: Theorem~\ref{thm:star}, Lemma~\ref{app:Se}, Corollary~\ref{app:gate},
Lemma~\ref{lem:infpoles}, the reduction \textup{(L)} (the Lean hypothesis \texttt{hred}), and the
finiteness of an algebraic function's
pole set in the disc. The file formalises \emph{no} analysis: no $q$-series, no steepest
descent, no numerics; its content is exactly that those hypotheses \emph{compose}, so its
hypothesis list is an audit of what Theorem~\ref{thm:U} rests on, not a verification of any of
them. In particular \texttt{hred}, the hypothesis ``the gate implies a non-zero residue'', is
assumed there and not derived; its content is Theorem~\ref{thm:RJ}, whose identities and
positivity step are checked in \texttt{RJClosedForm.lean} and \texttt{RJMain.lean} with the
eigen-relations and the gate as hypotheses (\S\ref{sec:RJlean}). The
quantitative step is isolated
as \texttt{gate\_ratio\_lt\_one}, whose hypothesis $K<c$ (numerator constant strictly below
denominator constant) is what the gate needs and what a numerator bound alone cannot supply:
with $K=0.619$ from Theorem~\ref{thm:star} and the constant $0.35$ of
Proposition~\ref{prop:selower} it is \emph{false}, which is precisely why Lemma~\ref{app:Se}'s
$0.63$ is required.
\end{remark}

\begin{remark}[the standing of Theorem~\ref{thm:U}]\label{rem:Ustatus}
The unconditional content of this section is the gate: at every travel pole,
$|P_{12}|\le0.619\,\tau^{3/2}$ (Theorem~\ref{thm:star}), $|S_e|\ge0.63\sqrt\tau$
(Lemma~\ref{app:Se}) and hence $|s|\le0.983<1$ (Corollary~\ref{app:gate}), together with the
infinitude and accumulation of the travel poles (Lemma~\ref{lem:infpoles}) and the fact that these
facts compose to non-algebraicity (Remark~\ref{rem:uassembly}).
The factorisation \eqref{eq:mobius}, the reciprocity $t_0=b_1$, the dictionary \eqref{eq:dict} and
the half $b_0>0$ of the gate are Theorem~\ref{thm:L}; the junction pairing is Theorem~\ref{thm:RJ}.

\textup{(T)} is now proved. Proposition~\ref{prop:travelinv} follows from
Corollary~\ref{cor:localzero}: a pure-travel element has no gap edge, so by
Theorem~\ref{thm:nogap} some relaxed-optimal realisation is a single open walk, which bounds
$\ell_T\le\ell_R$; the reverse inequality is immediate. The metric identity
$\ell_T=\ell_R+2c$ is not used.

\textup{(R-J)} is Theorem~\ref{thm:RJ}, and \textup{(M)} is Theorem~\ref{thm:model}. The standing of
Theorem~\ref{thm:U} is therefore that of its inputs: the gate (Appendix~\ref{app:star}, which
contains verified interval-arithmetic enclosures); the metric identity and the faithfulness of the
model (Lean); the assembly \textup{(M3)} (a hand proof, Appendix~\ref{app:M3prime}, with two
independent adversarial reviews and exact series agreement to $x^{26}$); and, for the statement at
every travel pole rather than infinitely many, the interval certificate at the first twelve poles.
The assembly is not formalised.
\end{remark}

\begin{remark}[what the turning-point atom does and does not give]\label{rem:Ugap}
Lemma~\ref{lem:atomN} discharges the amplitude \emph{through the turning point} $X=\tfrac32$, which
is the obstruction to any Liouville--Green treatment. It does \emph{not} supply the amplitude at
the \emph{pole} $X=w$ to the
$O(\tau^2)$-relative accuracy the numerator gate requires. There the two-Bessel form carries a
relative defect whose $O(1)$ part is exactly $\tfrac1{36}$; i.e.\ the classical
turning-point-to-pole connection factor $\tfrac{36}{35}$, which is itself derivable and, once
applied, matches $Y_3(1/q)$ to $O(\tau)$-relative; and whose \emph{residual} $O(\tau)$ part (the
amplitude drift from the turning point to the pole) is not supplied by this route. Because
$P_{12}=\tfrac{2q^3}{3(1-q^3)}Y_3(1/q)-\tfrac23S_e$ is an $O(\sqrt\tau)$ cancellation, that residual
is amplified to the size of $P_{12}$ itself. This residual is closed in
Appendix~\ref{app:star}, not by resumming the amplitude at all, but by the second-difference
identity of Proposition~\ref{prop:P12second}: $P_{12}$ is the symmetric second difference of the
$q$-cosine at its own zero, so the gate needs only $\max|\psi''|\le3.5w$ on a window of width
$\tau$, which Theorem~\ref{thm:star} proves; Lemma~\ref{app:Se} and Corollary~\ref{app:gate} then
close the remaining halves of the gate ($|S_e|\ge0.63\sqrt\tau$ and $b_0>0$). The two
off-diagonal cocycle entries obey the Abel relation $x_NY_N-X_Ny_N=-1$ (the discrete $q$-Wronskian,
conserved because $\det M_n\equiv1$), whence
\[
   P_{12}=\frac{X_N\,S_e-1}{x_N}\qquad\text{exactly},
\]
$x_N,X_N$ being the leading components of the same three-term recursion. Hence $P_{12}$ is fixed by the
\emph{same} family of Hahn--Exton blocks that controls $S_e$: one analytic class; that of
Remark~\ref{lem:cos}; not two. In this precise sense $U$ sits on $V$'s footing. What $U$ additionally
requires is the pole asymptotics of $x_N,X_N$; again Hahn--Exton, again governed by the exact
$q$-Wronskian of Olde Daalhuis \cite{Daalhuis1994},
$W_q(J_\nu,Y_\nu)=q^{\nu(\nu-1)}(1-q^2)/\Gamma_{q^2}(\tfrac12)^2$ for the numerically satisfactory pair,
whose $q\to1$ expansion is explicit via Moak's $q$-Stirling formula \cite{Moak1984} (e.g.\
$\Gamma_{q^2}(\tfrac12)=\sqrt\pi(1-\tfrac38\tau+O(\tau^2))$). The block-asymptotic assembly is carried
out \emph{elementarily} in Remark~\ref{rem:elemY3}: the numerator amplitude has a derived closed-form
expansion with exact rational constants, and the phase-matching of the travel-pole equation against the
amplitude's zero (both collapsing to the same $\sinh$-product family) derives the gate value
$1/(4\sqrt2)$. $U$'s analytic inputs thereby coincide with $V$'s: the further pieces ($(\star)$
and the closure of the gate) are proved in Appendix~\ref{app:star}. (The orthogonal arithmetic
route of
\S\ref{sec:modp} would settle $U$ without any of this analysis, but is itself open.)
\end{remark}

\section{Beyond algebraicity: \texorpdfstring{$D$}{D}-finiteness, functional equations and the
natural boundary}\label{sec:beyond}

Theorems~\ref{thm:V} and~\ref{thm:U} exclude algebraic equations. The pole set that proves them
excludes much more, because it is infinite and lies inside the disc $|x|<1$ on which $U$ and $V$ are
meromorphic (Lemma~\ref{lem:mero}), well beyond their radius of convergence $x_1=0.6704\ldots$. This
section records what it excludes. Every statement below about $U$ or $V$ has the standing of
Theorems~\ref{thm:V} and~\ref{thm:U} (Theorems~\ref{thm:model} and~\ref{thm:RJ}); every statement
about the travel resolvent $G^T_0$ alone is unconditional, as Theorem~\ref{thm:travelres} is.

\subsection{Poles of \texorpdfstring{$U$}{U} on the negative axis}

\begin{proposition}\label{prop:minusx}
At every travel pole $q_m$, $U$ has a pole at $x=-x_m$ as well as at $x=x_m$, both of order
$\pi$, the order of the pole of $(I-\mathcal T)^{-1}$ at $q_m$. At $x=-x_m$ the junction pairing is
negative:
\[
   \Pi_q\big|_{x=-x_m}=\frac{2q_mB_0(1-x_m)}{1+x_m}\Bigl(t_1\frac{1+x_m+q_m}{1+q_m}
   -\frac1{2x_m}\Bigr)<0,\qquad B_0=\frac1{1-\mathfrak g_Ut_1}.
\]
\textup{(}Standing: Theorems~\ref{thm:model} and~\ref{thm:RJ}; at the six poles with
$\tau_m>5\cdot10^{-3}$, the enumeration standing of Theorem~\ref{thm:star}.\textup{)}
\end{proposition}

\begin{proof}
The travel resolvent $(I-\mathcal T)^{-1}$, the bulk vector $\Psi$, $t_1$ and $B_0$ depend on $x$
only through $q=x^2$; $x$ itself enters $\mathcal W(x,q)$ only through the monomials of
\eqref{eq:app-mu}--\eqref{eq:app-E} and the prefactor $x^{-1}$. The proofs of
Proposition~\ref{prop:shape} and of \eqref{eq:RJclosedq} use only $q=x^2$, never the sign of $x$
(the reductions $x^{\max(2s+1,2b\pm1)}=x\,q^{\max(\cdot,\cdot)}$ and the final polynomial identities
hold for either root). Near $x=-x_m$ one has $q-q_m=(x+x_m)(x-x_m)$ with the second factor
non-zero, so $U=\mathcal N_q(-2x_m)^{-\pi}(x+x_m)^{-\pi}+(\text{lower order})$ with
$\mathcal N_q=c_\star(-2x_m)^{-1}\Pi_q|_{x=-x_m}^2$, and \eqref{eq:RJclosedq} at $x=-x_m$ is the
displayed formula. The prefactor is positive: $0<x_m<1$, and $1-\mathfrak g_Ut_1>0$ because
$|\mathfrak g_Ut_1|<|\mathfrak g_Vt_1|=|s|<1$ (Corollary~\ref{app:gate}). For the bracket write
$t_1=s(1-q)/q$ and use $(1+x+q)/(1+q)=1+x/(1+x^2)\le\tfrac32$:
\[
   2x_m\,|t_1|\,\frac{1+x_m+q_m}{1+q_m}\;\le\;3|s|\,\frac{1-q_m}{x_m}.
\]
For $\tau_m\le5\cdot10^{-3}$, $|s|\le0.983$ and $(1-q_m)/x_m\le\tau_me^{\tau_m/2}\le0.00502$, so the
right side is below $0.015$. At the six poles with $\tau_m>5\cdot10^{-3}$, $|s|\le0.296$
(Corollary~\ref{app:gate}) and $(1-q_m)/x_m\le(1-q_1)/x_1=0.8212$, since $q\mapsto(1-q)/\sqrt q$
decreases; the right side is at most $0.730$. In both cases $|t_1|(1+x_m+q_m)/(1+q_m)<1/(2x_m)$, the
bracket is negative, $\mathcal N_q\ne0$, and $-x_m$ is a pole of order exactly $\pi$.
\end{proof}

\noindent At the dominant pole the bracket is $-0.2155$, and it tends to $-\tfrac12$ along the
ladder. Only the uniform range is needed for infinitely many poles at $-x_m$; the six-pole list
upgrades this to every travel pole.

\begin{remark}[$V$ at $-x_m$: verified, not proved]\label{rem:Vminusx}
For $V$ the same substitution in \eqref{eq:RJclosed1} gives
\[
   (1-\mathfrak g_Vt_1)\,\Pi_1\big|_{x=-x_m}=\frac{2q_m(1-x_m)}{1-q_m}\Bigl(t_1-\frac{1-x_m}{2x_m}\Bigr),
\]
and here the gate does not decide the sign. With $t_1=s(e^{\tau}-1)$, $s\to\tfrac14$, and
$(1-x_m)/(2x_m)=(e^{\tau/2}-1)/2=\tfrac\tau4+O(\tau^2)$, the two terms cancel at leading order and
the bracket is $O(\tau^2)$; its sign is decided by the $O(\tau)$ correction to $s$, which is not
proved. Evaluated at $90$ digits at the first twelve travel poles, the bracket is positive, from
$0.117$ at $m=1$ to $2.9\cdot10^{-7}$ at $m=12$, and close to $\tau_m^2/8$. So $V$ has a pole at
$-x_m$ for $m\le12$; for all $m$ this is open. Nothing below uses it: every statement about $V$ uses
only the poles $x_m$.
\end{remark}

\subsection{\texorpdfstring{$D$}{D}-finiteness}

\begin{lemma}\label{lem:odeposes}
Let $\Omega\subset\C$ be a domain, let $a_0,\dots,a_r,b$ be holomorphic on $\Omega$ with
$a_r\not\equiv0$, and let $f$ be meromorphic on $\Omega$ with
$a_rf^{(r)}+\cdots+a_1f'+a_0f=b$ off the poles of $f$. Then every pole of $f$ in $\Omega$ is a zero
of $a_r$.
\end{lemma}

\begin{proof}
Let $z_0$ be a pole of $f$ of order $k\ge1$. Then $f^{(j)}$ has a pole of order exactly $k+j$ at
$z_0$, so for $j<r$ the term $a_jf^{(j)}$ has a pole of order at most $k+r-1$, and $b$ is
holomorphic. If $a_r(z_0)\ne0$, the term $a_rf^{(r)}$ has a pole of order exactly $k+r$, which
nothing else in the equation cancels. Hence $a_r(z_0)=0$.
\end{proof}

\begin{theorem}\label{thm:nonDfinite}
$U$ and $V$ are not $D$-finite: neither satisfies a linear differential equation
$\sum_{j=0}^ra_j(x)f^{(j)}(x)=b(x)$ with polynomial coefficients and $a_r\not\equiv0$, nor even one
whose coefficients $a_j$ and $b$ are holomorphic on a neighbourhood of the closed disc $|x|\le1$.
Equivalently, $(u_n)$ and $(v_n)$ are not P-recursive. The same holds for the travel resolvent
$G^T_0$, as a series in $q$ and as a series in $x$, and for $G^T_0$ it is unconditional.
\textup{(}Standing for $U$ and $V$: Theorems~\ref{thm:model} and~\ref{thm:RJ}.\textup{)}
\end{theorem}

\begin{proof}
Let $f$ be one of the series and suppose it satisfies such an equation. By Lemma~\ref{lem:mero}
(for $G^T_0$, Theorem~\ref{thm:blocks}, first clause) $f$ continues to a meromorphic function on the
open unit disc $D$. The expression $\sum_ja_jf^{(j)}-b$ is then meromorphic on $D$ and vanishes near
$0$, so by the identity theorem it vanishes on $D$ off the poles of $f$. By
Lemma~\ref{lem:odeposes} every pole of $f$ in $D$ is a zero of $a_r$. Since $a_r\not\equiv0$ is
holomorphic on a neighbourhood of the compact closed disc, its zeros there are isolated and hence
finite in number. But $f$ has infinitely many poles in $D$: $V$ at every $x_m$
(Theorem~\ref{thm:V}), $U$ at every $\pm x_m$ (Theorem~\ref{thm:U} and
Proposition~\ref{prop:minusx}), and $G^T_0$ at every $q_m$ in the variable $q$ and at every $\pm x_m$
in the variable $x$ (Theorem~\ref{thm:travelres}). This is a contradiction. $D$-finiteness of a
power series is equivalent to P-recursiveness of its coefficients \cite[\S6.4]{StanleyEC2}.
\end{proof}

\noindent The pole argument does not reach differential algebraicity: $1/\cos(\pi/(1-x))$ has
poles at $x=1-2/(2n+1)$, $n\ge1$, infinitely many in $|x|<1$ and accumulating at $x=1$, and satisfies
an algebraic differential equation. Whether $U$ and $V$ are $D$-algebraic is open; we expect not.

\begin{corollary}\label{cor:Wnotdfinite}
The bivariate series $\mathcal W(x,y)\in\Z[y][[x]]\subset\Z[[x,y]]$ is not $D$-finite in $(x,y)$.
\textup{(}Standing: Theorems~\ref{thm:model} and~\ref{thm:RJ}.\textup{)}
\end{corollary}

\begin{proof}
$D$-finiteness is preserved under substitution of algebraic power series without constant term
\cite{Lipshitz1989}. If $\mathcal W$ were $D$-finite, the substitution $y=x^2$ would make
$\mathcal W(x,x^2)=U$ $D$-finite (Corollary~\ref{cor:M3prime-growth}), contrary to
Theorem~\ref{thm:nonDfinite}.
\end{proof}

\subsection{Linear \texorpdfstring{$q$}{q}-difference and Mahler equations}

\begin{proposition}\label{prop:noqdiff}
Let $f$ be $U$ or $V$, as series in $x$, or $G^T_0$, as a series in $q$, and let $a_0,\dots,a_r,b$
be polynomials with $a_0a_r\not\equiv0$.
\begin{enumerate}[label=\textup{(\roman*)},leftmargin=2.2em,itemsep=0.1em,topsep=0.2em]
\item If $Q\in\C^\times$ is not a root of unity, then $\sum_{i=0}^ra_i(x)f(Q^ix)\ne b(x)$.
\item If $k\ge2$ is an integer, then $\sum_{i=0}^ra_i(x)f(x^{k^i})\ne b(x)$.
\end{enumerate}
For $G^T_0$ both statements are unconditional; for $U$ and $V$ the standing is that of
Theorems~\ref{thm:V} and~\ref{thm:U}. When $Q$ is a root of unity the case is open: the pole set can
be invariant under $x\mapsto Qx$, as that of $U$ is under $x\mapsto-x$
\textup{(}Proposition~\ref{prop:minusx}\textup{)}, and the argument below gives nothing.
\end{proposition}

\begin{proof}
Let $D$ be the open unit disc, $P$ the pole set of $f$ in $D$, and $\rho>0$ the radius of
convergence of $f$; $P$ lies in $\rho\le|x|<1$, has finitely many points in each disc $|x|\le R<1$,
and is infinite (proof of Theorem~\ref{thm:nonDfinite}). An equation as in (i) or (ii) holds as an
identity of power series near $0$.

(i) If $|Q|>1$, replacing $x$ by $Q^{-r}x$ and $i$ by $r-i$ gives an equation of the same form with
$Q^{-1}$ in place of $Q$ and new leading and trailing coefficients $a_r(Q^{-r}x)$, $a_0(Q^{-r}x)$; so
we may assume $|Q|\le1$. Then every $f(Q^ix)$ is meromorphic on $D$, and the equation holds on $D$
by the identity theorem. Let $E$ be the finite zero set of $a_0$ in $D$. At $p\in P\setminus E$ the
term $a_0f$ is singular and $b$ is not, so some $f(Q^ix)$ with $1\le i\le r$ is singular at $p$,
that is $Q^ip\in P$. Thus every $p\in P\setminus E$ has a \emph{successor} $Q^ip\in P$, $1\le i\le
r$, and every point of $P$ is the successor of at most $r$ points. Follow successors from any $p\in
P$; along the chain the exponent of $Q$ increases strictly.

If $|Q|<1$ the moduli decrease strictly, so the chain stays in the finite set $P\cap\{|x|\le|p|\}$
and ends, necessarily at a point of $E$. A chain of $n$ steps ends at modulus at most
$|Q|^n|p|<|Q|^n$, and that modulus is at least $\rho$, so $n<N:=\log\rho/\log|Q|$. Hence every
point of $P$ is reached backwards from $E$ in fewer than $N$ steps, and $|P|\le|E|\sum_{n<N}r^n$ is
finite.

If $|Q|=1$ the chain stays on the circle $|x|=|p|$, which meets $P$ in a finite set, and its points
are pairwise distinct because $Q$ is not a root of unity; so it ends in $E$. Hence $P$ lies on the
finitely many circles $|x|=|e|$, $e\in E$, each meeting $P$ in finitely many points, and $P$ is
finite. Both cases contradict the infinitude of $P$.

(ii) Here every $f(x^{k^i})$ is meromorphic on $D$ and the equation holds on $D$. Work on the
positive axis: let $P_+=P\cap(0,1)$ and let $E_+$ be the finite zero set of $a_0$ in $(0,1)$. As in
(i), every $p\in P_+\setminus E_+$ has a successor $p^{k^i}\in P_+$ with $1\le i\le r$; it is
smaller than $p$, so every chain ends in $E_+$, and reading a chain backwards,
$p=e^{k^{-n}}$ for some $e\in E_+$ and some integer $n\ge0$. Hence, for $0<R<1$ and
$\ell=\log(1/R)$,
\[
   \#\bigl(P_+\cap(0,R]\bigr)\;\le\;\sum_{e\in E_+}\#\bigl\{n\ge0:k^n\le\log(1/e)/\ell\bigr\}
   \;\le\;|E_+|\Bigl(1+\log_k\frac L\ell\Bigr),\qquad L=\max_{e\in E_+}\log(1/e),
\]
which grows like $\log(1/\ell)$ as $R\uparrow1$. On the other hand, by Lemma~\ref{lem:infpoles}
there is a travel pole with $w\in(m\pi,(m+1)\pi)$ for every $m\ge4$, at
$x=e^{-1/w^2}<e^{-1/((m+1)\pi)^2}$, equivalently $q=e^{-2/w^2}<e^{-2/((m+1)\pi)^2}$, and it lies in
$P_+$ (Theorems~\ref{thm:V}, \ref{thm:U}, \ref{thm:travelres}). Distinct windows give distinct
poles, so $P_+\cap(0,R]$ has at least $\lfloor c/(\pi\sqrt\ell)\rfloor-4$ points, with $c=1$ in the
variable $x$ and $c=\sqrt2$ in $q$. This grows like $\ell^{-1/2}$, faster than $\log(1/\ell)$, a
contradiction for $R$ close to $1$.
\end{proof}

\noindent The count in (ii) is made on the real axis on purpose. It uses only the travel poles that
are proved to exist, and it needs no statement about where on the circle the poles of $f$
accumulate; compare Remark~\ref{rem:polesdensity}. Mahler equations are those of
\cite{Nishioka1996}. The catalytic equation of \S\ref{sec:setup} is a $q$-difference equation in
the catalytic variable, not in $x$, and is not affected by (i).

\subsection{The natural boundary}

We use the following rationality theorem.

\begin{theorem}[P\'olya--Bertrandias]\label{thm:PB}
Let $\Omega\subset\C$ be a simply connected domain containing $0$ and stable under complex
conjugation, and let $\varphi$ be the Riemann map of the unit disc onto $\Omega$ with
$\varphi(0)=0$, $\varphi'(0)>0$. If $\varphi'(0)>1$ and $f\in\Z[[x]]$ is the Taylor series at $0$ of
a function meromorphic on $\Omega$, then $f$ is a rational function.
\end{theorem}

\noindent This is P\'olya's theorem \cite{Polya1928}, which replaced the disc of Borel's
rationality theorem by a domain measured by transfinite diameter; Bertrandias
\cite{Bertrandias1964} extended it to $p$-adic and adelic domains, whence the name. Bost and
Chambert-Loir prove a generalisation to curves over number fields
\cite[Thm.~7.9]{BostChambertLoir2009}, from which they recover the P\'olya--Bertrandias criterion
\cite[Example~7.10]{BostChambertLoir2009}, and Cantor's rationality theorem
\cite[Thm.~5.2.2]{Cantor1980} is a capacity-theoretic generalisation. Meromorphy suffices, but the strict
inequality $\varphi'(0)>1$ is needed: for $\Omega$ the unit disc, $\varphi'(0)=1$, and
$\sum_nx^{2^n}$ is not rational.

\begin{theorem}\label{thm:natboundary}
\begin{enumerate}[label=\textup{(\alph*)},leftmargin=2.2em,itemsep=0.1em,topsep=0.2em]
\item The travel resolvent $G^T_0=\Sigma_0/(1-\Sigma_1)\in\Z[[q]]$ has $|q|=1$ as a natural
boundary. This is unconditional.
\item $U$ and $V$ have $|x|=1$ as a natural boundary. \textup{(}Standing: Theorems~\ref{thm:model}
and~\ref{thm:RJ}.\textup{)}
\end{enumerate}
More precisely, in each case there is no point $\zeta$ of the unit circle and no disc $B$ about
$\zeta$ such that the function continues meromorphically to $D\cup B$, $D$ the open unit disc.
\end{theorem}

\begin{proof}
Let $f$ be $G^T_0$ (in $q$), $U$ or $V$ (in $x$). In each case $f$ has integer coefficients (the
travel-run counts, $u_n$, $v_n$), is meromorphic on $D$ (Lemma~\ref{lem:mero}), and is not rational,
having infinitely many poles in $D$ (Theorems~\ref{thm:travelres}, \ref{thm:V}, \ref{thm:U}).
Suppose $f$ continues meromorphically to $D\cup B$ with $B$ the disc of radius $\delta>0$ about
$\zeta$, $|\zeta|=1$. Shrinking $\delta$, we may assume $\delta<|\operatorname{Im}\zeta|$ when
$\zeta\notin\R$, so that $B$ and its conjugate $\bar B$ are disjoint; when $\zeta=\pm1$,
$\bar B=B$. Because the coefficients are real, $\overline{f(\bar x)}=f(x)$ near $0$, hence on $D$,
and $x\mapsto\overline{f(\bar x)}$ continues $f$ meromorphically to $\bar B$, agreeing with $f$ on
$\bar B\cap D$. So $f$ is meromorphic and single-valued on
\[
   \Omega=D\cup B\cup\bar B ,
\]
which is stable under conjugation and simply connected: it is the disc $D$ with one or two disjoint
discs attached, each meeting $D$ in a convex set. Let $\varphi$ be its Riemann map, $\varphi(0)=0$,
$\varphi'(0)>0$. By uniqueness of the Riemann map, $\overline{\varphi(\bar z)}=\varphi(z)$, so
$\varphi$ has real Taylor coefficients (Schwarz reflection). The map $\psi=\varphi^{-1}|_D$ sends $D$
into $D$ with $\psi(0)=0$ and is not onto, since $\Omega\ne D$; by the Schwarz lemma
$0<\psi'(0)<1$, that is $\varphi'(0)>1$. Theorem~\ref{thm:PB} now makes $f$ rational, a
contradiction.
\end{proof}

\noindent A $D$-finite series continues analytically along every path avoiding the finitely many
zeros of the leading coefficient of its equation, so Theorem~\ref{thm:natboundary} gives a second
proof of the polynomial-coefficient case of Theorem~\ref{thm:nonDfinite}. The first proof is
elementary and also covers coefficients holomorphic near the closed disc.

\begin{remark}[where on the circle the poles accumulate]\label{rem:polesdensity}
Theorem~\ref{thm:natboundary} does not locate the singularities on the circle. The poles proved
above accumulate only at $x=\pm1$ (at $q=1$ for $G^T_0$), yet every point of the circle is singular.
Numerically the pole set is much larger. The function $1-\Sigma_1=\cos(Z;q^2)$ has $64$ zeros in
$|q|<0.94$ by the argument principle: the travel poles $q_1,q_2$ and $62$ non-real zeros in $31$
conjugate pairs, the first at $q=-0.4075\pm0.5449i$. At each non-real zero $\Pi_q$ was found
non-zero at $40$ digits, so these are candidate poles of $G^T_0$ and of $U$ off the real axis. Plotted,
the zeros approach roots of unity along straight lines. We conjecture that the poles are dense on
the circle; a route is a $q$-cosine resummation at $q=\zeta e^{-t}$, $\zeta$ a root of unity. The
natural boundary does not depend on this.
\end{remark}

\section{An orthogonal arithmetic route}\label{sec:modp}

For each prime $p$, $U\bmod p\in\F_p[[q]]$ is algebraic over $\F_p(q)$ iff its coefficient sequence
is $p$-automatic (Christol \cite{Christol1979,CKMR1980}), iff its $p$-kernel is finite. We find the $p$-kernel
to be \emph{maximally} non-automatic: through the accessible depth its size tracks the full $p$-ary
tree with \emph{no} deficit. (An apparent deficit of $1$ at level $3$; $1,4,13,39$ vs.\ the full
$1,4,13,40$ at $p=3$; is a \emph{comparison-length artifact}: a level-$k$ decimation subsequence
retains only $\sim N/p^{k}$ terms, so at a short prefix two distinct decimations can coincide; the
deficit disappears once the comparison length reaches $9$.) Two instances are machine-checked in
Lean~4; both proofs use \texttt{native\_decide}, so besides the three standard axioms they depend
on the compiler-trust axiom \texttt{\_native.native\_decide.ax\_1\_1} of the declaration concerned.
For the true series itself, the $3$-kernel of $(u_n\bmod3)$ through level $3$ is the full
ternary tree, deficit $0$: all $1+3+9+27=40$ decimation words of length $\le3$ are pairwise distinct: at comparison lengths $9$ and $10$ alike; on the $271$-term prefix $n\le270$
(\texttt{UKernel.lean}, from the validated \texttt{tools/u\_modp\_rust} output, whose $u_n$ are
self-checked against the $43$ known terms and byte-identical to the earlier validated $n\le180$ run
on the overlap). For the sibling block $\Sigma_1$, transcendental over $\Q(q)$ by
Corollary~\ref{cor:polepath}, the same level-$3$ full tree is certified in \texttt{SigmaKernel.lean}
(built from the $A/C$ recursion over $\F_3$), with the short-prefix deficit-$1$ confirmed spurious. The Berlekamp--Massey \cite{Massey1969} linear complexity of $(u_n\bmod 3)$ is $\approx N/2$. This is the
strongest finite-data evidence for non-automaticity, and hence for transcendence over $\Q(x)$,
independent of all of \S\S\ref{sec:sd}--\ref{sec:U}. A proof for $U$ would require a non-automaticity
criterion valid for \emph{dense}-support series (the coefficients carry no lacunary gaps), which
the standard gap/Mahler methods do not provide.

\section{Conclusion}\label{sec:conclusion}

The arithmetic of A396406 separates cleanly. The P\'olya--Carlson dichotomy for the catalytic
building blocks is unconditional (Theorem~\ref{thm:blocks}), and so is the lower bound $\beta_2$
for the growth rate; the matching upper bound, and with it the identification of $\beta_2$ as the
rate, consumes the transfer model \textup{(M)}, Theorem~\ref{thm:model}, and nothing else
(Corollary~\ref{cor:beta}). For the two denominators $\Sigma_1,S_1$ the exclusion of the rational
alternative needs no growth hypothesis: it follows from the infinitude of their level sets in
$(0,1)$, hence from the same Lemma~\ref{lem:T2abs} that the rest of the paper consumes
(Corollary~\ref{cor:polepath}). For the two numerators it rests on Hypothesis~\ref{hyp:growth},
verified out to degree $1156$ and unproved, except for the joint statement that $\Sigma_0$ and
$S_0$ are not both rational, which the pole route supplies. \textbf{The travel resolvent
$\Sigma_0/(1-\Sigma_1)$ is transcendental over $\Q(x)$, unconditionally}
(Theorem~\ref{thm:travelres}), and it needs no growth hypothesis. Its non-cancellation inputs are
two and are separate. The denominator input is Proposition~\ref{prop:selower},
proved in Appendix~\ref{app:annulus} by an elementary chain with explicit numerical constants: an
annulus representation, an Euler bound, a Poisson theta envelope, and a $q$-Wronskian invariant. No
asymptotic theorem is cited on that path. The numerator input is
Proposition~\ref{prop:numnonvanish}, the identity $\Sigma_0(q_m)S_0(q_m)=2q_m/(1-q_m)$, which
follows from the closed forms $\Sigma_0=\tfrac{2q}Z\mathfrak s(Z)$, $S_0=\tfrac{2q}{z_0}\mathfrak
s(z_0)$ and the $q$-Pythagorean identity at a zero of the $q$-cosine; it is exact, it holds at
every pole with no exceptional set, and it consumes no asymptotics. Remark~\ref{lem:cos} gives the
denominator conclusion by a
different and shorter route (the saddles $z^\ast_\pm=\pm iW/2$ with all constants derived
explicitly, the two saddles supplying the global cancellation that no partial-summation method can), but that route is a sketch only (Remark~\ref{rem:cosgap}), and no theorem uses it. Both $V$ and
$U$ take their pole set from Lemma~\ref{lem:infpoles}, which consumes only the rectangle bound
$|T_2|\to0$ of Lemma~\ref{lem:T2abs}. The block assembly is developed in \S\ref{sec:assembly}, on
the observation that the travel and the bulk block are the same object, the resolvent of
$q^{\max(a,b)}$ against a geometric weight: this proves the travel telescoping, the bulk dressing,
the reciprocity $t_0=b_1$ with its cocycle form $\det P=1$, the identity $t_1=P_{12}/S_e$, the
M\"obius factorisation \eqref{eq:mobius}, and the shape of the local form at a travel pole, whose
leading numerator is the square of the junction pairing $\langle\lambda,R\rangle$. It
also refutes the product form of the numerator: the marker junction has rank at least $3$
(Proposition~\ref{prop:junction}); at a travel pole the pairing nevertheless collapses to an exact
closed form in $t_1$, positive at every pole (Theorem~\ref{thm:RJ}). \textbf{$V$ is transcendental
over $\Q(x)$} (Theorem~\ref{thm:V}). The true series $U$ stands on the same
\emph{analytic} footing: its further analytic inputs are proved in Appendix~\ref{app:star}: the gate block is the
symmetric second difference of the $q$-cosine at its own zero (Proposition~\ref{prop:P12second}),
the two-interval uniform certificate of Theorem~\ref{thm:star} gives
$|P_{12}|\le0.619\,\tau^{3/2}<\tau^{3/2}/\sqrt2$ for $\tau\le5\cdot10^{-3}$ with the six poles
above that threshold checked directly, and the exact identity $S_e=\tfrac{qZ}2\mathfrak s(Z)-P_{12}$ turns the
same certificate into the denominator bound that closes the gate
($|s|\le0.983$; Corollary~\ref{app:gate}), consistent with the phase-matched leading
values $|P_{12}|/\tau^{3/2}\to1/(4\sqrt2)$ and $s\to\tfrac14$.
\textbf{$U$ is transcendental over $\Q(x)$} (Theorem~\ref{thm:U}). The input
\textup{(T)} that earlier versions of this theorem carried is gone: Proposition~\ref{prop:travelinv}
is now proved (\S\ref{sec:inputT}), via Theorem~\ref{thm:nogap} and Corollary~\ref{cor:localzero},
and does not depend on the metric identity. \textup{(R-J)} is
Theorem~\ref{thm:RJ}, and the hypothesis \textup{(L)} of earlier versions is Theorem~\ref{thm:L}. Neither series is $D$-finite,
and $|x|=1$ is a natural boundary of both (Theorems~\ref{thm:nonDfinite} and~\ref{thm:natboundary},
by the P\'olya--Bertrandias theorem); for the travel resolvent the natural boundary $|q|=1$ is
unconditional. Two questions remain open: whether the poles are dense on the circle, and whether
$U$ and $V$ are $D$-algebraic.
Every statement in \S\ref{sec:assembly} beyond the operator algebra rests on the transfer model
\textup{(M)}, now Theorem~\ref{thm:model}, whose assembly step is a hand proof not yet formalised.
No input separates $U$ from $V$. Two finite ranges are computer-assisted: the six gate poles above
$\tau=5\cdot10^{-3}$ and the positivity of the pairing at the first twelve poles; neither is
needed for transcendence, only for the statements ``at every travel pole''. The series in question
is the growth series of the universal group $W_{\mathrm{univ}}$; for a particular rational triangle
it agrees with that of $W_T$ only to the effective universality depth of \cite{Bonfioli2026a}.
The turning-point amplitude atom on the $U$-side is the elementary rational bound
of Lemma~\ref{lem:atomN}, formally verified in Lean~4. The independent mod-$p$ route of
\S\ref{sec:modp} would settle $U$ by a completely different mechanism, if a dense-support
non-automaticity criterion can be found. Whether the \emph{number} $\beta_2$ is transcendental is
untouched by all of this and remains open; Appendix~\ref{app:beta2} reduces it to a value statement
for a Hahn--Exton $q$-Bessel function and shows the obstruction to be the arithmetic ratio
$\rho=\tfrac12$.

\appendix
\section{An explicit lower bound for $S_e$ at the travel poles}\label{app:annulus}

Throughout $z_0=\sqrt{2(1-q)}$, $Z=z_0/\sqrt q$, $w=\sqrt{2/\tau}$, $f(z)=\cos(z;q^2)=C(z^2)$,
$C(y)=\sum_{j\ge0}(-1)^jq^{j^2+j}y^j/(q;q)_{2j}$, and $s(z)=\sin(z;q^2)$; by
Proposition~\ref{prop:qtrig}, $S_e=f(z_0)$ and the travel poles are the solutions of $f(Z)=0$.

\begin{proposition}\label{prop:selower}
At every travel pole $q_m$ with $\tau_m\le1.29\cdot10^{-4}$ one has
$|S_e(q_m)|\ge0.35\sqrt{\tau_m}$. The finitely many (Lemma~\ref{lem:polefinite}) poles with
$\tau_m>1.29\cdot10^{-4}$ (the
dominant pole $q_1=0.449453\ldots$ ($|S_e|/\sqrt{\tau_1}=0.581$;
\texttt{star\_gate\_certificate.py}, Corollary~\ref{app:gate}) and the thirty-nine tabulated below
it, $\tau_{41}=1.2354\cdot10^{-4}$ being the first below the threshold) satisfy the same bound
by direct high-precision computation.
\end{proposition}

The proof occupies the rest of the appendix: an annulus representation turns all derivative bounds
into cancellation-free integrals, and the pole condition supplies the value-level input.

\begin{lemma}\label{app:rep}
For $0<r<1$,
$C(y)=\frac1{2\pi}\int_0^{2\pi}\Theta(y/x^2)\,(x;q)_\infty^{-1}\,d\theta$ with $x=re^{i\theta}$ and
$\Theta(t)=\sum_{j\in\Z}(-1)^jq^{j^2+j}t^j=(q^2;q^2)_\infty(q^2t;q^2)_\infty(1/t;q^2)_\infty$.
\end{lemma}

\begin{proof}
Euler's identity gives $1/(q;q)_{2j}=[x^{2j}]\,(x;q)_\infty^{-1}$; extract the coefficient by
Cauchy's formula on $|x|=r$. For $j\le-1$ the integrand $x^{-2j}(x;q)_\infty^{-1}$ is analytic in
$|x|<1$ (the zeros of $(x;q)_\infty$ lie at $x=q^{-k}$, $|q^{-k}|\ge1$), so the negative-index terms
vanish and the sum completes to the bilateral one, which is the triple product. Equivalently this is
the $p=q^2$ case of the $q$-Borel--$q$-Laplace inversion of \cite{Zhang2000} and
\cite[Lemma~1]{Morita2011}, the $q^2$-Borel transform of $C$ being the even part of
$(x;q)_\infty^{-1}$.
\end{proof}

\begin{lemma}\label{app:euler}
For $0<r\le\tfrac12$: $\ (x;q)_\infty^{-1}\le\exp(w_r\cos\theta+K_r)$ in modulus, with
$w_r=r/(1-q)$ and $K_r=r^2/\bigl(2(1-r)(1-q^2)\bigr)$.
\end{lemma}

\begin{proof}
$-\log|1-u|=\operatorname{Re}u+\operatorname{Re}\sum_{n\ge2}u^n/n
\le\operatorname{Re}u+\tfrac12|u|^2/(1-|u|)$ for $|u|<1$; summing over $u=xq^j$ gives
$\sum_j\operatorname{Re}(xq^j)=w_r\cos\theta$ and
$\sum_j r^2q^{2j}/(2(1-rq^j))\le r^2/\bigl(2(1-r)(1-q^2)\bigr)=K_r$.
\end{proof}

\begin{lemma}\label{app:theta}
Let $p=q^2$, $t=\rho e^{i\varphi}$, $a=(\log\rho-\tau)+i\tilde\varphi$ with
$\tilde\varphi\in(-\pi,\pi]$ representing $\varphi+\pi$. For $n\in\{0,1,2\}$,
\[
 |\Theta^{(n)}(t)|\le\sqrt{\pi/\tau}\;e^{(\log\rho-\tau)^2/(4\tau)}\,\rho^{-n}
 \bigl(|a|/(2\tau)+\tau^{-1/2}\bigr)^n\,G(\tilde\varphi),\quad
 G(\tilde\varphi)=\sum_{k\in\Z}e^{-\pi^2(k-\tilde\varphi/2\pi)^2/\tau}.
\]
\end{lemma}

\begin{proof}
$\Theta(t)=\theta_p(-pt)$ with $\theta_p(y)=\sum_np^{n(n-1)/2}y^n$. Poisson summation gives the
\emph{exact} identity
\begin{equation}\label{eq:poisson}
 S(a):=\sum_{j\in\Z}e^{-\tau j^2+aj}=\sqrt{\pi/\tau}\sum_{k\in\Z}e^{a_k^2/(4\tau)},
 \qquad a_k:=a-2\pi ik .
\end{equation}
Since $\operatorname{Re}(a_k^2)=(\log\rho-\tau)^2-(\tilde\varphi-2\pi k)^2$, the imaginary part of $a$ contributes
\emph{decay}, and $\sum_k e^{\operatorname{Re}(a_k^2)/(4\tau)}=e^{(\log\rho-\tau)^2/(4\tau)}G(\tilde\varphi)$; no
cancellation is involved, the term $-(\tilde\varphi-2\pi k)^2$ being precisely what constitutes $G$.
Differentiating \eqref{eq:poisson} in $a$, the $j$- and $j(j-1)$-weighted sums are $S'$ and
$S''-S'$, and $\partial_ae^{a_k^2/(4\tau)}=(a_k/2\tau)e^{a_k^2/(4\tau)}$, so the $k$-th image carries
the weight
\[
 W_1(a_k)=\frac{|a_k|}{2\tau},\qquad
 W_2(a_k)=\frac{|a_k|^2}{4\tau^2}+\frac1{2\tau}+\frac{|a_k|}{2\tau}.
\]
At the radius $\rho=e^\tau$ of Lemma~\ref{app:M2} one has $a_k=i(\tilde\varphi-2\pi k)$ exactly, so
$|a_k|=|\tilde\varphi-2\pi k|$. To pull the weight out of the $k$-sum, note that for $k\ne0$ the
excess $|a_k|-|a_0|$ is paid for by the Gaussian factor, which is smaller by
$e^{-((\tilde\varphi-2\pi k)^2-\tilde\varphi^2)/(4\tau)}$; the aggregate excess is maximised at the
nearest image and is at most $\max_{\delta>0}2\delta e^{-\pi\delta/\tau}=2\tau/(\pi e)\le0.2342\,\tau$,
which is below $\tau^{-1/2}$ for $\tau\le0.0198$. Hence
$\sum_kW_n(a_k)e^{\operatorname{Re}(a_k^2)/(4\tau)}\le\bigl(|a|/(2\tau)+\tau^{-1/2}\bigr)^n
\sum_ke^{\operatorname{Re}(a_k^2)/(4\tau)}$, which is the stated bound.
\end{proof}

\begin{lemma}\label{app:M2}
For $\tau\le0.0198$ and every $\xi\in[qZ,Z]$, with $r^2=\xi^2e^{-\tau}$:
$\ M_2:=\max_{[qZ,Z]}|f''|\le7w^4$.
\end{lemma}

\begin{proof}
$f''(z)=2C'(z^2)+4z^2C''(z^2)$, and differentiating Lemma~\ref{app:rep} under the integral,
$|C^{(n)}(y)|\le(2\pi r^{2n})^{-1}\int_0^{2\pi}|\Theta^{(n)}(t)|\,(x;q)_\infty^{-1}d\theta$ at
$\rho=y/r^2=e^\tau$, which zeroes $(\log\rho-\tau)^2$. The key step is the exact unfolding: as
$\theta$ runs over $[0,2\pi)$ the variable $u=\pi-2\theta$ runs over two periods, consecutive
branches carrying opposite signs of $\sin(u/2)$, so
\[
 \frac1{2\pi}\int_0^{2\pi}e^{w_r\cos\theta}G(\pi-2\theta)\,d\theta
 =\frac1{2\pi}\int_{\R}\cosh\bigl(w_r\sin\tfrac u2\bigr)e^{-u^2/(4\tau)}\,du ,
\]
an identity; the derivative weights unfold with it, becoming polynomials in $u$, and
$|\sin(u/2)|\le|u|/2$. The even weight gives an exact Gaussian moment,
$\int v^2e^{-(v-s)^2/(4\tau)}dv=\sqrt{4\pi\tau}(2\tau+s^2)$ with $s=\tau w_r\le\sqrt{2\tau}e^{\tau/2}$,
and $\sqrt{\pi/\tau}\cdot\sqrt{4\pi\tau}=2\pi$ cancels the prefactors; the odd weight does
\emph{not}, giving instead
$\int|v|e^{-(v-s)^2/(4\tau)}dv=\sqrt{4\pi\tau}\,s\,\mathrm{erf}\bigl(s/(2\sqrt\tau)\bigr)
+4\tau e^{-s^2/(4\tau)}$, and the error function and the exponential are kept, not discarded.
Writing $s=\tau w_r$, the two assembled factors are
\[
 Q_1=\tau^{-1/2}+\frac{s\,\mathrm{erf}\bigl(s/(2\sqrt\tau)\bigr)+2\sqrt{\tau/\pi}\,e^{-s^2/(4\tau)}}{2\tau},
 \quad
 Q_2=\frac{2\tau+s^2}{4\tau^2}+\tau^{-3/2}\Bigl[\text{same bracket}\Bigr]+\frac1\tau ,
\]
the prefactor $\sqrt{\pi/\tau}$ of Lemma~\ref{app:theta} cancelling against $\sqrt{4\pi\tau}$ in
every term. On the whole range, $\sup\sqrt\tau\,Q_1=1.8273$ and $\sup\tau\,Q_2=3.6596$, both
attained at $\tau=0.0198$, $\xi=Z$. With the weights of
Lemma~\ref{app:theta} this assembles, for $\tau\le0.0198$ (so $r\le0.1981<\tfrac12$, $K_r\le0.631$,
$e^{\tau w_r^2/4}\le e^{0.506}$), to
\[
 |C'(y)|\le e^{K_r-\tau+\tau w_r^2/4}\,\bigl(1.83\,\tau^{-1/2}\bigr)r^{-2},\qquad
 |C''(y)|\le e^{K_r-2\tau+\tau w_r^2/4}\,\bigl(3.670/\tau\bigr)r^{-4}.
\]
With $4\xi^2/r^4\le(2/\tau)e^{3\tau}(1+2\tau)$ the total $M_2\tau^2$ is at most $25.04$ at
$\tau=0.0198$, decreasing to $20.2$ at $\tau=1.29\cdot10^{-4}$, in both cases below the target
$7w^4\tau^2=28$. Those two numbers come from pairing the largest $K_r$ and the largest
$\tau w_r^2/4$, both attained at $\xi=Z$, with the largest $r^{-2}$, attained at $\xi=qZ$; taking
instead the supremum over $\xi\in[qZ,Z]$ of the assembled bound at a single $\xi$ gives $23.33$ at
$\tau=0.0198$, $20.06$ at $\tau=1.29\cdot10^{-4}$ and $19.84$ as $\tau\to0$, so the true margin
against $28$ is $16.7\%$ rather than the $10.6\%$ the factored form reports. Every constant in this
proof, and the two $\xi$-suprema, are regenerated by \texttt{annulus\_certificate.py} at two
precisions; see Remark~\ref{rem:M2cert}.
\end{proof}

\begin{remark}[what the certificate checks, and what it found]\label{rem:M2cert}
\texttt{annulus\_certificate.py} recomputes, at $30$ and again at $60$ digits with zero
disagreement: the window geometry ($\sup r=0.19802$, $\sup K_r=0.62963$, $\sup\tau w_r^2/4=0.50497$
against the quoted $\tfrac12$, $0.631$, $0.506$); the excess constant $2/(\pi e)=0.234199$ of
Lemma~\ref{app:theta}; Lemma~\ref{app:euler} against the true $|(x;q)_\infty^{-1}|$ on a
$(\tau,\xi,\theta)$ grid; the unfolding identity of this proof against numerical quadrature; the
two moment suprema above; the assembled $M_2\tau^2$ on a $400$-point $\tau$ grid with the
$\xi$-supremum taken by search rather than at an endpoint; and the chain of
Proposition~\ref{prop:selower} below. It also checks that the bound dominates the true
$\max|f''|$ where both are computable, the ratio $\max|f''|/(\text{bound})$ coming out between
$0.023$ and $0.045$ at $\tau=0.0198,\,0.0127,\,5\cdot10^{-3},\,10^{-3}$, consistent with
Remark~\ref{rem:M2sharp}. Two constants of an earlier version of this proof did not survive it and
are corrected above: $\sup\sqrt\tau\,Q_1=1.82733$ is not bounded by $1.825$, and
$2/(\pi e)=0.234199$ is not bounded by $0.234$. Both corrections move $M_2\tau^2$ by less than
$0.01$ and neither affects any conclusion. The arithmetic model is mpmath \texttt{mpf} at the two
stated precisions, not interval arithmetic.
\end{remark}

\begin{remark}[how lossy Lemma~\ref{app:M2} is, and why one sample cannot say]\label{rem:M2sharp}
The ratio $M_2/w^4$ does not converge: $f$ is the $q$-cosine, so
$f''(z)\approx-(W/z_0)^2\cos(zW/z_0)$ with $(W/z_0)^2=e^{-\tau}/(\tau(1-q))\sim w^4/4$, while the
window $[qZ,Z]$ spans a phase of only $(1-q)Z/\tau\sim\sqrt{2\tau}$ radians. As $\tau$ varies the
window slides through the cosine and $M_2/w^4$ sweeps $[0,\tfrac14)$; a single value of $\tau$
therefore reports an arbitrary point of that sweep and says nothing about the maximum. The maximum
is attained where the window covers an extremum of the phase, i.e.\ along the extreme-phase points
$w=m\pi$, and there $M_2/w^4\uparrow\tfrac14$: measured $0.2499993$ at $w=120\pi$, against
$0.22599$ at the top of the range $\tau=0.0198$ and $0.02337$ at $\tau=0.0163$. Over the
$165$ values of $\tau$ scanned, $M_2<\tfrac14w^4$ (\texttt{m2\_certificate.py}, mpmath at two
precisions, agreeing to ten digits at every point; a scan, not a supremum proof). Lemma~\ref{app:M2}
is therefore lossy by a factor of about $28$ at the top of the range, and its target $7w^4$ clears
the measured values by about $31$.
\end{remark}

\begin{proof}[Proof of Proposition~\ref{prop:selower}]
At a pole $f(Z)=0$, so the Koornwinder--Swarttouw rules give the exact facts $f(qZ)=qZ\,s(Z)$ and
$s(qZ)=s(Z)$; the mean value theorem on $[qZ,Z]$ then pins
$f'(\zeta)=-q\,s(Z)/(1-q)$ at some $\zeta\in(qZ,Z)$, and a second mean value on
$[z_0,Z]\subset[qZ,Z]$ gives
\[
 |S_e|=(Z-z_0)\,|f'(\tilde\xi)|\ \ge\ \frac{q(1-\sqrt q)}{1-q}\,Z\,|s(Z)|-(1-\sqrt q)(1-q)Z^2M_2 .
\]
The sine factor is bounded below by an invariant of Casoratian ($q$-Wronskian) type for the
Hahn--Exton pair; the closed-form $q$-Wronskian is due to Swarttouw \cite{Swarttouw1992} (
eq.~(4.3.8)); see also \cite[eq.~(4.21)]{Daalhuis1994}. Concretely, with $c_n=\cos(Zq^n;q^2)$, $s_n=\sin(Zq^n;q^2)$
and $F_n=c_n^2-Zq^{n+1}c_ns_n+q\,s_n^2$, the rules give the exact telescoping
$F_n-F_{n+1}=-Zq^{n+1}(1-q)c_{n+1}s_{n+1}$ with $F_n\to1$; the form has smallest eigenvalue
$\lambda_n\ge q-Z/2$ (by $\sqrt{a^2+b^2}\le a+b$), $\sum_{n\ge0}Zq^{n+1}(1-q)=qZ$, and the discrete
Gronwall gives $F_0\ge e^{-0.61Z}\ge1-0.61Z$ for $Z\le\min(0.2,q/3)$; at the pole $F_0=q\,s(Z)^2$,
so $|s(Z)|\ge\sqrt{1-0.61Z}$. (The Gronwall sum is $qZ/(2q-Z)$, whose supremum over the range is
$0.5568\,Z$, so $0.61$ has room; $Z\le\min(0.2,q/3)$ holds throughout, $Z=0.19999$ at
$\tau=0.0198$ and $Z=0.016063$ at $\tau=1.29\cdot10^{-4}$.) Inserting Lemma~\ref{app:M2}, the error fraction is
$28\sqrt2\,(1.02)\sqrt\tau\le0.46$ for $\tau\le1.29\cdot10^{-4}$, a range containing every pole with
$m\ge41$, whence $|S_e(q_m)|\ge(1-0.02-0.46)\sqrt{\tau_m/2}\ge0.35\sqrt{\tau_m}$ there. Evaluated
rather than rounded, the same chain gives $0.3856\sqrt{\tau_m}$ at the threshold and more below it,
the exact error fraction being $0.4520$ against the quoted $0.46$
(\texttt{annulus\_certificate.py}). The tabulated poles are verified numerically at $60+$ digits
(\texttt{qtrig\_verify.py}, mpmath, no interval arithmetic:
$|S_e|/\sqrt\tau\in[0.698,0.708]$ as asserted by the script; the tabulated set is poles
$2,\dots,60$ in the indexing of \S2 (the dominant pole is not among them, its value $0.581$
coming from \texttt{star\_gate\_certificate.py}) and the table omits poles $53$ and $60$ (rows
$52,59$), whose $\tau<7.7\cdot10^{-5}$ lies inside the proved range), and the two ranges overlap on
every tabulated pole with $m\ge41$; only poles $2,\dots,40$ (table rows $1,\dots,39$) are
load-bearing.
The polynomial identities underlying the telescoping and the pole reduction are machine-checked
(\texttt{GramCasoratian.lean}: \texttt{F\_shift}, \texttt{F\_drift}, \texttt{pole\_qsine},
\texttt{sigma1\_term}; each audits inside the standard three, the first two on
\texttt{propext},~\texttt{Quot.sound} alone; no \texttt{sorry}).
\end{proof}

\noindent On the intermediate window $1.29\cdot10^{-4}<\tau_m\le5\cdot10^{-3}$ the bound also
follows, with margin, from Lemma~\ref{app:Se} ($0.63>0.35$), independently of any enumeration of
the poles there.

\section{Proof of the amplitude estimate: the gate holds at every travel pole}\label{app:star}

This appendix proves the estimate $(\star)$ of Theorem~\ref{thm:U} (at every travel pole,
$|P_{12}(q_m)|\,w_m^3<2$, $w_m=\sqrt{2/\tau_m}$; quantitatively
$|P_{12}|\le\tfrac{3.5\sqrt2}8\,\tau^{3/2}<0.619\,\tau^{3/2}$ for $\tau\le5\cdot10^{-3}$, and
$|P_{12}|w^3\le0.75$ at the six poles above that threshold) and then closes the full gate
\eqref{eq:gate}: the same certificate yields the denominator bound $|S_e|\ge0.63\sqrt\tau$
(Lemma~\ref{app:Se}), hence $s<1$ at every travel pole, and, granting the transfer model
\textup{(M)} for the closed form of $b_0$ (Corollary~\ref{cor:dict}), also $b_0>0$ there
(Corollary~\ref{app:gate}). Nothing in this appendix depends on \textup{(M)} except the statement
just named as depending on it.
Throughout, $h=\tau/2$, $\sigma^2=\tau/2$, $\psi(t)=\cos(Ze^{t};q^2)$; by
Proposition~\ref{prop:P12second}, $P_{12}=-\tfrac12[\psi(h)+\psi(-h)]$ and $\psi(0)=0$, so
\begin{equation}\label{eq:taylorgate}
 |P_{12}|\;\le\;\tfrac{h^2}2\max_{|t|\le h}|\psi''(t)| .
\end{equation}

\subsection*{The gate block is a second difference at a zero}

The following exact identity is not needed for Theorem~\ref{thm:V}, but it is the sharpest
available description of the object the $U$-gate must bound, and it makes the size
$\tau^{3/2}$ structural rather than numerical.

\begin{lemma}[the gate block in closed form]\label{lem:P12closed}
Identically in $0<q<1$, with $c(u)=\cos(u;q^2)$ and $\mathfrak s(u)=q^{-1/2}\sin(q^{1/2}u;q^2)$:
\[
 P_{12}=\tfrac{qZ}2\,\mathfrak s(Z)-c(z_0),\qquad
 S_o:=\sum_{j\ge0}\frac{(-1)^j[2(1-q)]^j\,q^{j^2+2j}\,(1-q)}{(q;q)_{2j+1}}=\tfrac{z_0}2\,\mathfrak s(z_0),
\]
$S_o$ being the odd bulk block of the closed form $b_0=\tfrac{2q}{1-q}\,S_o/S_e$ of
Corollary~\ref{cor:dict}.
\end{lemma}

\begin{proof}
From \eqref{eq:p12closed}, $P_{12}=\tfrac{2q^3}{1-q^3}\sum_kd_k$ with
$d_k=(-2)^k(1-q)^kq^{k^2+3k}/[(q^2;q^2)_k(q^5;q^2)_k]$; since
$(q;q)_{2k+3}=(q;q^2)_{k+2}(q^2;q^2)_{k+1}$ and $(q;q^2)_{k+2}=(1-q)(1-q^3)(q^5;q^2)_k$, this is
\begin{equation}\label{eq:p12banked}
 P_{12}=\sum_{k\ge1}(-1)^{k-1}\,\frac{[2(1-q)]^k\,q^{k^2+k+1}\,(1-q^{2k})}{(q;q)_{2k+1}} .
\end{equation}
Since $q^{1/2}Z=z_0$ and $z_0^2=2(1-q)$, the $j$-th summand of
$\tfrac{qZ}2\mathfrak s(Z)=\tfrac{z_0}2\sin(z_0;q^2)$ is
$(-1)^j(1-q)[2(1-q)]^jq^{j^2+j}/(q;q)_{2j+1}$, and that of $c(z_0)$ is
$(-1)^j[2(1-q)]^jq^{j^2+j}/(q;q)_{2j}$. Their difference carries the factor
$(1-q)-(1-q^{2j+1})=q^{2j+1}-q$: the $j=0$ terms cancel, and what remains is
\eqref{eq:p12banked} term by term. The $S_o$ identity is the same computation one half-step down:
the $j$-th summand of $\tfrac{z_0}2\mathfrak s(z_0)$ is
$(-1)^j(1-q)[2(1-q)]^jq^{j^2+2j}/(q;q)_{2j+1}$, which is $S_o$'s summand verbatim. All three
matchings are exponent bookkeeping of the kind machine-checked in \texttt{QTrigIdentities.lean};
the identities are verified off-pole (hence as identities) and at the six exceptional poles to
$10^{-100}$ (\texttt{star\_gate\_certificate.py}); the collapse identities behind
Proposition~\ref{prop:P12second} are further verified at the $58$ tabulated poles
(\texttt{halfstep\_verify.py}).
\end{proof}

\begin{proposition}\label{prop:P12second}
At every travel pole, with $Z=z_0/\sqrt q$ the corresponding zero of $\cos(\,\cdot\,;q^2)$,
\[
 \boxed{\;P_{12}\;=\;-\tfrac12\bigl[\cos(z_0;q^2)+\cos(z_0/q;q^2)\bigr].\;}
\]
Since $z_0=Ze^{-\tau/2}$ and $z_0/q=Ze^{+\tau/2}$, and since $\cos(Z;q^2)=0$, the right-hand side is
exactly the symmetric second difference of $\cos(\,\cdot\,;q^2)$ at its own zero, taken in the
variable $\log u$ with step $\tau/2$.
\end{proposition}

\begin{proof}
Write $c(u)=\cos(u;q^2)$ and $\mathfrak s(u)=q^{-1/2}\sin(q^{1/2}u;q^2)$. The shift rules and the
pole invariant are Lemma~\ref{lem:halfstep}:
$c(qu)=c(u)+q^{3/2}u\,\mathfrak s(\sqrt q\,u)$, $\mathfrak s(qu)=\mathfrak s(u)-u\,c(\sqrt q\,u)$,
and $q^{3/2}\mathfrak s(z_0)\mathfrak s(Z)=1$ at every travel pole.
Taking $u=Z/\sqrt q$ in the first shift rule gives $c(z_0)-c(z_0/q)=qZ\,\mathfrak s(Z)$. Taking
$u=Z$ in the Hahn--Exton $q$-Wronskian $\mathfrak s(u)c(qu)-c(u)\mathfrak s(qu)=u$ and in the first
shift rule, and using $c(Z)=0$ twice, gives
$\mathfrak s(Z)\,c(qZ)=Z$ and $c(qZ)=q^{3/2}Z\,\mathfrak s(z_0)$, whence again the pole invariant
and $c(qZ)=Z/\mathfrak s(Z)$. Substituting the closed
form $P_{12}=\tfrac{qZ}2\mathfrak s(Z)-c(z_0)$ of Lemma~\ref{lem:P12closed} into
$c(z_0)+c(z_0/q)=2c(z_0)-qZ\,\mathfrak s(Z)$ yields the boxed identity.
\end{proof}

\noindent The two algebraic steps of the proof are machine-checked
(\texttt{SelowerSkeleton.lean}): \texttt{shift\_at\_pole} is the instantiation of the half-step rule
at $u=z_0/q$, written so that the index bookkeeping is explicit, and
\texttt{P12\_second\_difference} is the elimination; both are stated in doubled form over an
arbitrary commutative ring and depend only on \texttt{propext}. The identity is
\emph{verified} at all $58$ tabulated poles to a worst relative deviation of
$2\cdot10^{-44}$ (\texttt{halfstep\_verify.py}). The consequence, \eqref{eq:taylorgate} above, is that the gate is a bound on a second
derivative on a window of width $\tau$, not an asymptotic evaluation.

\begin{lemma}[The travel poles are locally finite]\label{lem:polefinite}
For every $\tau_\star>0$ there are only finitely many travel poles with $\tau_m>\tau_\star$.
\end{lemma}

\begin{proof}
By Proposition~\ref{prop:qtrig} the travel poles are the zeros in $q\in(0,1)$ of
$g(q):=\cos(Z;q^2)=C(Z^2)$ with $Z^2=2(1-q)/q$, i.e.\ of
\[
 g(q)=\sum_{j\ge0}(-1)^j\,\frac{2^j\,q^{j^2}(1-q)^j}{(q;q)_{2j}} .
\]
First, $g$ has no zero near $q=0$. For $0<q\le\tfrac1{10}$ one has
$(q;q)_{2j}\ge\prod_{i\ge1}(1-q^i)\ge1-\tfrac q{1-q}\ge\tfrac89$ and $(1-q)^j\le1$, so
$|g(q)-1|\le\tfrac98\sum_{j\ge1}2^jq^{j^2}\le\tfrac98\bigl(2q+4q^4/(1-q)\bigr)\le0.24<1$, whence
$g(q)>0$ there.

Second, $g$ is real-analytic on $(0,1)$: the series converges locally uniformly, since on
$q\le1-\delta$ the $j$-th term is bounded by $2^jq^{j^2}/(q;q)_\infty$ and $q^{j^2}$ dominates $2^j$.
Hence the zeros of $g$ are isolated in $(0,1)$, and any zero set with an accumulation point in the
\emph{open} interval would force $g\equiv0$, contradicting $g(\tfrac1{10})>0$.

Now fix $\tau_\star>0$. The poles with $\tau_m>\tau_\star$ lie in $q\in[\tfrac1{10},e^{-\tau_\star}]$,
a compact subinterval of $(0,1)$. An analytic function on a domain whose zeros are isolated has only
finitely many zeros on any compact subset; hence there are finitely many such poles.
\end{proof}

\begin{remark}
Only local finiteness is proved here; that the poles \emph{do} accumulate at $q=1$
(equivalently $\tau_m\downarrow0$) is Lemma~\ref{lem:infpoles}. The present lemma says merely that
they cannot accumulate anywhere inside $(0,1)$, so that any threshold $\tau_\star$ leaves a finite
set to be checked directly. It is used four times: to make the splice in
Proposition~\ref{prop:selower} well posed, to license the phrase ``for all large $m$'' in
Theorem~\ref{thm:V}, to make the exceptional set of Theorem~\ref{thm:star} finite, and in
Lemma~\ref{lem:infpoles}'s accumulation clause (its proof also shows $g>0$ on $(0,\tfrac1{10}]$,
so every travel pole has $q_m>\tfrac1{10}$).
\end{remark}

\begin{lemma}\label{app:rep2}
For $0<\rho<1$, with $\theta\sim N(\pi/2,\tau/2)$:
$\ \mathbb{E}[(\rho e^{i\theta};q)_\infty^{-1}]=\sum_{M\ge0}q^{M^2/4}(i\rho)^M/(q;q)_M=:\Phi(\rho)$, and
$\operatorname{Re}\Phi(\rho)=\cos(\rho/\sqrt q;q^2)$. The sum is the $q$-Airy function of
Kajiwara--Masuda--Noumi--Ohta--Yamada \cite{KMNOY2004} at base $\sqrt q$, and the representation is Zhang's
$q$-Laplace transform of the second kind at level $\tfrac12$ \cite{Zhang2000}; only the
probabilistic packaging is ours.
\end{lemma}

\begin{proof}
Euler's identity $(x;q)_\infty^{-1}=\sum_Mx^M/(q;q)_M$ integrated termwise (absolute uniform
convergence on $|x|=\rho<1$), with $\mathbb{E}[e^{iM\theta}]=i^Mq^{M^2/4}$; splitting $M$ by parity and matching exponents
($(2j+1)^2=4(j^2+j)+1$; Lean \texttt{parity\_exponent}) identifies $\operatorname{Re}\Phi$ with the
cosine series of Proposition~\ref{prop:qtrig}.
\end{proof}

\begin{lemma}\label{app:cum}
Write $A(\xi)=(\rho e^{i(\pi/2+\xi)};q)_\infty^{-1}$. Then
$\log A(\xi)=\sum_{k\ge1}(i\rho)^ke^{ik\xi}/(k(1-q^k))$; the cumulants
$\Lambda_n=i^n\sum_kk^{n-1}(i\rho)^k/(1-q^k)$ satisfy
$|\Lambda_n|\le(n-1)!\,\rho/((1-q)(1-\rho)^n)$; Gaussian integration by parts gives, at radius
$\rho(t)=z_0e^{t}$,
\[
 \psi'(t)=\tfrac2\tau\,\operatorname{Im}\,\mathbb{E}[\xi A],\qquad
 \psi''(t)=\tfrac4{\tau^2}\,\operatorname{Re}\,\mathbb{E}[(\tfrac\tau2-\xi^2)A];
\]
and for complex $a,b$ with $\operatorname{Re}(1-b\sigma^2)>0$,
$\mathbb{E}[\xi^me^{a\xi+b\xi^2/2}]=\beta^{-1/2}e^{a^2\sigma^2/(2\beta)}M_m$, $\beta=1-b\sigma^2$,
$M_0=1$, $M_1=\mu=a\sigma^2/\beta$, $M_m=\mu M_{m-1}+(m-1)(\sigma^2/\beta)M_{m-2}$.
\end{lemma}

\begin{proof}
$-\log(1-u)=\sum u^m/m$ summed over $u=\rho e^{i\theta}q^j$; $1-q^k\ge1-q$ and a binomial estimate;
$\rho\partial_\rho A=-i\partial_\theta A$ with $\mathbb{E}[f']=\mathbb{E}[f\xi/\sigma^2]$ once and twice; completion of
the square and analytic continuation in $(a,b)$. The exponent bookkeeping is machine-checked
(\texttt{QTrigIdentities.lean}).
\end{proof}

\begin{lemma}\label{app:psi1}
At every travel pole with $\tau\le5\cdot10^{-3}$ and every $t\in[-\tau,h]$:
$|\psi'(t)|\le1.10\,w$.
\end{lemma}

\begin{lemma}\label{app:sZ}
At every travel pole with $\tau\le5\cdot10^{-3}$, writing
$\mathfrak s(u)=q^{-1/2}\sin(q^{1/2}u;q^2)$: $\ \cos(qZ;q^2)=Z/\mathfrak s(Z)$, and
$|\mathfrak s(Z)|\ge Z/(\tau\max_{-\tau\le t\le0}|\psi'|)\ge0.90$. Moreover, for $t\in[-\tau,0]$,
$|\Phi(\rho(t))|\ge q^{3/4}\bigl(|\mathfrak s(Z)|-0.49\,\tau\bigr)$.
\end{lemma}

\begin{proof}[Proof of Lemmas \ref{app:psi1} and \ref{app:sZ}]
The identity is the pole reduction underlying Proposition~\ref{prop:P12second} (the Hahn--Exton
$q$-Wronskian at the zero, cited there). The bound in Lemma~\ref{app:sZ} is the mean value theorem
on $[-\tau,0]$, where $qZ=Ze^{-\tau}$ exactly and all radii satisfy $\rho(t)\le z_0$. For the last
bound: $\operatorname{Im}\Phi(\rho)=q^{3/4}\,\mathfrak s(\rho/\sqrt q)$ (the odd half of the parity
split in Lemma~\ref{app:rep2}), and the drift of $\mathfrak s(\rho(t)/\sqrt q)$ over $t\in[-\tau,0]$
is the even/odd $G_4$ term of the certificate, $\le0.49\,\tau$ (\texttt{budget.py},
\texttt{G4\_coef}). For
Lemma~\ref{app:psi1}: the principal term is the $m=1$ tilted moment of Lemma~\ref{app:cum} at
$a=\Lambda_1$, $b=\Lambda_2$, at most
$e^{c_1\sqrt\tau}|\beta|^{-1/2}L_1(\operatorname{Re}\beta)^{-1}\cdot w\le1.088\,w$, where $L_1$
is the interval constant bounding $|\Lambda_1|/w$ (Lemma~\ref{app:cum}, $n=1$) and
$c_1\le0.177$ bounds the exponent residual per $\sqrt\tau$ (the two leading halves of the exponent
cancel exactly; every residual carries a geometric tail); the tail (linear orders $3..10$ with
$\xi$-weighted moments, the product convolution, and the degree-$\ge11$ remainder) is at most
$0.131$ \emph{absolute}, i.e.\ $\le0.0066\,w$ at the top endpoint $w=20$. All constants are
produced by the committed certificate \texttt{code/zeta\_probe/budget.py} (function
\texttt{uniform\_stack}), which works in $30$-digit floating point with the interval constants
realised as endpoint evaluations. That evaluation is \emph{re-certified in directed-rounding
interval arithmetic} by \texttt{budget\_iv.py}: the same stack is evaluated with $\tau$ entering as
an \emph{interval}, over $160$ cells covering $[1.2\cdot10^{-4},5\cdot10^{-3}]$ and $160$ covering
$[10^{-8},1.2\cdot10^{-4}]$ (contiguous, with both outer endpoints clamped exactly), giving
enclosures valid for \emph{every} $\tau$ in each cell; so on
those ranges neither the floating-point arithmetic nor the monotonicity of any atom is assumed. The
certified maxima are $3.349$ and $2.650$, against the budget $3.5$. On the remaining range
$(0,10^{-8}]$, which no finite set of cells can cover, the power-counting argument below (every
term a positive power of $\tau$) is what proves the bound. The companion
\texttt{star\_gate\_certificate.py} works in the floating-point model at $45$--$260$ digits with
explicit truncation control.
\end{proof}

\begin{theorem}[the amplitude estimate $(\star)$]\label{thm:star}
At every travel pole, $|P_{12}(q_m)|\,w_m^3<2$. Quantitatively: for $\tau\le5\cdot10^{-3}$,
$\max_{|t|\le h}|\psi''(t)|\le3.5\,w$, so \eqref{eq:taylorgate} gives
$|P_{12}|\le\tfrac{3.5\sqrt2}8\tau^{3/2}<0.619\,\tau^{3/2}<\tau^{3/2}/\sqrt2$; and the six poles
with $\tau>5\cdot10^{-3}$,
\[
 \tau_1=0.79972,\quad \tau_2=9.049\cdot10^{-2},\quad \tau_3=3.248\cdot10^{-2},\quad
 \tau_4=1.656\cdot10^{-2},\quad \tau_5=1.001\cdot10^{-2},\quad \tau_6=6.701\cdot10^{-3},
\]
satisfy $|P_{12}|w^3=0.745,\,0.530,\,0.510,\,0.505,\,0.503,\,0.502$ by direct evaluation at
$120$-digit precision with explicit truncation control (\texttt{star\_gate\_certificate.py}; the
underlying identities are independently verified by \texttt{halfstep\_verify.py} at the $58$
tabulated poles, $47$--$54$ digits on the five it shares with this list: the dominant pole is
covered by \texttt{star\_gate\_certificate.py} alone). The first is the dominant pole $q_1=1/\beta_2^2$ of
Corollary~\ref{cor:beta}. Enumeration of the exceptional set: every travel pole has
$q_m>\tfrac1{10}$ and the set is finite (Lemma~\ref{lem:polefinite}); the count of six is
established by a derivative-majorant subdivision for $\tau\ge0.03$ (three poles; in the arithmetic
model disclosed above) and needs no enumeration for $\tau\le5\cdot10^{-3}$, where the uniform
certificate covers every pole; on the window $5\cdot10^{-3}<\tau<0.03$ the three listed poles are
located by a $\Delta w\approx5\cdot10^{-4}$ scan (they sit on the $(k+\tfrac12)\pi$ ladder to three
decimals, and $|\cos(Z;q^2)|$ swings to $1.00$ between them with minimum $0.019$ outside the root
cells). The scan is the weakest step of the enumeration and is stated as such; transcendence in
Theorem~\ref{thm:U} does not consume it.
\end{theorem}

\begin{proof}
By Lemma~\ref{app:cum},
$\psi''=\frac4{\tau^2}\operatorname{Re}\bigl[-e^{\Lambda_0}G_Q(\Lambda_2\sigma^4/\beta+\mu^2)
+e^{\Lambda_0}\mathbb{E}[(\tfrac\tau2-\xi^2)e^{Q}(T-1)]\bigr]$, three terms, with
$Q=\Lambda_1\xi+\Lambda_2\xi^2/2$, $G_Q=\mathbb{E}[e^Q]=\beta^{-1/2}e^{\Lambda_1^2\sigma^2/(2\beta)}$,
$\beta=1-\Lambda_2\sigma^2$, $\mu=\Lambda_1\sigma^2/\beta$ (Lemma~\ref{app:cum} at $m=0,1$), and
$T=e^\Omega$, $\Omega=\sum_{n\ge3}\Lambda_n\xi^n/n!$, the exponential of the higher cumulants. The $\mu^2$
term's phase is pinned by the pole:
$|\cos\arg\Phi(\rho(t))|=|\psi(t)|/|\Phi(\rho(t))|\le\max|\psi'|\,|t|/|\Phi|$ with $\psi(0)=0$
(Lemma~\ref{app:psi1}) and $|\Phi|\ge q^{3/4}(|\mathfrak s(Z)|-0.49\tau)$ (Lemma~\ref{app:sZ});
its argument's sine obeys $|\sin2\arg\mu|\le13.3\,\tau^{3/2}$, the two $O(1/w)$ tilts of
$\Lambda_1$ and $\beta$ cancelling to $O(\tau)$ because they are the same Lambert data shifted by
one index. Every remaining quantity is bounded by an interval constant in scaled variables, each
resting on a one-line monotonicity ($(1-e^{-\tau})/\tau$ decreasing; $e^{c\tau}$, $\rho^2=2(1-q)e^\tau$
increasing; $w,q$ decreasing), with tails via $1-q^k\ge(1-q)kq^{k-1}$ and the signed cancellations
via the alternating-series bound $\tau^2/6-\tau^3/12\le1+q-2(1-q)/\tau\le\tau^2/6$. The
degree-$\ge11$ remainder is handled on two intervals: on $[1.2\cdot10^{-4},5\cdot10^{-3}]$ by an
exponential majorant whose profile $e^{0.097/\sqrt\tau}\tau^{4.5}$ is increasing exactly there
(threshold $1.16\cdot10^{-4}$), total $\le3.03\,w$; on $(0,1.2\cdot10^{-4}]$ with no exponential
device (linear remainder geometric with ratio $\tfrac12$ on the cutoff region; all products by
$|e^\Omega-1-\Omega|\le\tfrac12|\Omega|^2e^{|\Omega|}$ with $|\Omega|\le0.45$ on $|\xi|\le t_0\kappa^{-1/3}$; the
intermediate zone suppressed by the Gaussian; here $t_0=1-\rho$ and $\kappa=\rho/(1-q)$),
total $\le2.10\,w$. Both are endpoint evaluations of quantities proved monotone; on
$[10^{-8},5\cdot10^{-3}]$ the same conclusion is obtained without any monotonicity assumption by
the interval-arithmetic certification \texttt{budget\_iv.py} described above, and the monotonicity
argument is what carries the remaining range $(0,10^{-8}]$. The assembly is the committed script
\texttt{budget.py}.
\end{proof}

\begin{lemma}[the denominator at the pole]\label{app:Se}
At every travel pole, exactly: $S_e=\tfrac{qZ}2\,\mathfrak s(Z)-P_{12}$. Consequently, for
$\tau\le5\cdot10^{-3}$,
\[
 |S_e|\;\ge\;0.90\cdot\tfrac{qZ}2-|P_{12}|\;\ge\;
 \Bigl[\tfrac{0.90}{\sqrt2}\sqrt{q\,v}-0.619\,\tau\Bigr]\sqrt\tau\;\ge\;0.63\,\sqrt\tau,
 \qquad v=\tfrac{1-q}\tau,
\]
and $S_e$ has the sign of $\mathfrak s(Z)$.
\end{lemma}

\begin{proof}
The identity is Lemma~\ref{lem:P12closed} rearranged, since $S_e=\cos(z_0;q^2)$
(Proposition~\ref{prop:qtrig}). For the bound: $|\mathfrak s(Z)|\ge0.90$ (Lemma~\ref{app:sZ}),
$\tfrac{qZ}2=\sqrt{q(1-q)/2}=\sqrt{\tau/2}\,\sqrt{qv}$, $|P_{12}|\le0.619\,\tau^{3/2}$
(Theorem~\ref{thm:star}), and $\sqrt{qv}\ge\sqrt{(1-\tau)(1-\tau/2)}\ge1-\tfrac45\tau$ for
$\tau\le\tfrac12$ (since $(1-\tfrac45\tau)^2\le(1-\tau)(1-\tfrac\tau2)$ there); at
$\tau=5\cdot10^{-3}$ the bracket is $\ge0.630$ (exact endpoint value $0.6309$) and it increases as
$\tau$ decreases. The sign
statement holds because $0.90\cdot\tfrac{qZ}2>|P_{12}|$ with the same margin.
\end{proof}

\begin{corollary}[the gate closes]\label{app:gate}
At every travel pole the gate \eqref{eq:gate} holds; with Theorem~\ref{thm:RJ} this gives
$\Pi_1(q_m),\Pi_q(q_m)>0$ for $\tau_m\le5\cdot10^{-3}$ (the enumeration
standing of the six-pole list is that of Theorem~\ref{thm:star}). The half $|s|<1$ is proved here
outright; the half $b_0>0$ consumes the closed form $b_0=\tfrac{2q}{1-q}S_o/S_e$, which is
Corollary~\ref{cor:dict} and therefore rests on the transfer model \textup{(M)},
Theorem~\ref{thm:model}. Quantitatively,
for $\tau\le5\cdot10^{-3}$:
\[
 |s|\;=\;\frac{q}{1-q}\,\frac{|P_{12}|}{|S_e|}\;\le\;\frac{\tau}{e^\tau-1}\cdot\frac{0.619}{0.63}
 \;\le\;0.983\;<\;1,
\]
and $b_0>0$: with the closed form $b_0=\tfrac{2q}{1-q}\,S_o/S_e$ of Corollary~\ref{cor:dict} and
$S_o=\tfrac{z_0}2\,\mathfrak s(z_0)$ (Lemma~\ref{lem:P12closed}), the pole invariant
$q^{3/2}\mathfrak s(z_0)\mathfrak s(Z)=1$ (Proposition~\ref{prop:P12second}'s proof) gives
$\operatorname{sign}S_o=\operatorname{sign}\mathfrak s(Z)=\operatorname{sign}S_e$
(Lemma~\ref{app:Se}), so $S_o/S_e>0$. At the six poles with $\tau>5\cdot10^{-3}$ the gate
quantities are evaluated directly (\texttt{star\_gate\_certificate.py}):
$|s|=0.296,\,0.256,\,0.252,\,0.251,\,0.251,\,0.250$ and
$b_0\tau=2.174,\,2.001,\,2.000,\,2.000,\,2.000,\,2.000$; moreover $|S_e|/\sqrt\tau\ge0.58$ at
all six, so the bound $0.35\sqrt\tau$ of Proposition~\ref{prop:selower} holds at the dominant pole,
as its statement records.
\end{corollary}

\subsection*{Machine verification of the $q$-trigonometric layer}

Every series identity of the $q$-trigonometric sections and of this appendix is proved by
matching the $j$-th summand; the matching is exponent bookkeeping, which is where such arguments
fail in practice. That bookkeeping, and the Casoratian constancy behind the half-step invariant,
are machine-checked in Lean~4 (\texttt{QTrigIdentities.lean}; the last four rows of the table
below live in \texttt{SelowerSkeleton.lean}); convergence and the interchange of summation are
analytic and are not formalised.

\begin{center}\small
\begin{tabular}{lll}
\hline
Paper step & Lean declaration & axioms \\
\hline
half-step rule for $c$, exponent match & \texttt{cshift\_exponent} & \texttt{propext}\\
half-step rule for $s$, exponent match & \texttt{sshift\_exponent} & \texttt{propext}\\
parity split of $\Phi$ & \texttt{parity\_exponent} & \texttt{propext}\\
functional equation of $F$ & \texttt{qairy\_exponent} & \texttt{propext}\\
$P_{12}$ $k$-sum split, low branch & \texttt{P12\_split\_lo} & \texttt{propext}\\
$P_{12}$ $k$-sum split, high branch & \texttt{P12\_split\_hi} & \texttt{propext}\\
Casoratian constancy & \texttt{casoratian\_re\_const} & standard\\
constancy $\Rightarrow$ every rung & \texttt{casoratian\_re\_eq\_zero} & standard\\
shift rule at the pole & \texttt{shift\_at\_pole} & \texttt{propext}\\
$P_{12}$ as a second difference & \texttt{P12\_second\_difference} & \texttt{propext}\\
scaling $\tau^2w=\sqrt2\,\tau^{3/2}$ & \texttt{tau\_sq\_mul\_w\_sq} & standard\\
gate threshold $2/w^3=\tau^{3/2}/\sqrt2$ & \texttt{gate\_threshold} & standard\\
\hline
\end{tabular}
\end{center}

The algebraic skeleton of Proposition~\ref{prop:selower} is likewise machine-checked
(\texttt{SelowerSkeleton.lean}): \texttt{pole\_facts} derives the two exact facts
$s(qZ)=s(Z)$ and $c(qZ)=qZ\,s(Z)$ from the two shift rules together with $c(Z)=0$;
\texttt{F\_at\_pole} collapses the invariant to $q\,s(Z)^2$ at a pole;
\texttt{form\_lower\_bound} is the quadratic-form estimate behind $\lambda_n\ge q-Z/2$; and
\texttt{prod\_one\_sub\_ge} is the Weierstrass bound $\prod(1-\delta_i)\ge1-\sum\delta_i$ used in the
discrete Gronwall step. The analytic inputs (the annulus representation, the Poisson envelope and
the mean value theorem) are hypotheses there, not formalised.

\noindent ``standard'' means \texttt{propext}, \texttt{Classical.choice}, \texttt{Quot.sound}. Every
file named in this paper lives in the Mathlib-based package of the accompanying development; none of
its $62$ files contains \texttt{sorry}, each ends with a \verb|#print axioms| audit, and every
declaration cited anywhere in this paper appears in one of those audits. (The audits are selective
in the other direction: a file may prove auxiliary lemmas that its audit block does not list.) Some
declarations report ``does not depend on any axioms'', which is the legitimate output
for a goal closed by \texttt{decide} or \texttt{rfl} (for instance
\texttt{OverDetermination.linear\_order\_19\_ok}) and is also, character for character, what an
audit prints for a declaration that failed to elaborate; the audits quoted here were therefore read
from a cold elaboration of each file one at a time (\texttt{lake env lean}, exit status checked),
not from a warm \texttt{lake build}, which does not replay the output for an up-to-date library.
Three declarations cited in this paper are proved by \texttt{native\_decide} and carry, besides the
standard axioms, the compiler-trust axiom named at their point of use
(Remark~\ref{rem:growthstatus}; \S\ref{sec:modp}). The realness of the ladder coefficient is what
makes the added term in \texttt{casoratian\_re\_const} purely imaginary, and is carried as an
explicit hypothesis.

\subsection*{The boundary of the formalisation: the analytic layer}\label{sec:leanboundary}

The combinatorial layer of this paper is formalised (\S\ref{sec:sitecost}), the analytic layer is
not, and the obstruction is one and the same for every statement in it: Mathlib carries no
$q$-Pochhammer symbol $(a;q)_n$ or $(a;q)_\infty$, no product form of a theta function, no Bessel
or $q$-Bessel function, no $q$-difference equations, and no steepest-descent, stationary-phase or
van der Corput machinery. It does carry the two-variable Jacobi theta series with its summability,
analyticity and modular functional equation, and the Mellin transform; it does not carry the
Jacobi triple product, which is the identity that connects the two. The list below records, for
each analytic statement of this paper, the object whose absence blocks a formalisation. The
statements are at PROVED, not at VERIFIED, and this table is the reason.

\begin{center}\small
\begin{tabular}{ll}
\hline
Statement & missing prerequisite \\
\hline
Lemma~\ref{app:euler} & $(x;q)_\infty$ and its logarithm \\
Lemma~\ref{lem:P12closed} & $(q;q)_n$; the $q$-cosine and $q$-sine \\
Proposition~\ref{prop:P12second} & $(q;q)_n$; the $q$-cosine \\
Lemma~\ref{app:Se} & $(q;q)_\infty$; the $q$-sine \\
Lemma~\ref{app:rep} & the Jacobi triple product; $(q^2;q^2)_\infty$ \\
Lemma~\ref{lem:polefinite} & analyticity of a $q$-series and its zero set \\
Lemma~\ref{lem:infpoles} & the same, plus a limit at $q\to1$ \\
Lemma~\ref{lem:infzeros} & the same \\
Proposition~\ref{prop:selower} & $(q;q)_\infty$; quantitative $q$-trigonometric bounds \\
Lemma~\ref{app:theta} & a Gaussian majorant for a theta series \\
Lemma~\ref{app:M2} & second derivatives of a $q$-series \\
Lemma~\ref{app:psi1} & derivative bounds for the $q$-sine \\
Lemma~\ref{app:sZ} & the $q$-sine and $q$-cosine at a zero \\
Theorem~\ref{thm:star} & a quantitative Taylor bound on a $q$-phase \\
Lemma~\ref{app:rep2} & the $q$-Airy function; a Gaussian expectation \\
Lemma~\ref{app:cum} & cumulant bounds; Gaussian integration by parts \\
Lemma~\ref{lem:Bbounded} & $\log\Gamma$ with an explicit error term; $\log(q;q)_\infty$ \\
Lemma~\ref{lem:T2abs} & Mellin--Barnes contour shifting; a verified contour-integral enclosure \\
Lemma~\ref{lem:atomN}, confluence & Bessel $J_\nu$ at half-integer order \\
\hline
\end{tabular}
\end{center}

\noindent The first five rows would be reached by a $q$-Pochhammer layer alone, together with the
Jacobi triple product; the next four would be reached by that layer together with Mathlib's
existing complex analysis; the remainder need asymptotic machinery that is a larger absence than
the $q$-series one. The half of Lemma~\ref{lem:atomN} that is an explicit rational inequality
\emph{is} verified (\texttt{AtomN.lean}); the confluence half is the row above
(Remark~\ref{rem:lean}).

\section{The number \texorpdfstring{$\beta_2$}{beta2}: value transcendence}\label{app:beta2}

This appendix concerns the \emph{number} $\beta_2=1/\sqrt{q^\ast}$ of Corollary~\ref{cor:beta},
not the function $V$. Nothing in \S\S\ref{sec:blocks}--\ref{sec:conclusion} or in
Appendices~\ref{app:annulus}--\ref{app:star} depends on any statement made here; the material is
recorded because it locates the obstruction to value transcendence precisely, at the arithmetic
ratio $\rho=\tfrac12$, one notch beyond the proven $\rho\ge1$ theory.

\begin{remark}[the number $\beta_2$: a value-transcendence question]\label{rem:betanumber}
Whether the \emph{number} $\beta_2$ is transcendental is open, but the question has a clean
reduction. Since $\beta_2=1/\sqrt{q^\ast}$ with $\Sigma_1^T(q^\ast)=1$, the number $\beta_2$ is
algebraic iff $q^\ast$ is; if $q^\ast$ were algebraic, the series $\Sigma_1^T$ would take the
\emph{rational} value $1$ at an algebraic interior point of its disc. For the Jacobi-theta class
this is impossible: by Nesterenko's theorem on Eisenstein values \cite{Nesterenko1996},
$\sum_n\alpha^{n^2}$ is transcendental at \emph{every} algebraic $\alpha$ with $0<|\alpha|<1$
\cite{DNNS1996}. The blocks here are theta-\emph{like}; their coefficients ride the squares
$q^{(j+1)^2}$, which is the source of the natural boundary of Theorem~\ref{thm:blocks}; but
their Euler-product denominators $(q;q)_{2k+1}$ spread the support densely, and no analogue of
\cite{Nesterenko1996} is known off the modular class.

Concretely the reduction localizes on a single confluent basic hypergeometric function. Writing
$\Sigma_1^T(q)=1-S(q)$,
\[
   S(q)\;=\;\sum_{k\ge0}\frac{\bigl(-2(1-q)\bigr)^k q^{k^2}}{(q;q)_{2k}}
   \;=\;{}_1\phi_1\!\bigl(0;q;q^2,\,2q(1-q)\bigr),
\]
a $q$-Bessel-class function (using $(q;q)_{2k}=(q;q^2)_k(q^2;q^2)_k$); $\beta_2$ is algebraic iff the
least positive zero $q^\ast=0.449453630\ldots$ of $S$ is. Certified exclusions at $430$ digits
(residual $10^{-442}$; PSLQ with post-verification of every candidate): neither $q^\ast$ nor
$\beta_2$ satisfies an integer polynomial of degree $\le8$ with height $10^{20}$, degree $\le12$
with height $10^{18}$, or degree $\le16$ with height $10^{14}$; no bilinear relation ties
$q^\ast$ to the fold point $q_c$ of the zero curve (height $10^{12}$). The continued fraction of
$q^\ast$ ($383$ certified partial quotients, two-precision guard) is Gauss--Kuzmin generic:
digit frequencies match, Khinchin mean $2.63$ ($2.685$), L\'evy exponent $3.21$ ($3.276$),
largest quotient $a_{170}=8250$ (probability $\approx6\%$ over $383$ draws, consistent with chance).
The number offers no arithmetic foothold of any tested kind. We record what the natural tools give here,
since it locates the difficulty precisely. First, $S$ is \emph{neither} $q$-holonomic \emph{nor} a
Mahler function in $q$ (verified to $O(q^{48})$): the diagonal specialization $z=2q(1-q)$ ties the
argument to the base and destroys the single-variable $q$-difference structure, so the methods of
Mahler--Nishioka \cite{Nishioka1996} and of holonomy do not apply. Second, two elementary refinements fail
\emph{decisively}, not for want of effort: an ultrametric argument is void; the $v$-adic and
archimedean sums of $S(\alpha)$ are unrelated numbers, so $|S(\alpha)|_v=1$ does not bound the real
value; and a Liouville/height bound is beaten by the $q$-Pochhammer denominators, whose growth over
the conjugates of modulus $>1$ is cubic ($\sim\!K^3$) against the merely quadratic ($\sim\!K^2$) tail,
Kronecker-sharply (equality would need Mahler measure $1$, i.e.\ a root of unity). A more promising reformulation uses two variables. Set $G(q,z)={}_1\phi_1(0;q;q^2,z)=\sum_k(-1)^k
q^{k(k-1)}z^k/(q;q)_{2k}$, so $S(q)=G(q,2q(1-q))$; unlike $S$, the function $G$ \emph{is} $q$-holonomic
in $z$, satisfying the second-order $q$-difference equation
\[
   q\,G(q,z)\;-\;(q+1-qz)\,G(q,q^2z)\;+\;G(q,q^4z)\;=\;0 .
\]
One checks that $z^\ast:=2q^\ast(1-q^\ast)=0.4948901\ldots$ is the \emph{smallest positive zero} of
$G(q^\ast,\cdot)$. Hence $\beta_2$ algebraic would force an \emph{algebraic zero of a $q$-Bessel-class
function at an algebraic base}; and the classical analogue is false: the zeros of the Bessel function
$J_\nu$ are transcendental (Siegel \cite{Siegel1929}). A $q$-analogue of Siegel's theorem for the zeros of ${}_1\phi_1$
would therefore settle $\beta_2$. The right framework is arithmetic. Writing the equation as a system $Y(q^2z)=A(z)Y(z)$ with
$Y=(G(z),G(q^2z))^{\!\top}$, the companion matrix
\[
   A(z)=\begin{pmatrix}0&1\\-q&\ q+1-qz\end{pmatrix}
\]
is \emph{polynomial} in $z$ with $\Z[q]$ entries and $\det A=q$, so the iterated matrices carry no
denominators and the Casoratian obeys $W(q^2z)=qW(z)$, i.e.\ $W(z)=c\,z^{1/2}$: the two solutions form a
$q$-Bessel pair with Frobenius exponents $0,\tfrac12$. The second solution is explicit,
$\widetilde G(z)=z^{1/2}H(z)$ with $H(z)=\sum_k(-1)^kq^{k^2}(1-q)z^k/(q;q)_{2k+1}$, and the pair satisfies
the exact $q$-Wronskian identity
\[
   q\,G(z)\,H(q^2z)-G(q^2z)\,H(z)\;=\;q-1 ,
\]
so at the zero $z^\ast$ one has the clean relation $G(q^2z^\ast)H(z^\ast)=1-q$. Moreover $G(q,\cdot)$ is
\emph{entire of order zero}; a $q$-$E$-function, not a $q$-$G$-function. Hence $\beta_2$ transcendental would follow from a
$q$-analogue of the Siegel--Shidlovskii theorem for this $q$-$E$-function: such a theorem gives
$G(q,\alpha)$ transcendental at every algebraic $\alpha\ne0$ off the singular locus, whereas
$G(q^\ast,z^\ast)=0$ is algebraic; forcing $q^\ast$, and with it $\beta_2$, transcendental (and by the
shared $q$-connection core, likewise $U$). All the structural hypotheses hold here (polynomial companion
matrix, constant determinant, entire order-zero solution); the sole open point is whether the available
$q$-Siegel--Shidlovskii results \textup{(Andr\'e \cite{Andre2000}; Di~Vizio--Hardouin
\cite{DiVizioHardouin2012}; and the Siegel-method theory of \cite{AMV2007})} extend to a ${}_1\phi_1$ that is
\emph{irregular} at $z=\infty$; a precise, expert-checkable question, not an absence of theory. (Note
the contrast with the \emph{function-level} results of this paper, which say nothing about individual
interior values.)

For \emph{irrationality} (weaker than transcendence) the classical architecture transports verbatim.
$G$ and $H$ are Hahn--Exton (third Jackson) $q$-Bessel functions: $G(q,z)\propto
J_{-1/2}(\sqrt z/q;q^2)$ and $H(q,z)\propto J_{+1/2}(\sqrt{z/q};q^2)$ (verified; the zeros correspond
exactly), i.e.\ the $q$-analogues of $\cos$ and $\sin$; and $z^\ast$ is the $q$-analogue of
$(\pi/2)^2$. Niven's irrationality proof for $\pi$ runs on the Lommel-polynomial ladder
$J_{m+1/2}(r)=R_{m,1/2}(r)J_{1/2}(r)-R_{m-1,3/2}(r)J_{-1/2}(r)$ evaluated at the zero of $J_{-1/2}=
\sqrt{2/\pi r}\cos$; the exact $q$-analogue of every ingredient exists (Koelink--Swarttouw \cite{KoelinkSwarttouw1994}): the ladder
$J_{\nu+1}(x;q^2)=\bigl((1-q^{2\nu})/x+x\bigr)J_\nu-J_{\nu-1}$, the $q$-Lommel connection
$J_{\nu+m}=R_{m,\nu}J_\nu-R_{m-1,\nu+1}J_{\nu-1}$ with $x^mR_{m,\nu}\in\Z[q,x^2]$, and the second-kind
functions carrying $q^{m^2/2}$ prefactors; all verified numerically for our pair. Moreover, since a
rational zero would make $G(q^\ast,z^\ast)=0$ a \emph{rational value} at rational base and argument,
$\beta_2$ is irrational as soon as this Hahn--Exton function omits the value $0$ at rational points: a plain \emph{value-irrationality} statement, strictly weaker than every transcendence target above.
What blocks even this is now a single quantified invariant: the \emph{arithmetic ratio} $\rho$ of the
series ($b$-adic decay exponent of the terms over $b$-adic growth exponent of the cleared denominators,
$q=s/t$ in lowest terms). Every proven $q$-irrationality theorem (Tschakaloff, $q$-exponential, the
B\'ezivin--Bundschuh--V\"a\"an\"anen classes) lives at $\rho\ge1$, where the Pochhammer denominator
itself supplies the decay; our $G$ has an independent $q^{k^2}$ numerator against the double Pochhammer
$(q;q)_{2k}$, giving $\rho=\tfrac12$. All three classical constructions fail by exactly this margin,
measured: the $z$-Pad\'e remainder $q^{3.5n^2}$ against height $q$-degree $\approx9n^2$; the
$z$-direction continued fraction (convergents in $\Z[q,z]$, Pincherle limit the $H$-ratio) converges
only geometrically against $t^{n^2}$ heights; and the $\nu$-ladder at the zero has
$J_{m+1/2}(x_0)=R_{m,1/2}(x_0)J_{1/2}(x_0)\sim Cx_0^m$; geometric, because the $q$-factorial is
geometric where Niven's $m!$ is super-exponential; against $t^{m^2}$. The $\nu$-ladder is in fact
fully costed: the $q$-Lommel coefficients satisfy $\deg_q p_m=m^2$ \emph{exactly} (immediate from the
recurrence, verified symbolically), so their denominators at $q=s/t$ are the pure power $t^{m^2}$: carrying \emph{no} cyclotomic structure, so the lcm-compression that powers the little-$q$-Legendre
proofs has nothing to act on; while by Poincar\'e--Perron \cite{Poincare1885,Perron1909} the $m$-recurrence, whose coefficient tends
to $x_0+1/x_0$ with \emph{distinct} characteristic roots $x_0,1/x_0$, admits \emph{only geometric}
solutions: no theta-type $q^{m^2}$-decaying solution exists in the $\nu$-direction for either kind (the
second kind's $q^{\nu^2/2}$ prefactor is exactly cancelled by the $(q^{-\nu+1};q)_\infty$ blow-up of its
shifted argument). Hence the cleared $\nu$-ladder forms grow like $t^{m^2}x_0^{-m}\to\infty$ for
\emph{every} rational $q=s/t$, including $s=1$. The $b\leftrightarrow z$ transformation
${}_1\phi_1(0;q;q^2,z)=\frac{(z;q^2)_\infty}{(q;q^2)_\infty}{}_1\phi_1(0;z;q^2,q)$ redistributes but
conserves $\rho$. The structural summary: the theta-type decay lives only in the $z$-direction, where
the equation is \emph{irregular} at $z=\infty$ (Newton polygon slopes $\{0,1\}$) and every integral
structure carries the double-Pochhammer weight $\rho=\tfrac12$; the $\nu$-direction is \emph{regular}
(Poincar\'e class) with clean $\Z[q,X]$ arithmetic but only geometric decay. Decay and arithmetic live
in different directions and never meet; the same irregularity that obstructs $q$-Siegel--Shidlovskii
creates the arithmetic deficit. The precise open door remains a $q$-irrationality criterion at
$\rho=\tfrac12$; one notch beyond the proven $\rho\ge1$ theory.
\end{remark}

\begin{remark}[the $\beta_2$ family]\label{rem:family}
$q^\ast$ is the first member of an infinite family. On $(0,1)$ the parabola $w(q)=2q(1-q)$ falls
like $1-q$ while every zero branch falls like $(1-q)^2$ (the confluent limit): the branches are
$z_j(q)=2q(1-q)\lambda_j(q)$ with $\lambda_j$ the eigenvalue branches of
Theorem~\ref{thm:fredholm}, and $w=z_j$ means $\lambda_j=1$. So $w$ overtakes each branch at least
once, distinct branches at distinct points, and $S$ has infinitely many zeros
$q^\ast_1<q^\ast_2<\cdots\to1$, which are the travel poles $q_j$
(Proposition~\ref{prop:infpolesalt}). That $w$ overtakes each branch \emph{exactly} once, so that
the $j$-th zero is the crossing of the $j$-th branch, $w(q^\ast_j)=z_j(q^\ast_j)$, is
Proposition~\ref{prop:simplepoles}\textup{(i)}, at the standing of Hypothesis~\ref{hyp:Tstar}
(verified at $j=2$ to $43$ digits: exactly one zero of $G(q^\ast_2,\cdot)$ lies below $w$). The accumulation obeys the classical cosine-zero law
\[
   1-q^\ast_j\ \sim\ \frac{8}{\pi^2(2j-1)^2},\qquad
   \text{measured }(1-q^\ast_j)\,\frac{(2j-1)^2\pi^2}{8}
   =0.679,\,0.961,\,0.986,\,0.993,\,0.9956,\,0.9970,\,0.9979,\,0.9984\ (j\le8).
\]
The second member is $q^\ast_2=0.913486638731047141696\ldots$, i.e.\
$\beta_2^{(2)}:=1/\sqrt{q^\ast_2}=1.046282352687417855\ldots$; neither satisfies an integer
polynomial of degree $\le8$ with height $10^{10}$, and no bilinear relation ties $q^\ast_1$ to
$q^\ast_2$ (PSLQ, $100$ digits). The family interpolates between the open head ($j=1$, the wild
$\beta_2$) and the classical tail (the $\pi^2$-law): the roots become ``provably transcendental
in the limit'' while the first, the one of combinatorial interest, remains untouched; the
confluent boundary of the step-ratio mechanism made visible inside a single object. (That the
higher roots are genuine singularities of $U$ itself is the non-vanishing $\Pi_q(q_m)\ne0$, which
is Theorem~\ref{thm:RJ}.)

Three refinements (all computations at $115$ digits, $30$ cosine and $10$ sine members). First,
the sine side has its own family, $H(q^{\mathrm s}_j,2q(1-q))=0$ with
$1-q^{\mathrm s}_j\sim2/(j\pi)^2$ (law ratios $0.913,\dots,0.9991$ for $j\le10$), and the two
families \emph{interlace}: $q^\ast_j<q^{\mathrm s}_j<q^\ast_{j+1}$ (verified $j\le9$); its
head is $q^{\mathrm s}_1=0.815032640290\ldots$ Since $H(q,2q(1-q))=(1-q)\Sigma_0(q)/(2q)$, the sine
family is exactly the zero set of $\Sigma_0$ in $(0,1)$, that is, of
$q\mapsto J^{(3)}_{1/2}(z_0(q);q^2)$ with $z_0=\sqrt{2(1-q)}$ (Lemma~\ref{lem:qsinenum} and the
normalisation after Proposition~\ref{prop:qtrig}). Half of the interlacing, at least one sine member
in each $(q^\ast_j,q^\ast_{j+1})$, is Proposition~\ref{prop:signlaws}\textup{(b)}, at the standing
of Hypothesis~\ref{hyp:Tstar}; that there is exactly one is verified only. Second, the accumulation law has an exact-looking
rational second order, identical for both families:
\[
   1-q^\ast_j=\frac{2}{\mu_j}\Bigl(1-\frac{8}{9\mu_j}+\frac{B}{\mu_j^2}
   +O(\mu_j^{-3})\Bigr),\qquad
   \mu_j=\Bigl(\tfrac{(2j-1)\pi}{2}\Bigr)^2\ \text{resp.}\ (j\pi)^2,
\]
with $A=-8/9$ matched to $12$ digits (fit residual $3\cdot10^{-18}$ out of sample) and
$B$ derived below. The constant $A$ is derived, and the derivation separates two regimes that
must not be conflated. Writing $\varepsilon=1-q$, the logarithm of the ratio of the $k$-th term
of $G(1-\varepsilon,\varepsilon^2y)$ to $(-1)^ky^k/(2k)!$ equals
$\tfrac k2\varepsilon+c_2(k)\varepsilon^2+O(\varepsilon^3)$ with
\[
   c_2(k)=-\frac{k^3}{9}-\frac{k^2}{12}+\frac{23k}{72};
\]
to first order $G$ is $\cos\!\sqrt{y\,e^{\varepsilon/2}}$, so at \emph{fixed} $j$ the zero
branch obeys $z_j=\varepsilon^2\mu_j(1-\tfrac12\varepsilon+O(\varepsilon^2))$: the fixed-$j$
confluent correction is $-\tfrac12$ (measured $-0.5004$ at $\varepsilon=0.002$, converging
linearly). The crossing, however, sits in the \emph{uniform} regime $\mu\varepsilon\approx2$,
where the $k^3$-term acts as the operator $-\tfrac{\varepsilon^2}9D^3$, $D=(u/2)\,d/du$ on
$\cos u$, $u=\sqrt{X}$; at the zeros $D^3\cos u=\tfrac{u^3}8\sin u+O(u^2)$, while even powers
of $D$ are $\cos$-weighted and drop out (parity protects the order). With $m=\varepsilon\mu$
the crossing equation $2(1-\varepsilon)e^{\varepsilon/2}=\varepsilon\mu(1-\varepsilon^2\mu/36)$
becomes $m=2-\varepsilon+\varepsilon m^2/36+O(\varepsilon^2)$, whence $m=2-\tfrac89\varepsilon$:
\[
   A=-1+\tfrac19=-\tfrac89 ,
\]
identically for the sine family (its $\varepsilon^2$-coefficient carries the same leading
$-k^3/9$). The two regimes carry different constants ($-\tfrac12$ versus the crossing's
$-\tfrac89$). One more order of the same calculus derives $B$ exactly. The ingredients:
the exact prefactor $2(1-\varepsilon)e^{\varepsilon/2}=2-\varepsilon-\tfrac34\varepsilon^2-\cdots$;
the cubic coefficient $c_3(k)=-\tfrac{k^3}9-\tfrac{k^2}{12}+\tfrac{17k}{72}$ (the $m^3$-terms
cancel in its defining sum, so $c_3$ shares the leading $-k^3/9$); the quartic
$c_4(k)=\tfrac{k^5}{450}+O(k^4)$ together with $\tfrac12c_2(D)^2=\tfrac{D^6}{162}+\tfrac{D^5}{108}+\cdots$;
and one subtlety: at order $\varepsilon^6u^9$ the cube $\tfrac16(c_2(D)\varepsilon^2)^3$
contributes $+\tfrac{u^9}{2239488}\sin u$ through the leading term of $D^9$, cancelling the
$\delta^3/6$ of the cosine expansion \emph{identically}; truncating the exponential before the
cube gives a spurious $m^5$-term. With that cancellation the crossing equation closes as
$m=2-\tfrac89\varepsilon-\tfrac{1891}{8100}\varepsilon^2+O(\varepsilon^3)$, and inversion gives
\[
   B=2m_2+\tfrac{64}{81}=\frac{1309}{4050}=0.3232098765432\ldots,
\]
confirmed by a five-order free fit to the computed roots: $B=0.3232098765429\ldots$ agrees to
$13$ digits (with $A=-8/9$ simultaneously reproduced to $17$), and the $(A,B)$-constrained fit has
out-of-sample residuals below $4\cdot10^{-16}$. The accumulation law of the $\beta_2$ family
therefore begins
\[
   1-q^\ast_j=\frac{2}{\mu_j}\Bigl(1-\frac{8}{9\mu_j}+\frac{1309}{4050\,\mu_j^2}
   +O(\mu_j^{-3})\Bigr),
\]
with both corrections rational and derived. Three closing observations. (a) \emph{Universality}:
the sine family obeys the \emph{same} law; a constrained fit to $22$ sine members gives
$B_{\sin}\cdot4050=1309.000000$ to the fit precision $2\cdot10^{-8}$; so through second order
the two interlaced families share one accumulation series, with only $\mu_j$ distinguishing them.
(b) Every ingredient of the crossing calculus is rational (the $c_n(k)$ are rational polynomials;
$D^n\cos u$ has odd $u$-powers on $\sin$ and even on $\cos$, an invariant of $D=(u/2)d/du$
preserved by the one-line induction; the $\delta$-elimination therefore stays in odd powers and
the $m$-equation in rational coefficients), so we conjecture the expansion is rational to all
orders. The third constant resists identification: $C=-0.4745357072870\ldots$ is stable to
$10^{-14}$ across model orders and members through $j=55$, excludes every denominator up to
$\sim10^6$ (PSLQ), and the fitted higher coefficients grow ($F\approx4.3$), marking the series as
asymptotic-divergent; which caps the extractable precision and leaves $C$'s (conjecturally
rational, large-denominator) value to the third-order calculus. (c) At $j=1$ the same series
becomes a rational-coefficient expansion in $4/\pi^2$ evaluated at the wild head: divergent, hence
a resummation question. If a Borel-type summation of the rational series could be shown to
reproduce $q^\ast$, the arithmetic of $\pi^2$ would enter the problem for the first time through
the family; we record this as a direction and have tested its carrying capacity twice. A Borel--Pad\'e
resummation of the seven known coefficients reproduces the four accessible samples (both family
heads and both second members) to residuals $7.9\cdot10^{-3}$, $1.7\cdot10^{-6}$,
$4.6\cdot10^{-9}$, $7.4\cdot10^{-11}$ at $\mu=2.47,\,9.87,\,22.2,\,39.5$. An attempted
transseries extraction from these residuals then \emph{refuted its own premise}: the log-slopes
between consecutive samples are $1.14,\,0.48,\,0.24$, matching the truncation signature $7/\mu
=1.17,\,0.47,\,0.24$ exactly (three points), not an exponential $e^{-\beta\mu}$; the
residuals are the $x^7$-tail of the divergent series, not a nonperturbative signal. At the head
this is the standard resurgence degeneracy: seven coefficients place the evaluation at optimal
truncation, where truncation floor and first transseries term coincide in scale
($e^{-t_0\mu_1}\approx0.01$, with $t_0\approx0.75\pm1.70i$ the Pad\'e--Borel singularities),
so the head residual $-0.008(4)$ \emph{cannot} be attributed to a nonperturbative sector on
present data. That cost has now been paid: the crossing calculus has been automated in exact rational
arithmetic (weight-graded formal Newton for the zero shift; parity and residual controls exact),
producing the law's coefficients through order thirty; \emph{all rational}, e.g.\
$A_3=-1695089/3572100$ (whose denominator $2^23^65^27^2$ explains the failed PSLQ) and
$A_4=10199641/535815000$. The coefficients are computed by running the calculus over $\F_p$ and
recovering each $A_n$ by rational reconstruction from twenty $62$-bit primes; the values are
independently confirmed by extracting $A_n$ from the accumulation law applied to the roots
themselves, computed to $105$ digits.

The Borel geometry is only partly determined by these coefficients. The signs of $A_n$ follow a
three-positive, three-negative pattern from $n\approx12$, indicating a conjugate pair of
singularities near $60^\circ$; the pattern is not exact over the full range. This period-$6$
oscillation dominates the naive growth diagnostics, and must be removed before the Gevrey order can
be read off. Setting
\[
   G_n:=\tfrac16\bigl(\log|A_n|-\log|A_{n-6}|\bigr),
\]
which is phase-matched and therefore free of the oscillation, the model $|A_n|\sim C^n(n!)^s$ gives
$G_n\sim\log C+s\log n$, so $s$ is the slope of $G_n$ against $\log n$. Least squares yields
$s=0.844$ on $16\le n\le30$, $0.897$ on $23\le n\le30$ and $1.033$ on $26\le n\le30$: the estimate
rises towards $1$ as the fitting window advances, and the data are consistent with Gevrey-$1$.
(The unnormalised ratio $\log|A_n|/(n\log n)$, which equals $s+(\log C-s)/\log n$, is $0.313$ at
$n=14$ and $0.411$ at $n=30$; its corrections are of order $1/\log n$ and it is therefore useless
for extrapolation at these lengths.)

For the singularity modulus three independent estimates are available and agree only in order of
magnitude: the phase-matched growth constant gives $|t_0|=1/C=3.59$--$3.98$ over $24\le n\le30$;
a three-term recurrence fit, appropriate to a conjugate pair, gives
$|t_0|=3.285,\,3.394,\,3.493,\,3.558$ at $n_0=22,25,28,30$, extrapolating in $1/n_0$ to $4.30$,
with $\arg t_0=60.8^\circ,\,61.7^\circ,\,62.4^\circ,\,62.3^\circ$; and the Pad\'e--Borel
approximants give $|t|=2.61,\,2.56,\,2.96,\,3.09,\,3.33$ for $[6/6]$ through $[14/14]$.
Thus $|t_0|$ lies between $2.9$ and $4.3$ and is not pinned. Gevrey-$1$ growth, and with it the
$n!$-Borel framework, is supported but not established; the location of the singularities is not
established; and the Borel representation below remains a conjecture whose analytic premise is
open. One structural constraint is available: by the weight lemma, only
Bernoulli numbers of index $\lesssim n$ reach $A_n$ (a $B_j$ occurring in $c_m$ sits at $k$-degree
$\deg c_m-j$ and reaches $A_n$ only if $j\le2n-m+1$), which excludes the Gevrey-$2$ growth that
$B_{2n}$ would otherwise force, but does not pin the class. Where testable the resummation succeeds: at the sine head the Pad\'e--Borel ladder
converges $1.7\cdot10^{-6}\to3.7\cdot10^{-11}$ (any transseries prefactor $\lesssim10^{-6}$ unless
phase-suppressed), while at the cosine head the plain ladder reaches its approximation wall at
$\sim10^{-3}$ with residuals oscillating about zero. We therefore pose the
\emph{Borel representation conjecture}: $1-q^\ast_j=(2/\mu_j)\,\mathcal S(1/\mu_j)$ for both
families and all $j$, where $\mathcal S$ is the Borel--Laplace sum of the explicit rational series
$\sum A_nx^n$; in particular
\[
   q^\ast\;=\;1-\frac{8}{\pi^2}\,\mathcal S\!\Bigl(\frac{4}{\pi^2}\Bigr),
\]
the first exact (conjectural) bridge between $q^\ast$ and $\pi^2$-arithmetic. Numerically the head
representation is confirmed to $\sim10^{-4}$ by conformal-map Borel resummation (the residual
crosses zero smoothly with no plateau above the estimator scatter), and the \emph{same} resummation
machinery reproduces the three lower-$x$ heads to $10^{-13}$--$10^{-26}$; the coefficients are now
computed through $A_{22}$, still rational (e.g.\ $A_{21}$, $A_{22}$ have denominators
$\sim10^{93}$), consistent with the rationality of the expansion at every order.

Two honest boundaries frame the bridge. \emph{Analytically}, its proof is \emph{not} a
cite-the-theorems exercise: it reduces to a single hard lemma; a $q$-uniform, turning-point
(Olver/Airy-grade) factorial remainder bound $|A_n|\le C\sigma^nn!$ for the Hahn--Exton crossing as
$q\to1$; which has no citable form in the literature and is plausibly the program's recurring
degenerate-saddle obstruction in Borel-summability costume. \emph{Arithmetically}, the Borel
transform $b(t)=\sum(A_n/n!)t^n$ carries a sharp, exceptionless prime law: for every odd prime
$p\le43$, $p\mid\operatorname{den}(b_n)$ exactly when $2n+1\ge p$, with $p$-adic valuation exactly
$2$ at the onset $n=(p-1)/2$, and thereafter non-vanishing with valuation growing linearly
\emph{on average} (least-squares slopes $2.38,1.17,0.73,\dots$ for $p=3,5,7$; the valuation is
non-monotone locally, e.g.\ $v_5$ dips $6\!\to\!5$ and $v_{17}$ dips $2\!\to\!1$). The clean-before-onset half is a theorem for all $p$. (i) The polynomials $c_n(k)$ built from the
calculus are $p$-integral for $n\le p-2$ (von Staudt--Clausen \cite{IrelandRosen1990} at the Faulhaber layer). (ii) $A_n$
depends only on the calculus data through order $2n$; this is now proved, in two steps. First, by
series reversion $A_n$ is a function of $m_1,\dots,m_n$ alone (the order-$n$ inverse coefficient uses
only the first $n$ forward coefficients; verified by perturbing $m_{>n}$). Second, $m_i$ depends only
on $c_1,\dots,c_{2i}$: in the turning-point grading, where $k\sim\varepsilon^{-1/2}$ assigns
$\varepsilon^au^b$ the weight $2a-b$ (physical order $=$ weight$/2$) and every operation is additive
on weights, the polynomial $c_m$ enters at minimum weight $2m-\deg_k c_m$, attained by its
top-degree term. Here $\deg_k c_m=m+[m\text{ even}]$: the degree-$m$ coefficient of the
$\varepsilon^m$ log-ratio $\ell_m(i)$ equals $(-1)^m[x^m]\log\!\bigl((e^x-1)/x\bigr)$, which vanishes
for every odd $m\ge3$ because $\log\!\bigl((e^x-1)/x\bigr)=\tfrac x2+\log\!\bigl(2\sinh(x/2)/x\bigr)$
and the last term is \emph{even} in $x$; the Faulhaber sum then raises the $k$-degree by one. Hence
$c_m$ has minimum weight $m-[m\text{ even}]$ and first reaches physical order $i$ only when $m\le2i$,
so $m_i$ (and therefore $A_n$) sees no $c_m$ with $m>2n$. (iii) The assembly introduces no
prime $p$ at output order $n<p$: the Bernoulli denominators of the Faulhaber sums acquire $p$ only
when $(p-1)\mid2m$, i.e.\ $2m\ge p-1$, and all factorial and division indices stay below $p$. For
$n<(p-1)/2$ one has $2n\le p-3$, so every ingredient feeding $A_n$ is $p$-integral and
$p\nmid\operatorname{den}(A_n)$.

The matching lower direction; that $p$ \emph{does} appear at $n=(p-1)/2$, with valuation
exactly $2$; reduces to a unit. At $n=(p-1)/2$ the sole prime-$p$ source is $c_{p-1}$ (all lower
$c$ are $p$-integral), and it enters $A_n$ \emph{linearly}, since $c_{p-1}^2$ has weight
$2(p-2)>2n$. Only the top three coefficients of $c_{p-1}$ (degrees $\ge p-2$) survive the weight
cutoff; of the two that carry $p^2$; the $k^1$ and $k^p$ monomials (von Staudt--Clausen at the
Faulhaber layer, $\deg_kc_{p-1}=p$); only $k^p$ reaches $A_n$, as $k^1$ sits at weight
$2p-3>2n+1$. The coefficient of that top monomial in $A_{(p-1)/2}$ is a \emph{unit},
$\partial A_{(p-1)/2}/\partial c_{p-1}[k^p]=\pm1$, and this now has a leading-symbol proof. The
monomial perturbs the confluent expansion of $G$ by $\varepsilon^{p-1}k^p$; the turning-point map
(the $D^p$-action, $D=(u/2)\,d/du$) sends $k^p$ to $a_pu^p\sin u+b_pu^{p-1}\cos u$ with
$a_p=(-1)^{n-1}2^{-p}$ the leading sine coefficient and $b_p=(-1)^nn(2n+1)2^{-p}$ the leading cosine
coefficient ($p=2n+1$; both closed forms follow from the recursion
$(A,B)\mapsto\bigl(\tfrac u2(A'-B),\tfrac u2(B'+A)\bigr)$). At a simple zero $u_*$ of $\cos$ the
cosine term vanishes and every sine power below $u^p$ contributes only at higher order in
$\varepsilon$, so the zero shift is $\delta=\varepsilon^{p-1}a_pu_*^{\,p}$ (the factor
$\sin u_*=\pm1$ cancels against $G'=-\sin$). With $m=\varepsilon u^2$ and the leading balance
$u_*^2=2/\varepsilon$ (from $\varepsilon u^2=2(1-\varepsilon)e^{\varepsilon/2}\to2$),
\[
   \delta m=2\varepsilon u_*\,\delta=2a_p\,\varepsilon^{p}u_*^{\,p+1}
   =a_p\,2^{\,n+2}\,\varepsilon^{n},\qquad\text{so}\quad
   \frac{\partial m_n}{\partial c_{p-1}[k^p]}=a_p2^{\,n+2}=\pm2^{\,1-n},
\]
and with the reversion Jacobian $\partial A_n/\partial m_n=2^{n-1}$ (from $\varepsilon\sim2/\mu$)
the product is $\pm1$. Only the top sine power $u^p$ reaches order $\varepsilon^n$, so the magnitude
is exact; the sign is $(-1)^n$ (verified for all $n\le10$, independently of $p$). Since
$c_{p-1}[k^p]$ has $p$-adic valuation $-2$ while the other reaching monomials ($k^{p-1},k^{p-2}$)
have valuation $\ge-1$ against $p$-integral multipliers, $v_p(\operatorname{den}A_{(p-1)/2})=2$
exactly. Thus the onset law; $p\mid\operatorname{den}(A_n)\iff 2n+1\ge p$, with valuation $2$ at
$n=(p-1)/2$; is \emph{proved for all odd $p$}. The average-linear
(Legendre/factorial) growth; not the logarithmic growth an $\mathrm{lcm}(1,\dots,2n)$ or
$\zeta$-type denominator would give; makes
$\operatorname{den}(A_n)$ super-geometric ($\log\operatorname{den}(A_n)/n$ rises $5.4\to10.2$ over
$n\le20$, unbounded). Consequently $b(t)$ is \emph{not} a $G$-function (super-geometric denominators
alone suffice, independent of holonomy), so $\sum A_nx^n$ is not of Andr\'e arithmetic-Gevrey type \cite{Andre2000}
and lies outside the $G$-function, $E$-function, and $E$-operator (Fischler--Rivoal
\cite{FischlerRivoal2016}) classes
entirely. The closest existing frame, the $E$-operator theory, has exactly one blocking gap; the
hypothesis that $b$ is a $G$-function; which the prime law refutes rather than supplies.
The missing tool is therefore genuinely new, not an incremental hypothesis-removal: a value theory
for Borel--Laplace sums whose Borel transform is non-holonomic with factorial denominators,
evaluated at a point ($4/\pi^2$) transcendental over $\Q$. This constant is a canonical first target
for such a theory. Third, the family sum
converges to
\[
   T:=\sum_{j\ge1}\bigl(1-q^\ast_j\bigr)=0.7357490877157094\ldots\ (\pm3\cdot10^{-17},
   \text{ tail summed via the fitted law and Hurwitz zeta}),
\]
for which no relation with $\{1,\pi^{-2},\pi^{-4},\log2,\gamma,\pi^2/6\}$ exists at height
$10^4$ within the stated uncertainty (and $T\ne2/e$): the classical tail is $\pi^2$-rational
piece by piece, but the wild head shifts the sum off every tested constant; a sum rule, if one
exists, must come from global structure not yet visible.
\end{remark}

The failure pattern above is not an artifact of particular constructions; it is forced by exact
arithmetic facts, which we record as a theorem-level ledger.

\begin{proposition}[Arithmetic no-go ledger at $\rho=\tfrac12$]
\label{prop:nogo}
Let $q=s/t\in(0,1)$ with $\gcd(s,t)=1$, and let $g_k=(-1)^kq^{k(k-1)}/(q;q)_{2k}$ be the Taylor
coefficients of $G(q,\cdot)$. Then:
\begin{itemize}
\item[\textup{(i)}] \textup{(exact denominators)} in lowest terms
$g_k=(-1)^k\,s^{k^2-k}\,t^{k^2+2k}\big/\prod_{i=1}^{2k}(t^i-s^i)$; in particular
$\operatorname{den}(g_k)=\prod_{i=1}^{2k}(t^i-s^i)=t^{k(2k+1)(1+o(1))}$ exactly, while
$|g_k|=(s/t)^{k^2(1+o(1))}$, so the arithmetic ratio is
$\rho=\log(t/s)/(2\log t)\le\tfrac12$, with equality iff $s=1$.
\item[\textup{(ii)}] \textup{($\nu$-ladder arithmetic)} the $q$-Lommel polynomials of
Koelink--Swarttouw \cite{KoelinkSwarttouw1994} satisfy $\deg_qp_m=m^2$ exactly, with leading $q$-coefficient $(-1)^m$; hence at
$q=s/t$ their denominators are the pure power $t^{m^2}$, carrying no cyclotomic structure.
\item[\textup{(iii)}] \textup{($\nu$-direction regularity)} the ladder recurrence
$u_{m+1}=\bigl((1-q^{2m+1})/x_0+x_0\bigr)u_m-u_{m-1}$ has coefficients converging exponentially to
$x_0+1/x_0$ with distinct characteristic roots $x_0^{\pm1}$; by Poincar\'e--Perron
\cite{Poincare1885,Perron1909} every nonzero
solution satisfies $|u_m|^{1/m}\to x_0^{\pm1}$. No theta-type ($q^{cm^2}$-decaying) solution exists
in the $\nu$-direction, for either kind. Consequently the cleared $\nu$-ladder forms are
$\gg t^{m^2}x_0^{-m}\to\infty$ for every rational $q$, including $s=1$.
\item[\textup{(iv)}] \textup{(cyclotomic floor)} any linear form built over the coefficients
$g_0,\dots,g_M$ has decay at best $q^{M^2(1+o(1))}$, while its cyclotomic support is
$\{\Phi_d(t,s):d\le2M\}$; even if every $\Phi_d$ appeared at most \emph{once}; the idealized best
case of any lcm/acceleration argument; the denominator exceeds
$t^{\sum_{d\le2M}\varphi(d)}=t^{(12/\pi^2+o(1))M^2}$, and $12/\pi^2=1.2159\ldots>1$. Hence no
cyclotomic-compression (Rhin--Viola-type \cite{RhinViola2001}) refinement of the Pad\'e, continued-fraction, or ladder
families can produce integer forms tending to $0$. Moreover the true solved Pad\'e denominators
contain growing \emph{non-cyclotomic} irreducible factors (at $n=3$, one of degree $26$ with
coefficients up to $4$; no integer $d$ has $\varphi(d)=26$), which no lattice-of-cyclotomes argument
addresses.
\end{itemize}
\end{proposition}

\begin{proof}
(i) $q^{k(k-1)}=s^{k(k-1)}/t^{k(k-1)}$ and $1/(q;q)_{2k}=t^{k(2k+1)}/\prod_{i\le2k}(t^i-s^i)$; the
numerator $s^{k^2-k}t^{k^2+2k}$ is coprime to each factor since
$\gcd(st,\,t^i-s^i)=1$. (ii) Induction: $p_{m+1}=(X+1-q^{2m+1})p_m-Xp_{m-1}$, and the top
$q$-degree term of $p_{m+1}$ is $-q^{2m+1}\cdot(\text{top of }p_m)$, of degree $m^2+2m+1=(m+1)^2$,
which cannot cancel against $Xp_{m-1}$ (degree $(m-1)^2$). (iii) The perturbation
$(1-q^{2m+1})/x_0-1/x_0$ decays geometrically, so the Poincar\'e--Perron theorem \cite{Poincare1885,Perron1909} applies with the
limit equation $u_{m+1}=(x_0+1/x_0)u_m-u_{m-1}$, whose roots $x_0\ne1/x_0$ are simple. (iv) A form
using $g_0,\dots,g_M$ vanishes to order at most $M$ at $z=0$, so its value at a fixed $z_0$ inside
the disc is $\gg|g_{M+1}z_0^{M+1}|=q^{M^2(1+o(1))}$ in the truncation-based families; its
denominators involve $(t^i-s^i)$ for $i\le2M$, i.e.\ $\Phi_d(t,s)$ for $d\le2M$, and
$\sum_{d\le N}\varphi(d)=\tfrac3{\pi^2}N^2+O(N\log N)$ gives the floor with $N=2M$. The degree-$26$
factor is exhibited by the exact $n=3$ computation
(\texttt{code/zeta\_probe/beta2\_qniven\_ladder.py}).
\end{proof}

\begin{remark}[scope]
\label{rem:nogo-scope}
Proposition~\ref{prop:nogo} closes every construction whose coefficients live in the natural
$\Z[q^{\pm1},z]$-lattices of the $q$-difference module; Taylor Pad\'e and Hermite--Pad\'e,
the argument-direction continued fraction, the $\nu$-ladders of both kinds; together with all
cyclotomic-compression refinements of them. It does \emph{not} exclude a fundamentally different
family of forms; producing one is exactly what a $\rho=\tfrac12$ criterion would mean. The ledger
thereby locates $\beta_2$'s arithmetic at a precise frontier of $q$-Diophantine approximation.
\end{remark}

The remaining freedom can be narrowed further: the finite places of $\Q$ are provably useless here.
The series $S(q)=\sum_k g_k z^{\ast k}$, $z^\ast=2q(1-q)$, has rational terms, so besides its real sum
it possesses $s$-adic and $t$-adic \emph{avatars}; the same series summed in $\Q_s$ and $\Q_t$,
where it also converges (the valuations of the terms tend to $+\infty$).

\begin{lemma}[adelic units]
\label{lem:adelic-units}
Let $q=s/t\in(0,1)$, $\gcd(s,t)=1$, and $z^\ast=2q(1-q)$. For every $k\ge1$ and every prime $p$:
if $p\mid t$ with $e=v_p(t)$, then
\[
   v_p\bigl(g_kz^{\ast k}\bigr)\;=\;k^2e+k\,[\,p=2\,]\;\ge\;k^2 ,
\]
(exact, verified symbolically), and if $p\mid s$ then $v_p(g_kz^{\ast k})\ge k^2$. Consequently the
series converges in $\Q_p$ for every $p\mid st$, and the finite-place avatars satisfy
$S_p\equiv1\pmod{p^{\,v_p(st)}}$; in particular $S_t\equiv1\pmod t$ and $S_s\equiv1\pmod s$: they
are \emph{units}, hence nonzero; regardless of whether the real avatar vanishes.
\end{lemma}

\begin{proof}
Each factor $t^i-s^i$ of $\operatorname{den}(g_k)$ is a $p$-unit for every $p\mid st$ (it is
$\equiv-s^i$ or $\equiv t^i\bmod p$), so by Proposition~\ref{prop:nogo}\textup{(i)},
$v_p(g_k)=(k^2+2k)e$ for $p\mid t$ and $v_p(g_k)=(k^2-k)\,v_p(s)$ for $p\mid s$. From
$z^\ast=2s(t-s)/t^2$: for $p\mid t$, $v_p(t-s)=0$ and $v_p(z^\ast)=[\,p=2\,]-2e$, giving
$v_p(g_kz^{\ast k})=(k^2+2k)e+k([\,p=2\,]-2e)=k^2e+k[\,p=2\,]$; for $p\mid s$,
$v_p(z^\ast)=v_p(2s)\ge v_p(s)\ge1$, giving $v_p(g_kz^{\ast k})\ge(k^2-k)+k=k^2$. All are $\ge1$ for
$k\ge1$, so each avatar is $1$ plus terms of positive valuation.
\end{proof}

\begin{remark}[exact adelic balance, and the shape of any future criterion]
\label{rem:adelic-balance}
Lemma~\ref{lem:adelic-units} has two consequences. First, it extends the no-go ledger from the
module's coefficient lattices to the \emph{values}: no divisibility or smallness can be sourced from
the finite places, since the finite-place avatars are units; a hypothetical rationality of
$\beta_2$ would be a maximally place-asymmetric event, the real avatar vanishing while every finite
avatar equals a unit $\equiv1$. Second, it makes the $\rho=\tfrac12$ balance exact in the sense of
the product formula: the term $g_kz^{\ast k}$ carries archimedean decay $(s/t)^{k^2(1+o(1))}$,
$s$-adic and $t$-adic numerator content $s^{k^2}t^{k^2}$, and denominator $t^{2k^2(1+o(1))}$; the
books balance with \emph{zero margin} at the $k^2$ scale, so any successful construction must be
lossless at all three places simultaneously and win at the geometric scale. Every mechanism surveyed
above loses a factor $t^{ck^2}$ somewhere; a $\rho=\tfrac12$ criterion is therefore forced to be a
purely \emph{archimedean cancellation} phenomenon. This suggests the precise principle that would
settle $\beta_2$, which we record as a conjecture: for the irregular rank-two $q$-difference module
with polynomial companion matrix and $\det A=q$, the vanishing locus is \emph{adelically rigid}: the entire function $G_v(q,\cdot)$ cannot vanish at a rational point $z_0$ at which all finite-place
avatars are units (as at $z_0=z^\ast$, by Lemma~\ref{lem:adelic-units}) at one place $v$ of $\Q$
while the avatars at all other convergent places are units. The finite-place half of this rigidity is exactly Lemma~\ref{lem:adelic-units};
its archimedean half implies $\beta_2\notin\Q$. This is a $q$-analogue, for irregular $q$-difference
modules, of the place-uniformity enjoyed by $E$- and $G$-function values, where the exceptional sets
of Siegel--Shidlovskii--Andr\'e theory are independent of the place.
\end{remark}

Further analysis of the archimedean obstruction yields several exact structural facts and a faithful
reduction. All are verified numerically to $40$--$60$
digits (\texttt{code/zeta\_probe/beta2\_*.py}).

\begin{proposition}[structure of the zero curve]
\label{prop:zerostructure}
Write $Q=q^2$ and let $0<z_1(q)<z_2(q)<\cdots$ be the positive zeros of $G(q,\cdot)$.
\begin{itemize}
\item[\textup{(i)}] \textup{(base change)} $G(q,z)=\dfrac{(q^2;q^2)_\infty}{\sqrt q\,(q;q^2)_\infty}\,
z^{1/4}\,J^{(3)}_{-1/2}\!\bigl(\sqrt z/q;\,q^2\bigr)$, where $J^{(3)}_\nu$ is the Hahn--Exton (third
Jackson) $q$-Bessel function; thus $G$ is, up to an explicit prefactor, the \emph{Hahn--Exton
$q$-cosine}, and its zeros are those of a classical, well-studied object.
\item[\textup{(ii)}] \textup{(zero lattice)} $z_k(q)\,Q^{k-1}\to1$ as $k\to\infty$, with
theta-suppressed dressing $\delta_k:=z_{k+1}Q^{k}-1=O(c^{k}Q^{k^2})$ (numerically $\delta_0=-0.5051,\ \delta_1=-1.17\cdot
10^{-2},\ \delta_2=-1.03\cdot10^{-5},\dots$); $z^\ast=z_1(q^\ast)$ is the maximally dressed member,
displaced from the bare lattice point $Q^{0}=1$ by an $O(1)$ correction; it is \emph{not} a
torsion/lattice point of the connection curve.
\item[\textup{(iii)}] \textup{(power sums are rational)} the symmetric functions
$\sigma_p(q):=\sum_{n\ge1}z_n(q)^{-p}$ lie in $\Q(q)$ for every $p\ge1$ (Newton's identities on the
$\Q(q)$-coefficients $g_k$); e.g.\ $\sigma_1=1/((1-q)(1-q^2))$. Consequently
$z^\ast(q)=\lim_{p\to\infty}\sigma_p(q)/\sigma_{p+1}(q)$ with geometric rate $(z_1/z_2)^p$, giving an
all-orders $\Q(q)$-approximation of the zero curve.
\end{itemize}
\end{proposition}

Two of the no-go's ingredients also sharpen. First, the \emph{cause} of $\rho=\tfrac12$ is exactly the
\emph{double} Pochhammer: $\rho=\log(t/s)/(2\log t)$, and the factor $2$ in the denominator is the
$2k$ in $(q;q)_{2k}$; a single Pochhammer would give $\rho>1$ (whence Liouville) when $t>s^2$. The
intractability is intrinsic to the base-$q^2$ Hahn--Exton structure, not to the choice of construction.
Second, the finite-place primes are \emph{disjoint} from $st$: $\prod_{i\le2N}(t^i-s^i)$ is coprime to
$st$ (only the prime $2$, via lifting-the-exponent when $s,t$ are both odd, contributes an $O(N)$ term),
so the arithmetic of every partial-sum denominator is supported off $\{s,t\}$; orthogonal to the
unit residues of Lemma~\ref{lem:adelic-units}. The Hankel/determinantal escape is likewise closed
exactly: $\det(g_{i+j})_{0\le i,j\le n}$ has denominator of \emph{cubic} archimedean size
($\log|\!\operatorname{den}|\sim c\,n^3$), one power worse than the quadratic remainder.

Finally, the Nesterenko route is blocked more precisely than by non-holonomicity alone. Nesterenko's
method needs a finitely generated $\delta$-closed ($\delta=q\,d/dq$) $\C(q)$-algebra containing $S$,
i.e.\ that $S$ be \emph{differentially algebraic}; non-holonomicity does not suffice (the quasimodular
$E_2$ is non-holonomic yet $\delta$-algebraic). One computes
$\delta S=\sum_k a_k\bigl(k^2-\tfrac{kq}{1-q}+L_k\bigr)$ with
$L_k=\sum_{i=1}^{2k} iq^i/(1-q^i)\to(1-E_2)/24$: $S$ carries \emph{every} finite Eisenstein truncation
$L_k$, not merely the limit, which is precisely why it escapes the modular differential ring. Whether
$S$ is differentially algebraic is itself open, but the exclusion is now proof-grade over a wide
window: $S$ satisfies no algebraic differential equation in the $\delta$-jets at any of the
profiles (order, degree, $\deg_q$) $=(2,12,3)$, $(3,8,3)$, $(4,5,6)$, $(5,4,8)$, $(6,3,15)$,
$(8,2,35)$, $(12,1,120)$; exact linear algebra on $2000$ coefficients modulo a prime, every
matrix of full column rank; since a primitive integer relation survives reduction, full rank modulo
one prime excludes a relation over $\Q$ at that profile outright. (The pipeline is validated on
$E_2$, for which it recovers the $q$-Chazy equation
$\delta^3E_2=E_2\delta^2E_2-\tfrac32(\delta E_2)^2$ at its exact profile.) The same exclusion
holds verbatim for the zero branch $z_1(q)$. The Nesterenko route is blocked \emph{modulo the
differential transcendence of $S$}, which is itself open; and, with no functional equation in
$q$ available for $S$ (no shift, $q$-difference, or Mahler equation exists), no current method can
even attack it.

\begin{remark}[reduction to a theta non-resonance]
\label{rem:sharpest}
All routes converge on the same faithful reformulation: $\beta_2$ is irrational iff
$G(q^\ast,z^\ast)\ne0$ is forced by the arithmetic, and (Ramis--Sauloy--Zhang \cite{RSZ2013}) this is a
\emph{theta non-resonance} statement; the connection ratio $R_q(z)=[P_{21}\,y_{\mathrm{irr}}]/
[P_{11}\,y_{\mathrm{reg}}]$, an explicit $\theta_{q^2}$-quotient, must satisfy $R_q(2q(1-q))\ne-1$ at
every algebraic $q$. The evidence points toward truth: in the confluent limit $q\to1$ this is exactly
Lindemann's theorem that $\cos$ has no algebraic zero. But the reduction is not yet shown \emph{easier}:
$R_q(z)=-1$ is a functional-equation-consistent value that a nonconstant theta-quotient must attain
somewhere, and on the diagonal $z=2q(1-q)$ the theta-arguments appear algebraically dependent, so no
current independence theorem (Bertrand; Nesterenko \cite{Nesterenko1996}; Schneider--Lang
\cite{Lang1966}) forbids it. The archimedean
obstruction is thus the exact mirror of the proven finite-place balance: units ($\equiv1$) at every
$p\mid st$, and a putative $-1$ at $\infty$, with neither forbidden. The required connection-entry
computation is carried out below.
\end{remark}

\begin{proposition}[the connection entries]
\label{prop:pinned}
Write $Q=q^2$, $\theta_b(x)=\sum_{n\in\Z}b^{n(n-1)/2}x^n$. The equation for $G$ has, at $z=\infty$, the
formal solution basis $f_1(z)=\theta_Q(-z)\,\hat g(z)$ and $f_2(z)=h(z)/\theta_Q(-z/q)$, where
$h(z)=\sum_kc_kz^{-k}$ is \emph{convergent} (entire in $1/z$, $c_k=O(Q^{k^2/2})$) while
$\hat g(z)=\sum_kb_kz^{-k}$ is $q$-Gevrey divergent, with the closed form
\[
   b_k=\frac{q^{\,k-k(k-1)/2}}{(q;q)_k}
\]
(equivalent to the recursion $b_k=[-(q+1)Q\,b_{k-1}-Q^{3-k}b_{k-2}]/(q(Q^k-1))$; verified
symbolically for $k\le8$). By Euler's theorem $\sum_jx^j/(q;q)_j=1/(x;q)_\infty$, the $q$-Borel
transform is therefore explicit:
\[
   \widehat B(\xi)=\sum_{k\ge0}b_k\,q^{k(k+1)/2}\xi^{-k}
   =\sum_{k\ge0}\frac{(q^2/\xi)^k}{(q;q)_k}
   =\frac1{(q^2/\xi;\,q)_\infty},
\]
meromorphic on $\C^*$ with poles exactly on $\xi\in q^{2+\mathbb{N}_0}$, and satisfying
$\widehat B(\xi)+\tfrac{q+1}{\xi}\widehat B(q\xi)+(\xi^{-2}-1)\widehat B(q^2\xi)=0$ identically
(direct computation from the product form; in particular $\widehat B(-1)=(q+1)\widehat B(-q)$). The
theta-kernel $q$-Laplace
\[
   g^{[\lambda]}(z)=\frac{\sum_{m\in\Z}q^{m(m-1)/2}(\lambda/z)^m\,\widehat B(\lambda q^m)}
        {\theta_q(\lambda/z)},\qquad
   f_1^{[\lambda]}=\theta_Q(-z)\,g^{[\lambda]}(z),
\]
is then singular only for $\lambda q^{\Z}$ meeting the pole set, i.e.\ for the single class
$[\lambda]\in q^{\Z}$ of $E_q=\C^*/q^{\Z}$ (the positive-real lattice class, near which the zeros of
$G$ lie); every other class is regular. For $\lambda\notin q^{\Z}$ the bilateral sum converges
absolutely on compacts of $\C^*\setminus(-\lambda q^{\Z})$ and $Lf_1^{[\lambda]}=0$ \emph{exactly}:
the theta shifts $\theta_q(q^{-2}x)=q^{-3}x^2\theta_q(x)$ and $\theta_q(q^{-4}x)=q^{-10}x^4\theta_q(x)$
turn $f_1^{[\lambda]}(Qz)$ and $f_1^{[\lambda]}(Q^2z)$ into the same kernel sum with
$\widehat B(\lambda q^m)$ replaced by $\widehat B(\lambda q^{m+2})$, $\widehat B(\lambda q^{m+4})$;
collecting the coefficient of $q^{m(m-1)/2}(\lambda/z)^m/\theta_q(\lambda/z)$ in
$Lf_1^{[\lambda]}/\theta_Q(-z)$ then yields
$q\bigl[\widehat B(\xi)+\tfrac{q+1}{\xi}\widehat B(q\xi)+(\xi^{-2}-1)\widehat B(q^2\xi)\bigr]$ at
$\xi=\lambda q^m$, identically zero. (Each step was verified numerically to $40$ digits; the kernel
normalization was validated on convergent input to $34$ digits, and $Lf_1^{[\lambda]}\approx10^{-61}$
numerically.) Writing
$G=C_1f_1^{[\lambda]}+C_2^{[\lambda]}f_2$, numerically to $59$--$60$ digits:
\begin{itemize}
\item[\textup{(i)}] since $\det A=q$, every Casoratian
$\operatorname{Cas}(u,v)(z):=u(z)v(Qz)-u(Qz)v(z)$ of solutions satisfies
$\operatorname{Cas}(Qz)=q\operatorname{Cas}(z)$ (immediate from $u(Q^2z)=a\,u(Qz)-q\,u(z)$), so
$C_1=\operatorname{Cas}(G,f_2)/\operatorname{Cas}(f_1^{[\lambda]},f_2)$ and
$C_2^{[\lambda]}=\operatorname{Cas}(f_1^{[\lambda]},G)/\operatorname{Cas}(f_1^{[\lambda]},f_2)$ are
$Q$-elliptic; proven, and verified to $10^{-60}$;
\item[\textup{(ii)}] $C_1$ is a \emph{constant}, independent of both $z$ and the resummation direction
$\lambda$: this is proven in Theorem~\ref{thm:zeroproduct} below (the limit argument applies along
every orbit class $\ne[1]_Q$, and ellipticity extends the value); the slope-$1$ component of $G$ is
Stokes-invariant, and all $q$-Stokes ambiguity is carried by $C_2^{[\lambda]}$, nonconstant elliptic
with $\lambda^2$ in its zero divisor;
\item[\textup{(iii)}] \textup{(dressing-product formula, verified to $59$ digits)}
\[
   C_1=\frac1{(Q;Q)_\infty\prod_{k\ge1}z_kQ^{\,k-1}}\;:
\]
the Stokes-invariant connection constant is the regularized product of the zero dressings;
\item[\textup{(iv)}] at $q=q^\ast$ the entire apparatus evaluates: $C_1(q^\ast)=2.69898382616391375\ldots$,
and the zero condition $G(q^\ast,z^\ast)=0$ takes the pinned form
\[
   C_2^{[\lambda]}(z^\ast)\;=\;-\,C_1\,\frac{f_1^{[\lambda]}(z^\ast)}{f_2(z^\ast)}
   \;=\;0.011592994120624266\ldots\qquad(\lambda=\tfrac32),
\]
verified to $30$ digits. This is the explicit instance of the theta non-resonance: every object in it
is now constructed.
\end{itemize}
\end{proposition}

The constant $C_1$ is in fact \emph{modular}, and this upgrades the dressing-product formula into an
exact evaluation of the zero product.

\begin{theorem}[the connection constant, and the zero product of the Hahn--Exton $q$-cosine]
\label{thm:zeroproduct}
For $0<q<1$,
\[
   C_1=\frac1{(q;q)_\infty},
   \qquad\text{and consequently}\qquad
   \prod_{k\ge1} z_k(q)\,Q^{\,k-1}\;=\;(q;q^2)_\infty ,
\]
where $z_1<z_2<\cdots$ are the positive zeros of $G(q,\cdot)$. Likewise, for the second solution's
underlying entire function $H$ with zeros $w_1<w_2<\cdots$,
$\prod_{k\ge1}w_k(q)\,q\,Q^{\,k-1}=(q^3;q^2)_\infty$. \textup{(Verified at three bases including
$q^\ast$, to $41$, $23$, $26$ digits for $C_1$ and $58$ digits for the product.)}
\end{theorem}

\begin{proof}
\emph{(1) An exact identity.} From $(q;q)_{2k}=(q;q)_\infty/(q^{2k+1};q)_\infty$ and Euler's
expansion $(x;q)_\infty=\sum_{j\ge0}(-1)^jq^{j(j-1)/2}x^j/(q;q)_j$ at $x=q^{2k+1}$, together with
$k(k-1)+2kj=(k+j)(k+j-1)-j(j-1)$ and the substitution $m=k+j$, absolute convergence of the double
sum permits the rearrangement
\[
   (q;q)_\infty\,G(z)\;=\;\sum_{m\ge0}(-1)^mq^{m(m-1)}z^m\,P_m(1/z),
   \qquad
   P_m(w):=\sum_{j=0}^m c_jw^j,\quad c_j=\frac{q^{\,j-j(j-1)/2}}{(q;q)_j}
\]
(the $c_j$ coincide with the $b_j$ of Proposition~\ref{prop:pinned}; identity checked symbolically
through $z^8$ and numerically to $41$ digits).

\emph{(2) The constant.} Fix $u<0$ and put $z=uQ^{-N}$. The moduli $q^{m(m-1)}|z|^m$ peak at
$m_0=N+O(1)$ and decay like $q^{(m-m_0)^2+O(1)}$ away from it. For $1\le j\le m\le 2N$,
$2Nj-j(j-1)/2\ge Nj$, so $|c_jz^{-j}|\le Cq^{Nj}|u|^{-j}$ and
$|P_m(1/z)-1|\le C'q^{N}$ uniformly; the terms with $m>2N$ contribute a relative
$O(q^{N^2})$. The negative-$m$ part of $\theta_Q(-z)$ is $O(|z|^{-1})$, negligible against the peak.
Hence $(q;q)_\infty G(z)/\theta_Q(-z)\to1$ along $z=uQ^{-N}$, for each fixed $u<0$.

\emph{(3) Geometric zero dressing.} On the circle $|z-Q^{1-k}|=\varepsilon q^kQ^{1-k}$ the identity
of step (1) and the bounds of step (2) give
$|(q;q)_\infty G-\theta_Q(-z)|\le Cq^k|\theta_Q(-z)|_{\text{peak}}$, while $\theta_Q(-z)$ has a
simple zero at $Q^{1-k}$ with derivative bounded below there; Rouch\'e gives
$z_k=Q^{1-k}(1+O(q^k))$, i.e.\ $\log(z_kQ^{k-1})=O(q^k)$, absolutely summable.

\emph{(4) The product.} $G(z)=\prod_k(1-z/z_k)$ (order $0$, genus $0$, $G(0)=1$) and
$\theta_Q(-z)=(Q;Q)_\infty(z;Q)_\infty(Q/z;Q)_\infty$ (Jacobi). Along $z=-R$, $R\to\infty$,
\[
   \frac{G(-R)}{\theta_Q(R)}
   =\frac1{(Q;Q)_\infty(-Q/R;Q)_\infty}\prod_{k\ge1}\frac{1+R/z_k}{1+RQ^{k-1}} ,
\]
and for $x,y>0$ the function $R\mapsto\log\frac{1+Rx}{1+Ry}$ is monotone from $0$ to $\log(x/y)$, so
each factor is dominated by $|\log(z_kQ^{k-1})|$, summable by step (3). Dominated convergence gives
$G(-R)/\theta_Q(R)\to1/\bigl[(Q;Q)_\infty\prod_kz_kQ^{k-1}\bigr]$ as $R\to\infty$; evaluating along
the subsequence $R=|u|Q^{-N}$ by step (2), the limit equals $1/(q;q)_\infty$, whence
$\prod_kz_kQ^{k-1}=(q;q)_\infty/(Q;Q)_\infty=(q;q^2)_\infty$ by Euler's
$(q;q)_\infty=(q;q^2)_\infty(q^2;q^2)_\infty$. Finally, in the basis of
Proposition~\ref{prop:pinned}: along the negative axis (a regular class, $[-1]\notin q^{\Z}$) one has
$g^{[\lambda]}(z)\to1$ and $f_2(z)/\theta_Q(-z)=O\bigl(1/(\theta_Q(R)\theta_Q(R/q))\bigr)\to0$, so
the limit just computed is the entry $C_1$; hence $C_1=1/(q;q)_\infty$. The $H$-statement is
identical with comparison function $\theta_Q(-qz)$ and ratio $(1-q)/(q;q)_\infty$.
Finally, steps (1)--(2) apply verbatim along $z=uQ^{-N}$ for any $u\in\C^*$ whose orbit avoids the
zero class $[1]_Q$ of $\theta_Q(-\,\cdot\,)$, using the standard theta lower bound
$|\theta_Q(-z)|\ge c([u])\cdot(\text{peak term})$ off that class; along each such orbit
$C_1(z)$ is constant (ellipticity) and $f_2/f_1\to0$, $g^{[\lambda]}\to1$, so
$C_1([u])=1/(q;q)_\infty$ on every class $\ne[1]_Q$, and by meromorphy $C_1\equiv1/(q;q)_\infty$: for every regular $\lambda$. In particular $C_1$ is $z$- and $\lambda$-independent.
\end{proof}

\begin{remark}[Diophantine consequences at $z^\ast$]
\label{rem:zeroproduct-dioph}
At $q=q^\ast$ the identity splits off the candidate rational zero:
\[
   \prod_{k\ge2}z_k(q^\ast)\,Q^{\,k-1}\;=\;\frac{(q;q^2)_\infty}{z^\ast}
   \;=\;0.988299884628746\ldots\quad(\text{verified to }30\text{ digits}).
\]
Were $\beta_2$ rational, the \emph{modular} value $(q;q^2)_\infty$ (transcendental at algebraic $q$ by
Nesterenko's theorem applied to the Euler products) would equal the rational $z^\ast$ times the tail
dressing product $\prod_{k\ge2}z_kQ^{k-1}$. Thus the entire modular sector of the problem is now
explicit; the Stokes-invariant connection constant is $1/(q;q)_\infty$, the zero product is
$(q;q^2)_\infty$; and \emph{all} remaining unknown arithmetic is provably confined to the tail
dressing product, equivalently to the $q$-Stokes cocycle $C_2^{[\lambda]}$: the non-modularity of the
$\beta_2$ problem \emph{is} its $q$-Stokes data. A transcendence statement for that tail (an algebraic
independence of the higher-zero dressings from the modular values, or a rigidity theorem for the
$q$-Stokes cocycle) would settle $\beta_2$; no current theorem provides it.

The cocycle itself is now explicit. Writing $g=g^{[\lambda]}$ for the resummed series and $h$ for
$f_2$'s convergent one,
\[
   C_2^{[\lambda]}(z)\;=\;\theta_Q(-z/q)\,
   \frac{g(z)\,G(Qz)+z^{-1}g(Qz)\,G(z)}
        {-(z/q)\,g(z)\,h(Qz)+z^{-1}g(Qz)\,h(z)}
\]
This is an identity: it is the ratio of the Casoratian evaluations
$\operatorname{Cas}(f_1,G)=\theta_Q(-z)\bigl[g(z)G(Qz)+z^{-1}g(Qz)G(z)\bigr]$ and
$\operatorname{Cas}(f_1,f_2)=\tfrac{\theta_Q(-z)}{\theta_Q(-z/q)}
\bigl[-(z/q)g(z)h(Qz)+z^{-1}g(Qz)h(z)\bigr]$, both instances of
$\theta_Q(Qw)=w^{-1}\theta_Q(w)$ (and verified numerically to $10^{-61}$); with divisor now fully
described: an
\emph{explicit} zero lattice $z\in q^{-1}Q^{\Z}$ from the theta factor (the computed zero at
$z=1/q$, exact to $10^{-63}$), a direction-dependent zero at $z=\lambda^2$ arising from the exact
\emph{resummation resonance}
$g^{[\lambda]}(\lambda^2)G(Q\lambda^2)+\lambda^{-2}g^{[\lambda]}(Q\lambda^2)G(\lambda^2)=0$
(verified to $10^{-62}$ at $\lambda\in\{\tfrac32,2,\tfrac{13}5\}$; a proof for general regular
$\lambda$ is open),
and poles at the zeros of the $h$-bracket. Substituting into the pinned zero condition at $z^\ast$
and cancelling the common theta factor reduces the entire $\beta_2$ question to \emph{one} relation
between the two modular values $\theta_Q(-z^\ast),\ (q;q)_\infty$ and five values of the non-modular
pair $(g^{[\lambda]},h)$ on the orbit of $z^\ast$. The
remaining input is a transcendence/independence statement for those five values; supplying it is one
of three research programs (a non-modular extension of Nesterenko's method via the deformation of this
connection, a $\rho<\tfrac12$-regime irrationality criterion, or adelic rigidity
of the vanishing locus); none is currently available.
\end{remark}

The first of those programs now has its foundation: the deformation of the module in $q$ closes
finitely.

\begin{proposition}[the deformation module]
\label{prop:deformation}
Let $L=q-a\sigma+\sigma^2$ with $a=q+1-qz$, $\sigma f(z)=f(Qz)$, and let $\theta=z\partial_z$,
$\delta=q\partial_q$ (acting on all $q$-dependence, including the base $Q=q^2$, so that
$\delta\sigma=\sigma\delta+2\sigma\theta$). Then
\[
   L(\theta G)=-qz\,\sigma G,
   \qquad
   L(\delta G)=-qG+q(1+3z)\,\sigma G+4q\,\theta G-2a\,\sigma\theta G
\]
(two commutator computations; verified to $51$ digits). Consequently the
$\Q(q,z)\langle\sigma\rangle$-module generated by $\{G,\theta G,\delta G\}$ has finite rank $\le6$,
with explicit $\sigma$-companion relations $\sigma^2G=a\sigma G-qG$,
$\sigma^2\theta G=a\sigma\theta G-q\theta G-qz\,\sigma G$, and
$\sigma^2\delta G=a\sigma\delta G-q\delta G-qG+q(1+3z)\sigma G+4q\theta G-2a\,\sigma\theta G$.

This finiteness is of the \emph{first jet only}: the module is \emph{not} closed under a second
application of $\theta$ or $\delta$. Both $\theta^2G$ and $\delta^2G$ escape
$\operatorname{span}_{\Q(q,z)}\{G,\sigma G,\theta G,\sigma\theta G,\delta G,\sigma\delta G\}$
(verified by exact modular linear algebra at $q=\tfrac13$, $z$-series to $70$ terms, coefficient
degrees $\le6$: the augmented rank strictly exceeds the plain rank). The two derivations are moreover
linked by the exact identity
\[
   \delta G=(\theta^2-\theta)G+C,\qquad
   C=\sum_{k\ge0}(-1)^k\Bigl(\textstyle\sum_{i=1}^{2k}\tfrac{iq^i}{1-q^i}\Bigr)
   \frac{q^{k(k-1)}z^k}{(q;q)_{2k}}
\]
($=z^2G_{zz}+C$; immediate from $g_k$'s logarithmic $q$-derivative), whose weight $C$ carries the
truncated Eisenstein sums $\sum_{i\le2k}iq^i/(1-q^i)$ term by term; the same off-modular data that
obstructs the differential closure. In particular, differentiating the pinned relation of
Proposition~\ref{prop:pinned} in $q$ does \emph{not} close: the second $q$-derivative leaves the
module, consistent with Proposition~\ref{prop:nofinitering}. The value of the rank-$6$ structure is at
the \emph{function} level; it is the module whose difference Galois group is computed in
Theorem~\ref{thm:sl2}; not as a Nesterenko precondition in the descent variable, which
Proposition~\ref{prop:nofinitering} shows does not exist.
\end{proposition}

\begin{corollary}[equation of the zero curve]
\label{cor:zerocurve}
Implicit differentiation of $G(q,z_1(q))=0$ gives
\[
   z_1'(q)\;=\;-\,\frac{z_1\,(\delta G)(q,z_1)}{q\,(\theta G)(q,z_1)},
\]
with $\theta G,\delta G$ the first-jet generators of Proposition~\ref{prop:deformation} (formula
checked against a finite-difference derivative to $20$ digits at $q=0.3$). The right-hand side is not a
closed vector field; iterating $d/dq$ leaves the rank-$6$ module (Proposition~\ref{prop:deformation}); so this is one exact relation, not an autonomous system. The number $q^\ast$ is the crossing of the
zero curve with the diagonal: $z_1(q^\ast)=2q^\ast(1-q^\ast)$ (residual $<10^{-27}$).
\end{corollary}

\begin{lemma}[irreducibility of the $\sigma$-module]
\label{lem:irreducible}
For $0<q<1$ the rank-$2$ $\sigma$-module of $G$ (in the sense of \cite{PutSinger1997}) is
irreducible over $\C(z)$: the Riccati equation
\[
   u(z)\,u(Qz)-a(z)\,u(z)+q=0,\qquad a=q+1-qz,
\]
has no solution $u\in\C(z)$.
\end{lemma}

\begin{proof}
Suppose $u=P/R$ with $P,R\in\C[z]$ coprime, of degrees $p,r$, leading coefficients $\alpha,\beta$.
Clearing denominators,
\[
   F(z)\;:=\;P(Qz)P(z)-a(z)P(z)R(Qz)+q\,R(z)R(Qz)\;\equiv\;0 .
\]
\emph{At $z=0$:} $F(0)=P(0)^2-(q+1)P(0)R(0)+qR(0)^2=0$; if $R(0)=0$ then $P(0)=0$, contradicting
coprimality, and symmetrically; so $P(0)R(0)\ne0$ and $u(0)\in\{1,q\}$.
\emph{Top degree:} the three terms have degrees $2p$, $p+r+1$, $2r$. If $p=r$, the top degree
$2p+1$ is attained by the middle term alone, with coefficient $q\alpha\beta Q^{r}\ne0$;
if $|p-r|\ge2$, the top degree is attained by the first or third term alone, with coefficient
$\alpha^2Q^{p}$ or $q\beta^2Q^{r}$. Both are impossible, so $|p-r|=1$.
\emph{Root pairing:} let $\zeta\ne0$ be a root of $R$ of multiplicity $m$, so $P(\zeta)\ne0$. At
$z=\zeta/Q$ the three terms of $F$ have orders $k$, $k+m+[\,a(\zeta/Q)=0\,]$, $m_0+m$, where
$k=\operatorname{ord}_{\zeta/Q}P$ and $m_0=\operatorname{ord}_{\zeta/Q}R$; since $F\equiv0$ the
minimum order cannot be attained by a single term, forcing $k\ge m_0+m\ge m$. Symmetrically, a root
$\xi\ne0$ of $P$ of multiplicity $k$ forces $\operatorname{ord}_{Q\xi}R\ge k$ (coprimality gives
$\operatorname{ord}_{Q\xi}P=0$, so the first term has order $k$ at $\xi$, the third
$\operatorname{ord}_{Q\xi}R$, and the middle at least $k+\operatorname{ord}_{Q\xi}R$; again no
single minimum). The two maps $\zeta\mapsto\zeta/Q$, $\xi\mapsto Q\xi$ are mutually inverse on the
nonzero root sets, and chains close at length two by coprimality. Summing multiplicities,
$p\le r\le p$, so $p=r$; contradicting $|p-r|=1$.
\end{proof}

\begin{proposition}[the symmetric square]
\label{prop:sym2}
Products $w=y\tilde y$ of solutions of $L$ satisfy the third-order equation
\[
   w(Q^3z)\;=\;\frac{A_1(A_0A_1-q)}{A_0}\,w(Q^2z)\;+\;q\,(q-A_0A_1)\,w(Qz)\;+\;\frac{q^3A_1}{A_0}\,w(z),
\]
with $A_0=a(z)$, $A_1=a(Qz)$, obtained by eliminating the mixed product
$m=y\,\tilde y(Q\,\cdot)+\tilde y\,y(Q\,\cdot)$ through $m(Qz)=2A_0\,w(Qz)-q\,m(z)$. Its
characteristic roots at $z=0$ are $\{1,q,q^2\}$; verified on the solutions $G^2$, $z^{1/2}GH$,
$zH^2$ to $41$--$43$ digits.
\end{proposition}

\begin{lemma}[divisor orientation]
\label{lem:divisor}
Let $r\in\C(z)^\times$ satisfy the rational-quotient relation of Proposition~\ref{prop:sym2},
$r_0r_1r_2=c_2r_0r_1+c_1r_0+c_0$ with $r_k=r(Q^kz)$, $c_2=A_1(A_0A_1-q)/A_0$, $c_1=q(q-A_0A_1)$,
$c_0=q^3A_1/A_0$. On every $Q$-orbit in $\C^*$ disjoint from the singular orbit $z_aQ^{\Z}$
($z_a=(q+1)/q$, the zero of $a$), the maximal-modulus zero-or-pole of $r$ cannot be a pole, and the
minimal-modulus zero-or-pole cannot be a zero.
\end{lemma}

\begin{proof}
Off $z_aQ^{\Z}$ the $c_i$ are finite and $c_0\ne0$. If the maximal-modulus zero-or-pole $\zeta$ on
such an orbit were a pole, evaluate at $z=Q^{-2}\zeta$: there $r_0,r_1$ are finite and nonzero (all
higher-modulus orbit points carry no zero or pole), $r_2$ has a pole, so the left side has a pole
while the right side is finite. If the minimal-modulus zero-or-pole $\zeta'$ were a zero, rewrite
the relation for $s=1/r$ as $c_0s_0s_1s_2+c_1s_1s_2+c_2s_2=1$ and evaluate at $z=\zeta'$: $s_0$ has
a pole, $s_1,s_2$ are finite and nonzero, $c_0(\zeta')\ne0$, so the left side has a pole while the
right side is $1$.
\end{proof}

\begin{theorem}[the difference Galois group contains $\mathrm{SL}_2$]
\label{thm:sl2}
For every $0<q<1$, the difference Galois group over $\C(z)$ of the module of $G$ contains
$\mathrm{SL}_2$.
\end{theorem}

\begin{proof}
For $N\ge1$ suppose $\operatorname{Sym}^N\!M$ has a rank-$1$ submodule, i.e.\ there is a nonzero
solution $w$ of the $N$-th symmetric-power equation with $r:=w(Q\,\cdot)/w\in\C(z)$. The products
$f_1^{[\lambda]\,N-k}f_2^{\,k}$ ($0\le k\le N$) form a solution basis (their Casoratian is a power
of $\operatorname{Cas}(f_1,f_2)\ne0$; checked numerically for $N=2$), so
$w=\sum_kc_kf_1^{N-k}f_2^k$ with constants $c_k$, and the case split is $\lambda$-robust since
Stokes jumps preserve $k_0:=\min\{k:c_k\ne0\}$.

\emph{Case $k_0<N$.} Along a ray in a regular class, $f_2/f_1$ decays super-polynomially, so the
rational function $r$ has the Poincar\'e asymptotic expansion of the quotient of the dominant
component:
\[
   r\;\sim\;(-1/z)^{N-k_0}(-z/q)^{k_0}\,
   \Bigl(\tfrac{\hat g(Qz)}{\hat g(z)}\Bigr)^{N-k_0}
   \Bigl(\tfrac{h(Qz)}{h(z)}\Bigr)^{k_0}.
\]
The order-$j$ coefficient $d_j$ of $\hat g(Qz)/\hat g(z)$ satisfies
$d_j\sim-b_j$ with $b_j=q^{\,(3j-j^2)/2}/(q;q)_j$ (verified: $d_j/(-b_j)=0.99986,\,0.99998,\,
0.999998$ at $j=10,12,14$, $q=\tfrac13$, converging to $1$). Since this factor enters $r$ to the
power $N-k_0\ge1$ against rational, hence geometrically bounded, prefactors, the coefficients of $r$
inherit the growth $b_j\sim q^{-j^2/2}$, which is super-geometric, whereas a rational function's
expansion at $\infty$ has geometrically bounded coefficients. Since matching asymptotic
expansions are unique, no rational $r$ exists.

\emph{Case $k_0=N$.} Then $r=(-z/q)^N\rho^N$ with $\rho:=h(Q\,\cdot)/h$ meromorphic on $\C^*$;
$\rho^N\in\C(z)$ and single-valuedness force $\rho\in\C(z)$ (an $N$-th root of a rational function
is single-valued on $\C^*$ only if every divisor multiplicity, including at $0$ and $\infty$, is
divisible by $N$). Then $f_2(Qz)/f_2(z)=-(z/q)\rho(z)$ is rational, giving a rational solution of
the rank-$2$ Riccati equation; contradicting Lemma~\ref{lem:irreducible}.

Hence no $\operatorname{Sym}^N\!M$ has a rank-$1$ submodule. $N=1$ is irreducibility; $N=2$
excludes the imprimitive (dihedral) case; and a projectively finite irreducible subgroup
($A_4,S_4,A_5$) fixes a line in some $\operatorname{Sym}^N$ with $N\le12$ (the classical binary
polyhedral invariants), which is likewise excluded. By the classification of algebraic subgroups of
$\mathrm{GL}_2$, the group contains $\mathrm{SL}_2$.
\end{proof}

\begin{proposition}[no $\sigma\delta$-integrability]
\label{prop:nonintegrable}
For every $0<q<1$ there is no $B\in\mathrm{gl}_2(\C(z))$ and no $c\in\C(z)$ with
$\theta(A)=\sigma(B)A-AB+cA$, where $A=\bigl(\begin{smallmatrix}0&1\\-q&a\end{smallmatrix}\bigr)$ and
$\theta=z\partial_z$.
\end{proposition}

\begin{proof}
Writing $B=\bigl(\begin{smallmatrix}\alpha&\beta\\\gamma&\delta\end{smallmatrix}\bigr)$, the entries give
$\gamma=-q\beta_1$, $\delta=\alpha_1+a\beta_1-c$, and eliminating $\alpha$ (the central terms $c$
cancel identically; verified symbolically) leaves the necessary condition
\[
   qa\,\beta_3\;=\;a_1(aa_1-q)\,\beta_2\;-\;a(aa_1-q)\,\beta_1\;+\;qa_1\,\beta\;+\;q\,z\bigl(a_1+q^2a\bigr),
\]
$\beta_k=\beta(Q^kz)$, $a=a(z)$, $a_1=a(Qz)$. A polynomial part of degree $m\ge0$ contributes top
coefficient $\propto\beta_m(1-q^{2+2m})\ne0$ at degree $m+3$, above the inhomogeneity: so $\beta$
has no polynomial part. If $\beta$ has a pole off the orbit $z_aQ^{\Z}$, the maximal-modulus pole
appears in exactly one term (the slot $z=\zeta Q^{-3}$, coefficient $qa\ne0$ there), and the
minimal-modulus pole likewise (slot $z=\zeta'$, coefficient $qa_1$): both impossible; so the poles
lie in the window $\{z_aQ^{-1},\dots,z_aQ^{3}\}$ ($a(z_aQ^k)=(q+1)(1-Q^k)$ vanishes only at $k=0$).
Sweeping the slots $z=z_aQ^{-4},\dots,z_aQ^{-1}$ kills the orders at $z_aQ^{-1},\dots,z_aQ^{2}$ one
by one (each pole appears in a single term with nonvanishing coefficient), and the order at
$z_aQ^{3}$ is killed at the slot $z=z_aQ$; coefficient $a_1(aa_1-q)$ with
$(aa_1-q)(z_aQ)=(q+1)^2(1-q^2)(1-q^4)-q$; or, where that vanishes (one root
$q\approx0.78027$), at the slot $z=z_aQ^{2}$; coefficient $a(aa_1-q)$ with
$(aa_1-q)(z_aQ^2)=(q+1)^2(1-q^4)(1-q^6)-q$ (one root $q\approx0.86753$); the two polynomials are
coprime, so for every $q\in(0,1)$ at least one slot applies. Hence $\beta\equiv0$, contradicting
$z(a_1+q^2a)=z\bigl[(q+1)(1+q^2)-2q^3z\bigr]\not\equiv0$. (Cross-check: at $q=\tfrac13$ the exact
linear system for poles of order $\le3$ in the window has $\operatorname{rank}(A)=15$,
$\operatorname{rank}[A|b]=16$.)
\end{proof}

\begin{corollary}[differential transcendence of $G$]
\label{cor:hypertranscendence}
For every $0<q<1$ the function $G(q,\cdot)$ satisfies no algebraic differential equation over
$\C(z)$: the Hahn--Exton $q$-cosine is differentially transcendental.
\end{corollary}

\begin{proof}[Proof, via parametrized difference-Galois theory]
We apply the parametrized ($\sigma,\theta$)-difference Galois theory with derivation
$\theta=z\partial_z$, over the difference-Galois formalism of \cite{PutSinger1997} and in the
form of \cite{DHR2018,ADR2021}. The determinant equation $\sigma(y)=\det(A)\,y=q\,y$ has the solution
$y=z^{1/2}$ (indeed $\sigma(z^{1/2})=(Qz)^{1/2}=q\,z^{1/2}$, as $Q=q^2$), which is algebraic over $\C(z)$; hence the parametrized difference Galois group of $\sigma(y)=\det(A)y$
lies in $\C^\times$. This is exactly the hypothesis of Dreyfus--Hardouin--Roques
\cite[Prop.~2.9]{DHR2018}, which for a rank-$2$ system whose difference Galois group contains
$\mathrm{SL}_2$ (here Theorem~\ref{thm:sl2}) gives the dichotomy: the parametrized group
$\mathrm{Gal}^\theta$ either is conjugate into $\mathrm{GL}_2(\C)$ (all solutions differentially
algebraic) or contains $\mathrm{SL}_2(\widetilde{\mathcal C})$ (all nonzero solutions
differentially transcendental), and the \emph{first} alternative holds if and only if there exists
$B\in\mathrm{gl}_2(\C(z))$ with $\theta(A)=\sigma(B)A-AB$ (their integrability condition~(7);
$\theta(\det A)=0$ since $\det A=q$ is constant in $z$, so no trace term appears). By
Proposition~\ref{prop:nonintegrable} no such $B$ exists; indeed none exists even allowing a
scalar twist $cA$. Therefore the second alternative holds:
$\mathrm{Gal}^\theta\supseteq\mathrm{SL}_2(\widetilde{\mathcal C})$, and every nonzero solution,
in particular $G(q,\cdot)$, is differentially transcendental over $\C(z)$.
\end{proof}

The transcendence of the \emph{number} $q^\ast$ is a value-level statement, and the natural route (Nesterenko's multiplicity method fed by Corollary~\ref{cor:hypertranscendence}) meets a precise
structural obstruction, which we isolate.

\begin{proposition}[where the differential closure fails in the descent variable]
\label{prop:nofinitering}
Write $\delta=q\partial_q$, $\theta=z\partial_z$, and let $u_j(q):=G(q,Q^jz^\ast(q))$ denote the
orbit values along the diagonal $z^\ast(q)=2q(1-q)$.
\begin{itemize}
\item[\textup{(i)}] \textup{(orbit collapse)} The $\sigma$-orbit is \emph{not} an obstruction to
finiteness: the $z$-recurrence evaluated on the diagonal has coefficients
$a(Q^jz^\ast)=q+1-qQ^jz^\ast\in\Q(q)$, so $u_{j+2}=a(Q^jz^\ast)\,u_{j+1}-q\,u_j$ and every orbit
value lies in $\operatorname{span}_{\Q(q)}\{u_0,u_1\}$ \textup{(verified for $j\le3$ to $32$
digits)}.
\item[\textup{(ii)}] \textup{(the $E_2$-coupled jet structure)} The $q$-derivative of $G$ obeys the
exact identity
\[
   \delta G=(\theta^2-\theta)G+\varphi(q)\,G-T,\qquad
   \varphi(q)=\sum_{i\ge1}\frac{iq^i}{1-q^i},\qquad
   T=\sum_{k\ge0}(-1)^k\Bigl(\sum_{i>2k}\tfrac{iq^i}{1-q^i}\Bigr)\frac{q^{k(k-1)}z^k}{(q;q)_{2k}}
\]
\textup{(verified to $41$ digits)}, where $\varphi$ is the quasimodular Eisenstein weight
($24\varphi=1-E_2$ up to normalization) and the correction $T$ has coefficients smaller than $G$'s by
the factor $\sum_{i>2k}iq^i/(1-q^i)=O(kq^{2k})$. Expanding these weights as Lambert series and
interchanging the (absolutely convergent) sums exhibits $T$ as an Eisenstein-weighted trace of the
module over the forward $\sigma$-orbit:
\[
   T(q,z)\;=\;\sum_{d\ge1}\frac{2q^d}{1-q^d}\,(\theta G)(q,q^{2d}z)
   \;+\;\sum_{d\ge1}\frac{q^d}{(1-q^d)^2}\,G(q,q^{2d}z)
\]
\textup{(verified to $41$ digits)}: the entire non-modular content of $\delta G$ is an orbit
convolution of $(G,\theta G)$ against $E_2$-type weights. On the diagonal the orbit values collapse
by \textup{(i)} into four generators $u_0,u_1,\theta$-analogues $v_0,v_1$ (the $v$'s obey the
inhomogeneous transported recurrence $v_{j+2}=a_jv_{j+1}-qv_j-qQ^jz^\ast u_{j+1}$), but the collapse
coefficients accumulate into transfer-product series $A,B,C,D\in\Z[[q]]$ (integrality is immediate:
transfer data and Lambert weights lie in $\Z[[q]]$), which are \emph{not} of any classical type: they
satisfy no relation over $\operatorname{span}\{1,\varphi,\sum\sigma_3(n)q^n\}$ with rational-function
coefficients of degree $\le8$, and no linear $q$-ODE of order $\le3$ and degree $\le6$ (exact linear
algebra on $45$ coefficients). The identity relocates, and does not remove, the obstruction. Since $\delta$ and $\theta$ commute, each
application of $\delta$ raises the $z$-jet order by two: $\delta^nS$ involves
$\theta^{2n}G|_{z=z^\ast}$ with unit leading coefficient.
\item[\textup{(iii)}] \textup{(the actual obstruction)} The $z$-jets do not close: $\theta^2G$ and
$\delta^2G$ escape the rank-$6$ module of Proposition~\ref{prop:deformation}, and by
Corollary~\ref{cor:hypertranscendence} the tower $\{\theta^mG\}_m$ satisfies no algebraic relation
over $\C(z)$. A finitely generated $\delta$-closed $\C(q)$-algebra containing $S$ would force an
algebraic relation over $\C(q)$ among the diagonal restrictions
$\{\theta^mG(q,z^\ast(q))\}_{m\le M}$ for some finite $M$; a statement equivalent to $S$ being
differentially algebraic in $q$, for which there is no supporting structure and proof-grade
negative evidence across a wide window: no algebraic differential equation at profiles up to
$(2,12,3)$, $(3,8,3)$, $(4,5,6)$, $(5,4,8)$, $(6,3,15)$, $(8,2,35)$, $(12,1,120)$ in
(order, degree, $\deg_q$), by full column rank on $2000$ coefficients modulo a prime
(\texttt{beta2\_diffalg\_search.py}, \texttt{beta2\_ade\_sweep.py}; $E_2$/$q$-Chazy control
recovered).
\end{itemize}
Thus the Nesterenko route closes \emph{if and only if} $S$ is differentially algebraic in $q$; the
obstruction is carried by the $z$-jet tower; whose function-level independence is exactly
Corollary~\ref{cor:hypertranscendence}; and not by the $\sigma$-orbit, which collapses to rank
two.
\end{proposition}

\begin{remark}[the value-level descent as a joint multiplicity estimate]
\label{rem:valuedescent}
Nesterenko's method derives transcendence of a value from three ingredients: a finite ring closed
under the relevant operator, a \emph{multiplicity estimate} bounding
$\operatorname{ord}P(\text{generators})\le c\,(\deg P)^\kappa$, and routine analytic size/Siegel-lemma
constructions. Theorem~\ref{thm:sl2} and Proposition~\ref{prop:nonintegrable} furnish precisely the
nondegeneracy a multiplicity estimate consumes (no invariant subvariety; no first integral). But
Proposition~\ref{prop:nofinitering} shows the first ingredient fails in the descent variable $q$: the
finiteness of the $\beta_2$ module is intrinsically \emph{joint} in $(\sigma_z,\delta_q)$, and
$\sigma_z$ is a non-infinitesimal $q$-lattice shift, outside the Pfaffian/differential multiplicity
theories of Nesterenko and Binyamini. Thus the transcendence of $q^\ast=\beta_2^{-2}$ reduces to a
\emph{joint $q$-difference--differential multiplicity estimate} for the rank-$\le6$ module at
$(q^\ast,z^\ast)$: a bound $\operatorname{ord}_{(q^\ast,z^\ast)}P(G,\theta G,\delta G)\le c\,N^\kappa$
($N=\deg P$) in the local ring, uniform along the $\sigma$-orbit. Such an estimate for finite
$(\sigma,\delta)$-modules is not available in the literature; it is the sole non-routine input
remaining, and the precise open problem on which the (ir)rationality of $\beta_2$ now turns.
\end{remark}

\begin{remark}[the obstruction inside $q$-Siegel--Shidlovskii theory]
\label{rem:qSS}
The same wall appears, quantified, in the value-transcendence theory built on Siegel's method. The
approach of Amou--Matala-aho--V\"a\"an\"anen \cite{AMV2007} (and the Matala-aho-school Baker-type
refinements) proves non-vanishing of values of Poincar\'e-type $q$-difference solutions at algebraic
points under two hypotheses: the system is Poincar\'e-type at $0$ (met here: $A(0)$ is invertible,
$\det A=q$), and a \emph{Siegel margin} $\gamma<\Gamma$ holds, where $\gamma$ is the ratio of the
logarithmic height of the cleared denominators to that of the analytic decay. For the Hahn--Exton
class that ratio is exactly $\gamma=\rho=\log(t/s)/(2\log t)\le\tfrac12$ of Proposition~\ref{prop:nogo}: the double Pochhammer $(q;q)_{2k}$ is the factor $2$; so the margin \emph{fails}, and the
theory does not apply. Thus Conjecture~Z (transcendence of the Hahn--Exton zeros at algebraic base,
which by $z^\ast=z_1(q^\ast)$ implies $\beta_2$ transcendental) reduces to extending $q$-Siegel--
Shidlovskii past the margin for a class that is additionally irregular at $\infty$. The connection
theory of this paper supplies the irregular-$\infty$ data explicitly; the margin itself is the
value-theoretic face of the multiplicity estimate of Remark~\ref{rem:valuedescent}, and is equal in
strength to $\beta_2$-transcendence, not weaker. A detailed program is recorded in the repository
(\texttt{paper/journal/qSS\_program.md}).

One further mechanism deserves record, because it isolates which half of the Siegel scheme is
actually blocked. Under the rationality hypothesis all tail power sums
$\tau_p=\sum_{k\ge2}z_k^{-p}=\sigma_p-z^{\ast-p}$ are rational, and they are the moments of a
positive measure with infinite support (unit atoms at $1/z_k$); hence every Hankel determinant
$\det[\tau_{i+j+2}]_{i,j\le n}$ is \emph{strictly positive for free}; no auxiliary nonvanishing
argument is needed; and, being a positive rational, is bounded below by the reciprocal of its
denominator ($\le t^{8n^3/3+O(n^2)}$). Analytically these determinants decay like
$Q^{n^3/3(1+o(1))}$ (atoms at $\approx Q^{k-1}$; verified: strictly positive with super-quadratic
exponent growth at $q^\ast$). The comparison $2\log(t/s)>8\log t$ fails for every rational $q$: the
positivity route conserves the ledger at a factor $4$. The instructive point is that
\emph{nonvanishing was never the bottleneck} of any construction in this ledger; positivity
supplies it gratis here; the entire difficulty of the problem resides in the size-versus-height
comparison. Nor is there an Ap\'ery-type denominator miracle hiding in that comparison: the
\emph{measured} reduced denominators of the partial sums $S_N(s/t)$ equal $96$--$100\%$ of the naive
clearing $\prod_{i\le2N}(t^i-s^i)$ in logarithmic size, with the ratio tending to $1$ as $N$ grows
(exact rational arithmetic, $N\le12$, at $q=\tfrac12$ and $\tfrac25$); the cancellation is an
$O(N)$ boundary effect in an $O(N^2)$ exponent. The arithmetic of the ledger is exactly its
worst-case costing.
\end{remark}

\begin{remark}[backward-orbit normal form of the zero condition]
\label{rem:backward}
The zero condition admits a compression into a single explicit rational sequence. If $G(q,z_1)=0$,
the three-term recurrence along the backward $\sigma$-orbit, $u_{-m}=\bigl(a(Q^{-m}z_1)u_{-m+1}
-u_{-m+2}\bigr)/q$ with $u_n=G(q,Q^nz_1)$, loses one solution constant: $u_{-m}=u_1\,r_m(q)$ with
$r_0=0$, $r_1=-1/q$, and $r_m\in\Q(q,z_1)$ computable by the same recurrence. The asymptotics of
Theorem~\ref{thm:zeroproduct} then give the exact evaluation
\[
   \lim_{m\to\infty} r_m(q)\,Q^{m(m+1)/2}(-z_1)^{-m}
   \;=\;\frac{\theta_Q(-z_1)}{(q;q)_\infty\,G(q,Qz_1)},
\]
with geometric rate $Q^m$ and $q$-Gevrey correction series $\sum_jb_jz_1^{-j}Q^{mj}$ (verified to
$17$ digits at $q^\ast$ by $m=24$). Under the hypothesis $q^\ast=s/t$ the normalized terms are
rationals of height $(st)^{m^2+O(m)}$ approximating the limit to $O(Q^m)$; accelerated forms that
cancel $J$ correction terms gain decay $Q^{mJ}$ but pay the quadratic heights of the $b_j$, and the
arithmetic ratio stays $\le\tfrac12$; the balance of Proposition~\ref{prop:nogo} once more. The
normal form is recorded as the sharpest single-sequence statement of what any future method must
exploit: the entire rationality hypothesis is equivalent to the arithmetic of this one explicit
sequence. The $\theta$-jet admits the same normal form one level up: along the backward orbit,
\[
   \frac{(\theta G)(q,Q^{-m}z_1)}{G(q,Q^{-m}z_1)}-m
   \;\longrightarrow\;
   -\,\frac{z_1\,\theta_Q'(-z_1)}{\theta_Q(-z_1)}
   \qquad(\text{rate }Q^m;\ \text{verified to }10\text{ digits at }q=0.3),
\]
an exact theta log-derivative: each backward step increments $\theta\log\theta_Q(-z)$ by exactly one
(quasi-periodicity), and the $\hat g$-correction vanishes in the limit.

The \emph{forward} orbit admits the complementary normal form, in the language of rational dynamics.
The Riccati relation gives the explicit continued fraction
$q/(a(z)-q/(a(Qz)-\cdots))$, which by Pincherle's theorem \cite{LorentzenWaadeland1992} converges to the minimal-solution ratio
$q\,H(Qz)/H(z)$; the zero condition $G(q,z_1)=0$ forces the \emph{dominant} ratios
$w_n:=G(q,Q^{n+1}z_1)/G(q,Q^nz_1)$ to obey $w_1=a(z_1)$ exactly and thereafter the rational map
\[
   w_{n+1}=a(Q^nz_1)-\frac{q}{w_n},
\]
so that under the hypothesis $q,z_1\in\Q$ the entire forward orbit is an explicit $\Q$-orbit
converging to the attracting fixed point $1$ with $w_n-1=\Theta(Q^n)$ (verified to $25$ digits at
$q^\ast$). Its heights compound quadratically ($t^{O(n^2)}$), so the Liouville comparison
$Q^n$ vs.\ $t^{-n^2}$ is again consistent: both normal forms; backward theta-growth and forward
rational dynamics; conserve the balance of Proposition~\ref{prop:nogo}, and together they exhaust
the one-orbit reformulations of the zero condition.

Combining the two orbits with the $q$-Wronskian makes the relation-count of the rationality
hypothesis explicit. Under $G(q,z_1)=0$ the unknown transcendental data reduces to three numbers ($g_1:=G(q,Qz_1)$, $H(q,z_1)$, and the backward limit $\rho_\infty$), bound by exactly two
relations: $g_1H(q,z_1)=1-q$ (the $q$-Wronskian at the zero) and
$\rho_\infty g_1=\theta_Q(-z_1)/(q;q)_\infty$ (the backward normal form). Eliminating $g_1$,
\[
   \frac{\rho_\infty}{H(q,z_1)}=\frac{\theta_Q(-z_1)}{(1-q)\,(q;q)_\infty}
\]
(verified to $32$ digits at $q^\ast$), with modular right-hand side. Three unknowns, two relations:
the rationality hypothesis is consistent at this level, and a proof requires precisely one further
independent relation; the multiplicity statement of Remark~\ref{rem:valuedescent}, in yet another
equivalent form.
\end{remark}

\begin{remark}[truncation roots, and the exact deficit]
\label{rem:truncroot}
A scheme not covered by the lattice ledger prices the wall in its simplest form. Clearing the
truncation of $S$ gives integer polynomials
$P_N=(q;q)_{2N}\sum_{k\le N}(-2)^k q^{k^2}(1-q)^k/(q;q)_{2k}\in\Z[q]$ of degree $2N^2+N$, whose
heights are only $e^{O(N)}$ (the pentagonal sparsity of $q$-Pochhammer coefficients; measured
$\log_2H=2.6,4.6,\dots,12.1$ for $N=2,\dots,7$), and whose root nearest $q^\ast$ satisfies
$|q_N-q^\ast|\asymp q^{\ast\,(N+1)^2}$ (verified $N\le7$). If $q^\ast=s/t$ then
$|P_N(q^\ast)|\ge t^{-(2N^2+N)}$, while analytically $|P_N(q^\ast)|=e^{O(N)}q^{\ast\,(N+1)^2}$;
contradiction would require $q^\ast<t^{-2}$, impossible for $q^\ast=s/t\ge1/t$. The deficit is
exactly the factor $2$ in the clearing degree; the double Pochhammer once more, now in
$s$-independent form. The same computation shows that an equivalent representation of the root
condition with \emph{single}-Pochhammer denominators (clearing degree $\sim N^2/2$) would yield
$q^\ast\ne s/t$ for all $s<\sqrt t$. No such representation exists within the
Rogers--Ramanujan/Bailey framework \cite{GasperRahman2004}: those transformations require the argument to be a $q$-power,
and the argument here is the catalytic point $z=2q(1-q)$. In this precise sense the obstruction to
$\beta_2$'s irrationality is the non-$q$-power (catalytic) argument itself.

The catalytic locus is moreover \emph{seat-independent}. Writing $S(q)={}_1\phi_1(0;q;Q,z^\ast)$
($Q=q^2$), the parameter--argument transformation
${}_1\phi_1(0;b;Q,x)=\tfrac{(x;Q)_\infty}{(b;Q)_\infty}\,{}_1\phi_1(0;x;Q,b)$ gives the
equivalent zero condition
\[
   T(q)\;:=\;{}_1\phi_1\bigl(0;\,2q(1-q);\,q^2,\,q\bigr)
   \;=\;\sum_{k\ge0}\frac{(-1)^kq^{k^2}}{(q^2;q^2)_k\,(2q(1-q);q^2)_k}\;=\;0
\]
(verified to $51$ digits; the prefactor cannot vanish since $z^\ast<1$), in which the argument is
the $q$-power $q$ and the catalytic value sits in the \emph{parameter}; equivalently, $q^\ast$ is
the point where the Hahn--Exton $q$-Bessel of $q$-dependent order $\nu(q)=\log_{q^2}(2q(1-q))-1$
vanishes at the fixed argument $q^{-1/2}$. The clearing degree of $T$'s truncations is again
$\sim2N^2$: the deficit is invariant under the swap, and in either seat it is the non-special value
$2q(1-q)$ that lies outside every Rogers--Ramanujan/Bailey lattice. The well-poised generalization does
not help: a WP-Bailey specialization producing a single-Pochhammer form would require $2q(1-q)$ to be
a $q$-power (the only special-value type such summations admit), whereas at $q=q^\ast$ one has
$q^\ast<2q^\ast(1-q^\ast)<1$, strictly between consecutive $q$-powers (the two loci where
$2q(1-q)$ is a $q$-power, $q=\tfrac12$ and $q=\tfrac23$, both miss $q^\ast$). The catalytic argument
is thus $q$-power--free in every seat, which is the transformation-theoretic content of the
$\rho=\tfrac12$ wall.
\end{remark}

\begin{proposition}[effective non-rationality of $q^\ast$]
\label{prop:effective}
The least positive root $q^\ast$ of $S$ admits no rational representation $s/t$ with
$t\le10^{1044}$. Consequently, if $\beta_2=u/v$ in lowest terms then $u>10^{522}$.
\end{proposition}

\begin{proof}
All steps are certified interval arithmetic (outward rounding; \texttt{mpmath.iv}; script
\texttt{beta2\_effective\_irrationality.py}), with the series tail bounded by
$(q;q)_{2k}\ge(1-2q)/(1-q)>0$ for $q<\tfrac12$ and term ratios $<\tfrac12$ beyond the truncation.
\textup{(i)} $S>0$ on $[0,0.40]$ ($400$ certified chunks). \textup{(ii)} $S'\le-1$ on $[0.40,0.46]$
($600$ certified chunks; $S'$ evaluated term-by-term with the logarithmic-derivative weights), so
$S$ has at most one zero in $[0.40,0.46]$. \textup{(iii)} with $c$ the
$2090$-digit truncation of $q^\ast$, $a=c-2\cdot10^{-2090}$ and $b=c+2\cdot10^{-2090}$, certified
evaluation gives $S(a)>0>S(b)$ (working precision $2130$ digits, series tail enclosed in
$\pm10^{-2125}$); by \textup{(i)}--\textup{(ii)}, $q^\ast\in[a,b]$. \textup{(iv)} the interval
$[a,b]$ contains the rational $s_0/t_0$ (Stern--Brocot construction) with
$9.5\cdot10^{1044}<t_0<10^{1045}$, verified by exact rational comparison, and certified
$S(s_0/t_0)\ne0$. If $m/n\in[a,b]$ with $n\le10^{1044}$ and $m/n\ne s_0/t_0$, then
$|m/n-s_0/t_0|\ge1/(nt_0)>10^{-2089}>b-a=4\cdot10^{-2090}$, impossible; so $m/n=s_0/t_0$, whence
$n=t_0>10^{1044}$, a contradiction. Thus no rational with denominator $\le10^{1044}$ lies in
$[a,b]$ at all, in particular none equals $q^\ast$. For $\beta_2=u/v$ in lowest terms,
$q^\ast=v^2/u^2$ in lowest terms, so $u^2>10^{1044}$.
\end{proof}


% ============================================================
\section{Proof of the block assembly \textup{(M3$'$)}}\label{app:M3prime}

\paragraph{Status.}
This appendix proves the assembly step \textup{(M3)} of Theorem~\ref{thm:model} in the corrected form
\eqref{eq:assembly}. It is a \emph{hand proof}. Two independent adversarial reviews reconstructed it
and found it correct; the second also compared an independent tuple enumerator against the
model, with no mismatch. The proof is \emph{not} formalised. The Lean development formalises the
objects it is about: the site costs, the cut predicate and
$\ell_R^{\mathrm{true}},c^{\mathrm{true}}$ (\texttt{Realisation.lean},
\texttt{CorrectedSpan.lean}), the junction-cut characterisation (\texttt{RJLemmaJ.lean}), and
the metric identity $|g|=\ell_R^{\mathrm{true}}(g)+2c^{\mathrm{true}}(g)$
(\texttt{metricAll}). The proof is supported by the computational checks recorded in
\S\ref{sec:bridge}: with $y$ symbolic, all $816$ coefficients with $\ell_R+2c\le31$ agree with
exhaustive enumeration in all three sectors ($5.03\times10^6$ elements), and a deliberately
corrupted model is detected; at $y=q$ the model agrees to $x^{26}$ and at $y=q^2$ to $x^{22}$. The
checks are evidence, not part of the proof.

In this appendix we write $k,\varepsilon,\delta$ for the marker data $k^\ast,\varepsilon^\ast,\delta^\ast$
of the main text, and \texttt{Elt}, $\mathrm{SameElt}$ for the element type of the Lean development
and its equality of the four normal-form fields
\cite[Theorem ``Normal form and lamp model'', \textup{thm:normal-form}]{Bonfioli2026a}.

\subsection{Setting and statement}

\paragraph{Elements.}
Edges and sites are indexed by $\mathbb Z$. Site $s$ lies between edge $s-1$ and edge $s$.
For $k\in\mathbb Z$, the \emph{travel indicator} is
\[
  f^{(k)}_j=\begin{cases}+1,&0\le j<k,\\-1,&k\le j<0,\\0,&\text{otherwise.}\end{cases}
\]
Let $\mathcal G$ be the set of tuples $(k,\varepsilon,\delta,d)$ such that
\begin{itemize}
\item $k\in\mathbb Z$, $\varepsilon\in\{\pm1\}$ and $\delta\in\{0,1\}$;
\item $d:\mathbb Z\to\mathbb Z$ has finite support;
\item $d_j\equiv f^{(k)}_j\pmod 2$ for all $j$.
\end{itemize}

\begin{lemma}[classes]\label{M3:lem:classes}
The map $g\mapsto(g.\mathtt{kstar},g.\mathtt{eps},g.\mathtt{delta},g.\mathtt d)$ induces a
bijection from $\mathrm{SameElt}$-classes of \texttt{Elt} onto $\mathcal G$. The functions
$\ell_R^{\mathrm{true}}$ and $c^{\mathrm{true}}$ of \texttt{CorrectedSpan.lean} are
constant on classes, so they are functions on $\mathcal G$.
\end{lemma}

\begin{proof}
$\mathrm{SameElt}$ is by definition equality of the four fields, so the induced map is well
defined and injective. The fields of an \texttt{Elt} satisfy the conditions defining
$\mathcal G$: \texttt{hpar} gives the parity condition, and \texttt{hsupp} gives finite support.

For surjectivity, take $(k,\varepsilon,\delta,d)\in\mathcal G$. Put
$\mathtt{supp}=\operatorname{supp}d\cup\{j:f^{(k)}_j\ne0\}$. This set is finite, and it gives an
\texttt{Elt} with the prescribed fields.

The only field not determined by the class is \texttt{supp}. It enters
$\ell_R^{\mathrm{true}}$ and $c^{\mathrm{true}}$ only through
$\mathtt{occTrue}=\mathtt{supp}\cap\{j:d_j\ne0\vee f_j\ne0\}$. By \texttt{hsupp}, this
intersection equals $\{j:d_j\ne0\vee f_j\ne0\}$, which is independent of \texttt{supp}. All
the other ingredients are functions of $(k,\varepsilon,\delta,d)$:
\texttt{mu}, \texttt{siteCost}, \texttt{cut} and \texttt{ShieldFires}.
\end{proof}

From now on we write $g=(k,\varepsilon,\delta,d)\in\mathcal G$ and $f=f^{(k)}$.

\paragraph{Local data.}
These are transcribed from \texttt{Realisation.lean}. Write $[\cdot]$ for the Iverson bracket.
\begin{align*}
  \alpha_s&=d_{s-1}-[s=0]+\varepsilon[s=k][\delta=0],&
  \beta_s&=d_s-\varepsilon[s=k][\delta=1],\\
  \Phi_s&=f_{s-1}+[s=0]-[s=k][\delta=0],&
  \operatorname{site}(s)&=\max(|\alpha_s|,|\beta_s|),\\
  \mu_j&=\begin{cases}2,&d_j=f_j=0,\\ \max(|d_j|,|f_j|),&\text{otherwise,}\end{cases}&
  \operatorname{cut}(s)&\iff\alpha_s=\beta_s=\Phi_s=0 .
\end{align*}
The term $-[s=0]$ is the virtual arrival \texttt{vArr}. It is the ``$\pm1$ correction'' at
the junction.

\paragraph{Corrected span.}
These are transcribed from \texttt{CorrectedSpan.lean}. Let
$\mathrm{occ}=\{j:d_j\ne0\vee f_j\ne0\}$. Then
\[
  A=\begin{cases}\min(0,\min\mathrm{occ}),&\mathrm{occ}\neq\varnothing,\\ 0,&\text{otherwise,}\end{cases}
  \qquad
  B=\begin{cases}\max(-1,\max\mathrm{occ}),&\mathrm{occ}\neq\varnothing,\\ -1,&\text{otherwise,}\end{cases}
\]
\[
  \ell_R^{\mathrm{true}}(g)=\sum_{j=A}^{B}\mu_j+\sum_{s=A}^{B+1}\operatorname{site}(s),\qquad
  c^{\mathrm{true}}(g)=\#\{s:A<s<B+1,\ \operatorname{cut}(s)\}
     +\bigl[\mathrm{SF}(g)\wedge\operatorname{cut}(0)\bigr],
\]
where the \emph{boundary shield} is
\[
  \mathrm{SF}(g)\iff k=0\ \wedge\ \delta=0\ \wedge\ (\forall j<0:\ d_j=0)\ \wedge\ (\exists j\ge0:\ d_j\ne0).
\]

\paragraph{Series.}
Put $q=x^2$ and
\[
  \mathfrak g=\sum_{L\ge1}q^Ly^{L-1}.
\]
The bulk vector $P=(P_b)_{b\ge1}$ is the solution of
\begin{equation}\label{eq:app-bulk}
  P_b=2q^b\Bigl(q^b+\sum_{a\ge1}K_{\mathfrak g}(a,b)P_a\Bigr),\qquad
  K_{\mathfrak g}(a,b)=q^{\max(a,b)}+q^{a+b}\mathfrak g ,
\end{equation}
which exists and is unique by Lemma~\ref{M3:lem:conv}. From it we form
\[
  S=\sum_{b\ge1}q^bP_b,\qquad B_0=1+\mathfrak gS,\qquad B_{0z}=1+y\,\mathfrak gS,
\]
and, for $s\ge0$,
\begin{align}
  \mu^{(y)}_s&=B_0\,x^{2s+1}+\sum_{b\ge1}\frac{P_b}{2}\bigl(x^{\max(2b-1,2s+1)}+x^{\max(2b+1,2s+1)}\bigr),
  \label{eq:app-mu}\\
  E^{(y)}_s&=B_0^{[s]}x^{2s}+B_0\,x^{2s+2}+\sum_{b\ge1}P_b\bigl(x^{\max(2s,2b)}+x^{\max(2s+2,2b)}\bigr),
  \qquad B_0^{[0]}=B_{0z},\ B_0^{[s]}=B_0\ (s\ge1),\label{eq:app-E}\\
  \lambda^{(y)}_s&=E^{(y)}_s+2\mu^{(y)}_s .\label{eq:app-lambda}
\end{align}
Let $D_e=\operatorname{diag}(e_s)_{s\ge0}$ with $e_s=2q^{1+s}$, let $K(s,t)=q^{\max(s,t)}$,
and let $\mathcal T=D_eK$. Define
\[
  \mathcal W_+=x^{-1}\bigl\langle\mu^{(y)},(I-\mathcal T)^{-1}D_e\lambda^{(y)}\bigr\rangle,\qquad
  \mathcal W_-=x^{-1}\bigl\langle\tfrac12E^{(y)},(I-\mathcal T)^{-1}D_e\lambda^{(y)}\bigr\rangle,
\]
and let $\mathcal W_0$ be the branch sum \eqref{eq:app-W0} of \S\ref{app:k0}, a double sum over the
left and right chain kinds (summable, not finite).

\begin{theorem}[\textup{(M3$'$)}]\label{thm:M3prime}
The following identity holds in $\mathbb Z[y][[x]]$:
\[
  \sum_{[g]\in\mathrm{Elt}/\mathrm{SameElt}}x^{\ell_R^{\mathrm{true}}(g)}y^{c^{\mathrm{true}}(g)}
  \;=\;\sum_{g\in\mathcal G}x^{\ell_R^{\mathrm{true}}(g)}y^{c^{\mathrm{true}}(g)}
  \;=\;\mathcal W_++\mathcal W_-+\mathcal W_0 .
\]
The three terms are the contributions of the sectors $k>0$, $k<0$ and $k=0$. Moreover
\[
  \mathcal W_++\mathcal W_-=\frac1{2x}\bigl\langle\lambda^{(y)},(I-\mathcal T)^{-1}D_e\lambda^{(y)}\bigr\rangle ,
\]
and each of $\mathcal W_+$, $\mathcal W_-$, $\mathcal W_0$ lies in $\mathbb Z[y][[x]]$.
\end{theorem}

\begin{corollary}\label{cor:M3prime-growth}
By \texttt{metricAll} (Lean: $|g|=\ell_R^{\mathrm{true}}(g)+2c^{\mathrm{true}}(g)$ for every
\texttt{Elt}, all of which are reachable), the specialisation $y=x^2$ of
Theorem~\ref{thm:M3prime} is the growth series $\sum_{[g]}x^{|g|}$ of the group modelled by
\texttt{Elt}. This group is $W_{\mathrm{univ}}$. For a particular triangle group $W_T$, the
growth series agrees with this one only up to the depth at which $\rho_T$ stops being
faithful on $W_{\mathrm{univ}}$
\cite[Theorem ``Effective universality'', \textup{thm:effective-universality}]{Bonfioli2026a}. For
the $(1,2)$ triangle the first deviation is at depth~$33$
\cite[\textup{thm:universality-sharp}]{Bonfioli2026a}.
\end{corollary}

The first equality of Theorem~\ref{thm:M3prime} is Lemma~\ref{M3:lem:classes}. The rest of this
appendix proves the second.

\subsection{Convergence and integrality}\label{app:conv}

We use the $x$-adic topology on $\mathbb Z[y][[x]]$. A family of series is \emph{summable}
if only finitely many of its members have a nonzero coefficient of any given $x^n$. We write
$\operatorname{ord}$ for the $x$-adic order.

\begin{lemma}\label{M3:lem:conv}
\begin{enumerate}
\item[(a)] For each $n$, only finitely many $g\in\mathcal G$ have $\ell_R^{\mathrm{true}}(g)=n$,
  and for these $c^{\mathrm{true}}(g)\le n+1$. So the left side of Theorem~\ref{thm:M3prime} is a
  well-defined element of $\mathbb Z_{\ge0}[y][[x]]$.
\item[(b)] Equation \eqref{eq:app-bulk} has a unique solution with $P_b\in\mathbb Z[y][[x]]$
  and $\operatorname{ord}P_b\ge4b$. Every $P_b$ is divisible by $2$.
\item[(c)] $(I-\mathcal T)^{-1}=\sum_{n\ge0}\mathcal T^n$ converges entrywise, and every pairing
  in the definitions of $\mathcal W_\pm$ is a summable sum.
\item[(d)] $\mu^{(y)}$, $E^{(y)}$, $\lambda^{(y)}$, $\mathcal W_+$, $\mathcal W_-$ and
  $\mathcal W_0$ all have coefficients in $\mathbb Z[y]$.
\end{enumerate}
\end{lemma}

\begin{proof}
(a) Every edge of $[A,B]$ has $\mu_j\ge1$, so $B-A+1\le n$. Moreover $|d_j|\le\mu_j\le n$
and $|k|\le B-A+1$. This leaves finitely many choices of $(k,\varepsilon,\delta,d)$. Also
$c^{\mathrm{true}}\le\#\{s:A<s<B+1\}+1=(B-A+1)+1\le n+1$.

(b) Let $\Psi(P)_b=2q^b(q^b+\sum_aK_{\mathfrak g}(a,b)P_a)$. The entry $K_{\mathfrak g}(a,b)$
has order $\ge2\max(a,b)\ge 2a$, so $\sum_a$ is summable whenever $\operatorname{ord}P_a\ge0$.
If two inputs satisfy $\operatorname{ord}(P_a-P'_a)\ge m$ for all $a$, then
$\operatorname{ord}(\Psi(P)_b-\Psi(P')_b)\ge 2b+2+m\ge m+4$. So $\Psi$ is a contraction.
Starting from $P=0$, the iterates therefore converge to the unique fixed point. Every
iterate has $\operatorname{ord}P_b\ge4b$ and carries the factor $2$. The coefficients lie
in $\mathbb Z[y]$, since the coefficient of $x^n$ in $\mathfrak g$ is $0$ or $y^{n/2-1}$.

(c) $\mathcal T(s,t)=2x^{2s+2+2\max(s,t)}$ has order $\ge 2s+2t+2$. So $(\mathcal T^n)(s,t)$
has order $\ge 2n+2t$, the Neumann series converges, and for fixed $s$ the sums over $t$ are
summable. By (b), $\operatorname{ord}\mu^{(y)}_s,\operatorname{ord}E^{(y)}_s\ge 2s$, so the
outer pairings are summable as well.

(d) $\mu^{(y)}_s$ is integral because $P_b/2$ is integral, by (b). For $\mathcal W_\pm$,
factor $D_e=2D'$ with $D'=\operatorname{diag}(q^{1+s})$. The push-through identity gives
\[
  (I-D_eK)^{-1}D_e=D_e(I-KD_e)^{-1}=2D'(I-KD_e)^{-1}.
\]
Hence
\[
  \mathcal W_-=x^{-1}\bigl\langle\tfrac12E^{(y)},\,2D'(I-KD_e)^{-1}\lambda^{(y)}\bigr\rangle
  =\bigl\langle E^{(y)},\,(x^{-1}D')(I-KD_e)^{-1}\lambda^{(y)}\bigr\rangle ,
\]
and $x^{-1}D'=\operatorname{diag}(x^{2s+1})$ is integral. \textbf{So $\mathcal W_-$ is
integral because $D_e$ carries the factor $2$.} The $\tfrac12$ on $E^{(y)}$ is cancelled by
that factor, and $E^{(y)}$ itself is not divisible by $2$: its $s=0$ entry has constant
term $1$. The same factorisation shows $\mathcal W_+\in\mathbb Z[y][[x]]$, with the factor
$2$ to spare. $\mathcal W_0$ is a summable double sum of products $x^jw_{\mathrm L}w_{\mathrm R}$
with weights from $1$, $\mathfrak gS$, $y\mathfrak gS$ and $P_b/2$, as \eqref{eq:app-W0} shows.
\end{proof}

All rearrangements below are of families of monomials with nonnegative integer coefficients,
and each family is summable by Lemma~\ref{M3:lem:conv}(a). So they are legitimate in
$\mathbb Z[y][[x]]$.

\subsection{Decomposition of an element: chains and the clamp analysis}\label{app:decomp}

\paragraph{Chains.}
A \emph{chain} is a finitely supported $c:\mathbb Z_{\ge1}\to2\mathbb Z$. Its \emph{length}
is $N(c)=\max\{i:c_i\ne0\}$, with $N(0)=0$. There are three kinds of chain:
\begin{itemize}
\item \emph{empty} (kind $\mathsf E$): $N=0$;
\item \emph{gap-opening} (kind $\mathsf Z$): $N\ge1$ and $c_1=0$;
\item \emph{deposit-opening} (kind $\mathsf N_{\sigma,b}$): $c_1=2\sigma b$ with $\sigma=\pm1$ and $b\ge1$.
\end{itemize}

Given $g=(k,\varepsilon,\delta,d)$, put $k_-=\min(0,k)$ and $k_+=\max(0,k)$. The
\emph{left chain} and the \emph{right chain} of $g$ are
\[
  c^{\mathrm L}_i=d_{k_--i},\qquad c^{\mathrm R}_i=d_{k_+-1+i}\qquad(i\ge1).
\]
Both consist of even deposits, since $f=0$ outside $[k_-,k_+)$. The \emph{travel deposits}
are $d_j$ for $k_-\le j<k_+$. There are $|k|$ of them, all odd. Write $\ell=N(c^{\mathrm L})$
and $r=N(c^{\mathrm R})$.

\begin{lemma}[bijection]\label{M3:lem:bij}
The map
\[
  g\mapsto\bigl(k,\varepsilon,\delta,(d_j)_{k_-\le j<k_+},c^{\mathrm L},c^{\mathrm R}\bigr)
\]
is a bijection from $\mathcal G$ onto the set of tuples consisting of $k\in\mathbb Z$,
$\varepsilon$, $\delta$, a word of $|k|$ odd integers, and two chains.
\end{lemma}

\begin{proof}
The three index sets $\{j<k_-\}$, $[k_-,k_+)$ and $\{j\ge k_+\}$ partition $\mathbb Z$, and
the parity condition says exactly that the deposits are odd on the middle set and even off
it.
\end{proof}

\begin{lemma}[span; clamp analysis]\label{M3:lem:span}
For every $g\in\mathcal G$,
\[
  A=k_--\ell,\qquad B=k_+-1+r .
\]
Hence:
\begin{enumerate}
\item[(i)] The edges of $[A,B]$ are the $\ell$ left-chain edges, the $|k|$ travel edges and
  the $r$ right-chain edges.
\item[(ii)] The site range $[A,B+1]$ consists of the following sites:
  \begin{itemize}
  \item the \emph{junction sites} $k_-$ and $k_+$, which coincide when $k=0$;
  \item the $|k|-1$ \emph{travel-interior} sites $k_-<s<k_+$;
  \item the \emph{chain-interior} sites: $A<s<k_-$ on the left and $k_+<s<B+1$ on the right;
  \item the \emph{outer sites} $A$ and $B+1$, but only when $\ell>0$, resp.\ $r>0$.
  \end{itemize}
\item[(iii)] The junction site $k_-$ lies in the open interval $(A,B+1)$ if and only if
  $\ell>0$ or $k<0$, i.e.\ $k_-<k_+$. Symmetrically, $k_+$ lies in it if and only if $r>0$
  or $k>0$. The exception is $k=0$: there site $0$ lies in the open interval if and only if
  $\ell>0$ \emph{and} $r>0$.
\end{enumerate}
\end{lemma}

\begin{proof}
$\mathrm{occ}$ consists of the travel range $[k_-,k_+)$, the edges $k_--i$ with
$c^{\mathrm L}_i\ne0$, and the edges $k_+-1+i$ with $c^{\mathrm R}_i\ne0$.

\emph{Case $k>0$.} Here $0\in\mathrm{occ}$. So $\min\mathrm{occ}=-\ell\le0$ (it is $0$ when
$\ell=0$), and the clamp $\min(0,\cdot)$ is inactive. Likewise
$\max\mathrm{occ}=k-1+r\ge0>-1$, and the clamp $\max(-1,\cdot)$ is inactive.

\emph{Case $k<0$.} The argument is symmetric, using $-1\in\mathrm{occ}$.

\emph{Case $k=0$.} Now $\mathrm{occ}=\{-i:c^{\mathrm L}_i\ne0\}\cup\{i-1:c^{\mathrm R}_i\ne0\}$
and the clamps are active:
\begin{itemize}
\item If $\ell=0$, then $\mathrm{occ}\subseteq[0,\infty)$ or $\mathrm{occ}=\varnothing$, and
  either way $A=0=-\ell$.
\item If $r=0$, then $B=-1=r-1$.
\item If $\ell>0$, resp.\ $r>0$, the unclamped values $-\ell$, resp.\ $r-1$, already lie on
  the correct side of the clamp.
\end{itemize}
So the formula for $A$ and $B$ holds in all cases.

Items (i) and (ii) now follow by listing $[A,B]$ and $[A,B+1]$. When $\ell=0$, the outer site
$A$ coincides with the junction site $k_-$, and similarly on the right.

For (iii), $A<k_-$ if and only if $\ell>0$, and $k_-<B+1=k_++r$ if and only if $k_-<k_+$ or
$r>0$. The symmetric statement for $k_+$ is proved the same way.
\end{proof}

So the clamps in \texttt{ATrue} and \texttt{BTrue} change nothing when $k\ne0$. When $k=0$
they do exactly two things. First, they keep site $0$ in the site range even when both
chains are empty. Second, they put site $0$ \emph{outside} the open interval unless both
chains are nonempty. That second effect is what $\mathrm{SF}$ partly undoes
(\S\ref{app:k0}).

\begin{lemma}[locality]\label{M3:lem:local}
Let $s$ be a site of $[A,B+1]$ other than a junction site. Then
$\alpha_s=d_{s-1}$, $\beta_s=d_s$ and $\operatorname{site}(s)=\max(|d_{s-1}|,|d_s|)$. Moreover:
\begin{enumerate}
\item[(a)] A travel-interior site has odd cost $\max(|d_{s-1}|,|d_s|)$, and it is never a cut,
  because $\Phi_s=f_{s-1}=\pm1$.
\item[(b)] A chain-interior site is a cut if and only if $d_{s-1}=d_s=0$. Every
  chain-interior site lies in $(A,B+1)$.
\item[(c)] An outer site has cost $|c_N|$, where $c$ is the corresponding chain and $N$ its
  length. It is not a cut, and it does not lie in $(A,B+1)$.
\end{enumerate}
In particular $\ell_R^{\mathrm{true}}$ splits as a sum of the following parts:
\begin{itemize}
\item left-chain edges and left chain-interior and outer sites;
\item travel edges and travel-interior sites;
\item right-chain edges and right chain-interior and outer sites;
\item the junction sites.
\end{itemize}
The cut count $c^{\mathrm{true}}$ splits the same way, except that the junction sites are
governed by Lemma~\ref{M3:lem:span}(iii) and by $\mathrm{SF}$.
\end{lemma}

\begin{proof}
For $s\notin\{0,k\}$ all the brackets vanish, so $\alpha_s=d_{s-1}$, $\beta_s=d_s$ and
$\Phi_s=f_{s-1}$. For the cut criterion see also \texttt{RJLemmaJ.bulk\_cut}.

For (a), if $k_-<s<k_+$ then $s-1\in[k_-,k_+)$, so $d_{s-1}$ and $d_s$ are odd and
$f_{s-1}=\pm1$.

For (b), if $s$ is chain-interior then $s-1$ and $s$ are chain edges with $f=0$. Such $s$
lie strictly between $A$ and $B+1$ by Lemma~\ref{M3:lem:span}.

For (c), at the outer site $B+1$ (when $r>0$) we have $\alpha=d_B=c^{\mathrm R}_r\ne0$ and
$\beta=0$. The site $A$ is symmetric.

The cost and cut conditions of every non-junction site involve only the deposits of one
block, which gives the splitting.
\end{proof}

\subsection{Chain generating functions}\label{app:chains}

For a chain $c$ of length $N\ge1$, define its \emph{weight} $\omega(c)$ as $x^{\ell(c)}y^{\gamma(c)}$,
where
\begin{align*}
  \ell(c)&=\underbrace{\sum_{i=1}^{N}\bigl(|c_i|+2[c_i=0]\bigr)}_{\text{edges}}
          +\underbrace{\sum_{i=2}^{N}\max(|c_{i-1}|,|c_i|)}_{\text{interior sites}}
          +\underbrace{|c_N|}_{\text{outer site}},\\
  \gamma(c)&=\#\{2\le i\le N:\ c_{i-1}=c_i=0\}.
\end{align*}
The empty chain has weight $1$. By Lemmas~\ref{M3:lem:span} and~\ref{M3:lem:local}, the left and
right chains of $g$ contribute exactly $\omega(c^{\mathrm L})$ and $\omega(c^{\mathrm R})$ to
$x^{\ell_R^{\mathrm{true}}}y^{c^{\mathrm{true}}}$. Everything else comes from the junction
sites and the travel block. Note that the junction site is \emph{not} part of $\omega$.

\begin{lemma}[bulk]\label{M3:lem:bulk}
For $b\ge1$,
\[
  \sum_{c:\ |c_1|=2b}\omega(c)=P_b,\qquad
  \sum_{c:\ c_1=2\sigma b}\omega(c)=\tfrac12P_b\quad(\sigma=\pm1),\qquad
  \sum_{c\text{ of kind }\mathsf Z}\omega(c)=\mathfrak gS .
\]
Consequently
\[
  \sum_{c\text{ of kind }\mathsf E\text{ or }\mathsf Z}\omega(c)=B_0 .
\]
\end{lemma}

\begin{proof}
Let $Q_b=\sum_{|c_1|=2b}\omega(c)$. We split according to what follows edge $1$.
\begin{itemize}
\item \emph{$N=1$.} The weight is $x^{2b}$ for the edge times $x^{2b}$ for the outer site,
  i.e.\ $q^{2b}$, for each of the two signs.
\item \emph{Next deposit $c_2=\pm2a\ne0$.} Site $2$ costs $x^{2\max(a,b)}$.
\item \emph{A gap run $c_2=\dots=c_{L+1}=0$ with $L\ge1$, followed by $c_{L+2}=\pm2a$.} Site
  $2$ costs $x^{2b}$. The $L$ gap edges give $q^L$. The $L-1$ sites strictly inside the gap
  cost $1$ and are cuts, giving $y^{L-1}$. The site before $c_{L+2}$ costs $x^{2a}$.
\end{itemize}
The tail $(c_{i+1})_{i\ge1}$ in the second case, and $(c_{L+1+i})_{i\ge1}$ in the third, is
again a chain with first entry of magnitude $2a$. Its interior sites, outer site and cuts
are exactly those counted in its own $\omega$. So
\[
  Q_b=2q^b\Bigl(q^b+\sum_{a}\bigl(q^{\max(a,b)}+q^{b}\,\mathfrak g\,q^{a}\bigr)Q_a\Bigr),
\]
which is \eqref{eq:app-bulk}. We have $\operatorname{ord}Q_b\ge4b$, so $Q=P$ by the
uniqueness in Lemma~\ref{M3:lem:conv}(b).

The negation $c\mapsto-c$ preserves $\omega$, which gives the signed statement.

A chain of kind $\mathsf Z$ is a gap run of length $L\ge1$, followed by a chain with
$|c_1|=2a$. Its weight is $q^Ly^{L-1}\cdot x^{2a}\cdot\omega(\text{tail})$. Here the sites
strictly inside the gap are the $L-1$ cuts, and $x^{2a}$ is the site between the gap and the
deposit. Summing gives $\mathfrak g\sum_aq^aP_a=\mathfrak gS$.
\end{proof}

\subsection{Travel block}\label{app:travel}

\begin{lemma}[travel]\label{M3:lem:travel}
Let $k\ne0$. Fix everything except the signs of the travel deposits, and write
$|d_j|=2s_j+1$. The travel edges and travel-interior sites contribute
\[
  \prod_{j}x^{2s_j+1}\prod_{\text{interior}}x^{2\max(s_{j},s_{j+1})+1}
  =x^{-1}\prod_j\bigl(x^{2s_j+2}\bigr)\prod_{\text{interior}}x^{2\max(s_j,s_{j+1})}.
\]
This does not depend on the signs.
\end{lemma}

\begin{proof}
The run has $|k|$ edges and $|k|-1$ interior sites. Moving one $x$ from each interior site to
the edge that follows it leaves one extra $x$, which is where the factor $x^{-1}$ comes from.
Sign independence follows from Lemma~\ref{M3:lem:local}(a).
\end{proof}

Summing over the travel words with the two end magnitudes fixed, and counting the $2^{|k|}$
sign choices as the factor $2$ in each $e_s$, gives
\begin{equation}\label{eq:app-transfer}
  \sum_{|k|\ge1}\ \sum_{\text{words}}(\cdots)=x^{-1}\sum_{n\ge0}\bigl\langle v,(D_eK)^nD_ew\bigr\rangle
  =x^{-1}\bigl\langle v,(I-\mathcal T)^{-1}D_ew\bigr\rangle .
\end{equation}
This formula needs a condition on the ends: the end weights $v_{s_{\mathrm{near}}}$ and
$w_{s_{\mathrm{far}}}$ must be the \emph{sign-averaged} end weights. That is, the ends
together must contribute the factor $\tfrac12\cdot2$ per end sign, which the next lemma makes
precise. Here $n=|k|-1$, and $K$ is symmetric, so the reading direction does not matter.

\begin{lemma}[sign decoupling]\label{M3:lem:sign}
Let $k\ne0$. Let $\sigma_{\mathrm n}$ and $\sigma_{\mathrm f}$ be the signs of the travel
deposits adjacent to the near junction (site $0$) and the far junction (site $k$). Write
the contribution of the two junction sites, together with the adjacent chains, as
$N(\sigma_{\mathrm n})\cdot F(\sigma_{\mathrm f},\varepsilon,\delta)$. Here $N$ is summed
over the chain at site $0$ and $F$ over the chain at site $k$. Suppose that
\begin{equation}\label{eq:app-Finv}
  \Lambda:=\sum_{\varepsilon,\delta}F(\sigma,\varepsilon,\delta)\quad\text{does not depend on }\sigma .
\end{equation}
Then the sum of $N\cdot F$ over $\varepsilon$, $\delta$ and the end signs equals
$2\cdot\overline N\cdot\Lambda$ when $|k|=1$ (so $\sigma_{\mathrm n}=\sigma_{\mathrm f}$), and
$4\cdot\overline N\cdot\Lambda$ when $|k|\ge2$. Here
$\overline N=\tfrac12\sum_\sigma N(\sigma)$. In both cases this is what
\eqref{eq:app-transfer} produces with $v=\overline N$ and $w=\Lambda$.
\end{lemma}

\begin{proof}
If $|k|\ge2$ the two signs are independent, and
$\sum_{\sigma_{\mathrm n},\sigma_{\mathrm f},\varepsilon,\delta}NF
=(\sum_\sigma N)(\sum_{\sigma_{\mathrm f}}\Lambda)=2\overline N\cdot2\Lambda$.

If $|k|=1$ there is one travel edge and one sign $\sigma$, so
$\sum_{\sigma,\varepsilon,\delta}N(\sigma)F(\sigma,\varepsilon,\delta)=\sum_\sigma N(\sigma)\Lambda
=2\overline N\Lambda$, using \eqref{eq:app-Finv}.

In both cases these are the edge factors $2$ carried by $D_e$ at the end edges.
\end{proof}

\begin{remark}[decoupling holds only after the sum over $\varepsilon$]\label{M3:rem:eps}
Hypothesis \eqref{eq:app-Finv} is essential when $|k|=1$, and it fails before the sum over
$\varepsilon$.

Take $k=-1$, so that $d_{-1}=\sigma(2s+1)$ is both the near and the far travel deposit. Then
$N(\sigma)$ involves $|d_{-1}-1|$, which is $2s$ or $2s+2$ according to $\sigma$. For fixed
$\varepsilon$ and $\delta=1$, the far site involves $|d_{-1}-\varepsilon|$, which is also
$2s$ or $2s+2$ according to $\sigma\varepsilon$. So for fixed $\varepsilon$ the sum
$\sum_\sigma N(\sigma)F(\sigma,\varepsilon,1)$ pairs the cheap near cost with the cheap or
the expensive far cost depending on $\varepsilon$. It is \emph{not} equal to
$\overline N\sum_\sigma F(\sigma,\varepsilon,1)$.

Only after summing over $\varepsilon$ does the far weight lose its $\sigma$-dependence,
because the substitution $\varepsilon\mapsto-\varepsilon$ exchanges the two values. Then the
product factorises.

For $k=+1$, the near weight depends only on $|d_0|$ (see \eqref{eq:app-k+near}), so there
the factorisation is harmless. The model nevertheless uses the $\varepsilon$-summed far
vector $\lambda^{(y)}$ in both sectors, and that is legitimate only because of
\eqref{eq:app-Finv}.
\end{remark}

Condition \eqref{eq:app-Finv} is verified in \eqref{eq:app-k+far} and \eqref{eq:app-k-far}
below. In both, the $\sigma$-independence comes from the involution
$\varepsilon\mapsto-\varepsilon$, combined with negating the far chain where necessary.
Negating the far chain preserves $\omega$ and all cut conditions.

\subsection{Sector $k>0$}\label{app:kpos}

The near junction is site $0$. There $\alpha_0=d_{-1}-1$, $\beta_0=d_0$ and
$\Phi_0=f_{-1}+1=1$. So \textbf{site $0$ is never a cut} (\texttt{RJLemmaJ.J\_site0\_pos}).
Its chain is $c^{\mathrm L}$, with $u:=c^{\mathrm L}_1=d_{-1}$. Let $d_0=\sigma(2s+1)$. The
cost is $\max(|u-1|,2s+1)$, which does not depend on $\sigma$. Summing over $c^{\mathrm L}$
with Lemma~\ref{M3:lem:bulk}:
\begin{itemize}
\item $u=0$ (kinds $\mathsf E$, $\mathsf Z$): the cost is $2s+1$ and the weight is $B_0$;
\item $u=+2b$: the cost is $\max(2b-1,2s+1)$ and the weight is $\tfrac12P_b$;
\item $u=-2b$: the cost is $\max(2b+1,2s+1)$ and the weight is $\tfrac12P_b$.
\end{itemize}
Hence
\begin{equation}\label{eq:app-k+near}
  N(\sigma)=\overline N=\mu^{(y)}_s .
\end{equation}
No $y$ occurs here: site $0$ is never a cut, and in any case it lies in $(A,B+1)$ only when
$\ell>0$.

The far junction is site $k$. Put $d_{k-1}=\sigma(2s+1)$ and $v:=c^{\mathrm R}_1=d_k$. Then
\[
  \alpha_k=\sigma(2s+1)+\varepsilon[\delta=0],\qquad \beta_k=v-\varepsilon[\delta=1],\qquad
  \Phi_k=f_{k-1}-[\delta=0]=[\delta=1].
\]
By Lemma~\ref{M3:lem:span}(iii), site $k$ lies in $(A,B+1)$ if and only if $r>0$.

\emph{$\delta=1$.} Here $\Phi_k=1$, so there is no cut. The cost is
$\max(2s+1,|v-\varepsilon|)$.
\begin{itemize}
\item $v=0$: the cost is $2s+1$ for both $\varepsilon$, contributing $2B_0x^{2s+1}$.
\item $v=\pm2b$: as $\varepsilon$ runs over $\pm1$, $|v-\varepsilon|$ runs over
  $\{2b-1,2b+1\}$. Each sign of $v$ has weight $\tfrac12P_b$, so the contribution is
  $P_b(x^{\max(2b-1,2s+1)}+x^{\max(2b+1,2s+1)})$.
\end{itemize}
The total is $2\mu^{(y)}_s$.

\emph{$\delta=0$.} Here $\Phi_k=0$ and $\beta_k=v$. As $\varepsilon$ runs over $\pm1$,
$|\alpha_k|$ takes the value $2s+2$ (for $\varepsilon=\sigma$) and $2s$ (for
$\varepsilon=-\sigma$).
\begin{itemize}
\item $\varepsilon=\sigma$: the cost is $\max(2s+2,|v|)$ and there is no cut. The
  contribution is $B_0x^{2s+2}+\sum_bP_bx^{\max(2s+2,2b)}$.
\item $\varepsilon=-\sigma$: the cost is $\max(2s,|v|)$. Site $k$ is a cut if and only if
  $s=0$ and $v=0$ (\texttt{RJLemmaJ.J\_sitek\_pos}: $\delta=0$, $d_{k-1}=-\varepsilon$,
  $d_k=0$). The cut is \emph{counted} if and only if moreover $r>0$, i.e.\ the right chain
  has kind $\mathsf Z$; for kind $\mathsf E$, site $k=B+1$ is the outer end. So for $v=0$
  the weight is $1+y^{[s=0]}\mathfrak gS=B_0^{[s]}$, and the contribution is
  $B_0^{[s]}x^{2s}+\sum_bP_bx^{\max(2s,2b)}$.
\end{itemize}
The total is $E^{(y)}_s$.

Therefore
\begin{equation}\label{eq:app-k+far}
  \sum_{\varepsilon,\delta}F(\sigma,\varepsilon,\delta)=E^{(y)}_s+2\mu^{(y)}_s=\lambda^{(y)}_s
  \quad\text{for both }\sigma,
\end{equation}
which verifies \eqref{eq:app-Finv}.

Combining Lemmas~\ref{M3:lem:local}, \ref{M3:lem:bulk}, \ref{M3:lem:travel} and~\ref{M3:lem:sign}
with \eqref{eq:app-transfer}:
\[
  \sum_{g:\,k>0}x^{\ell_R^{\mathrm{true}}}y^{c^{\mathrm{true}}}
  =x^{-1}\bigl\langle\mu^{(y)},(I-\mathcal T)^{-1}D_e\lambda^{(y)}\bigr\rangle=\mathcal W_+ .
\]
There is no shield term, because $\mathrm{SF}$ requires $k=0$.

\subsection{Sector $k<0$, and the factor $\tfrac12$ on $E$}\label{app:kneg}

The near junction is site $0$. Its chain is now $c^{\mathrm R}$, with $v=c^{\mathrm R}_1=d_0$,
and the adjacent travel deposit is $d_{-1}=\sigma(2s+1)$. Then
\[
  \alpha_0=\sigma(2s+1)-1,\qquad\beta_0=v,\qquad\Phi_0=f_{-1}+1=0 .
\]
So $|\alpha_0|=2s$ for $\sigma=+1$, and $|\alpha_0|=2s+2$ for $\sigma=-1$. Site $0$ is a cut
if and only if $d_{-1}=1$ and $d_0=0$ (\texttt{RJLemmaJ.J\_site0\_neg}). It lies in
$(A,B+1)$ if and only if $r>0$, by Lemma~\ref{M3:lem:span}(iii). So:
\begin{itemize}
\item $\sigma=+1$: the contribution is $B_0^{[s]}x^{2s}+\sum_bP_bx^{\max(2s,2b)}$. The $y$
  appears exactly for $s=0$ with $c^{\mathrm R}$ of kind $\mathsf Z$.
\item $\sigma=-1$: the contribution is $B_0x^{2s+2}+\sum_bP_bx^{\max(2s+2,2b)}$, with no cut.
\end{itemize}
Thus $N(+1)+N(-1)=E^{(y)}_s$, and
\begin{equation}\label{eq:app-k-near}
  \overline N=\tfrac12E^{(y)}_s .
\end{equation}

\paragraph{The $E/2$ argument.}
Here, unlike the case $k>0$, the near weight depends on $\sigma$. The sum over $\sigma$ of
the near end produces the whole of $E^{(y)}_s$. But the edge carrying $\sigma$ is a travel
edge, and its two signs are already counted by the factor $2$ in $e_s$. Pairing
$E^{(y)}$ with $D_e$ would therefore count that sign twice. The correct near vector is the
sign average, $\tfrac12E^{(y)}$. This is exactly what Lemma~\ref{M3:lem:sign} says. Its
integrality was recorded in Lemma~\ref{M3:lem:conv}(d).

The far junction is site $k$. Its chain is $c^{\mathrm L}$, with $u=c^{\mathrm L}_1=d_{k-1}$,
and $d_k=\sigma(2s+1)$. Then
\[
  \alpha_k=u+\varepsilon[\delta=0],\qquad\beta_k=\sigma(2s+1)-\varepsilon[\delta=1],\qquad
  \Phi_k=f_{k-1}-[\delta=0]=-[\delta=0].
\]
By Lemma~\ref{M3:lem:span}(iii), site $k$ lies in $(A,B+1)$ if and only if $\ell>0$.

\emph{$\delta=0$.} Here $\Phi_k=-1$, so there is no cut. The cost is
$\max(|u+\varepsilon|,2s+1)$. By the same computation as for $\delta=1$ in \S\ref{app:kpos},
with the chains mirrored, the sum over $\varepsilon$ and $c^{\mathrm L}$ is $2\mu^{(y)}_s$.

\emph{$\delta=1$.} Here $\Phi_k=0$ and $\alpha_k=u$, while $|\beta_k|$ is $2s$ for
$\varepsilon=\sigma$ and $2s+2$ for $\varepsilon=-\sigma$. Site $k$ is a cut if and only if
$u=0$, $s=0$ and $\varepsilon=\sigma$ (\texttt{RJLemmaJ.J\_sitek\_neg}). The cut is counted
if and only if $\ell>0$, i.e.\ the left chain has kind $\mathsf Z$. The sum over
$\varepsilon$ and $c^{\mathrm L}$ is $E^{(y)}_s$, with the slot $B_{0z}$ at $s=0$.

Hence
\begin{equation}\label{eq:app-k-far}
  \sum_{\varepsilon,\delta}F(\sigma,\varepsilon,\delta)=\lambda^{(y)}_s\quad\text{for both }\sigma .
\end{equation}
By Lemma~\ref{M3:lem:sign}, which is now needed in earnest at $k=-1$ (Remark~\ref{M3:rem:eps}),
\[
  \sum_{g:\,k<0}x^{\ell_R^{\mathrm{true}}}y^{c^{\mathrm{true}}}
  =x^{-1}\bigl\langle\tfrac12E^{(y)},(I-\mathcal T)^{-1}D_e\lambda^{(y)}\bigr\rangle=\mathcal W_- .
\]
Finally, $\mu^{(y)}+\tfrac12E^{(y)}=\tfrac12\lambda^{(y)}$ by \eqref{eq:app-lambda}, which gives
$\mathcal W_++\mathcal W_-=(2x)^{-1}\langle\lambda^{(y)},(I-\mathcal T)^{-1}D_e\lambda^{(y)}\rangle$.

\subsection{Sector $k=0$: the complete case analysis}\label{app:k0}

When $k=0$ there is no travel, every $d_j$ is even, and site $0$ is the only junction site.
Put $u=c^{\mathrm L}_1=d_{-1}$ and $v=c^{\mathrm R}_1=d_0$. Then
\[
  \alpha_0=u-1+\varepsilon[\delta=0],\qquad \beta_0=v-\varepsilon[\delta=1],\qquad
  \Phi_0=1-[\delta=0]=[\delta=1].
\]
The term $-1$ is \texttt{vArr}. By Lemmas~\ref{M3:lem:span}, \ref{M3:lem:local} and~\ref{M3:lem:bulk},
\begin{equation}\label{eq:app-W0}
  \mathcal W_0=\sum_{\varepsilon=\pm1}\sum_{\delta\in\{0,1\}}\ \sum_{(u,w_{\mathrm L})}\ \sum_{(v,w_{\mathrm R})}
  x^{\,j(\varepsilon,\delta,u,v)}\,w_{\mathrm L}\,w_{\mathrm R}\,y^{\chi},
  \qquad j=\max\bigl(|u-1+\varepsilon[\delta=0]|,\ |v-\varepsilon[\delta=1]|\bigr),
\end{equation}
where:
\begin{itemize}
\item $(u,w_{\mathrm L})$ ranges over the left kinds: $\mathsf E:(0,1)$,
  $\mathsf Z:(0,\mathfrak gS)$ and $\mathsf N_{\pm,b}:(\pm2b,\tfrac12P_b)$ for $b\ge1$;
\item $(v,w_{\mathrm R})$ ranges over the same list for the right chain;
\item $\chi\in\{0,1\}$ records whether site $0$ is counted in $c^{\mathrm{true}}$.
\end{itemize}
Only $w_{\mathrm R}$ of kind $\mathsf Z$ can carry the extra $y$: $y^\chi w_{\mathrm R}$ is
$y\mathfrak gS$ when $\chi=1$. The chain interiors are accounted for inside $w_{\mathrm L}$ and $w_{\mathrm R}$.

\paragraph{Site-$0$ cost, by branch.}
\begin{center}
\begin{tabular}{c|c|l}
$\delta$ & $\varepsilon$ & $j(\varepsilon,\delta,u,v)$\\\hline
$0$ & $+1$ & $\max(|u|,|v|)$\\
$0$ & $-1$ & $\max(|u-2|,|v|)$\\
$1$ & $+1$ & $\max(|u-1|,|v-1|)$\\
$1$ & $-1$ & $\max(|u-1|,|v+1|)$
\end{tabular}
\end{center}
For example, $u=v=0$ gives $0$, $2$, $1$ and $1$ in the four rows. These are the identity
and three elements of length $\le2$.

\paragraph{When is site $0$ a cut?}
$\operatorname{cut}(0)$ holds if and only if all three of the following hold:
\begin{itemize}
\item $\delta=0$, since $\Phi_0=[\delta=1]$;
\item $v=0$;
\item $u=1-\varepsilon$, that is, $u=0$ if $\varepsilon=+1$ and $u=2$ if $\varepsilon=-1$.
\end{itemize}

\paragraph{When is site $0$ counted?}
Site $0$ is counted if and only if $\operatorname{cut}(0)$ holds and one of the following
holds:
\begin{itemize}
\item[(I)] $0\in(A,B+1)$, which by Lemma~\ref{M3:lem:span}(iii) means $\ell>0$ and $r>0$;
\item[(S)] $\mathrm{SF}(g)$ holds. This means $\delta=0$, $\ell=0$ (the left chain has kind
  $\mathsf E$; recall $d_j=0$ for all $j<0$ says exactly that), and $r>0$.
\end{itemize}
Since $v=0$ and $r>0$ force the right chain to have kind $\mathsf Z$, we get $\chi=1$ only if
$\delta=0$ and $c^{\mathrm R}$ has kind $\mathsf Z$. The full table follows. Rows are the
left kind, and the $\mathsf N$ kinds are split by $u$.

\begin{center}
\small
\begin{tabular}{c|c|c|c|c|c|c|c}
$\delta$&$\varepsilon$&left kind ($u$)&right kind&$\operatorname{cut}(0)$&(I)&(S)&$\chi$\\\hline
$0$&$+1$&$\mathsf E$ ($u=0$)&$\mathsf Z$&yes&no ($\ell=0$)&\textbf{yes}&$\mathbf 1$\\
$0$&$+1$&$\mathsf Z$ ($u=0$)&$\mathsf Z$&yes&yes&no&$\mathbf 1$\\
$0$&$+1$&$\mathsf N$ ($u\ne0$)&$\mathsf Z$&no ($u\ne0$)&yes&no&$0$\\
$0$&$-1$&$\mathsf E$ ($u=0$)&$\mathsf Z$&no ($u\ne2$)&no&yes&$0$\\
$0$&$-1$&$\mathsf Z$ ($u=0$)&$\mathsf Z$&no ($u\ne2$)&yes&no&$0$\\
$0$&$-1$&$\mathsf N_{+,1}$ ($u=2$)&$\mathsf Z$&yes&yes&no&$\mathbf 1$\\
$0$&$-1$&$\mathsf N$, $u\ne2$&$\mathsf Z$&no&yes&no&$0$\\\hline
$0$&$\pm1$&$\mathsf E$&$\mathsf N$&no ($v\ne0$)&no&yes&$0$\\
$0$&$\pm1$&$\mathsf Z,\mathsf N$&$\mathsf N$&no ($v\ne0$)&yes&no&$0$\\
$0$&$\pm1$&$\mathsf E$&$\mathsf E$&$[\varepsilon=+1]$&no ($\ell=r=0$)&no ($r=0$)&$0$\\
$0$&$\pm1$&$\mathsf Z,\mathsf N$&$\mathsf E$&possible&no ($r=0$)&no ($\ell>0$, $r=0$)&$0$\\\hline
$1$&$\pm1$&any&any&no ($\Phi_0=1$)&--&no ($\delta=1$)&$0$
\end{tabular}
\end{center}

The table exhausts all branches. There are $2\times2$ marker choices, the left kinds
$\mathsf E,\mathsf Z,\mathsf N_{\pm,b}$, and the right kinds $\mathsf E,\mathsf Z,\mathsf N_{\pm,b}$.
Every marker choice with $\delta=0$ appears in the first eleven rows, split by right kind, and
$\delta=1$ is the last row. The rows marked (S)\,=\,yes with $\chi=0$ are the cases in which
the shield \emph{fires but site $0$ is not a cut}. They contribute nothing, which is why
\texttt{cTrue} conjoins $\mathrm{SF}$ with $\operatorname{cut}(0)$. The first row is the only
branch in which the shield changes the count. There site $0$ is a genuine cut, but it lies
at the clamped left end $A=0$ of the span and so falls outside the open interval.

\paragraph{Summary.}
The right $\mathsf Z$ weight is $y\mathfrak gS$ exactly in the following cases:
\begin{itemize}
\item $\delta=0$, $\varepsilon=+1$ and the left chain has kind $\mathsf E$ or $\mathsf Z$;
\item $\delta=0$, $\varepsilon=-1$ and $u=+2$.
\end{itemize}
In every other case it is $\mathfrak gS$. This is branch for branch the rule used by the
bivariate enumeration check of \S\ref{sec:bridge}. This proves
\[
  \sum_{g:\,k=0}x^{\ell_R^{\mathrm{true}}}y^{c^{\mathrm{true}}}=\mathcal W_0 .
\]
Integrality is clear from \eqref{eq:app-W0}: $P_b/2\in\mathbb Z[y][[x]]$, and the sum is
summable because $\operatorname{ord}P_b\ge4b$.

\begin{remark}[the regrouping of $\mathcal W_0$]\label{M3:rem:W0regroup}
Let $\Psi=(\Psi_b)_{b\ge1}$ be the ungapped bulk vector, the solution of \eqref{eq:app-bulk} with
$\mathfrak g=0$, which depends on $q$ alone. By the Sherman--Morrison formula
\eqref{eq:smresolvent}, $P_b=B_0\Psi_b$, and $\mathfrak gS=B_0-1$. So every weight in
\eqref{eq:app-W0} is affine in $B_0$: $w\in\{1,\ B_0-1,\ \tfrac12B_0\Psi_b\}$, and $y^\chi w_{\mathrm R}=y(B_0-1)$
when $\chi=1$. Each term $x^jw_{\mathrm L}w_{\mathrm R}y^\chi$ is therefore a polynomial of degree
at most $2$ in $B_0$, and collecting powers gives
\[
   \mathcal W_0=A_0+B_0A_1+B_0^2A_2,
\]
with $A_0,A_1,A_2$ absolutely convergent sums of monomials in $x$ and $y$ times at most two entries
of $\Psi$. The $A_i$ contain no travel resolvent. As functions of $q$ they
are singular only where $\Psi$ is, i.e.\ where $I-M_0$ is not invertible ($S_e=0$,
Corollary~\ref{cor:sing}), and $B_0=1/(1-\mathfrak gt_1)$ only where $1-\mathfrak gt_1=0$. Neither
occurs at a travel pole (Proposition~\ref{prop:selower}, Corollary~\ref{app:gate}), so
$\mathcal W_0$ has no travel pole.
\end{remark}

\subsection{Conclusion of the proof}

The sectors $k>0$, $k<0$ and $k=0$ partition $\mathcal G$. Lemma~\ref{M3:lem:bij} identifies each
sector with the product set of its blocks, and Lemma~\ref{M3:lem:local} shows that
$x^{\ell_R^{\mathrm{true}}}y^{c^{\mathrm{true}}}$ is the product of the block weights and the
junction-site factors. The three sector computations above then give
Theorem~\ref{thm:M3prime}. The rearrangements are justified by Lemma~\ref{M3:lem:conv}. \qed

\paragraph{What is and is not established here.}
The following are established:
\begin{itemize}
\item the identity of Theorem~\ref{thm:M3prime}, as a hand proof;
\item integrality;
\item $x$-adic convergence.
\end{itemize}
The following are not established here:
\begin{itemize}
\item any statement about a specific rational triangle beyond the universality depth
  (Corollary~\ref{cor:M3prime-growth});
\item a Lean formalisation of the assembly.
\end{itemize}
The inputs taken from Lean are the definitions of \texttt{siteCost}, \texttt{cut},
\texttt{ATrue}, \texttt{BTrue}, \texttt{lRTrue}, \texttt{cTrue} and \texttt{ShieldFires}, the
junction lemmas of \texttt{RJLemmaJ.lean}, and \texttt{metricAll} (for the corollary only).

\begin{thebibliography}{99}
\bibitem{AMV2007} M.~Amou, T.~Matala-aho, K.~V\"a\"an\"anen, \emph{On Siegel--Shidlovskii's theory
for $q$-difference equations}, Acta Arith.\ \textbf{127} (2007), 309--335.

\bibitem{Andre2000} Y.~Andr\'e, \emph{S\'eries Gevrey de type arithm\'etique I, II}, Ann.\ of Math.\
(2) \textbf{151} (2000), 705--740 and 741--756.

\bibitem{ADR2021} C.~E.~Arreche, T.~Dreyfus, J.~Roques, \emph{Differential transcendence criteria
for second-order linear difference equations and elliptic hypergeometric functions},
J.~\'Ec.\ polytech.\ Math.\ \textbf{8} (2021), 147--168.

\bibitem{BellMishna2020} J.~P.~Bell and M.~Mishna, \emph{On the complexity of the cogrowth
sequence}, J.~Comb.\ Algebra \textbf{4} (2020), 73--85.

\bibitem{Bertrandias1964} F.~Bertrandias, \emph{Diam\`etre transfini dans un corps valu\'e.
Application au prolongement analytique}, S\'eminaire Delange--Pisot--Poitou, Th\'eorie des
nombres \textbf{5} (1963--1964), expos\'e no.~3, 13~pp.

\bibitem{Bonfioli2026a} V.~Bonfioli, \emph{The universal right-triangle reflection sequence:
metric, universality, and the lamplighter structure}, preprint (2026),
\texttt{doi:10.5281/zenodo.20370090}.

\bibitem{Bonfioli2026c} V.~Bonfioli, \emph{Universality and a dimensional transcendence dichotomy
for orthoscheme reflection groups}, preprint (2026), \texttt{doi:10.5281/zenodo.20370090}.

\bibitem{supplement} V.~Bonfioli, \emph{Route-B lifting and transcendence supplement}, unpublished
working notes (2026), deposited with the project source release on Zenodo,
\texttt{doi:10.5281/zenodo.20370090} (concept DOI, resolving to the current version). The notes
record the block assembly and the lifting reduction and prove neither; the mathematics this paper
takes from them is reproduced and proved in full in \S\ref{sec:assembly} and
Appendix~\ref{app:M3prime}. No
theorem of this paper depends on the notes.

\bibitem{BostChambertLoir2009} J.-B.~Bost and A.~Chambert-Loir, \emph{Analytic curves in algebraic
varieties over number fields}, in: Algebra, Arithmetic, and Geometry (in honor of Yu.~I.~Manin),
Vol.~I, Progr.\ Math.~\textbf{269}, Birkh\"auser, 2009, 69--124; Thm.~7.9 and Example~7.10.

\bibitem{Cannon1984} J.~W.~Cannon, \emph{The combinatorial structure of cocompact discrete
hyperbolic groups}, Geom.\ Dedicata \textbf{16} (1984), 123--148.

\bibitem{Cantor1980} D.~G.~Cantor, \emph{On an extension of the definition of transfinite
diameter and some applications}, J.~Reine Angew.\ Math.\ \textbf{316} (1980), 160--207;
Thm.~5.2.2.

\bibitem{Carlson1921} F.~Carlson, \emph{\"Uber Potenzreihen mit ganzzahligen Koeffizienten},
Math. Z. \textbf{9} (1921), 1--13.

\bibitem{Christol1979} G.~Christol, \emph{Ensembles presque p\'eriodiques $k$-reconnaissables},
Theoret. Comput. Sci. \textbf{9} (1979), 141--145.

\bibitem{CKMR1980} G.~Christol, T.~Kamae, M.~Mend\`es France, G.~Rauzy, \emph{Suites alg\'ebriques,
automates et substitutions}, Bull.\ Soc.\ Math.\ France \textbf{108} (1980), 401--419.

\bibitem{DiVizioHardouin2012} L.~Di~Vizio, C.~Hardouin, \emph{Descent for differential Galois theory
of difference equations: confluence and $q$-dependence}, Pacific J. Math.\ \textbf{256} (2012),
79--104.

\bibitem{DHR2018} T.~Dreyfus, C.~Hardouin, J.~Roques, \emph{Hypertranscendence of solutions of
Mahler equations}, J.~Eur.\ Math.\ Soc.\ \textbf{20} (2018), 2209--2238.

\bibitem{DuchinShapiro2019} M.~Duchin and M.~Shapiro, \emph{The Heisenberg group is pan-rational},
Adv.\ Math.\ \textbf{346} (2019), 219--263.

\bibitem{DNNS1996} D.~Duverney, Ke.~Nishioka, Ku.~Nishioka and I.~Shiokawa, \emph{Transcendence
of Jacobi's theta series}, Proc.\ Japan Acad.\ Ser.~A \textbf{72} (1996), 202--203.

\bibitem{Exton1978} H.~Exton, \emph{A basic analogue of the Bessel--Clifford equation},
J\~n\=an\=abha \textbf{8} (1978), 49--56.

\bibitem{FischlerRivoal2016} S.~Fischler, T.~Rivoal, \emph{Arithmetic theory of $E$-operators},
J.~\'Ec.\ polytech.\ Math.\ \textbf{3} (2016), 31--65.

\bibitem{GasperRahman2004} G.~Gasper and M.~Rahman, \emph{Basic Hypergeometric Series}, 2nd ed.,
Encyclopedia of Mathematics and its Applications~96, Cambridge Univ.\ Press, 2004.

\bibitem{Grigorchuk1984} R.~I.~Grigorchuk, \emph{Degrees of growth of finitely generated groups
and the theory of invariant means}, Izv.\ Akad.\ Nauk SSSR Ser.\ Mat.\ \textbf{48} (1984),
939--985.

\bibitem{Hahn1953} W.~Hahn, \emph{Die mechanische Deutung einer geometrischen
Differenzengleichung}, Z. Angew.\ Math.\ Mech.\ \textbf{33} (1953), 270--272.

\bibitem{Humphreys1990} J.~E.~Humphreys, \emph{Reflection Groups and Coxeter Groups}, Cambridge
Studies in Advanced Mathematics~29, Cambridge Univ.\ Press, 1990; \S5.12 (Poincar\'e series of a
Coxeter group).

\bibitem{IrelandRosen1990} K.~Ireland and M.~Rosen, \emph{A Classical Introduction to Modern Number
Theory}, 2nd ed., Graduate Texts in Mathematics~84, Springer, 1990; Ch.~15 (the von
Staudt--Clausen theorem).

\bibitem{IsolaPiergallini} S.~Isola and R.~Piergallini, \emph{On the generic triangle group},
arXiv:1405.1881 [math.MG], 2014.

\bibitem{KMNOY2004} K.~Kajiwara, T.~Masuda, M.~Noumi, Y.~Ohta, Y.~Yamada, \emph{Hypergeometric
solutions to the $q$-Painlev\'e equations}, Int.\ Math.\ Res.\ Not.\ \textbf{2004}, no.~47,
2497--2521.

\bibitem{Kato} T.~Kato, \emph{Perturbation Theory for Linear Operators}, 2nd ed., Grundlehren
Math.\ Wiss.~132, Springer, Berlin, 1976; reprinted in Classics in Mathematics, 1995.

\bibitem{KoelinkSwarttouw1994} H.~T.~Koelink and R.~F.~Swarttouw, \emph{On the zeros of the
Hahn--Exton $q$-Bessel function and associated $q$-Lommel polynomials}, J. Math. Anal. Appl.
\textbf{186} (1994), 690--710.

\bibitem{KS1992} T.~H.~Koornwinder and R.~F.~Swarttouw, \emph{On $q$-analogues of the Fourier and
Hankel transforms}, Trans.\ Amer.\ Math.\ Soc.\ \textbf{333} (1992), 445--461.

\bibitem{Lang1966} S.~Lang, \emph{Introduction to Transcendental Numbers}, Addison--Wesley, 1966;
Ch.~III (the Schneider--Lang theorem).

\bibitem{Lipshitz1989} L.~Lipshitz, \emph{$D$-finite power series}, J.~Algebra \textbf{122} (1989),
353--373.

\bibitem{LorentzenWaadeland1992} L.~Lorentzen and H.~Waadeland, \emph{Continued Fractions with
Applications}, Studies in Computational Mathematics~3, North-Holland, 1992; \S IV.2 (Pincherle's
theorem).

\bibitem{Massey1969} J.~L.~Massey, \emph{Shift-register synthesis and BCH decoding}, IEEE Trans.\
Inform.\ Theory \textbf{15} (1969), 122--127.

\bibitem{MathlibCommunity} The mathlib Community, \emph{The Lean mathematical library},
CPP~2020, 367--381; \url{https://github.com/leanprover-community/mathlib4}.

\bibitem{Moak1984} D.~S.~Moak, \emph{The $q$-analogue of Stirling's formula}, Rocky Mountain J.
Math.\ \textbf{14} (1984), 403--414.

\bibitem{Morita2011} T.~Morita, \emph{A connection formula of the Hahn--Exton $q$-Bessel function},
SIGMA \textbf{7} (2011), 115, 11~pp.; arXiv:1105.1998.

\bibitem{Nesterenko1996} Yu.~V.~Nesterenko, \emph{Modular functions and transcendence questions},
Sb.\ Math.\ \textbf{187} (1996), 1319--1348.

\bibitem{Nishioka1996} Ku.~Nishioka, \emph{Mahler Functions and Transcendence}, Lecture Notes in
Mathematics~1631, Springer, 1996.

\bibitem{Niven1956} I.~Niven, \emph{Irrational Numbers}, Carus Mathematical Monographs~11,
Mathematical Association of America, 1956; Cor.~3.12 (the rational values of $\cos$ at rational
multiples of $\pi$).

\bibitem{Daalhuis1994} A.~B.~Olde Daalhuis, \emph{Asymptotic expansions for $q$-gamma, $q$-exponential,
and $q$-Bessel functions}, J. Math. Anal. Appl. \textbf{186} (1994), 896--913; \S4 (the numerically
satisfactory pair and exact $q$-Wronskians for the Hahn--Exton equation).

\bibitem{Olver} F.~W.~J.~Olver, \emph{Asymptotics and Special Functions}, Academic Press, 1974 (reprinted A K Peters, 1997); Ch.~4, Thm.~7.1 (steepest descent with error bound).

\bibitem{Parry1992} W.~Parry, \emph{Growth series of some wreath products}, Trans.\ Amer.\ Math.\
Soc.\ \textbf{331} (1992), 751--759.

\bibitem{Perron1909} O.~Perron, \emph{\"Uber einen Satz des Herrn Poincar\'e}, J. Reine Angew.
Math.\ \textbf{136} (1909), 17--37.

\bibitem{Poincare1885} H.~Poincar\'e, \emph{Sur les \'equations lin\'eaires aux diff\'erentielles
ordinaires et aux diff\'erences finies}, Amer.\ J. Math.\ \textbf{7} (1885), 203--258.

\bibitem{Polya1928} G.~P\'olya, \emph{\"Uber gewisse notwendige Determinantenkriterien f\"ur die
Fortsetzbarkeit einer Potenzreihe}, Math.\ Ann.\ \textbf{99} (1928), 687--706.

\bibitem{PutSinger1997} M.~van der Put and M.~F.~Singer, \emph{Galois Theory of Difference
Equations}, Lecture Notes in Mathematics~1666, Springer, 1997.

\bibitem{RSZ2013} J.-P.~Ramis, J.~Sauloy, C.~Zhang, \emph{Local analytic classification of
$q$-difference equations}, Ast\'erisque \textbf{355} (2013).

\bibitem{RhinViola2001} G.~Rhin and C.~Viola, \emph{The group structure for $\zeta(3)$}, Acta Arith.
\textbf{97} (2001), 269--293.

\bibitem{Siegel1929} C.~L.~Siegel, \emph{\"Uber einige Anwendungen diophantischer Approximationen},
Abh.\ Preuss.\ Akad.\ Wiss.\ Phys.-Math.\ Kl.\ \textbf{1929}, no.~1, 1--70.

\bibitem{Simon} B.~Simon, \emph{Trace Ideals and Their Applications}, 2nd ed., Math.\ Surveys
Monogr.~120, Amer.\ Math.\ Soc., Providence, RI, 2005.

\bibitem{OEIS} N.~J.~A.~Sloane et al., \emph{The On-Line Encyclopedia of Integer Sequences},
\url{https://oeis.org}, sequence A396406.

\bibitem{StampachStovicek2014} F.~\v{S}tampach and P.~\v{S}\v{t}ov\'{\i}\v{c}ek, \emph{The
Hahn--Exton $q$-Bessel function as the characteristic function of a Jacobi matrix}, Spec.\
Matrices \textbf{2} (2014), 131--147.

\bibitem{StanleyEC2} R.~P.~Stanley, \emph{Enumerative Combinatorics, Volume~2}, Cambridge Univ.
Press, 1999.

\bibitem{Stoll1996} M.~Stoll, \emph{Rational and transcendental growth series for the higher
Heisenberg groups}, Invent.\ Math.\ \textbf{126} (1996), 85--109.

\bibitem{Swarttouw1992} R.~F.~Swarttouw, \emph{The Hahn--Exton $q$-Bessel function}, PhD thesis,
Delft University of Technology, 1992; eq.~(4.3.8) (the closed-form $q$-Wronskian).

\bibitem{Wong} R.~Wong, \emph{Asymptotic Approximations of Integrals}, Academic Press, 1989
(reprinted SIAM, 2001); Ch.~II (Laplace's method and stationary phase, with error bounds).

\bibitem{Zhang2000} C.~Zhang, \emph{Transformations de $q$-Borel--Laplace au moyen de la fonction
th\^eta de Jacobi}, C.\ R.\ Acad.\ Sci.\ Paris S\'er.\ I \textbf{331} (2000), 31--34.
\end{thebibliography}

\end{document}
